% Height-source replacement: see work/remove-vkk-20261104/REPORT.md for scope and review status.
\documentclass[11pt]{article}
\usepackage[T1]{fontenc}
\usepackage[utf8]{inputenc}
\usepackage{lmodern,amsmath,amssymb,amsthm,mathtools}
\usepackage[margin=1in]{geometry}
\usepackage{booktabs,longtable,array,microtype,enumitem}
\usepackage[hidelinks]{hyperref}
\setlist{itemsep=3pt,topsep=5pt}
\setlength{\emergencystretch}{2em}
\newtheorem{theorem}{Theorem}[section]
\newtheorem{proposition}[theorem]{Proposition}
\newtheorem{lemma}[theorem]{Lemma}
\newtheorem{corollary}[theorem]{Corollary}
\theoremstyle{definition}
\newtheorem{definition}[theorem]{Definition}
\theoremstyle{remark}
\newtheorem{remark}[theorem]{Remark}
\newcommand{\Q}{\mathbb Q}
\newcommand{\Z}{\mathbb Z}
\newcommand{\N}{\mathbb N}
\newcommand{\F}{\mathbb F}
\newcommand{\Cset}{\mathcal C}
\providecommand{\mathscr}[1]{\mathcal{#1}}
\newcommand{\Sel}{\operatorname{Sel}}
\newcommand{\Sha}{\operatorname{Sha}}
\newcommand{\ord}{\operatorname{ord}}
\newcommand{\Norm}{\operatorname{N}}
\newcommand{\Tr}{\operatorname{Tr}}
\newcommand{\Trd}{\operatorname{Trd}}
\newcommand{\Close}{\operatorname{Close}}
\newcommand{\Test}{\operatorname{Test}}
\newcommand{\Val}{\operatorname{Val}}
\newcommand{\tw}[1]{2\mathbin{\uparrow\!\uparrow}#1}
\title{$\mathbb{Z}$ is Diophantine in $\mathbb{Q}$\\[4pt]\large Version 3.}
\author{}
\date{November 6, 2026}
\begin{document}
\maketitle
\begin{abstract}
We construct an integer-coefficient polynomial in
$292\,890\,231\,905$ variables whose rational zero projection is exactly
$\Z$. Its degree is less than $2\mathbin{\uparrow\!\uparrow}7$, and the
absolute value of every coefficient is less than
$2\mathbin{\uparrow\!\uparrow}8$.

We also specify a quartic with $9\,625$ variables, $82\,476$ nonzero monomials, and coefficients of magnitude below $2^{80}$,
with rational zero projection exactly $\Z$.

Well-known consequences of $\Z$ being Diophantine in $\Q$ include a refutation of Mazur's conjecture that the topological closure of the rational points of a variety over $\Q$ has finitely many components, and existence of a universal rational Diophantine polynomial. Moreover, the set $\operatorname{HTP}(\Q)$ of integer polynomials with a rational zero is many-one equivalent to the halting set $0'$.
\end{abstract}
\tableofcontents

% Begin 01-specification-model.tex
\section{The main theorem and proof structure}\label{sec:main}
This construction builds on OpenAI's preprint
\emph{Hilbert's tenth problem over the rational numbers}~\cite{openaihtp}.

We use $\N=\{0,1,2,\ldots\}$. For a reduced rational number $a/b$, with
$b>0$, put $h(a/b)=\log\max(|a|,b)$. Define $\tw1=2$ and
$\tw{(r+1)}=2^{\tw r}$; every tower is associated to the right.

\begin{theorem}[A polynomial definition of the integers]\label{thm:target}
There exists one polynomial
\[
 P\in\Z[X_0,X_1,\ldots,X_{292890231904}]
\]
such that
\begin{align*}
 \deg P&<\tw7,\\
 |[X^\alpha]P|&<\tw8\qquad\text{for every multi-index }\alpha,\\
 q\in\Z&\ \Longleftrightarrow\
 \exists(w_1,\ldots,w_{292890231904})\in\Q^{292890231904},\quad
 P(q,w_1,\ldots,w_{292890231904})=0\quad(q\in\Q).
\end{align*}
All conclusions concern the same polynomial. Here $q\in\Z$ means the
ordinary embedded integers. The ambient variable count does not assert that
every variable occurs essentially.
\end{theorem}

The file \nolinkurl{ordinary-integer-definition.json} specifies the
polynomial $P_*$. The argument first represents integer indices by points
on an elliptic curve over $\Q(\sqrt2)$. A separate value relation identifies
the ordinary rational value of an index. The point coordinates and the
integer index have distinct roles throughout.

Theorem~\ref{thm:target} refutes Mazur's conjecture that the closure of
the rational points of a variety in its real locus has finitely many
connected components. It also gives a universal rational Diophantine
polynomial, with
$\operatorname{HTP}(\Q)$ many-one equivalent to $0'$.
Section~\ref{sec:four-corollaries} proves these three consequences.

% Begin quartic introduction addition
Section~\ref{sec:quartics} gives the explicit $9\,625$-variable quartic.
Its reciprocal comparison source revises the base residual compiler of
Section~\ref{sec:degree8-original}, with a separate arithmetic proof.
The complete coefficient table and exact algebraic conversion certificate
are included with this manuscript.
% End quartic introduction addition

The final composition uses the following precise statements.
\begin{description}[style=nextline,leftmargin=1.8cm]
\item[$\mathsf G$] Canonical point codes form exactly the infinite cyclic
group $2E(\Q(\sqrt2))=\Z P_0$. The membership polynomial types arbitrary inputs;
raw addition describes the group law on canonical codes.
\item[$\mathsf M$] The actual guarded multiplication template has rational
witnesses exactly for triples with indices $a,b,ab$.
\item[$\mathsf H$] The six-contact height theorem holds with the constants
used by the precision computation and, separately, by multiplication.
\item[$\mathsf{PC}$] The explicit pole polynomial is soluble for every
input in Theorem~\ref{hyp:PC}.
\item[$\mathsf C$] The actual five-input natural polynomial and its
rational translation have the exact computation and substitution semantics
of Section~\ref{sec:compiler}, including the power template.
\item[$\mathsf B$] The templates, fixed constructors and private witness
substitutions specify the same saved expression $P_*$. The quantitative
syntax calculation applies to that expression.
\end{description}

\begin{lemma}[Assembly]\label{lem:assembly}
Assume $\mathsf G,\mathsf M,\mathsf H,\mathsf{PC},\mathsf C,\mathsf B$.
Then $P_*$ satisfies Theorem~\ref{thm:target}.
\end{lemma}
\begin{proof}
Section~\ref{sec:bridge} proves the exact ordinary-value relation using
one shared shift and disjoint private witness blocks. The literal
substitutions in $\mathsf B$ identify its existential polynomial with
$P_*$. Section~\ref{sec:size} proves the stated bounds for that same
expression.
\end{proof}

\begin{proof}[Proof of Theorem~\ref{thm:target}]
Section~\ref{sec:model}, including Proposition~\ref{prop:generator},
proves $\mathsf G$ with its canonical membership and addition formulas.
Theorem~\ref{thm:height} proves $\mathsf H$ for each of the two sets of
constants used below. Sections~\ref{sec:rc}--\ref{sec:limits}, with
Appendices~\ref{app:even-source}--\ref{app:endpoints}, prove the restricted
converse in Theorem~\ref{hyp:RC}. The simultaneous descent in
Section~\ref{sec:pole} then proves $\mathsf{PC}$. Using these results,
Section~\ref{sec:mult} proves $\mathsf M$ on all rational input codes,
including the finite false-triple guard.

Section~\ref{sec:compiler} proves $\mathsf C$ for the fixed five-input
natural polynomial and its rational model translation. Its witness
variables remain natural before that translation. Sections~\ref{sec:syntax}
and~\ref{sec:size} and Appendix~\ref{app:literal} prove $\mathsf B$, including
the fixed scalars and private witness substitutions. All six statements
therefore hold, and Lemma~\ref{lem:assembly} applies to the specified $P_*$.
\end{proof}

\section{A finite specification of the polynomial}\label{sec:syntax}
The delivered specification is
\nolinkurl{construction/ordinary-integer-definition.json}, a
$23\,179\,783$-byte JSON expression with SHA256
\begin{center}\small\ttfamily
4d639d4c44e69ea360d36d4591b8fdc82269d68b9363af0157d1e5c1174c7af6.
\end{center}
It contains 23 library templates and seven named recipes. The root recipe
is \texttt{integers}; its only parameter is the rational variable at index
zero. Every other variable is existentially quantified over $\Q$.
Sources S01--S03 fix the construction and its semantics. Their historical
specification had $23\,180\,023$ bytes and SHA256
\begin{center}\small\ttfamily
a86790bbf1095d6b1cea9c66067ad8432dd35adc320254eecde21f9f98574714.
\end{center}
The delivered expression replaces two diagnostic \texttt{sourceFile}
strings and the path key in \texttt{sourceHashes} by portable relative
paths. Direct comparison shows that every polynomial-defining field is
unchanged. Since these three provenance fields contribute no equations,
both byte representations denote exactly the same integer polynomial,
with the same variables and bounds. The shipped
\nolinkurl{construction/SOURCE_MANIFEST.json} records both identities and
the normalization; the companion \texttt{SOURCE\_MAP.json} retains the
historical source references.

\subsection{Literal operations and witness scope}
A scalar node is a variable, integer constant, sum, product, or a call to
an earlier template. A call substitutes its parameter expressions and a
private consecutive witness block into the callee's polynomial. Shared
nodes represent repeated subexpressions, not new equations. A recipe
sums squares of constraint polynomials or multiplies branch polynomials.
Over $\Q$,
\[
 \sum_i f_i^2=0\ \Longleftrightarrow\ \bigwedge_i f_i=0,
 \qquad\prod_i f_i=0\ \Longleftrightarrow\ \bigvee_i f_i=0.
\]
Private witness blocks can be moved to the front because the rational
domain is nonempty. A whole inactive product-call block may receive
arbitrary values. This does not discard fixed equations inside an active
callee. A rational constant $p/q$, $q\ne0$, is represented by a fresh
$c$ and the globally conjoined equation $qc=p$. In particular the cached
coordinates of $P_0$ satisfy
\[
 8g_0=1641,\qquad4g_1=-15465,\qquad32g_2=87483.
\]
These equations remain active outside internal disjunctions in the
multiplication, signed-encoding and exception templates. Inverses are
equations $uv=1$, never partial evaluation operations.

Strict positivity has the polynomial expression
\[
 s>0\ \Longleftrightarrow\
 \exists z_1,z_2,z_3,z_4\in\Q,\quad s\sum_i z_i^2=1.
\]
Indeed, if $1/s=a/b>0$, express $ab$ as four integer squares and divide
their coordinates by $b$. The reverse direction follows from the order
on $\Q$. Absolute value has the exact expression
\[
 a=|x|\quad\Longleftrightarrow\quad
 \exists z_1,\ldots,z_4\in\Q,\quad a=\sum_i z_i^2,\quad a^2=x^2.
\]

\subsection{Fixed finite constructors}
Three node types abbreviate finite polynomials. A
\texttt{fixed\_integer} inserts a fixed integer expression;
\texttt{psi\_component} inserts one coordinate of a division polynomial;
\texttt{finite\_guard\_product} inserts a product over a fixed box of
integers. Their bounds never depend on the free rational input.

For example, on $y^2=x^3+Ax$, write $\psi_n=f_n$ for odd $n$ and
$\psi_n=2yf_n$ for even $n$. The bases are
\begin{align*}
 f_0&=0,&f_1=f_2&=1,\\
 f_3&=3x^4+6Ax^2-A^2,&
 f_4&=2(x^6+5Ax^4-5A^2x^2-A^3).
\end{align*}
For $m\ge3$,
$f_{2m}=f_m(f_{m+2}f_{m-1}^2-f_{m-2}f_{m+1}^2)$.
For $m\ge2$, put $L=f_{m+2}f_m^3$, $R=f_{m-1}f_{m+1}^3$, and
$D_0=16(x^3+Ax)^2$. Then $f_{2m+1}=D_0L-R$ for even $m$ and
$L-D_0R$ for odd $m$. Negative indices satisfy $\psi_{-n}=-\psi_n$.
This is a terminating recursion. These are the short-Weierstrass
normalizations of \cite[\S\S5.5--5.6, Lemmas 5.20 and 5.24, Theorem 5.21]{sutherland}.
In its Jacobian-coordinate convention,
$x([n]P)=\phi_n(P)/\psi_n(P)^2$ and
$y([n]P)=\omega_n(P)/\psi_n(P)^3$.
The numerator formula has a minus sign:
\[
 \omega_n=\frac{\psi_{n+2}\psi_{n-1}^2-
                   \psi_{n-2}\psi_{n+1}^2}{4y},
 \qquad \psi_{2n}=2\psi_n\omega_n.
\]
Thus the fixed-scalar ordinate equation is
$2y([n]P)\psi_n(P)^4=\psi_{2n}(P)$.
For $n\ne0$ and a nonidentity point in the torsion-free carrier,
$\psi_n(P)\ne0$, so clearing these denominators is reversible.

The fixed prime catalogue consists of the first 22 primes above 324 in
each unit residue class modulo $W=4\prod_{p\le324}p$. Dirichlet's theorem
ensures termination. Put $N_s=(\prod_{p\le324}p)H_{\rm cat}$, where
$H_{\rm cat}$ is the catalogue product. Assign inverse abscissa zero
at the identity. For the normalized representation
$1/x([m]P_0)=(a+bs)/c$ with $a,b,c\in\Z$, $c>0$, and
$\gcd(a,b,c)=1$, define $M_0$ to be the least $m\ge1$ such that
$N_s^4\mid a,b$ and $\gcd(c,N_s)=1$ (S31).
This is a fixed integer, with a successful index bounded in
Section~\ref{sec:size}.
The fixed point coordinates in the guard have nonzero denominators by
the infinite order of $P_0$. Thus finite expansion gives one integer
polynomial without requiring a monomial expansion. Its arithmetic zero
set remains a separate issue. In particular, bounding the integer $T$
does not prove that it excludes all false multiplication triples.

\section{The elliptic model of integer indices}\label{sec:model}
Let
\[
 K_0=\Q(s),\quad s^2=2,\quad d=75-53s,\quad E:y^2=x^3-d^2x,
\]
and fix
\[
 P_0=\left(\frac{1641-1160s}{8},\frac{-123720+87483s}{32}\right).
\]
The identity has code $(1,0,0,0,0)$. An affine point has code
$(0,x_0,x_1,y_0,y_1)$ with $x=x_0+x_1s$, $y=y_0+y_1s$.
Linear independence of $1,s$ makes this unique. Write $\iota(n)$ for
the code of $[n]P_0$.

\begin{proposition}[Generator certificate]\label{prop:generator}
The point $P_0$ has infinite order and $2E(K_0)=\Z P_0$.
\end{proposition}
\begin{proof}
We give the argument of S05--S07, making its finite certificate explicit.
The ring of integers $\Z[s]$ has class number one because its Minkowski
bound is $\sqrt2<2$. It has four unit squareclasses. The sign vectors
of $-1,1+s$ are independent, so every totally positive unit is a square.

The $x$-squareclass map in the two-isogeny descent on $E$ has image
contained in the eight classes generated by $-1,1+s,d$. Away from $(d)$,
odd $v(x)$ makes the minimum of $3v(x)$ and $v(x)+2v(d)$ odd, impossible
for $y^2$. The torsion points give $1,-1,d,-d$. Put $\epsilon=1+s$,
$\delta=3-s$. The change $(x,y)\mapsto(\epsilon^4x,\epsilon^6y)$ takes
$E$ to $C:y^2=x^3-\delta^2x$, and $R=(1+s,2s)$ provides the remaining
unit class. Thus the first image has order eight.

On $E':y^2=x^3+4d^2x$, a nonzero $x$ is positive in both real embeddings.
Its valuations are even. At $(s)$ the two terms have valuations $3m$
and $m+4$, with equality only at $m=2$. At $(d)$ the only odd case
would require $u^2=-4$ in $\F_7$. Thus the second image is trivial.
The exact two-isogeny sequence gives $[E(K_0):2E(K_0)]=8$; its possible
extra kernel term vanishes since a nonkernel rational two-torsion point
maps to the dual kernel.

The good reductions at $(17,s-6)$ and $(23,s-5)$ have 20 and 24 points.
Prime-to-residue-characteristic torsion injects, so the torsion group
has order at most four, hence is exactly the displayed full two-torsion. Mordell--Weil now gives rank one.
The $x$-class of $R$ has opposite real signs, whereas all torsion classes
have equal signs. Its free index is therefore odd.

The duplication forms on $C$ are
\[
 F_2(a,b)=(a^2+\delta^2b^2)^2,\qquad
 G_2(a,b)=4ab(a^2-\delta^2b^2).
\]
At either real embedding their maximum lies between
$\max(|a|,|b|)^4$ and $676\max(|a|,|b|)^4$. For primitive $a,b$ their
common ideal has valuations at most six at $(s)$ and four at $(\delta)$,
and none elsewhere. For the latter bound use that $-1$ is nonsquare
modulo seven. For the former, a local unit is $A+Bs$ with $A$ odd, and
$v_s((A+Bs)^2+1)$ is two or three according to the parity of $B$.
Hence
\[
 4h_x(T)-\log392\le h_x(2T)\le4h_x(T)+\log676.
\]
With $2\widehat h(T)=\lim_m4^{-m}h_x(2^mT)$ this gives
$h_x(T)\le2\widehat h(T)+\log392/3$. Direct doubling gives
$x(2R)=(57-28s)/8$, with height $\log41$. If $R=mS+T$ for torsion
$T$ and $|m|\ge3$, it follows that
\[
 h_x(S)\le\frac{\log(41^3\cdot676)}{108}+\frac{\log392}{3}<\log9.
\]
The strict inequality is the exact integer comparison
$41^3\cdot676\cdot392^{36}<9^{108}$.

A nonrational $x\in K_0$ in this height range has primitive polynomial
$aX^2+bX+c$ with $1\le a\le81$, $|b|\le162$, $|c|\le81$, and
discriminant $2k^2$, $1\le k\le162$. Both roots are included.
Rational possibilities have coprime numerator of absolute value at most
nine and denominator from one to nine. The exact finite square test of
$x^3-\delta^2x$ gives precisely
\[
 0,\quad\delta,\quad-\delta,\quad1+s,\quad23-17s,
 \quad2+4s,\quad(5-4s)/2.
\]
The certificate covers 7,004 primitive quadratics, their 14,008 roots,
and 111 rational candidates. To test whether $u+vs$ is a square, take
both signs of a rational square root $t$ of $u^2-2v^2$, test
$y_0^2=(u+t)/2$, then set $y_1=v/(2y_0)$ if $y_0\ne0$.
The $y_0=0$ case requires $v=0$ and $y_1^2=u/2$. Substitution verifies
each result. This is a necessary and sufficient test, including zero.
The last four listed coordinates are those of $R+T$ for $T\in C[2]$.
Every surviving nontorsion point therefore has height $\widehat h(R)$,
contradicting $\widehat h(R)=m^2\widehat h(S)$. This finite computation
is an explicit part of the proof, with a recorded independent replay;
it is not replaced by a rank oracle.

Thus $2C(K_0)=\Z(2R)$. The inverse coordinate change sends $2R$ to
$P_0$, as asserted.
\end{proof}

\subsection{Polynomial membership and addition}
Apart from the identity branch, membership in $2E(K_0)$ is expressed by
choosing $X,Y,H\in K_0$, $u\in\Q$, and imposing
\[
 Y^2=X^3-d^2X,\quad u\Norm(Y)=1,\quad2YH=3X^2-d^2,
 \quad x=H^2-2X,\quad y=H(X-x)-Y.
\]
Every nonidentity double has a preimage with nonzero ordinate.
Each field equation is two rational equations. Membership is an all-input
statement. Raw addition has the narrower exact contract
\[
 \operatorname{Add}_{\rm raw}(\iota(a),\iota(b),R)
             \quad\Longleftrightarrow\quad R=\iota(a+b).
\]
Its branches cover both identities, opposite points, distinct abscissas
and doubling with nonzero ordinate. Every later use supplies canonical
input codes. Multiplication first imposes membership of all three roots
and is sound on arbitrary rational tuples. In the natural predicate,
four initially untyped tuples enter multiplication calls, so are typed
before their squared values are added. Lagrange's theorem then gives
exactly $\iota(\N)$. This uses multiplication completeness; the generator
proposition alone does not supply the natural predicate.

% End 01-specification-model.tex


% Begin 02-height-pole-multiplication.tex
\section{The six-contact height estimate}\label{sec:height}
Set
\[
 \Cset=\{0,-15,-20,-24,-30,-60\},\qquad
 F(u)=u(u+15)(u+20)(u+24)(u+30)(u+60),
\]
and $G(u)=7u^2+300u+3600$. For a finite prime set $S$ containing all
primes at most 324 and $u\notin\Cset$, define
\[
 M_S(u)=\prod_{\substack{p\notin S:\ v_p(u-c)>0\text{ odd}\\
                          \text{for some }c\in\Cset}}p.
\]
Negative valuations of $u$ do not enter the product. Put
\begin{align*}
 F_0&=324!,\qquad D=F_0\,3^{256F_0^4},\qquad G_0=116064000000,\\
 E_0&=10(DG_0)^2,\qquad Q_R=R^{RE_0},\\
 B_{\rm cm}&=D^2\,3^{21000D^2}+1,\\
 C_{\rm cm}&=2B_{\rm cm}+4096(1028+5D+84000D^2),\\
 H_6(S;R)&=2\cdot10^{13}\max(1,C_{\rm cm},Q_R^{17}),\\
 &\hspace{1em}R\ge\max(D,\max S,10^{12}).
\end{align*}
When the second argument is omitted, we use
$H_6(S)=H_6(S;\max(D,\max S,10^{12}))$.
The geometric construction is followed by a direct modularity argument
at coefficient prime five and a separate effective CM argument. The
constants in the displayed definition of $H_6(S;R)$ are unchanged.

\begin{theorem}[Six-contact estimate]\label{thm:height}
For every such $S,R$ and every $u\in\Q\setminus\Cset$,
\[
 h(u)\le H_6(S;R)M_S(u)^{24}.
\]
\end{theorem}
\begin{proof}
Take $B=(2,-3)$, the maximal order $\mathcal O=\Z[i,(1+j)/2]$, and
$\mu=2j+ij$. With its positive sign choice, the form
$\Trd(\mu^{-1}x\overline y)$ has matrix
\[
 \begin{pmatrix}0&0&-1&1\\0&0&1&-2\\1&-1&0&0\\-1&2&0&0\end{pmatrix}
\]
and determinant one, in the ordered basis
$(1,i,(1+j)/2,(i+ij)/2)$. The Riemann form is the convention of
\cite[equation (3) and Lemma 3]{linyang}; its Rosati involution is
$a\mapsto\mu^{-1}\overline a\mu$. This is the principal-polarization datum whose
normalizer quotient is $\mathbf P^1_j$, with branch values
$0,\infty,-16/27$ and indices $4,6,2$ \cite{babagranath,linyang}.
Put $K_0=\widehat{\mathcal O}^{\,*}$ and
$K_3=\ker(K_0\longrightarrow(\mathcal O/3\mathcal O)^*)$.
Let $N$ be the finite-adelic normalizer of $\widehat{\mathcal O}$.
It normalizes $K_3$. At a split prime the local normalizer is the
center times the maximal-order units; at each of the two division
primes its quotient by these groups has order two. Thus
$|N/(\mathbb A_f^*K_0)|=4$.
A central finite idele acts trivially on the level-three curve:
multiply it by a rational scalar to remove all finite valuations,
and choose the scalar's sign so that the residue at three is one.
The remaining idele lies in $K_3$, and the rational central scalar
acts trivially on the symmetric domain. Consequently the entire
level-three curve $Y_3$, with all its determinant components, admits
the action of $N/(\mathbb A_f^*K_3)$ with quotient $\mathbf P^1_j$,
and
\[
 \deg(Y_3\longrightarrow\mathbf P^1_j)
 \le4|(\mathcal O/3\mathcal O)^*|\le4\cdot81=324.
\]

We construct its integral model without changing this degree.
Apply \cite[Theorem 2.1 and Corollary 2.3]{buzzard} to $V_1(5)$
over $\mathbb Z[1/30]$. Its determinant image is full, its level
is maximal at the division primes, and $5\ge4$ is prime to six.
First add a full five-marking refining the $V_1(5)$-marking.
This gives the full-five moduli scheme, including every determinant
component, as a finite \'{e}tale cover. Next add the finite \'{e}tale sheaf
\[
 \operatorname{Isom}_{\mathcal O}
       ((\mathcal O/3\mathcal O),A[3]).
\]
This sheaf is nonempty on every geometric fiber. Indeed $B_3$ is
a division algebra and the rational Tate module is a rank-one
left $B_3$-module. Every $\mathcal O_3$-stable lattice in that module
is principal: choose a nonzero lattice element of least division
valuation and divide every other element by it on the right.
Thus the Tate lattice is free of rank one over $\mathcal O_3$.
One can also check this in characteristic zero and extend it to
all fibers by local constancy of the finite \'{e}tale $A[3]$ with its
$\mathcal O$-action.

The resulting smooth projective scheme $X_{35}$ carries full
markings at three and five. The group
$G_5=(\mathcal O/5\mathcal O)^*$ changes only the five-marking.
Its action is free: a stabilizer would give a polarized automorphism
preserving the full three-marking, hence a finite-order element of
$1+3M_4(\mathbb Z_3)$. That group is torsion-free, as follows by
expanding a hypothetical prime-order element at its first nonzero
$3$-adic term. The automorphism is therefore the identity, and
the full five-marking forces the group element to be the identity.
The quotient $X_{35}/G_5$ is smooth projective, its covering map is
finite \'{e}tale, and the universal abelian scheme, its canonical
principal polarization, its $\mathcal O$-action and its three-marking
descend. This quotient represents exactly the full-three moduli
problem, so its generic fiber is $Y_3$. The auxiliary five-marking
has contributed no degree factor to $Y_3\to\mathbf P^1_j$.

Each geometric component of $Y_3$ is a torsion-free cover of the
norm-one orbifold of Euler characteristic $-1/3$, and hence has genus
at least two. The normalizer action extends to its smooth proper
model by uniqueness of the stable model, including permutations of
components. Fix $p>324$ and work after strict henselization.
The action is tame. On each component use its effective action.
A finite-order automorphism of order prime to $p$ cannot specialize
to the identity, by considering its first nonzero term modulo a
power of the uniformizer. At a fixed point the same power-series
argument shows that the tangent character of the stabilizer is
faithful. Thus the stabilizer is cyclic.

For a nonidentity generator the tangent eigenvalue minus one is a
unit. Its fixed-point equation therefore has a unique lift by
Hensel's lemma. All powers of the generator fix that lift. It follows
that no new branch orbit appears and that distinct branch orbits
cannot collide: their stabilizers at a purported common specialization
would belong to one cyclic group, whose nontrivial elements have the
same unique fixed lift. Averaging a parameter gives a local invariant
ring $R[[z^e]]$. Hence the quotient is smooth and retains the three
disjoint branch sections and their indices $4,6,2$.

Its genus-zero model with these three ordered sections is isomorphic
to $\mathbf P^1_{\mathbb Z_p}$ carrying the ordered triple
$(0,\infty,-16/27)$, which is disjoint since $2$ and $3$ are units.
The isomorphism agrees with $j$ on the generic fiber, by uniqueness
of a projective coordinate on three points. Thus it is the integral
model in the stated $j$-coordinate; no further denominator prime
has been discarded.

Pull back by $J(u)=-57600F(u)^2/G(u)^6$. If
\[
 H(u)=(u^2-60u-1800)(13u^2+480u+3600)(11u^2+600u+7200),
\]
direct calculation gives $G^6-97200F^2=H^2$. The pullback indices are
$2,6,2$. Tame normalization replaces $e$ by $e/\gcd(e,a)$ and leaves
$2,1,1$. The discriminants and resultants have prime support at most 13,
so this persists at $p>324$. The normalized cover of $\mathbf P^1_u$
has degree at most 324, not $12\cdot324$, and is branched only at the
six contacts. Locally it has equation $u-c=z^2$ after unramified
extension. Positive even contact valuation is unramified; negative
valuation specializes to the unbranched point at infinity. This is
exactly the support specified by $M_S(u)$.

Let $K$ be the splitting field of the fiber over $u$ of the normalized
cover. The fiber is finite \'{e}tale of degree at most $324$, so $K/\mathbb Q$
is normal with $[K:\mathbb Q]\le F_0=324!$. The preceding local
calculation makes it unramified outside
$T=S\cup\operatorname{supp}M_S(u)$. Choose a fiber $A/K$.
The proper model and its universal abelian scheme give good reduction
outside the places over $S$.

All conjugates of $A$ have the same geometric polarized moduli point.
Indeed Baba--Granath's polarized equivalence relation gives
$V_{\mathbb Z}=V/W$; for discriminant six its two twisting involutions
generate the whole Atkin--Lehner group $W$. On the finite nonzero
$j$-locus the polarized point is equivalently specified by
\[
 [J_2,J_4,J_6,J_{10}]
 =[12(j+1),6(j^2+j+1),4(j^3-2j^2+1),j^3].
\]
Here $j=J(u)$ is rational, finite and nonzero: $F(u)\ne0$ and
$\operatorname{disc}(G)=-10800<0$.

To define the Hom field without choosing a Galois action on a
labelled collection, form $W_0=\operatorname{Res}_{K/\mathbb Q}A$.
Its geometric endomorphism lattice is the direct sum of the ordered
Hom groups among the conjugates. Its rank $r$ is at most
$16F_0^2$, by its faithful integral-homology realization.
Its Galois image is finite, and every finite subgroup of
$\operatorname{GL}_r(\mathbb Z)$ injects modulo three. If $E$ is
the kernel field of this action and $L=EK$, then $L/\mathbb Q$ is
normal and
\[
 [L:\mathbb Q]\le F_0\,3^{r^2}
 \le F_0\,3^{256F_0^4}=D.
\]
At a prime outside $T$, the field $K$ is unramified and every
conjugate has good reduction. Choose a different Tate prime.
Inertia acts trivially on the Tate module of $W_0$, and therefore
on its endomorphism lattice. Thus $L$ is unramified outside $T$.
It defines every geometric endomorphism and every ordered
conjugacy homomorphism of $A$.

By \cite[Part I, Theorems 2.8--2.9]{milne-av}, polarizable complex
tori are abelian varieties and their analytic homomorphisms are algebraic.
After clearing lattice denominators, the same assertion identifies
rational Hodge endomorphisms with algebraic quasi-endomorphisms.
These endomorphisms are defined over $\overline{\Q}$: every automorphism
of $\mathbb C/\overline{\Q}$ fixes their action on all torsion points, and
torsion is Zariski dense.
On rational homology the complex structure is right multiplication
by an element $J\in B\otimes\mathbb R$ with $J^2=-1$. In
$\operatorname{End}_{\mathbb Q}(B)=B\otimes B^{\mathrm{op}}$,
commutation with $J$ requires each right coefficient to commute
with $J$. Its rational centralizer is either $\mathbb Q$ or an
imaginary quadratic field $k$. Hence the geometric endomorphism
algebra is either $B$ or $M_2(k)$.
We treat these two cases separately. Throughout, put
$M=M_S(u)$, $d=[L:\Q]\le D$, and $Q=Q_R$.

\paragraph{CM fibers.}
Suppose $\operatorname{End}^0(A_{\overline\Q})=M_2(k)$.
All its elements are defined over $L$. Multiples of matrix units
therefore give an elliptic curve $E/L$ with complex multiplication
and an $L$-isogeny $E^2\to A$. Let $\Delta<0$ be the discriminant
of the full geometric endomorphism order of $E$. The degree of
$j(E)$ is its order class number, so $h(\Delta)\le d\le D$.
The effective bound of \cite[Proposition 2.7]{fowler} gives
\[
 \sqrt{|\Delta|}\le2D^2\exp(21000D^2).
\]
This bound includes every imaginary quadratic order and has no
exceptional-field exclusion.

Here is the height comparison with its normalization explicit.
After a finite extension $E$ has everywhere good reduction, since
its $j$-invariant is integral. This extension does not change its
stable Faltings height. For a reduced CM period $\tau=x+iy$ we have
$\sqrt3/2\le y\le\sqrt{|\Delta|}/2$. The archimedean summand in
the original Faltings normalization is
\[
 -\log(2\pi)+\frac{\pi y}{6}-\frac12\log y
       -2\sum_{n\ge1}\log|1-e^{2\pi in\tau}|,
\]
by the discriminant formula \cite[Eq. (6.2)]{cmheight}.
For $\rho=e^{-2\pi y}<1/100$, the final sum is at most
$2\rho/(1-\rho)^2<200/9801$, and
$-\tfrac12\log y<1/10$. Averaging and discarding the negative
$-\log(2\pi)$ gives
\[
 h_F(E)\le\frac{\pi\sqrt{|\Delta|}}{12}+1
       \le D^2 3^{21000D^2}+1=B_{\rm cm}.
\]
Apply the published isogeny estimate
\cite[Theorem 1.4]{gr2014} over the original field $L$, with
source $E^2$. It gives an isogeny to $A$ of degree at most
\[
 \bigl(28^{256}d\max(h_F(E^2),\log d,1)^2\bigr)^{8192}.
\]
The maximum is at most $2B_{\rm cm}$. Height variation under
isogeny \cite{faltings1983} and additivity give
\begin{align*}
 h_F(A)&\le2B_{\rm cm}
   +4096\bigl(256\log28+\log d+2\log(2B_{\rm cm})\bigr)\\
 &\le2B_{\rm cm}+4096(1028+5D+84000D^2)=C_{\rm cm}.
\end{align*}
For the last line use $\log28<4$, $\log d\le D$, and
$\log(2B_{\rm cm})\le2+2D+42000D^2$.
No bound for the degree of the auxiliary good-reduction extension
enters the isogeny estimate.

\paragraph{Non-CM fibers and scalar descent.}
Now suppose $\operatorname{End}^0(A_{\overline\Q})=B$.
In particular $A$ is absolutely simple. For each
$\sigma\in\operatorname{Gal}(L/\Q)$ choose a geometric isomorphism
$\psi_\sigma:{}^\sigma A\to A$. The automorphism
$b\mapsto\psi_\sigma{}^\sigma b\psi_\sigma^{-1}$ of $B$ is inner.
Correcting it by Skolem--Noether and clearing denominators gives
a $B$-linear isogeny $\mu_\sigma:{}^\sigma A\to A$, defined over
$L$ because all the relevant Hom groups are defined there. Take
$\mu_1=1$ and inflate this finite choice to $G_\Q$.
The discrepancy
\[
 c(\sigma,\tau)=
 \mu_\sigma{}^\sigma\mu_\tau\mu_{\sigma\tau}^{-1}
\]
commutes with $B$, hence belongs to $\Q^\times$.
Set $a_\sigma=(\deg\mu_\sigma)^{-1/4}>0$ and
$\nu_\sigma=a_\sigma\mu_\sigma$ in the category with algebraic
coefficients. Throughout these coefficient extensions Galois acts on
the varieties and morphisms and acts trivially on the added scalars;
the induced homology maps are extended linearly. In particular, the
$a_\sigma$ are constant coefficients in the products below.
Multiplicativity of degree and
$\deg[c]=c^4$ show that
\[
 \nu_\sigma{}^\sigma\nu_\tau
       =\varepsilon(\sigma,\tau)\nu_{\sigma\tau},
 \qquad \varepsilon(\sigma,\tau)\in\{1,-1\}.
\]
Thus $\varepsilon$ is a normalized sign cocycle inflated from
$\operatorname{Gal}(L/\Q)$. Only finitely many algebraic
coefficients $a_\sigma$ occur.

We need a splitting with controlled ramification, not merely
existence of a projective lift. Identify
$H^2(G_\Q,\{1,-1\})$ with $\operatorname{Br}(\Q)[2]$.
At a finite prime outside $T$ the class of $\varepsilon$ is zero:
its restriction is unramified and the unramified local quotient
has cohomological dimension one. Let $T_-$ be the finite primes
where its invariant is $1/2$. For each odd $p\in T_-$ choose an
odd character of $(\Z/p\Z)^\times$; at $2\in T_-$ use the
nontrivial character modulo four. Their product is a Dirichlet
character $\chi$, unramified outside $T$, whose order $m$ satisfies
\[
 m\le2\prod_{\substack{p\in T\\p\ne2}}(p-1)
       \le2R^R M.
\]
Here the product over an empty set is one.

To check that $\chi$ has the required square-root obstruction,
use the exact sequence of constant discrete Galois modules
$1\to\{\pm1\}\to\mu_\infty\xrightarrow{z\mapsto z^2}
\mu_\infty\to1$. By local reciprocity, a finite-order character
of $G_{\Q_p}$ has a character square root exactly when its value
on $-1\in\Q_p^\times$ is $1$. Indeed this is the only obstruction
on the torsion part of $\Q_p^\times$; the valuation factor and
the procyclic principal-unit factors admit such roots. Hence the
boundary of $\chi$ and $[\varepsilon]$ agree at every finite
place. Brauer reciprocity makes them agree at the real place,
and the injective global-to-local Brauer map makes them equal
globally \cite[Chapter I]{milne}.
Choosing a continuous pointwise square root of $\chi$ and
correcting it by a sign cochain consequently gives a continuous
cochain $\beta:G_\Q\to\mu_{2m}$ such that
\[
 \beta^2=\chi,\qquad
 \beta(\sigma)\beta(\tau)\beta(\sigma\tau)^{-1}
        =\varepsilon(\sigma,\tau).
\]
Its restriction to $G_L$ is a character invariant under $G_\Q$.
Outside $T$, its inertia restrictions are quadratic and nontrivial
at only finitely many primes. Twisting by the quadratic character
of the product of those odd primes kills all these restrictions:
at an odd prime there is a unique nontrivial quadratic character
on inertia. This introduces ramification only at those primes
and possibly at $2\in T$, and changes neither displayed identity.
We henceforth use this choice of $\beta$.

Let $L'$ be the kernel field of $\beta|_{G_L}$.
Conjugacy invariance shows that $L'/\Q$ is Galois, and
\[
 [L':\Q]\le2md\le4DR^R M\le QM.
\]
The last inequality follows from $R\ge\max(D,10^{12})$ and
$E_0=10(DG_0)^2$.

\paragraph{Modularity and a uniform conductor bound.}
Choose one number field $C$ containing all $a_\sigma$, the values of
$\beta$, and a splitting field of $B$. For each coefficient prime
$\ell$ choose an embedding of $C$ into $\overline{\mathbb Q}_\ell$
and extend scalars on $V_\ell(A)$ to a finite extension containing
its image. With the constant-coefficient convention above, define
\[
 \rho_\ell(\sigma)
   =\beta(\sigma)^{-1}a_\sigma\mu_\sigma\circ\sigma.
\]
The cocycle identities make this a representation of $G_\Q$.
It commutes with $B$, so matrix units give
$\rho_\ell\simeq r_\ell\oplus r_\ell$ with $\dim r_\ell=2$.
On $G_{L'}$ it is the actual Tate action of $A$.
Faltings' semisimplicity and endomorphism theorems
\cite{faltings1983} show that $r_\ell|_{G_{L'}}$ is absolutely
irreducible: after splitting $B$, its commutant is the scalars.
Lemma~\ref{lem:de-rham-descent} below shows that the representation
is de Rham at $\ell$, by proper-smooth comparison and finite descent
\cite{scholze2013,scholze-erratum}, and that its
two Hodge--Tate weights are $0,-1$ in the Tate convention.
Berger's $p$-adic monodromy theorem \cite{berger2002} makes this
de Rham representation potentially semistable, as required by
Pan's theorem.
It is odd. To see this without an assumption on a residual
representation, compare with Betti homology at a complex
embedding. Take the map on rational Betti homology induced by
antiholomorphic conjugation followed by the holomorphic
quasi-isogeny $\mu_c$, and extend it $\mathbb C$-linearly.
It exchanges the two Hodge summands. Each
summand is a simple module for $B\otimes\mathbb C=M_2(\mathbb C)$.
Consequently its action on the two-dimensional multiplicity
space is not scalar; since $r_\ell(c)^2=1$, its eigenvalues
are $1,-1$.
It is unramified outside $T\cup\{\ell\}$, since there $L$ is
unramified, $A$ has good reduction, and $\beta$ is unramified.

Use the Tate convention in which the cyclotomic character has
Hodge--Tate weight $-1$. Apply \cite[Theorem 1.0.4]{pan2022} at
$\ell=5$; there is no residual-image proviso at this prime.
It gives a cuspidal eigenform up to character twist. The twisting
character is a de Rham subrepresentation of the Hom representation
and is unramified outside finitely many primes. Local algebraicity
\cite[Appendix B, pp.~32--33]{conrad-local-characters} and class field
theory over $\mathbb Q$ make it a finite character times an integral
power of the cyclotomic character: removing its integral local
Hodge--Tate weight leaves finite inertia at five and at every other
prime, hence a finite ray-class character. Comparing $\{0,-1\}$
with the modular weights $\{0,1-k\}$ after that power forces
$k=2$ and power zero. Absorb the finite Dirichlet character into
the eigenform and take its primitive newform $f$. Li's Fricke formula
\cite[p.~296, before Lemma 6]{li-newforms} makes the span of its
coefficient conjugates Fricke-stable. Shimura's quotient construction
and Euler-factor identity \cite[Theorem 1]{shimura-factors} therefore
show that $r_5$ is a constituent of the Tate module of the modular abelian variety
$A_f$, or equivalently $r_5^\vee$ is its cohomological constituent.
This duality preserves the conductor.

The representations at five and three belong to this same system.
Enlarge $C$ to contain the coefficients of $f$, retaining the chosen
embedding at five. For $p\notin T\cup\{3,5\}$, choose a place $v$
of $L$ and a lift $\sigma\in G_{\mathbb Q_p}$ of the residue
field map $x\mapsto x^p$ in its decomposition group. Since
$L_v/\mathbb Q_p$ is unramified and $A$ has good reduction,
relative $p$-Frobenius maps the reduction $\overline A_v$ to
the reduction of ${}^\sigma A$ at $v$. Composing with
$\overline\mu_\sigma$ gives an endomorphism of $\overline A_v$.
The characteristic polynomial of an abelian-variety endomorphism
on rational Tate modules is independent of the coefficient prime.
Thus half its trace, multiplied by $\beta(\sigma)^{-1}a_\sigma$,
is one element of $C$ whose images are the traces of every $r_\ell$.
Away also from the finite set of primes dividing the level of $f$,
the equality with its modular trace at five is an equality of
algebraic numbers. Transporting it through an embedding at three,
Chebotarev and semisimplicity identify $r_3$ with the corresponding
modular Tate representation. We may therefore use $r_3$ to compute
the conductor at five and $r_5$ at all other primes.

Here are uniform local bounds, including the small primes.
For a finite extension of $\Q_p$ of degree at most $a$, its
normal closure has degree at most $a!$. For a finite Galois
local extension with ramification index $e$ and different
exponent $\delta$, let $f$ be the Eisenstein polynomial of a
uniformizer $\pi$ over the maximal unramified subfield. Then
$\delta\le e-1+e v_p(e)$. The derivative terms have distinct
valuations modulo $e$, so there is no cancellation at the
smallest valuation. The identity
$f'(\pi)=\prod_{\sigma\ne1}(\pi-\sigma\pi)$ over the inertia
conjugates gives $\delta=\sum_{i\ge0}(|G_i|-1)$. Changing variables
by the Herbrand function yields
$\delta/e=1-1/e+\int_0^\infty(1-|G^u|^{-1})\,du$.
Before the last positive upper break the integrand is at least
$1/2$, so that break is bounded by
$2\delta/e\le2(1+\log_2(a!))\le4a^2$
\cite[Sections 4.1--4.3]{serrelocal}.

At $p$, take $F=L_v$ and adjoin $A[3]$, using $A[5]$ if
$p=3$. This extension has degree at most $DG_0$ over $\Q_p$.
Raynaud's criterion \cite[Theorem 3.5]{silverbergzarhin}
gives semistable reduction there. Its wild inertia acts
trivially on each $V_\ell(A)$ with $\ell\ne p$.
The preceding break bound and the subgroup rule for upper
numbering therefore show that all absolute ramification
groups above $U_0=4(DG_0)^2$ act trivially on that Tate module.
On these groups $a_\sigma\mu_\sigma=1$ because they lie in $G_L$.
For odd $p$, the character $\beta|_{G_F}$ is trivial on wild
inertia: its square is the tame character $\chi|_{G_F}$, and a
pro-$p$ group has no quotient of order two. At $p=2$, local
reciprocity and the norm map give $\beta^4=1$ on the units of
$F$, since the $2$-primary local factor of $\chi$ has conductor
at most four. If $e=v_F(2)\le D$, the logarithm for $\widehat{\mathbf G}_m$ identifies
$U_F^{e+1}$ with $\mathfrak m_F^{e+1}$ since $e+1>e$
\cite[IV.6.4, p.~132]{silverman}. Multiplication by four gives
$(U_F^{e+1})^4=U_F^{3e+1}$. Thus the relative upper break of
$\beta$ is at most $3e$. To make the change of numbering explicit,
the same normal-closure estimate bounds the last upper break of
$F/\mathbb Q_2$ by $4D^2$. Above it, the absolute upper group
$G_{\mathbb Q_2}^{u}$ lies in $G_F$, and the subgroup formula is
\[
 G_{\mathbb Q_2}^{u}\cap G_F
       =G_F^{\psi_{F/\mathbb Q_2}(u)},
 \qquad \psi_{F/\mathbb Q_2}(u)\ge u\quad(u\ge0).
\]
Consequently $\beta$ is trivial on every absolute upper group with
$u>\max(4D^2,3e)$, which is at most $U_0$.
This proves the asserted absolute bound without a factor involving
$[L':\mathbb Q]$.
The Artin conductor exponent of $r_\ell$ is consequently
at most $2+2U_0\le E_0$.

For $p>R$, the extension $L_v/\Q_p$ is tame, since its
degree is at most $D<p$. The surface has good reduction,
and the same argument kills the wild part of $\beta$.
Hence the conductor exponent is at most two there.
Outside $T$ it is zero. Local--global compatibility for the
weight-two modular representation therefore bounds the
primitive level $N$ of $f$ by
\[
 N\le\Bigl(\prod_{p\le R}p^{E_0}\Bigr)M^2
      \le R^{RE_0}M^2=QM^2.
\]

\paragraph{Transfer from a modular Jacobian.}
Over $L'$ the Tate modules of $A$ and $J=J_1(N)$ have a common
constituent. Invariants in their finite-dimensional Hom space
commute with extension of the coefficient field. Faltings'
Hom theorem thus gives a nonzero $L'$-homomorphism $A\to J$.
Absolute simplicity and Poincar\'e reducibility give an
$L'$-isogeny $A\times Z\to J$ for an abelian variety $Z/L'$.
In particular $g=\dim J\ge2$.
The index of $\Gamma_1(N)$ is at most $N^2$, so both the
genus and the degree of its standard Belyi map are at most
$N^2$. By \cite[Theorem 1.1.1 and Section 2.3]{javanpeykar},
\[
 h_F(J)\le13000000N^{12}\le Q^{13}M^{24},\qquad
 g\le Q^3M^4,\qquad d'=[L':\Q]\le QM.
\]
The normalization in that theorem is the original stable
Faltings height of the Jacobian.

Apply \cite[Theorem 1.4]{gr2014} with source $J$ and target
$A\times Z$, both over $L'$. There is such an isogeny $\varphi$
with
\[
 \log\deg\varphi\le1024g^3
 \bigl(64g^2\log(14g)+\log d'
    +2\log\max(h_F(J),\log d',1)\bigr).
\]
For $Q\ge10^{12}$ and $M\ge1$, the three logarithms in order
are at most $5QM^4$, $2QM^4$, and $19QM^4$.
These follow from $\log Q\le Q$ and $\log M\le M^4/4$.
It follows that
\[
 \log\deg\varphi\le368640Q^{16}M^{24}
                       \le Q^{17}M^{24}.
\]
The general lower bound $h_F(Z)\ge-\tfrac32\dim Z\ge-2g$
is recorded in \cite[Section 9]{gr2014}.
Additivity and height variation now give
\[
 h_F(A)\le Q^{13}M^{24}+\tfrac12Q^{17}M^{24}+2Q^3M^4
          \le Q^{17}M^{24}.
\]
Thus both cases prove the exact estimate needed below:
\[
 \max(1,h_F(A))
      \le\max(1,C_{\rm cm},Q_R^{17})M_S(u)^{24}.
\]


Use the Hauptmodul normalized by the values $0,\infty,-16/27$
at the CM points of discriminants $-4,-3,-24$.
In a common weighted-projective representative,
\cite[equations (2) and (14), pp.~3 and 17]{linyang} gives
\[
 J_2=12(j+1),\quad J_4=6(j^2+j+1),\quad J_{10}=j^3,
 \qquad \psi_4=J_2^2-24J_4=144j,\quad \chi_{10}=-j^3/4.
\]
The conventions of \cite[\S2.2, p.~4]{lny} are
$\psi_4=2^{-2}\sum_{m\ {\rm even}}z_m^8$ and
$\chi_{10}=-2^{-14}\prod_{m\ {\rm even}}z_m^2$.
The quotient $\chi_{10}^2/\psi_4^5$ is unchanged by weighted
rescaling. Since our $j$ is finite and nonzero, its denominator
is nonzero, and substitution gives
\[
 j=16\cdot144^5\frac{\chi_{10}^2}{\psi_4^5}
  =3779136\frac{(\prod_{m\ {\rm even}}z_m)^4}
                 {(\sum_{m\ {\rm even}}z_m^8)^5}.
\]
The two homogeneous polynomials have degree $40$ and coefficient
$\ell^1$ norms $3779136$ and $10^5$. Hence
$h(j)\le40h([z_m])+22$.
The product and sum use the ten even theta nulls of the same
principally polarized surface. Adjoining the six zero odd
coordinates does not change their projective height.

In \cite{pazuki}, the level $r=2$ theta coordinates are the sixteen
half-characteristic values used here; the six odd values vanish.
The permitted changes of theta structure permute coordinates and
multiply them by roots of unity. His Euclidean projective height
dominates our supremum projective height. Put
$X=\max(1,h_F(A))$ and
$X_D=\max(1,h_F(A)+\log\pi)\le3X$.
His Corollary~1.3(2), in the Deligne metric, gives
\[
 h_\Theta(A)\le\tfrac12 X_D+C_2\log(X_D+2),
 \qquad C_2=1000\cdot2^4(\log(2^4))^5\le1000\cdot2^{24}.
\]
Using $\log(X_D+2)\le2X_D$ proves
$h([z_m])\le10000\cdot2^{24}X$.

Write $j=a/b$ in lowest terms. The nonzero integer polynomial
$-57600bF(X)^2-aG(X)^6$ has $u$ as a root and coefficient norm at
most $2\cdot4000^6\max(|a|,|b|)$. Remove any initial power of $X$
and use the actual highest nonzero coefficient. The rational-root
theorem bounds both numerator and denominator of $u$, giving
$h(u)\le h(j)+\log(2\cdot4000^6)<h(j)+74$.

A field extension for the theta structure changes neither absolute
height. Since $40\cdot10000\cdot2^{24}+96<10^{13}$, the preceding
inequalities give
\[
 h(u)\le10^{13}\max(1,h_F(A))
 \le10^{13}\max(1,C_{\rm cm},Q_R^{17})M_S(u)^{24}
 =\tfrac12 H_6(S;R)M_S(u)^{24}.
\]
Every estimate used only the displayed lower bound on $R$, so the
conclusion holds for every permitted larger $R$ as well.
\end{proof}

The final value argument uses $S_{\rm ht}=\{p:p\le324\}$. Multiplication
uses $S_{\rm mul}=S_{\rm ht}\cup\mathcal H\cup\operatorname{supp}(c_*)$,
where $\mathcal H$ is the fixed catalogue. A common numerical $R$ may
bound both sets. It does not identify them.

\section{The pole predicate and its converse input}\label{sec:pole}
Let $S_0=\{2,3,5,7,13\}$. With input $B$ and eleven witnesses
$b_0,\ldots,b_5,e,z,r_1,r_2,w$, define
\begin{align*}
 m&=(1-2b_0)(1+b_1)(1+2b_2)(1+4b_3)(1+6b_4)(1+12b_5),\\
 A&=mB,\qquad H=A(A-B^2),\\
 T_1&=A(A-e)(A+3e),\qquad T_2=A(Az^2-eH)(Az^2+3eH),\\
 Q_{\rm pole}&=B\Bigl(\sum_{i=0}^5(b_i^2-b_i)^2+(eHzw-1)^2\\
 &\hspace{4em}+(r_1^2-T_1)^2+(r_2^2-T_2)^2\Bigr).
\end{align*}
Write $\Phi(B)$ for rational solubility. For $B\ne0$, $m$ is a signed
divisor of 2730, and the equations give points $(A,r_1)\in E_e$ and
$(Az^2,zr_2)\in E_{eH}$ with $eHz\ne0$, where
$E_\ell:y^2=x(x-\ell)(x+3\ell)$. The polynomial has total degree 75.

\begin{lemma}[One-sided obstruction]\label{lem:pole}
If $B\ne0$ and $\Phi(B)$, then $v_p(B)$ is not negative and odd for
any $p\notin S_0$.
\end{lemma}
\begin{proof}
If $k=v_p(B)<0$ were odd, then $v_p(A)=k$, $v_p(H)=3k$.
Exactly one of $e,eH$ has even valuation. Square scaling makes its
parameter a unit, preserving the parity of the abscissa's valuation.
For a unit parameter the cubic has valuation $v_p(x)$ if $v_p(x)>0$
and $3v_p(x)$ if $v_p(x)<0$. Its being a square forces even valuation.
But both points have abscissa in the squareclass $A$, a contradiction.
\end{proof}

\begin{theorem}[Pole completeness, $\mathsf{PC}$]\label{hyp:PC}
For every rational $B$ with square reduced denominator,
$0<B<1/2$ and $v_p(B)\ge4$ for every $p\in S_0$ imply $\Phi(B)$.
\end{theorem}
\begin{theorem}[Restricted converse, $\mathsf{RC}$]\label{hyp:RC}
For every positive squarefree $\ell\equiv11\pmod{12}$,
\[
 \operatorname{corank}_{\Z_2}\Sel_{2^\infty}(E_\ell/\Q)=1
 \quad\Longrightarrow\quad\ord_{s=1}L(E_\ell,s)=1.
\]
\end{theorem}
Here $c_4=208\ell^2$, $\Delta=2304\ell^6$, and $j=35152/9$.
The nonintegral rational $j$ shows these curves are non-CM. They are
different from the CM generator curve. Theorem~\ref{hyp:RC} is proved in
Sections~\ref{sec:rc}--\ref{sec:limits}, with its coefficient and
normalization arguments in Appendices~\ref{app:even-source}--\ref{app:endpoints}.
We now prove Theorem~\ref{hyp:PC} from this restricted converse.

\subsection{The exact descent reduction}
The signed divisors of 2730 realize all real, $2$-adic and $3$-adic
squareclass signatures. After $-1,2,3$ fix sign and valuation parities,
the remaining unit vectors of $5,7,13$ are
$(0,1,1),(1,0,0),(0,1,0)$ and form a basis. Choose positive $m$ making
$A=mB$ square at two and a nonsquare unit squareclass at three.
Let $\beta,\theta$ be the positive squarefree representatives of $A,H$.
Then $\beta\equiv17\pmod{24}$, $\theta\equiv1\pmod{24}$, with disjoint
supports, $\beta>1$, and $\theta=1$ permitted. Each is a local square
at the support of the other. At a prime of $\beta$, the relevant odd
valuation of $B$ is positive, so $1-B/m$ is a square unit. At a prime
of $\theta$, $v(B)$ must be zero, since otherwise $v(H)$ is even;
there $B/m\equiv1$, so $A$ is a square.

The simultaneous construction in S11--S12 gives $n,n\theta$ positive
squarefree and congruent to eleven modulo twelve, both with two-Selmer
dimension three and with extra Kummer class $(\beta,1)$. Its matrix
conditions can be stated explicitly. For support primes $q_i$, set
$s_i=[-1/q_i]$, $\lambda_i=[3/q_i]$, and
\[
 (Kv)_i=\sum_{j\ne i}[q_j/q_i](v_j+v_i).
\]
One arranges $\ker K=\langle\mathbf1,\mathbf1_{\operatorname{supp}\beta}\rangle$
and $\lambda\notin\operatorname{im}K$ on both supports. Add $d-1$
prime slots of type $(s,\lambda)=(0,0)$ and one slot $q_0$ of type
$(1,1)$, where $d=\omega(\beta)$. The cross-block has columns
$e_i+e_d$. After the rank-one correction $K^\sharp=K+\lambda a^t$,
deleting $q_0$ leaves
\[
 \begin{pmatrix}D_1&C_1\\C_1^t&0\end{pmatrix},
 \qquad\mathbf1^tD_1\mathbf1=1.
\]
This is nonsingular: $C_1^tx=0$ makes $x$ constant, and the first
equation forces that constant and then $y$ to vanish. Adding the
support of $\theta$ preserves this block. Free $q_0$-symbols toggle
the new diagonal entries; the determinant's full-product coefficient
is one. Thus a single Boolean assignment works on both supports.
Reciprocity, CRT and Dirichlet realize it by distinct auxiliary primes.

In Kummer coordinates $(x+3\ell,x-\ell)$, a positive-second-coordinate
Selmer representative is $(3^\delta\prod q_i^{u_i},\prod q_i^{v_i})$.
The necessary local equations are
\[
 Ku+\operatorname{diag}(\lambda)v=\delta\lambda,\quad
 (K+\operatorname{diag}(s))v=0,\quad(\lambda+s)^tv=0.
\]
The final equation is the condition at three. The matrix properties
force $v=0$, $\delta=0$, $u\in\ker K$, so there are at most eight
classes in total. The extra class is actually local at every place,
including two, and independent of the two torsion classes, giving
dimension three. This sufficiency check for the particular class is
essential; arbitrary parity conditions at two would not suffice.

\begin{lemma}[Dimension three gives corank one]\label{lem:corank}
For an elliptic curve with full rational two-torsion,
$\dim_{\F_2}\Sel_2=3$ implies $\operatorname{corank}_{\Z_2}\Sel_{2^\infty}=1$.
\end{lemma}
\begin{proof}
The finite Kummer sequence gives $r+\dim\Sha[2]=1$.
Write $\Sha[2^\infty]=(\Q_2/\Z_2)^t\oplus F$ with $F$ finite.
The Cassels--Tate pairing on the quotient by the maximal divisible
subgroup is nondegenerate and alternating for an elliptic curve
\cite{poonenstoll}. Its finite primary cyclic factors are paired,
so $\dim F[2]$ is even. Thus $F=0$ and $r+t=1$. The restriction of
the pairing to $F[2]$ need not be nondegenerate and is not used.
\end{proof}

\begin{lemma}[Forward analytic rank zero and one]\label{lem:forward-rank}
Let $A/\Q$ be an elliptic curve. If $\ord_{s=1}L(A,s)=r\leq1$,
then $\operatorname{rank}A(\Q)=r$ and $\Sha(A/\Q)$ is finite.
Consequently $\operatorname{corank}_{\Z_2}\Sel_{2^\infty}(A/\Q)=r$.
\end{lemma}
\begin{proof}
Modularity follows from \cite[Theorem A]{bcdt}.
Apply \cite[Theorem 1, pp.~4--5]{kolyvagin}, whose statement
includes both $r=0$ and $r=1$. Its Theorems 3--5 provide the
auxiliary imaginary quadratic field, the nonzero Heegner height,
and finiteness over that field; the restriction kernel for $\Sha$
is contained in the finite group $\Sha(A/\Q)[2]$.
The ordinary Kummer sequence
\[
0\longrightarrow A(\Q)\otimes\Q_2/\Z_2
\longrightarrow\Sel_{2^\infty}(A/\Q)
\longrightarrow\Sha(A/\Q)[2^\infty]\longrightarrow0
\]
gives the corank assertion. Independently, the $r=0$ finiteness
assertion follows directly from \cite[Corollary 14.3]{kato},
applied to the modular quotient $J_1(N)\twoheadrightarrow A$
over $\Q$ and the trivial character. This result includes the
prime two and imposes no residual irreducibility hypothesis.
\end{proof}

Applying Theorem~\ref{hyp:RC} to both $n$ and $n\theta$,
Lemma~\ref{lem:forward-rank} gives rank one and finite $\Sha$, hence
$\Sha[2]=0$.
The Kummer class is then represented by actual rational points. Their
abscissas have squareclass $\beta$, by the cubic equation. Write
$H=\theta c^2$ and rescale the second point to $E_{nH}$. If the first
abscissa is $Aq^2$, divide both parameters by $q^2$ and both coordinates
by $q^2,q^3$. The abscissas become $A,Az^2$, giving precisely the
eleven witnesses, with $w=(eHz)^{-1}$. If $\theta=1$, one can reuse
a point on the same curve. This proves Theorem~\ref{hyp:PC}, with both
restricted-converse applications in their stated domain.

\section{Multiplication and the finite false-triple guard}\label{sec:mult}
The actual multiplication template first forces all three input codes
to be members of $2E(K_0)$. These membership calls are conjuncts inside
the height certificate; they are not assumed caller preconditions.
Thus even untyped rational inputs, if accepted, have unique indices.
The template has 35,134 rational witnesses. Its construction is S03
and its mathematical specification is S13.

We use the following local conventions for the displayed integral
Weierstrass equation. Put $z=-x/y$, and put $w=1/x$ with $w(O)=0$.
This $w$ differs from the auxiliary variable $-1/y$ used in
\cite[IV.1]{silverman}. The formal expansions give
$w=z^2$ times an integral unit, so $v(w)=2v(z)$.
For a complete discrete valuation with $e=v(p)$, the formal
logarithm converges on $v(z)>0$ and is an isometry with inverse
exponential on $v(z)>e/(p-1)$
\cite[IV.5.5, IV.6.3--6.4, pp.~129--132]{silverman}.
Indeed its degree-$r$ coefficient has denominator dividing $r$,
so each relative error has valuation at least
$(r-1)v(z)-e v_p(r)>0$ in this range.
Thus $\log([n]Q)=n\log(Q)$ and
$v(z([n]Q))=v(z(Q))+v(n)$ for nonzero $Q$ and $n$ there.
These statements concern the formal neighborhood of our integral
equation. When identifying it with the full reduction kernel of
a minimal model, we use
\cite[VII.2.1--2.2, pp.~188--191]{silverman}.

For this particular equation the invariant differential is
\[
 \frac{dx}{2y}=\frac{dz}{\sqrt{1+4d^2z^4}},\qquad
 \log(z)=\sum_{j\ge0}
 \frac{(-1)^j\binom{2j}{j}d^{2j}}{4j+1}z^{4j+1}.
\]
This proves the displayed coefficients of $L$ below. At a contact
prime $p>324$, with $v(z)\ge1$ and $v(p)=1$, every omitted
term relative to $z$ has valuation at least
$r-1-v_p(r)\ge28$ for $r\ge29$.
Consequently $L(z)/16575=\log(z)(1+\epsilon_z)$ with
$v(\epsilon_z)\ge28$.

For indices $n_1,n_2,n_3$, each test group chooses $H=[h]P_0$,
$U_i=[z_i]P_0$, and sets $V_0=H-[c_*]P_0$,
$V_i=U_i-[c_*n_i]P_0$, where $c_*=4M_0$ is fixed. Use
\begin{align*}
 L(z)={}&16575z-6630d^2z^5+11050d^4z^9-25500d^6z^{13}\\
        &+68250d^8z^{17}-198900d^{10}z^{21}+612612d^{12}z^{25}.
\end{align*}
This is 16575 times the degree-25 truncation of the normalized formal
logarithm. The twelve tests are four inverse abscissas at $H,U_i$,
their four conjugate counterparts at $V_0,V_i$, three differences of
the corresponding $L$-ratios, and the product of the first two unprimed
ratios minus the third. Identity inverse abscissas are zero; all
denominators are guarded.

For a slope $t\notin\Cset\cup\{1\}$, put
\[
 R_0=\eta F(t)/(t-1)^8,\quad\eta^2=1,\quad
 a=1-R_0q^2\ne0,\quad R=\ell_0^2a^2R_0.
\]
Denote the twelve unscaled test values in group $g$ by $D_{gj}$.
Their normalized values are
\[
 C_{gj}=\frac{D_{gj}}{R_0^{k_j}(t-1)^{k_j+1}},\qquad
 k_j=\begin{cases}1,&1\le j\le8,\\25,&9\le j\le12.\end{cases}
\]
For six independent groups set
$W_g=\prod_j(1+RC_{gj}^2)$ and
$\Theta=\Tr\sum_gW_g^{-1}$, require them nonzero, and impose
$\Phi(R/\Theta^2)$. The pole input is nonzero, so the prefactor $B$
in $Q_{\rm pole}$ cannot make these constraints vacuous.

At a positive odd contact prime outside $S_{\rm mul}$, each factor
valuation is the minimum of zero and an odd integer. No cancellation
occurs. The trace and Lemma~\ref{lem:pole} force one common group and
place where the first eight tests have valuation at least one and
the last four at least 25. The first tests put the points in the
formal group. At least one of $h,h-c_*$ is a unit at the prime, so its full-log
ratio is integral. The truncated logarithm is the full logarithm times
a unit congruent to one in valuation at least 28. Its ratio is therefore
integral. The difference test transfers integrality to the other truncated
ratio, then to the other full-log ratio. Only now do the relative errors
give absolute congruences modulo $p^{28}$. Write $A_i=L(z(U_i))/L(z(H))$. The difference test and the
full-log ratios give
\[
 z_i-hA_i\equiv0,\qquad
 z_i-c_*n_i-(h-c_*)A_i\equiv0\pmod{p^{25}}.
\]
Their difference is $c_*(n_i-A_i)\equiv0$.
Since $p\notin S_{\rm mul}$ implies $p\nmid c_*$, this yields
$n_i\equiv A_i\pmod{p^{25}}$.
The product test then implies $p^{25}\mid n_1n_2-n_3$.

To specify the height certificate, write the abscissa of each input as
$x_{i0}+x_{i1}s$, using $x_{i0}=x_{i1}=0$ at the identity. Choose
$a_{i0},a_{i1},b_i\in\Q$ with $b_i\ne0$ and
$b_i x_{ij}=a_{ij}$. Express each of these nine rational quantities
as a sum of three slopes. Impose the preceding slope predicate on
all 27 slopes, with private witnesses, and impose membership on the
three input codes. This conjunction is the height certificate.
The multiplication predicate conjoins it with the fixed scalar
graphs and the inverse of the finite guard below.

\begin{proposition}[Finite-box deduction]\label{prop:box}
Assuming these local test identities and the height estimate for
$S_{\rm mul}$ with constant $A_{\rm ht}$, every false triple accepted
by the unguarded certificate has $\max_i|n_i|<T=(48A_{\rm ht})^{13}$.
\end{proposition}
\begin{proof}
Put $f=n_1n_2-n_3\ne0$. For each of the 27 slopes,
$M_{S_{\rm mul}}(t)^{25}\mid f$, so $h(t)\le A_{\rm ht}|f|^{24/25}$.
The nine rational homogeneous quantities for the three abscissas
are sums of three slopes. The sum and quotient height bounds give
$h(x_j)\le6A_{\rm ht}|f|^{24/25}+2\log3$ for each rational component.
The fixed coordinate change and the duplication bounds yield
\[
 2n_i^2\widehat h(P_0)\le12A_{\rm ht}|f|^{24/25}+B_0,
\]
where $2\widehat h(P_0)>1$ and
$B_0=4\log3+\tfrac32\log2+2\log(1+\sqrt2)+\tfrac13\log676<10$.
For $N=\max|n_i|$, $|f|\le2N^2$, so
$N^2<24A_{\rm ht}N^{48/25}+10$.
If $N\ge5$, this implies $N^{2/25}<48A_{\rm ht}$, hence
$N<(48A_{\rm ht})^{25/2}\le T$. The case $N\le4$ is immediate.
\end{proof}

For each false triple $(a,b,c)\in[-T,T]^3$, write its fifteen code
coordinates as $A_j/B_j$ with $B_j>0$, put $D_{abc}=\prod B_j$, and
form
\[
 g_{abc}(X)=\sum_{j=1}^{15}
       (D_{abc}X_j-(D_{abc}/B_j)A_j)^2.
\]
The guard requires the product of all these factors to have an inverse.
A false accepted triple makes its own factor zero. At a true triple
every factor is positive by injectivity of codes, and the inverse exists.
Thus the sound direction uses the generator and corrected all-$S$
height theorem but does not use pole completeness.

% Replacement for the prime-pattern paragraph in manuscript Section 6,
% baseline lines 740--759. The existing application beginning
% "Apply this construction to the nine integral homogeneous quantities"
% follows unchanged. This fragment uses the existing greentao, moebius,
% and gtz bibliography keys. No fixed polynomial constant is changed.

\begin{proposition}[The fixed twenty-two-prime pattern]\label{prop:fixed-prime-pattern}
Put
\[
 \mathcal B=\{0,-15,-20,-24,-30,-60,1\},\qquad
 P_s=\prod_{p\le324}p,\qquad W=4P_s.
\]
Let $\mathcal H$ consist of the first twenty-two primes greater than
$324$ in each unit residue class modulo $W$, as in the fixed constant
prescription, and set
\[
 \mathcal K=\{dh:d\mid P_s^4,\ d>0,\ h\in\mathcal H\}.
\]
For every $A\in\mathbb Z$ and every real $B_0$ there are integers
$u_1,u_2,v$, with $v\ne0$, such that, on putting $u_3=Av-u_1-u_2$,
each of the twenty-two values
\[
 v,\qquad u_j-bv\quad(1\le j\le3,\ b\in\mathcal B)
\]
equals $\epsilon k p$, where $\epsilon\in\{1,-1\}$,
$k\in\mathcal K$, and $p>B_0$ is prime. Every such $p$ splits in
$\mathbb Q(\sqrt2)$, the two reductions of
\[
 E:y^2=x^3-(75-53\sqrt2)^2x
\]
are good, and their orders $d_1,d_2$ satisfy
\[
 \gcd(d_1,d_2)=4,\qquad p\nmid d_1d_2.
\]
The prime values need not be distinct. The assertion concerns the two
orders at the same prime only.
\end{proposition}

\begin{proof}
Write $T=(u_1,u_2,v)$ and denote the twenty-two homogeneous forms by
$F_\nu$. Their coefficient vectors are
\[
 (0,0,1),\quad (1,0,-b),\quad(0,1,-b),\quad(-1,-1,A-b).
\]
Each form has a coefficient equal to $1$ or $-1$. Within each of the
three seven-form families, a nonzero two-by-two minor is a contact
difference, of absolute value at most $61$. Between two different
families the first two coordinates give a minor $1$ or $-1$; a pair
involving $v$ also has such a minor. Thus every pair is independent
over $\mathbb Q$ and over $\mathbb F_q$ for every prime $q>324$,
for every integer $A$. No condition on three-form dependencies is
required below.

Choose a bounded open rectangular box $\Omega\subset\mathbb R^3$
of positive volume whose closure avoids all the zero planes. For each
$\nu$ choose $\epsilon_\nu$ such that $\epsilon_\nu F_\nu>0$ on
$\overline\Omega$. These signs allow all integers $A$, including zero
and negative integers.

\paragraph{The fixed multipliers and initial progression.}
For each prime $q\le324$, the proportion of $T$ modulo $q^5$ for
which at least one $F_\nu(T)$ vanishes is at most $22/q^5<1$.
Choose a residue outside this union and lift it to a residue modulo
$q^8$. On this residue each $F_\nu$ has a fixed valuation
$e_{\nu,q}\in\{0,1,2,3,4\}$. Set
\[
 d_\nu=\prod_{q\le324}q^{e_{\nu,q}}.
\]
The unit $\epsilon_\nu F_\nu/d_\nu$ is now fixed modulo $8$ at
$2$ and modulo $q$ at each odd $q\le324$. Prescribe the prime
residues
\[
 c_2=1\pmod8,\qquad c_7=3\pmod7,\qquad
 c_q=1\pmod q\quad(q\le324,\ q\ne2,7).
\]
Choose $h_\nu\in\mathcal H$ in the unit class specified by
$h_\nu\equiv(\epsilon_\nu F_\nu/d_\nu)c_q^{-1}$ at these
moduli. There are at most twenty-two demands in any class, so the
$h_\nu$ can be chosen pairwise distinct. Put $k_\nu=d_\nu h_\nu$.
Dirichlet's theorem proves that this is the same finite catalogue as
the prescribed one; no extra catalogue prime is introduced.

At $h=h_\nu$, impose $F_\nu(T)=0$. This plane has $h^2$ points.
For every $\mu\ne\nu$, the restriction of $F_\mu$ is nonconstant,
so excluding
\[
 F_\mu(T)=0,\qquad \epsilon_\mu F_\mu(T)=-k_\mu
\]
removes at most $42h<h^2$ points. Choose a remaining point.
Among its lifts modulo $h^2$, the unit
$\epsilon_\nu F_\nu/(d_\nu h)$ runs through every residue
modulo $h$: changing a coordinate with coefficient $1$ or $-1$
by $h$ proves surjectivity. Choose this residue different from $0$
and $-1$. The restrictions for $\mu\ne\nu$ remain valid.

The Chinese remainder theorem now gives one vector residue $r_0$
modulo
\[
 W_0=\prod_{q\le324}q^8\prod_\nu h_\nu^2.
\]
On $r_0+W_0\mathbb Z^3$ the functions
$\Psi_\nu=\epsilon_\nu F_\nu/k_\nu$ are integral. At every prime
dividing $W_0$ they are nonzero, and at every odd such prime they
are different from $-1$. In the affine coordinates on this
progression every variable coefficient is divisible by each prime
dividing $W_0$. Indeed, $q^8/d_\nu$ is divisible by $q$ even when
$e_{\nu,q}=4$, and $h^2/k_\nu$ is divisible by $h$.

All these choices depend on $A$ and the fixed catalogue, but not on
$B_0$ or the cutoff $L$ used next. Fix also
\[
 Q_*=\max(324,\max\mathcal H).
\]
In both upper sieves below the excluded fixed primes are exactly
those at most $Q_*$; in particular this exclusion does not grow
with $L$. There are constants $0<\alpha_A<\beta_A$ such that
\[
 \alpha_A X\le\Psi_\nu(T)\le\beta_A X
 \quad(T\in X\Omega,\ 1\le\nu\le22).
\]

\paragraph{The lower count with its cutoff dependence.}
Take $L>\max(Q_*,B_0,3)$. At each prime $q\le L$ not dividing
$W_0$, exclude the forty-four affine planes
\[
 F_\nu(T)=0,\qquad \epsilon_\nu F_\nu(T)=-k_\nu.
\]
Each has $q^2$ points and all these primes exceed $324$. With
\[
 W_L=W_0\prod_{\substack{q\le L\\q\nmid W_0}}q,
\]
let $\mathcal R_L$ be the resulting vector residues modulo $W_L$
that reduce to $r_0$ modulo $W_0$. Then
\begin{equation}\label{eq:prime-density44}
 \frac{|\mathcal R_L|}{W_L^3}
 \ge W_0^{-3}\prod_{\substack{q\le L\\q\nmid W_0}}
       (1-44/q)
 \ge c_A(\log L)^{-44}.
\end{equation}
Here and below a positive constant with subscript $A$ may depend
on the choices already made, including the fixed catalogue, but
never on $L$, $X$, or a varying fiber prime. The last inequality
follows from $\log(1-44/q)=-44/q+O(q^{-2})$ and Mertens' theorem \cite[Theorem 1, with modulus $2$]{williams-mertens}.

Fix $L$. For $r\in\mathcal R_L$ apply the finite-complexity
prime asymptotic of \cite[Main Theorem, p.~1763]{greentao},
with its hypotheses supplied by \cite[Theorem 1.1]{moebius}
and \cite[Theorem 1.3]{gtz}, to
\[
 n\longmapsto\frac{\epsilon_\nu F_\nu(r+W_Ln)}{k_\nu},
 \qquad r+W_Ln\in X\Omega.
\]
These are integral affine forms with pairwise independent linear
parts. Singleton classes in \cite[Definition 1.5]{greentao}
give complexity at most $20$, so the inputs $\mathrm{GI}(20)$
and $\mathrm{MN}(20)$ suffice. The inverse theorem is read in
arXiv version 5 together with the separate erratum, in particular
its Proposition 3.1 and modified deduction of Proposition 7.3.
All coefficient bounds and moduli are fixed before $X$ tends
to infinity. At $q\mid W_L$ the local factor is
$(1-1/q)^{-22}$. At $q\nmid W_L$, scaling and changing variables
preserve the twenty-two homogeneous zero planes, so
\[
 1-22/q\le
 \mathbb P(F_\nu\ne0\text{ for all }\nu)
 \le1-22/q+\binom{22}{2}q^{-2}.
\]
Consequently the local factor is $1+O(q^{-2})$, and
\[
 \log\beta_q\ge\log(1-22/q)-22\log(1-1/q)
 \ge-484/q^2.
\]
The singular product is therefore bounded below by an absolute
positive constant, uniformly in $L$ and $r$. The real factor is
$\operatorname{vol}(\Omega)X^3/W_L^3$.

There are finitely many $r$ for fixed $L$, so their asymptotic
errors can be made smaller than half their main terms with one
threshold for $X$, depending on $A,L$. This uses no uniform
prime-asymptotic error as $L$ varies. Summing over $r$ and using
\eqref{eq:prime-density44} gives a weighted lower bound
$c_A(\log L)^{-44}X^3$. Proper prime powers contribute at most
$O_A(X^{5/2}(\log X)^{23})$: there are
$O_A(X^{1/2}\log X)$ possible proper powers in each form's range,
each equation fixing a form value has $O_A(X^2)$ integer solutions
in the box, and each weight is $O_A((\log X)^{22})$.
Thus, for all sufficiently large $X$ depending on $A,L$,
\begin{equation}\label{eq:prime-lower}
 \#\{T\in X\Omega:\ T\bmod W_L\in\mathcal R_L,
                  \ \Psi_\nu(T)\text{ all prime}\}
 \ge a_A\frac{X^3}{(\log L)^{44}(\log X)^{22}},
 \qquad a_A>0.
\end{equation}

\paragraph{An upper sieve with uniform constants.}
We give the estimates needed for both upper bounds. Consider an
axis-parallel square of side $N\ge1$ in $\mathbb R^2$. For each sieving prime
$q$ let $\Omega_q\subset\mathbb F_q^2$ have size $\rho(q)$,
where $0<\rho(q)\le21q$ and $q>Q_*$. Put
\[
 g(q)=\rho(q)/q^2,\qquad
 h(q)=g(q)/(1-g(q)),
\]
and extend both multiplicatively on squarefree integers supported
on the sieving primes. For such $d$, CRT and the elementary count
of lattice points in each residue square give
\begin{equation}\label{eq:prime-residue-error}
 A_d=N^2g(d)+R_d,\qquad
 |R_d|\ll\rho(d)(N/d+1)\le Nd+d^2.
\end{equation}
The implied absolute constant is uniform in the square's location
and in the sets $\Omega_q$. The last inequality uses
$\rho(q)\le q^2$.

For $Z\ge2$ define
\[
 G(Z)=\sum_{d\le Z}h(d),\qquad
 \lambda_d=\frac{\mu(d)}{G(Z)\prod_{q\mid d}(1-g(q))}
       \sum_{\substack{m\le Z/d\\(m,d)=1}}h(m)
 \quad(d\le Z),
\]
and set $\lambda_d=0$ otherwise. All sums here are over the stated
squarefree support, with $g(1)=h(1)=1$. These are Selberg's weights.
They satisfy $\lambda_1=1$, $|\lambda_d|\le d$, and
\[
 \sum_{e,f\le Z}\lambda_e\lambda_f g([e,f])=G(Z)^{-1}.
\]
For completeness, set $y_r=\sum_{r\mid d}\lambda_dg(d)$.
The displayed quadratic form equals $\sum_r y_r^2/h(r)$,
because $1/g(m)=\sum_{r\mid m}1/h(r)$; the stated weights give
$y_r=\mu(r)h(r)/G(Z)$. The bound on the weights follows from
the inner sum being at most $G(Z)$ and
$(1-g(q))^{-1}\le q$ for $q>324$.

Squaring the divisor sum majorizes the indicator of avoiding
all $\Omega_q$ for $q\le Z$. Inserting
\eqref{eq:prime-residue-error}, and using $[e,f]\le Z^2$,
there are at most $Z^2$ pairs with $|\lambda_e\lambda_f|\le Z^2$.
It follows that
\begin{equation}\label{eq:prime-selberg}
 S(N,Z)\ll \frac{N^2}{G(Z)}+NZ^6+Z^8.
\end{equation}

Here are the two denominator estimates, including truncation.
For $z=Z^{1/100}$ put
\[
 H(z)=\prod_{\substack{q\le z\\q\text{ sieved}}}(1+h(q)).
\]
Give each squarefree divisor $d$ of this prime product probability
$h(d)/H(z)$. The events $q\mid d$ are independent and have
probability $g(q)$. Hence
\[
 \mathbb E\log d=\sum_{\substack{q\le z\\q\text{ sieved}}}
                  g(q)\log q
 \le21\sum_{q\le z}\frac{\log q}{q}
 \le42\log z=0.42\log Z
\]
for sufficiently large $z$, using
$\sum_{q\le z}(\log q)/q=\log z+O(\log\log z)$.
Indeed, \cite[Theorem 1, with modulus $2$]{williams-mertens}
gives $S(z):=\sum_{q\le z}q^{-1}
=\log\log z+C+O(1/\log z)$, including the prime $2$.
Partial summation in
$S(z)\log z-\int_2^z S(t)\,dt/t$ gives the asserted bound.
Markov's inequality gives probability at least $0.58$, and thus
at least $1/2$, that $d\le Z$. Every such divisor is included in
$G(Z)$, so
\begin{equation}\label{eq:prime-truncated-denominator}
 G(Z)\ge\tfrac12 H(Z^{1/100}).
\end{equation}
This proves a lower bound for the truncated sum, rather than
merely for an untruncated Euler product.

In the first application the sieved primes are
$q>Q_*$ with $q\equiv1\pmod4$, with at most one prime omitted,
and $g(q)=2/q-1/q^2$. Mertens' theorem in the fixed progression
$1\pmod4$ \cite[Theorem 1]{williams-mertens} gives
\[
 \log H(z)=\log\log z+O_{Q_*}(1).
\]
Omitting a prime $\ell$ changes $H(z)$ by $1-g(\ell)\ge1/2$
when $\ell\le z$, and has no effect otherwise. Therefore
$G(Z)\gg_{Q_*}\log Z$ uniformly in the omitted prime.
In the second application all primes $q>Q_*$ are sieved and
\[
 \frac{21}{q}-\frac{210}{q^2}\le g(q)\le\frac{21}{q}.
\]
Then $\log H(z)=21\log\log z+O_{Q_*}(1)$, with an error
independent of the individual sets $\Omega_q$, and hence
$G(Z)\gg_{Q_*}(\log Z)^{21}$. These statements and
\eqref{eq:prime-selberg} establish every uniform upper sieve
estimate used below.

\paragraph{The conjugate orders.}
Every prime counted by \eqref{eq:prime-lower} satisfies
$p\equiv1\pmod8$ and $p\equiv3\pmod7$. Thus it splits in
$\mathbb Q(\sqrt2)$ and $(7/p)=-1$. The two reductions of
$d=75-53\sqrt2$ have product $7$ and opposite square classes.
The two reduced curves are the quadratic twists of $y^2=x^3-x$
by these two classes. Their traces are opposite. The discriminant
$64d^6$ is a unit at these primes, so both reductions are good.

For $p\equiv1\pmod4$, choose $i\in\F_p$ with $i^2=-1$.
The automorphism $\iota:(x,y)\mapsto(-x,iy)$ satisfies
$\iota^2=[-1]$ and commutes with the arithmetic $p$-power
Frobenius $\varphi$, directly from its coordinate formula.
The injection
\[
 \operatorname{End}(\overline E)\otimes\Q_3
 \hookrightarrow\operatorname{End}_{\Q_3}(V_3\overline E)
\]
is \cite[III.7.4, pp.~89--91]{silverman}; here $p>324$.
Since $T^2+1$ is irreducible over $\Q_3$, the centralizer of
$\iota$ on this two-dimensional Tate space is $\Q_3(i)$.
The rational centralizer therefore has dimension at most two
and already contains $\Q(i)$, so it is exactly $\Q(i)$.
The Frobenius endomorphism is integral over $\Z$ by
\cite[V.2.3.1(b), p.~142]{silverman}. Hence it lies in $\Z[i]$;
write it as $a+bi$ with $a,b\in\Z$.
The trace and determinant formulas
\cite[III.8.6, p.~99]{silverman} give
\[
 a^2+b^2=p,\qquad d_1=p+1-2a,\qquad d_2=p+1+2a.
\]
Here $d_1,d_2$ are the cardinalities of the two reduced elliptic
curve groups; the opposite twist changes the Frobenius trace's
sign. In particular $a\ne0$, since $p$ is not a square.
Both orders are divisible by $4$ because all the two-torsion is
rational. Their sum is $2(p+1)$, of $2$-adic valuation exactly
$2$, so the $2$-part of their gcd is exactly $4$. Hasse's bound
\cite[V.1.1, p.~138]{silverman}
gives $0<d_i<2p$ for $p>324$, and $d_i\ne p$ since $4\mid d_i$.
Thus $p\nmid d_1d_2$.

If an odd prime $\ell$ divides both orders, it divides $a$ and
$p+1$, and $b^2\equiv-1\pmod\ell$. The progression conditions
already rule out every odd $\ell\le L$. It remains to discard
prime values with such an $\ell>L$.

\paragraph{The number of bad prime values.}
For a fixed form its prime value lies between $\alpha_A X$ and
$\beta_A X$. Representations $p=a^2+b^2$ lie in
$[-C_A\sqrt X,C_A\sqrt X]^2$. Overcount them to bound the
number of bad primes. If $L<\ell\le X^{1/8}$, there are at most
two residues $b_0$ with $b_0^2\equiv-1\pmod\ell$. Write
\[
 a=\ell U,\qquad b=b_0+\ell V.
\]
The required $(U,V)$ lie in a square of side
$N\asymp_A\sqrt X/\ell\gg_A X^{3/8}$.
For every $q>Q_*$ with $q\equiv1\pmod4$ and $q\ne\ell$,
the condition $a^2+b^2=0\pmod q$ is the union of two distinct
lines through the origin, of size $2q-1$. The affine substitution
is bijective modulo $q$. Take $Z=X^{1/100}$. Since the prime norm
is at least $\alpha_A X>Z$ for large $X$, every required pair
avoids these residues. Equations \eqref{eq:prime-selberg} and
\eqref{eq:prime-truncated-denominator} give
\[
 \#(U,V)\ll_A \frac{N^2}{\log X}+NZ^6+Z^8
 \ll_A\frac{X}{\ell^2\log X}.
\]
The absorption is uniform: $Z^6/N\ll_A X^{0.06-3/8}$ and
$Z^8/N^2\ll_A X^{0.08-3/4}$, both tending to zero faster
than $1/\log X$. Summing over $\ell>L$ gives
$O_A(X/(L\log X))$.

For $\ell>X^{1/8}$ use the lattice count directly. Since
$a\ne0$, there are $O_A(\sqrt X/\ell)$ possible multiples $a$,
with no additive constant. For each of the at most two root
classes there are $O_A(\sqrt X/\ell+1)$ choices of $b$.
Also $\ell\le|a|\le C_A\sqrt X$. Bounding the sum over primes
by the corresponding sum over positive integers yields
\[
 \sum_{X^{1/8}<\ell\le C_A\sqrt X}
       O_A\left(\frac{X}{\ell^2}+\frac{\sqrt X}{\ell}\right)
 \ll_A X^{7/8}+\sqrt X\log X.
\]
We have therefore proved the bound, uniformly in $L$,
\begin{equation}\label{eq:prime-bad-values}
 \#\{\text{bad prime values for a fixed form}\}
 \ll_A\frac{X}{L\log X}+X^{7/8}+\sqrt X\log X.
\end{equation}

\paragraph{The number of patterns on each nonzero fiber.}
Fix $\nu$ and $p\in[\alpha_A X,\beta_A X]$, and put
$c=\epsilon_\nu k_\nu p$. The integral solutions of $F_\nu(T)=c$
are parametrized by two original coordinates: solve for a
coordinate whose coefficient is $1$ or $-1$. For example, for
$F_\nu=-u_1-u_2+(A-b)v$, use $(u_1,v)$ and set
$u_2=(A-b)v-c-u_1$. This parametrization is valid over every
$\mathbb F_q$, and the points from $X\Omega$ lie in a square
of side $O_A(X)$ in these coordinates.

For $Q_*<q\le Z=X^{1/100}$, one has $q\nmid k_\nu$ and
$q<p$ when $X$ is large. Thus $c\ne0$ modulo $q$. Each remaining
form restricts to a nonconstant affine form. Their twenty-one
zero lines are distinct: if the restrictions of $F_j,F_k$ had
the same zero line, then for some $\lambda\ne0,\mu$,
\[
 F_j-\lambda F_k=\mu(F_\nu-c)
\]
as affine forms over $\mathbb F_q$. Comparing constants gives
$\mu c=0$, hence $\mu=0$, contradicting pairwise independence
of $F_j,F_k$. Parallel lines and multiple intersections cause
no difficulty. The union consequently has between
$21q-\binom{21}{2}$ and $21q$ points.

Every target pattern on the fiber avoids this union: its other
prime values are also at least $\alpha_A X>Z$, and every
$k_j$ is a unit at $q>Q_*$. Discarding the CRT progression
conditions only increases the count. The second denominator
estimate in \eqref{eq:prime-truncated-denominator} gives
\begin{equation}\label{eq:prime-fiber}
 \#\{\text{target patterns on }F_\nu=c\}
 \ll_A \frac{X^2}{(\log X)^{21}}+X^{1.06}+X^{0.08}
 \ll_A \frac{X^2}{(\log X)^{21}}.
\end{equation}
This is uniform in the fiber prime $p$ and independent of $L$.
In particular $q=p$ is outside the sieving range; no varying
exceptional-prime product has been suppressed.

\paragraph{The order of choices.}
Multiply \eqref{eq:prime-bad-values} by \eqref{eq:prime-fiber}
and sum over the twenty-two forms. The number of bad patterns
is at most
\[
 b_A\frac{X^3}{L(\log X)^{22}}
 +O_A\left(\frac{X^{23/8}}{(\log X)^{21}}
          +\frac{X^{5/2}}{(\log X)^{20}}\right),
\]
where $b_A$ is independent of $L$. First choose
$L>\max(Q_*,B_0,3)$ so large that
\[
 \frac{b_A}{L}<\frac{a_A}{2(\log L)^{44}}.
\]
Such an $L$ exists because $(\log L)^{44}/L\to0$.
Then keep $A,L$ and every preceding choice fixed, and take $X$
large enough for \eqref{eq:prime-lower}, the upper bounds, and
the remaining error to be less than
$a_A X^3/(2(\log L)^{44}(\log X)^{22})$.
A pattern survives. Its signed prime values tend to infinity
with $X$, $v\ne0$, and the preceding order calculation proves
all the assertions.
\end{proof}
Apply this construction to the nine integral homogeneous quantities
of a true triple. Each slope has at most six odd-contact primes.
At one such prime with reduced group cardinalities $d_1,d_2$,
which annihilate the reduced point, solve
\[
 h\equiv0\pmod{\operatorname{lcm}(M_0,d_1)},\qquad
 h\equiv c_*\pmod{d_2}.
\]
Compatibility follows from
\[
 \gcd(\operatorname{lcm}(M_0,d_1),d_2)
                  \mid M_0\gcd(d_1,d_2)=4M_0=c_*.
\]
Choose a solution larger than $c_*$. Each contact prime uses its own
private group within its own slope. No common shift across different
primes is imposed, and selected prime values need not be distinct.
The saved library allocates six private groups per slope and twenty-seven
private slope calls. Set $U_i=[n_i]H$ and assign one group to each prime. Fixed local
clearing at the old primes can be checked as follows.
Put $S_b=S_{\rm ht}\cup\mathcal H$, $B_p=4$ on $S_{\rm ht}$ and
$B_p=1$ on the remaining catalogue primes, $C_p=14B_p$, and
$D_p=51C_p$. The selected integral forms give
$|v_p(R_0)|,|v_p(t-1)|\le C_p$, while the small formal subgroup
in the definition of $M_0$ gives $v_p(D_{gj})\ge0$ at $S_b$.
Consequently $v_p(C_{gj})\ge-D_p$. The fixed normalization has
$v_p(\ell_0)=E_p=52C_p+v_p(12)+4$.

Weak approximation chooses $q$ small at $S_b$, integral at the
selected primes, and arbitrarily close to a real root of
$1-R_0q^2$, without equality; choose $\eta$ to make $R_0>0$.
Then $a$ is a unit at the specified finite primes. At $p\in S_b$,
\[
 v_p(RC_{gj}^2)\ge 2E_p-C_p-2D_p>v_p(12),\qquad
 v_p(\Theta-12)>v_p(12),
\]
so $v_p(R/\Theta^2)\ge103C_p+8\ge4$.
At a selected odd-contact prime, $v_p(R)=1$. Each of the twelve
summands of the trace of the six $W_g^{-1}$ has residue zero or one;
the assigned group and place contributes one. Since $p>324$, their
sum is a unit and the pole input has valuation one. At every other
prime outside $S_b$, its valuation is even. Thus its reduced
denominator is square. In both real embeddings each $W_g$ is positive;
as $a\to0$, $W_g\to1$, $\Theta\to12$, and $R/\Theta^2\to0$
through positive values. This verifies exactly $\mathsf{PC}$'s
hypotheses, including all nonzero guards. Negative indices, zero indices and
coincident points are handled by the group-law branches. Therefore
Theorem~\ref{hyp:PC} supplies the required pole witnesses. This proves
the full multiplication statement $\mathsf M$, including completeness
for every integer triple $(n_1,n_2,n_1n_2)$ and soundness on arbitrary
rational input codes.

% End 02-height-pole-multiplication.tex


% Begin 03-restricted-converse.tex
\section{The restricted converse and its dependency order}\label{sec:rc}
For a positive squarefree $\ell\equiv11\pmod {12}$ put
\[
 E_\ell:y^2=x(x-\ell)(x+3\ell).
\]
Its invariants are $c_4=208\ell^2$, $\Delta=2304\ell^6$ and
$j=35152/9$. Thus it has full rational two-torsion and is not a
CM curve, since a rational CM $j$-invariant is an integer.
The curve in this section is different from the fixed CM curve used
to represent integer indices.

The assertion needed in Theorem~\ref{hyp:RC} is
\[
 \operatorname{corank}_{\Z_2}\Sel_{2^\infty}(E_\ell/\Q)=1
 \quad\Longrightarrow\quad \ord_{s=1}L(E_\ell,s)=1.                 \tag{RC}
\]
We prove this assertion in this section and the next two sections.
Appendices~\ref{app:even-source}--\ref{app:endpoints} give the fixed-form,
theta, nonvanishing and normalization constructions used in this deduction.
By \cite[Theorem A and Introduction, conditions (1)--(6)]{bcdt},
$E=E_\ell$ is modular. Here $N$ denotes its conductor.
The modular realization gives a rational weight-two newform $f$
of level $N$ and an elliptic quotient $q_f:J_0(N)\twoheadrightarrow A_f$
in the rational isogeny class of $E$; see
\cite[Theorem 11.19 and Remark 11.18]{milnemf}.
Choose a $\Q$-isogeny $\alpha:A_f\to E$ and put
$\varphi=\alpha q_f$ and $\pi(x)=\varphi([x]-[\infty])$.
Thus $\pi(\infty)=0$ and, for $q\nmid N$,
$\varphi T_q=[a_q(f)]\varphi$.
Fix the integral Tate lattice
and the local squareclasses allowed at the primes of $2N$.
Constants depending on this data are chosen before any conductor,
companion field or binary dimension varies. Quotienting by a rational
point of order two gives a rational two-isogeny. The Selmer parity
theorem \cite[Theorem 1.8]{dd} therefore gives root number $-1$
under the left side of (RC). This application does not replace Selmer
corank by Mordell--Weil rank and does not assume finite $\Sha$.
\begin{remark}[Why this family is the input to the converse]\label{ex:sec7-rc-scope}
The congruence $\ell\equiv11\pmod {12}$ comes from the pole-completeness
construction. In the proof of Theorem~\ref{hyp:PC}, the two parameters
$n,n\theta$ are positive squarefree numbers in this congruence class;
their two-Selmer groups have dimension three. Lemma~\ref{lem:corank}
gives full two-primary Selmer corank one, so these are exactly the
two applications required of Theorem~\ref{hyp:RC}.

In the converse argument we fix one such $\ell$. In particular $\ell$
is odd, is prime to three, and is greater than one, hence is not a
square. We use the displayed full rational two-torsion and non-CM
properties, the resulting rational two-isogeny, and the explicit roots
$0,\ell,-3\ell$. The last roots enter the $K$-split seed argument below.
The constants may depend on this fixed curve; there is no bound on
$\omega(\ell)$ and no assertion that they are uniform as $\ell$ varies.
The statements about coefficient graphs below are the rational,
full-two-torsion case used by $E_\ell$ and its quadratic twists.
\end{remark}

Here is the dependency order of the proof. A fixed-form Euler-system
calculation gives the even coefficient bound. A varying-companion
arithmetic estimate, proved by the two limits below, then gives the
genus-point lower bound. Only after that lower bound do we normalize
the odd coefficients. Their finite-stage forms and the even forms give
one common binary family. Finally a fixed-field determinant argument
excludes a zero derivative at the base. The even estimate does not use
that final determinant argument on Heegner classes.

\subsection{The even coefficient calculation}
Use normalized central values $L_*(h)$ with the real and imaginary
periods fixed by $\pi$ and the integral lattice. Only their valuations
are used; every subsequent comparison uses this same normalization.
For every twist whose two-primary Tate--Shafarevich group is finite, let
$s(h)=\operatorname{length}_{\Z_2}\Sha(E^{(h)}/\Q)[2^\infty]$.

\begin{proposition}[The fixed-form even scalar]\label{ex:sec7-even-scalar}
Use the fixed $E,\pi$, integral lattice and allowed old local types
specified above. Put $\Lambda=\Z_2[[t]]$. For a finite elementary
abelian two-group $G$, let $\chi_x=\chi_0(x\circ\varepsilon)$
$(x\in\widehat G)$ be a family of quadratic characters with a common
real sign and the same restrictions at every place dividing $2N$.
There is one scalar $u^{\rm ev}\in\Lambda[1/2][G]$ satisfying (3.2).
Its clearing exponent and the bounded error in that formula depend
only on the fixed data, not on $G$ or the good-prime conductor support.
The central identity is used only at corank-zero constituents with
nonzero central value; at a corank-zero constituent the same central
coordinate is zero exactly when the central value is zero.
\end{proposition}
\begin{proof}\leavevmode
The following construction is expanded in
Propositions~\ref{ex:app-kato-system} and~\ref{ex:app-even-scalar},
with its central calculation in Proposition~\ref{ex:app-central-calibration}.
The fixed-form construction uses one integral symbol combination with
nonzero projections to both real signs. Kato's integral maps, with
$j_*$ cohomology taken modulo each power of two before passage to the
limit, give classes whose inertia restrictions vanish integrally.
Indeed the low-degree Leray edge map gives this vanishing before
rationalization. Restriction of transfer is the Mackey sum of the
corresponding upper inertia restrictions. Hence unnormalized transfer
preserves the vanishing even when the intervening extension is ramified.
No character projector divided by a growing group order is used.

For a finite binary character group, let $F$ be the least common
multiple of the odd parts of all its character conductors, including
the base character. Write $Z_F$ for the compatible Iwasawa class
obtained by pushing the retained-smoothed Kato classes at levels
$F2^n$ through the integral coefficient map.
Here and throughout the common-scalar construction, the binary family
is required to have one real sign and fixed local squareclasses at
$2N$. Writing its characters as
$\chi_x=\chi_0(x\circ\varepsilon)$, this means
$\varepsilon(c)=1$ for complex conjugation $c$, and
$\varepsilon|_{G_{\Q_v}}=1$ at every $v\mid2N$.
Indeed every character of $G$ evaluates these elements to one, and
the characters separate $G$. The lower bound uses $G=1$; the
upper-transfer cubes lie in the fixed signed unit filter of their
coefficient source. The final cubes have the same property.
At a character $\chi$ its unused-prime factors are
\[
 E_q(\chi)=1-\frac{a_q\chi(q)}q\sigma_q^{-1}
                 +\frac1q\sigma_q^{-2}.                           \tag{3.1}
\]
They are the actual determinants of the unused good local terms.
Active ramified good primes have rational inertia $-1$, so have no such
rational local term. The fixed bad and dyadic corrections form a factor
$F_E$ with nonzero central value.

The rational divisibility input is Kato's Theorem 13.4(2),
together with the nonzero-class input in 13.5(2)
\cite{kato}. We use the height-one primes not containing two.
Rational irreducibility is preserved on a quadratic subgroup of the
non-CM open image. That subgroup contains a nontrivial unipotent,
which supplies the required one-dimensional fixed space. Oddness
supplies the two real eigenspaces, and purity supplies the common
Frobenius weight. This application does not use residual irreducibility
at two or the stronger hypotheses of Theorem 13.4(3).

The smoothing constants $c,d_s$ can be chosen to satisfy
\[
 c,d_s\equiv1\pmod {3N_{\rm odd}F},\qquad
 c,d_s\equiv1+2^e\pmod {2^{e+1}},
\]
where $e\ge\max\{3,v_2(N)\}$ is fixed by the original dyadic data.
In weight two and
cohomological degree one the two smoothing exponents are one.
Kato's explicit reciprocity formula therefore contributes valuation
$2e$, independent of $F$. For the central normalization, first descend
to the primitive character conductor and then use the invariant basis
$\sqrt h\,e_{\chi_h}$ of the quadratic character line. The resulting
varying scalar $1/h$ is a two-adic unit. The complete calculation in
Appendix~\ref{app:even-source} includes discriminant $4h$; a Gauss basis
there has the extra factor $\sqrt{\Delta_h/h}$, which must also be
included in its change of basis. The numerical sizes of $c,d_s$ may
grow, while their relevant valuation remains fixed.

For completeness, the integral calculation at two uses the positive
complex. Finite continuous cochains are filtered unions of finite free
modules and hence flat. Splitting the residual complex into cohomology
and disks, then using the nilpotent perturbation calculation of
Section~\ref{sec:limits}, gives finite free models with their maps.
In the fixed real cyclotomic tower the ordinary two-primary class group
vanishes. One proves this inductively: only the prime above two ramifies
in each quadratic step; the abelianization of a possible Hilbert
two-class tower is generated by that inertia of order two. A finite
two-group with cyclic abelianization is cyclic, and inertia excludes a
larger cyclic extension. The unit-signature calculation is separate.
A norm-minus-one unit in positive layers, and $-1$ at the base, has
signature vector a unit in the characteristic-two cyclic group algebra.
For a scalar coefficient the resulting residual dimensions are
\[
 \dim H^1=2^n+s_n-1,\qquad \dim H^2=s_n-1.
\]
For each fixed support $s_n$ is bounded. Universal coefficients give
generic residual rank one in degree one; the two trivial elliptic
factors give rank two. Unit cancellations thus identify the positive
complex over $\mathscr O[G]$ with $\mathscr O[G]^2[-1]$.
This last rank is independent of support, although preliminary models
need not have a uniform size.

A generator of the fixed real invariant lattice gives a connecting class $a$.
Its central image need not be saturated.
Twice an ordinary class lifts because the relevant real $H^1$ is killed
by two. Wedge $a$ with a lift of $2F_E Z_F$ in the original determinant
line. This produces a common scalar $u^{\rm ev}$.
The first assertion below holds for the full family. The central valuation
identity is asserted only at a constituent of Selmer corank zero with
$L_*(h)\ne0$:
\[
 2^{C_{\rm ev}}u^{\rm ev}\in\Lambda[G],\qquad
 v_2u^{\rm ev}_h(0)=v_2L_*(h)-2\omega(h)-s(h)+O_E(1).               \tag{3.2}
\]
Here $\Lambda=\Z_2[[t]]$ and
$\mathscr O=\widehat{\Lambda_{(2)}}$. The error and $C_{\rm ev}$ are
fixed before the binary group and support. To justify membership in
$\Lambda[G]$, first use rational characterwise divisibility and Fourier
inversion for membership in $\Lambda[1/2][G]$, without estimating the
Fourier denominator. Intersect coefficientwise with the integral
$\mathscr O[G]$ result. The explicit coefficient construction also covers field discriminant $4h$,
with its dyadic character retained before the cofinal norm limit, as in
Appendix~\ref{app:even-source}. At a corank-zero constituent the same coordinate
vanishes centrally if and only if the central value vanishes.
At such a nonzero corank-zero constituent, keeping the real saturation
and torsion terms gives the local formula
\[
 v_2u(0)=v_2\exp^*z-s-\sum_v g_v+2g+1-r_1-e_2.
\]
At an active good odd prime $g_v=2$, since $A^{I_v}=E[2]$.
Unused terms cancel (3.1); all other terms belong to the fixed local
list. This proves (3.2) with the indicated active-prime coefficient.
\end{proof}

\subsection{Nonvanishing with all support restrictions}\label{sec:nonvan}
\begin{proposition}[Nonvanishing in a prescribed unit filter]\label{ex:sec7-nonvanishing}
The curve $E$ is the fixed curve of this section. The following
assertion allows any nonempty signed unit filter and any additional
finite avoidance set; its bound on the number of prime factors does
not depend on either set.
A signed unit filter prescribes a real sign and unit squareclasses at
a finite set $S$ containing $2,3$ and the bad primes. Let $B$ be any
additional finite set of forbidden primes. If its root number is even,
there is a squarefree parameter in this filter, coprime to $B$, with
at most five prime factors and nonzero central value. In an odd filter
there is a simple-zero parameter, with no asserted prime-count bound.

\end{proposition}
\begin{proof}\leavevmode
Here is the reduction to the precise classical statements.
Choose a signed prime $q_0$ in the prescribed unit classes by
Dirichlet's theorem, avoiding $S\cup B$. Apply Hoffstein--Luo
\cite{hl} in the even case, or Bump--Friedberg--Hoffstein
\cite{bfh} in the odd case, to the actual curve $E^{(q_0)}$.
Require their positive discriminant $d$ to split at
$S\cup B\cup\{|q_0|\}$ and at the conductor primes. Then $q_0d$ has
the prescribed original classes. Hoffstein--Luo gives at most four
prime factors of $d$, so at most five in total.
The BFH clause used in the odd case is the simple-zero clause.

The root-number check includes ramified dyadic quadratic characters.
At the primes in $S$ equal squareclasses give isomorphic local curves.
At a good odd prime newly ramified in the parameter, the root factor is
$(-1/p)$ \cite[Proposition 3.3]{desjardins}. Its product outside $S$
is fixed by the real sign and the dyadic unit class. Thus unit classes
$3$ or $7$ modulo eight are allowed; the field discriminant may be four
times the parameter. We do not invoke an odd-discriminant root formula
on those parameters.
See also Proposition~\ref{ex:app-five-prime} for the same
five-prime bound and its subsequent use in the companion construction.
\end{proof}

\subsection{Coefficient symbols and the finite address lemmas}
We state the finite interface before using it for the even upper bound.
The graph language records quadratic residue data of the coefficient
indices. It does not replace a Selmer group by an abstract graph.
Throughout this subsection all vertices are in the rational
full-two-torsion case. Thus there is one independent residue bit per
pair; the split/simple distinction needed for other two-torsion fields
does not occur here.

\begin{definition}[Compatible typed graphs]\label{ex:sec7-typed-graph}
Fix the finite support $S$ of a coefficient source, containing two,
and the finite quadratic characters used to read its coefficients.
The \emph{type} of a prime $p\notin S$ records its values under these
characters and its unit squareclasses at $S$, including its class
modulo eight. We write the resulting finite elementary two-group
additively as $\mathcal L$. We retain all relations among these
coordinates; a tuple of arbitrary character values need not be a type.
Additional fixed splitting conditions restrict us to their actual
realizable types. For an odd theta stage, its reduction prime is part
of this fixed support during that stage.
Let $S_{\rm lab}$ also contain the ramification of any additional
stationary label tests. Enlarging $S_{\rm lab}$ records labels and
prime exclusions; it does not change the coefficient source or add
relaxed places to a Selmer complex.

A \emph{compatible typed graph} consists of a finite ordered vertex
set $V$, an individually realizable type $\ell_v\in\mathcal L$ for
each vertex, and quadratic residue bits. For prime representatives,
write $z_{vw}=0$ or $1$ according as $(p_v/p_w)=1$ or $-1$.
For each unordered pair just one direction is freely specified; the
other satisfies
\[
 z_{vw}+z_{wv}=\epsilon_v\epsilon_w,\qquad
 \epsilon_v=(p_v-1)/2\pmod2,
\]
where $\epsilon_v$ is part of the type. There are no self-residue bits.
Only these compatible, realizable data are used below. A fixed filter
prescribes the total type $\sum_v\ell_v$ needed for the coefficient
index, including the signed parameter convention of the source.
If that convention is $h=\sigma D$, then $D=\prod_v p_v>0$ is the
Fourier index and $\sigma$ is fixed by the source. The signed factors
$p_v^*=(-1/p_v)p_v$ will be used when the same graph defines a family
of quadratic characters.

Contracting two vertices replaces their type by its sum and every
incidence to an external vertex by the sum of the two incidences.
Equivalently its template is the product $p_vp_w$; it is represented
by a fresh prime with those data. Reverse incidences are determined
by reciprocity. No bit internal to the contracted cluster is retained.
A vertex is \emph{primary} if its generator is a square unit at the
fixed support and its type is neutral in every retained quadratic
and stationary coordinate. It is \emph{neutral relative to a graph}
if in addition all its incidences to the other vertices are zero.

Two symbols may read projections of one common type group. The
projections retain the relevant coordinates, preserve primary types,
and commute with the product and external-incidence rules just given.
We require each indicated projection to be surjective onto its own
realizable alphabet; we do not require the map to their product to
be surjective.
\end{definition}

\begin{definition}[Coefficient symbols]\label{ex:sec7-symbol}
A \emph{symbol} is the function $A$ with values in $\F_2$ obtained
from one fixed coefficient source and one fixed normalized depth
$\delta$. On an allowed positive squarefree index $D$ it is
\[
 A(D)=
 \begin{cases}
  1,&v_2(a(D))=\delta+\omega(D),\\
  0,&v_2(a(D))>\delta+\omega(D),
 \end{cases}
\]
and it is zero off that source's filter. For the even source, exact
phase removal makes the relevant coefficient rational. For the odd
source, (3.8) gives the required phase removal modulo the retained
power $2^{M_\theta}$; use its integer representative at that precision.
A coefficient zero modulo $2^{M_\theta}$ is invisible. Whenever a
unit is tested, the precision exceeds $\delta+\omega(D)$, so the test
is independent of the representative. No exact rationality of an
arbitrary lifted odd theta coefficient is asserted. The lower bound making these alternatives exhaustive
is part of the coefficient construction, not part of this definition.
For the even source it follows from (3.2) and the coefficient-square
formula of Proposition~\ref{ex:app-even-phase}.

For an odd source this definition at one finite stage is used only
through a specified finite squarefree-weight cutoff and at a stage
where the lower bound holds through that cutoff. The chosen cutoff
and all required divisions must be included in the operator budget.
The fixed minimum and the compatible finite tables are constructed
in Propositions~\ref{ex:app-odd-minimum} and~\ref{ex:app-odd-stabilization}.
The resulting limit symbol is a function on finite compatible typed
graphs. It is not a valuation rule at all weights of one finite form.
The same source and depth are kept when a vertex is omitted.
\end{definition}

\begin{lemma}[The coefficient rules]\label{ex:sec7-coefficient-rules}
After the stated coefficient lower bounds have been proved, the even
symbol and the stabilized odd symbol satisfy the following rules on
compatible typed graphs. At an odd stage a rule is asserted only when
all its inputs, outputs and coefficient divisions lie in the chosen
finite budget. Configurations in these rules, and outputs of
contractions, are nonempty.
\begin{enumerate}
\item \emph{Graph dependence and whole-prime omission.}
For a new vertex $q$ adjoining a list $U$, its reading factors through
the elementary quadratic extension unramified outside $S\cup U$.
It is determined by its type and its incidences with $U$.
The reading after omission of $p\in U$ uses the same original form,
phase characters and depth, with varying support
$U\setminus\{p\}$. That extension is unramified at the whole
rational prime $p$. Changing inertia at the omitted prime has no
effect on this reading.
\item \emph{Contraction.}
For the independent bit of a pair $v,w$, write
$\Delta_{vw}A=A|_{z_{vw}=0}+A|_{z_{vw}=1}$, changing the reverse bit
as required by reciprocity. Then
\[
                 \Delta_{vw}A(\Gamma)=A(\Gamma/(v,w)).
\]
The graph on the right has the summed type and \emph{every} summed
external incidence from Definition~\ref{ex:sec7-typed-graph}. Its
fresh prime representative is read using the omitted-prime source
in the first rule.
\item \emph{Primary isolation.}
If a configuration contains a primary vertex neutral relative to
all its other vertices, its symbol is zero.
\end{enumerate}
These are identities of the actual coefficient sources, rather than
axioms imposed on an arbitrary binary function.
\end{lemma}
\begin{proof}\leavevmode
Proposition~\ref{ex:app-coefficient-rules} proves the rules from the
integral shell recurrence, the unary subtraction and the weight-two
trace and inertia identities. Divisibility by the integer two precedes
reduction modulo the residue field. In the present full-two-torsion
case every new vertex is split in $\Q(E[2])=\Q$; the contraction
coefficient is one. The omitted form is constructed from the same
original source, so its unramifiedness removes the entire prime and
does not choose a new minimum. The final fresh-copy comparison in that
proof shows dependence on the whole typed graph, including overlapping
prime realizations. The same argument proves the rules for the
even source; no odd normalization is needed in that application.
\end{proof}

\begin{lemma}[Finite realization]\label{ex:sec7-realization}
Keep finitely many old prime vertices and all their mutual data fixed.
Given individually realizable types for finitely many new vertices,
every prescription of the independent forward bits involving a new
vertex can be realized by distinct fresh primes outside any prescribed
finite avoidance set. This includes primary vertices and extra fixed
splitting conditions. A product template can likewise be replaced by
a fresh prime with the same type and all external incidences whenever
that output type is one of the permitted realizable types.
\end{lemma}
\begin{proof}\leavevmode
For rational quadratic characters the fixed type is a reduced residue
condition at a modulus supported on $S_{\rm lab}$, including eight. Choose the
new primes successively. A desired residue bit at an old prime is a
nonzero residue class there. The Chinese remainder theorem combines
these classes with the fixed type, and Dirichlet's theorem supplies
infinitely many primes in the resulting reduced residue class.
At each step avoid all previously chosen primes and the finite
avoidance set. Reciprocity gives the required reverse bits. Apply
the same construction to the external residues and support type of
a product template. The graph dependence in
Lemma~\ref{ex:sec7-coefficient-rules} makes replacement by this fresh
representative legitimate for the symbol.
\end{proof}

\begin{definition}[Unit seeds and partitionable lists]\label{ex:sec7-seed-partition}
A \emph{unit seed} is a finite compatible typed graph with symbol
one. A \emph{minimum nonempty seed} has the least positive vertex
count among such seeds. This terminology does not exclude the
possibility that the empty configuration is also a unit.
A nonempty prime list is \emph{partitionable} if its vertices admit
a partition into nonempty clusters, with a contraction tree in each
cluster, whose contracted graph is a unit seed. These trees and the
resulting types and incidences constitute a contraction certificate.
This condition makes no assertion about a Selmer group.
\end{definition}

\begin{lemma}[Minimum seeds and contraction certificates]\label{ex:sec7-seed-certificate}
At the least positive vertex count of a unit seed, its value is
independent of all internal pair bits and depends only on its vertex
types. A contraction certificate for a nonempty list makes its
polynomial in the independent pair bits nonzero. In particular, a
list with the same total type as a minimum nonempty seed is
partitionable if it contains anchors for the first $r-1$ seed types
and a nonempty remainder, where $r$ is the seed's size.
\end{lemma}
\begin{proof}\leavevmode
If $r\ge2$, a nonzero derivative in an internal bit would, by
contraction, give a unit on $r-1$ nonempty vertices. This contradicts
minimality. If $r=1$ there is no internal bit. Iterating the contraction
identity along a forest gives its contracted symbol; an output equal
to one therefore proves that the original polynomial is nonzero.
For the last assertion keep the $r-1$ anchors and merge the whole
remainder into the last cluster. Equality of total types forces the
last seed type. The internal bits of the minimum seed do not matter.
All contractions have nonempty output.
\end{proof}

\begin{definition}[Addresses and terminal polynomials]\label{ex:sec7-address}
For an integer $b\ge1$, an \emph{address} is $x\in\F_2^b$.
A vertex with activation $\lambda\in(\F_2^b)^*\setminus\{0\}$ is
present at $x$ exactly when $\lambda(x)=1$. Stationary vertices are
present at every address. At a fixed address, retain all vertex
types and any fixed external incidences, and regard a symbol as a
function of the independent pair bits between the active terminals.
Its unique multilinear polynomial over $\F_2$ is the
\emph{terminal polynomial}. Its variables describe possible changes
of those bits; the initial prime realization already has fixed bit
values. Finite differences in these variables are formal derivatives.
\end{definition}

\begin{lemma}[One symbol with a stationary base]\label{ex:sec7-one-symbol-address}
Fix a coefficient symbol satisfying
Lemma~\ref{ex:sec7-coefficient-rules}, a nonempty stationary terminal
list, its types and any stationary external data. Suppose its terminal
polynomial is nonzero and its total data lie in the source's filter.
For every $b\ge1$ and every fixed assignment of the list's pair bits, finitely
many new primary vertices with nonzero linear activations can be
adjoined, with only their new incidences chosen, so that the symbol is
one at every $x\ne0$. The original vertices and all their mutual data
are unchanged. No new vertex is present at zero. Fixed extra splitting
and neutral external-incidence requirements may be included in the
primary types.
For a fixed even source the original prime representatives can be
kept throughout. For a stabilized odd symbol, the fixed data in this
statement are the abstract types and bits; its actual theta stage is
chosen before its prime representatives, as specified in the proof.
\end{lemma}

\begin{lemma}[Two symbols on one address family]\label{ex:sec7-two-symbol-address}
Fix one or two coefficient symbols satisfying
Lemma~\ref{ex:sec7-coefficient-rules}, reading compatible projections
of one common realizable alphabet. Fix $b\ge1$ and terminals with
nonzero linear activations, including at least one terminal for each
nonzero linear form. Stationary external data may also be fixed.
Assume that at every $x\ne0$ the active configuration is nonempty,
its total data satisfy every required filter, and each symbol's
terminal polynomial is separately nonzero.
Then, starting with any fixed assignment of these compatible types
and bits, finitely many new primary vertices with nonzero linear activations can be
adjoined so that all the symbols are one at every $x\ne0$.
Only incidences involving the new vertices are chosen. All original
data and the zero address are preserved, and the same fixed splitting
and neutral external-incidence conditions may be imposed as above.
If an odd limit symbol is present, this is first a statement about
the finite typed graphs. It does not fix rational prime representatives
before the eventual theta stage.
\end{lemma}
\begin{proof}[Proof of the two address lemmas]
The forest and network calculation in
Proposition~\ref{ex:app-simultaneous-address} proves both statements.
For the stationary case its grounded-network columns attach to a
fixed terminal, so no terminal activation anchors are required.
For variable terminals the stated anchors give the row-sum identity
on every active restriction. Singular networks isolate prescribed
derivatives at one nonzero address; varying networks supply the
corresponding difference when two symbols share a variable. The
two-symbol proof takes the product of the top forms in the ordinary
polynomial ring before Boolean reduction. Its nonzero coefficient
forces a single assignment that works at all nonzero addresses.
It does not require the two symbols to have a common nonzero
evaluation before the auxiliary vertices are added.

Each construction uses a finite number of vertices and coefficient
readings once $b$ and the terminal polynomials are fixed. For an odd
symbol, include all of them in one finite weight and operator budget,
then use Proposition~\ref{ex:app-odd-stabilization} to choose one actual
theta stage carrying those tables and divisions. Only afterward use
Lemma~\ref{ex:sec7-realization} to choose its primes. Any old
prime representatives already chosen at that stage may be retained
when they have its full prescribed types and bits, by the same
realization lemma. In the even stationary application all original
primes can be retained from the outset. The number of
auxiliaries and the required theta precision may depend on $b$ and
the terminal data. The source's depth and the valuation constants
attached to a unit do not change. No common finite stage for all
values of $b$ is asserted.
\end{proof}

\paragraph{The extra upper-bound argument.}
Nonvanishing alone does not bound the valuation of $L_*$, and (3.2)
by itself gives a lower bound. The needed upper bound uses a separate
even-symbol selection and transfer argument, which we now give.

\begin{proposition}[The partitionable even upper estimate]\label{ex:sec7-even-upper}
Fix one even coefficient source, its signed unit filter and its depth.
If a parameter $d$ in that filter has a nonempty partitionable prime
list and $E^{(d)}$ has full two-primary Selmer corank zero, then (3.3b)
holds. Its constant depends only on the fixed source, not on the
number or sizes of the primes of $d$ or on later finite exclusions.
\end{proposition}
\begin{proof}\leavevmode
For the even coefficient system described below, a unit symbol at a
parameter $d$ gives
\[
       L_*(d)\ne0,\qquad v_2L_*(d)\le2\omega(d)+C_0.              \tag{3.3a}
\]
This follows from its same-filter coefficient-square formula and its
fixed normalization. That even normalization uses the lower bound
from (3.2), not the as-yet undefined odd minimum.
Partitionability is the nonempty contraction-certificate condition in
Definition~\ref{ex:sec7-seed-partition}. By
Lemma~\ref{ex:sec7-seed-certificate}, its iterated finite difference
is nonzero and hence so is the terminal polynomial. An empty unit,
when present, satisfies its coefficient bound directly and does not
require the stationary address lemma.

For a partitionable parameter $d$ already known to have Selmer
corank zero, keep all its prime factors stationary.
Lemma~\ref{ex:sec7-one-symbol-address} gives arbitrarily large binary
cubes based at
$d$, with no old prime removed, whose other addresses satisfy (3.3a).
The fixed-form common scalar in (3.2) has one uniformly cleared
denominator on every such cube. Choose a two-adic precision larger
than the nonzero-address scalar valuation bound and its fixed
clearing exponent. Its sign-evaluation coefficient is a Boolean
function with the bounded bit degrees proved in
Section~\ref{sec:limits}. A cube dimension larger than their sum
therefore supplies a nonzero address congruent to the base at this
precision. If the base scalar were zero or had larger valuation than
the bound, so would that nonzero address, a contradiction.
The corank-zero central comparison in (3.2) consequently gives
\[
 L_*(d)\ne0,\qquad
             v_2L_*(d)-2\omega(d)-s(d)\le C_1.                  \tag{3.3b}
\]
The constants are fixed by the even system. The dimension needed
may depend on $d$; the valuation bound does not. This is the
cyclotomic missing-address argument on the fixed-form Kato scalar,
not the later Heegner transfer.
\end{proof}

\begin{proposition}[A bounded companion in the prepared filter]\label{ex:sec7-rough-companion}
For (3.3), restrict $h$ to positive squarefree parameters in the
prepared odd-sign unit filter containing one, prime to its fixed
support. Thus $h\equiv1\pmod8$ and $w(E^{(h)})=-1$.
The simple-zero parameters used in (3.4)--(3.6) lie in this filter.
There are constants $A,C_E$ fixed before $h$ varies such that,
outside any prescribed additional finite avoidance set, one can
choose a negative odd fundamental discriminant $k'$ which is a
square at every place dividing $2Nh$, avoids the extra-unit fields,
and has $L_*(hk')\ne0$. It satisfies (3.3), in particular
$\omega(k')\le A$. The valuation estimate uses the fixed even
sources of the original $E$, although the analytic nonvanishing
theorem will be applied to an auxiliary twist of $E$.
\end{proposition}
\begin{proof}\leavevmode

Here is why the actual rough companion is partitionable.
Prepare an even symbol for every possible total type at the finite
fixed support. There are finitely many. In our full rational
two-torsion case the alphabet is all split. A least nonempty unit
seed has no dependence on internal edge bits: an edge derivative
would be a smaller nonempty unit configuration. Its individual
labels determine it. The negative-discriminant parity at two makes
the total label nonzero, so this seed is nonempty.
A list with the same total label and enough copies of the seed
labels contracts to it: keep the first $r-1$ seed anchors and merge
all remaining vertices into one nonempty cluster. Its label is
forced to be the last seed label. Thus the list is partitionable,
regardless of the number or labels of the remaining vertices.

Before $h$ varies, fix a bounded multiset of dummy labels supplying
the required anchors and reserves for every prepared total type.
For each $h$, realize these labels by distinct primes avoiding the
fixed support, the support of $h$, and all already fixed exclusions.
Put in $B$ the entire fixed support, the support of $h$, the chosen
dummy primes, and every other already fixed forbidden prime.
Include any already fixed reduction prime in $B$ as well.
Prescribe the total $k'$ to be a negative odd fundamental
discriminant, square at every place of $2Nh$ and outside the
extra-unit cases. Absorb $h$ and the dummy product into the curve
and use the preceding nonvanishing theorem in its even filter.
It adds at most five new primes. For the actual resulting total
type, the already chosen dummy supply makes $hk'$ partitionable.
Its analytic nonvanishing gives Selmer corank zero by the forward
rank-zero theorem. Apply (3.3b), obtaining
\[
 \omega(k')\le A,\qquad
 v_2L_*(hk')\le 2\omega(h)+2\omega(k')+s(hk')+C_E.                \tag{3.3}
\]
Here $A$ is the fixed dummy count plus five. This is the argument
of the partitionable-even and rough-companion lemmas, with their
stationary primes and order of choices retained.
Neither prime sizes nor a uniform analytic threshold are asserted.
Excluding a later coefficient prime does not enlarge the old
Selmer support.
The finite source preparation and the distinction between fixed dummy
labels and their later prime realizations are also detailed in
Proposition~\ref{ex:app-rough-companion}.
\end{proof}

\subsection{The genus lower bound precedes normalization}
Fix a negative genus partner $h_*$ with $L_*(h_*)\ne0$.
Here $h_*$ is chosen odd fundamental with $|h_*|>4$ and in the
same squareclasses at $2N$ as the prepared positive filter containing
one. Throughout (3.4)--(3.6), $h$ belongs to that filter and is
prime to the fixed support, including $h_*$; the companion $k'$
is supplied by the preceding construction with the primes of $h_*$
also in its avoidance set.
For a known simple-zero $h$, let $P_h^{\rm gen}$ be the unaveraged
conductor-one genus sum over $\Q(\sqrt{hh_*})$.
Let $P_{h,k'}^{\rm rc}$ be the unaveraged conductor-$|h|$ ring-class
sum over $\Q(\sqrt{k'})$, weighted by $\chi_h$.
They are different points in different native lattices.
Write $j_{\rm gen}(h)$ and $j_{\rm rc}(h,k')$ for their index valuations.

\begin{proposition}[The genus bound before odd normalization]\label{ex:sec7-genus-lower}
Use $h_*,h,k'$ and the two native index valuations defined in the
preceding paragraph, with $h$ a known simple-zero parameter and $k'$
as in Proposition~\ref{ex:sec7-rough-companion}. Then (3.4)--(3.6)
hold. In particular their constant in (3.6) is independent of $h$.
The arithmetic lower estimate (3.5) holds uniformly for the admissible
companions with $\omega(k')\le A$; it is only asserted here at these
known simple-zero parameters.
\end{proposition}
\begin{proof}\leavevmode
The explicit height formula \cite[Theorem 1.1]{cst} applies twice.
All level primes split, conductors are prime to $N$, extra-unit fields
are excluded, and the parametrization is the same. Its conductor factor
is $c\sqrt{|\operatorname{disc}K|}$ and its sum is unaveraged.
The two Rankin products have the common factor $L(E^{(h)},s)$.
Each rational free line has index at most two in its native rank-one
lattice: conjugation fixes the free class, and the norm makes twice the
class rational. Thus
\[
 2j_{\rm rc}(h,k')-2j_{\rm gen}(h)
   =v_2L_*(hk')-v_2L_*(h_*)+\epsilon,\qquad |\epsilon|\le2,        \tag{3.4}
\]
up to the one fixed initial period convention.

The arithmetic estimate of Section~\ref{sec:limits}, for each fixed $A$,
is
\[
 2j_{\rm rc}(h,k')\ge
 4\omega(h)+s(h)+s(hk')-C_R(E,A),\qquad \omega(k')\le A.           \tag{3.5}
\]
Combining (3.3)--(3.5) cancels $s(hk')$ and gives
\[
 2j_{\rm gen}(h)\ge 2\omega(h)+s(h)-C.                            \tag{3.6}
\]
This calculation uses no odd normalization.
For the forward reference in (3.5), combine
Proposition~\ref{ex:limits-uniform-lower-bound} with the same-class
central identity in Proposition~\ref{ex:arith-central-calibration}.
Those arguments use the known nonzero point in the present setup;
they do not use an odd coefficient minimum or the final family.
\end{proof}

\subsection{A fixed nonzero witness and actual odd forms}
Choose one positive squarefree $D_w$ in the odd filter containing one,
after adjoining the primes of $h_*$ and their compatible unit tests.
Section~\ref{sec:nonvan} supplies $L'(E^{(D_w)},1)\ne0$.
The explicit height formula makes $P_{D_w}^{\rm gen}$ nontorsion.
Its unwanted $h_*$ component has rank zero. One fixed multiplier $L$
kills the bounded-degree torsion and the projection denominator, so
$Q=LP_{D_w}^{\rm gen}$ is a rational nontorsion point on $E^{(D_w)}$.

Let $C$ be complex conjugation on its Tate module and choose a nonzero
integral row $b$ with $bC=b$. The Kummer translation image of $Q$ is a
full lattice. Indeed, if it vanished, its rational cocycle would descend
to the Tate image. A nonidentity central scalar $z$ gives
\[
 (z-1)c(g)=(g-1)c(z),
\]
so that descended cocycle is a coboundary, contrary to Kummer
injectivity. Irreducibility makes a nonzero translation span the Tate
space, and intersection with any fixed finite kernel preserves full
rank. Twisting preserves open image by restriction to an index-at-most-two
subgroup.

Choose one such translation $v$ with $2b(v)\ne0$.
Multiplying conjugation by this translation gives square translation
$(1+C)v$, detected by $2b(v)$ with one fixed valuation $c_w$.
The element is a single actual Galois element on the compositum.
It is trivial on $\sqrt{D_w}$ and nontrivial on $\sqrt{h_*}$ because
$D_w>0>h_*$. Its restrictions therefore already agree on intersections.
Finite Chebotarev, with a place chosen above the resulting prime, gives
arbitrarily advanced stages $(M_\theta,\ell_0)$ such that
\[
 \rho(\operatorname{Fr}_{\ell_0})\equiv C\pmod {2^{M_\theta}},
 \quad (D_w/\ell_0)=1,\quad(h_*/\ell_0)=-1,
 \quad v_2\lambda_{M_\theta}(\overline Q)=c_w<M_\theta.            \tag{3.7}
\]
Here $\lambda_{M_\theta}=b\kappa_{M_\theta}$ is the actual finite-field
Kummer map on $E(\F_{\ell_0^2})$. It has Frobenius sign $+1$.
All reductions use that same place. This is a fixed-witness statement,
not a bound for arbitrary points or arbitrary reduction primes.

The weighted ternary theta construction gives actual holomorphic
forms at these stages. At original level primes the root weight uses
the summand selected by the cyclic level structure. At two, the
integral spectral projectors of $(1+x)/2$ specify that summand; the
operator itself is not in general idempotent. The nilpotent weight at
$h_*$ and the division-algebra residue-line weight at $\ell_0$ have the
prescribed conjugation and uniformizer signs. Their product is even.
The root-character field is fixed before $\ell_0$ and $M_\theta$.

Appendix~\ref{app:odd-source} gives these functions, their exact
transformation rules, and the direct extraction using the form $Q/m$.
Poisson summation for the resulting compactly supported local functions
gives the scalar transformation law and holomorphy at every cusp.
A fixed stabilizer multiple, for example one divisible by
$|\operatorname{GL}_4(\F_3)|$, clears the finite automorphism groups.
For the Hecke recurrence, a split or irreducible nonzero reduction has
$1+(-n/p)$ invariant lines, a nonzero nilpotent has one, and zero has
$p+1$. The unique integral primitive nilpotent neighbor supplies the
remaining square term. Tracking the scaling characters gives the
recurrence with the same character convention as the coefficient
comparison.

For an admissible $D$ with $(D/\ell_0)=1$, the oriented embedding
count \cite[Theorem 3.1]{cornutjetchev}, with changed set $\{\ell_0\}$, gives
\[
 a_{M_\theta}(D)
   =c_0\,2^{\omega(N)}\zeta_{M_\theta}(D)
       \lambda_{M_\theta}(\overline{LP_D^{\rm gen}})
       \pmod {2^{M_\theta}},                                    \tag{3.8}
\]
where $c_0\ne0$ is fixed and $\zeta$ is a permitted root-of-unity phase.
For admissible $D$ with $(D/\ell_0)\ne1$, including $\ell_0\mid D$,
the local Schwartz support instead gives $a_{M_\theta}(D)=0$.
Formula (3.8) is asserted only on the first support.
There is no division by a class number. Formula (3.7) makes $D_w$
visible with fixed upper valuation in (3.8). It remains on every
stage's support, including the extra test $(D_w/\ell_0)=1$.

After phase removal, the coefficients used for valuation are integers.
Equation (3.6) gives a uniform lower bound for
$v_2a_{M_\theta}(D)-\omega(D)$ on visible bounded-weight sequences.
Take the finite integer values constant on a set belonging to the
chosen nonprincipal ultrafilter. This set of values is nonempty by
$D_w$, is bounded below, and hence has an attained minimum $\delta$.
If the resulting lower inequality failed for a fixed weight bound on
an ultrafilter-large set, choosing a violating index and then one of
finitely many smaller integer values would contradict minimality.
For any fixed finite operator budget, the coarse lower bound handles
the large weights and stabilization handles the bounded weights.
Consequently every required division is valid at one sufficiently
advanced actual stage. No infinite-precision modular form or
arbitrary-ring base-change assertion is needed.

\subsection{One common family and the final order of choices}
The even source is a fixed algebraic vector in the full good-Hecke
eigenspace. Waldspurger's Corollary 2, including its same-dyadic-squareclass
removal of the auxiliary condition, compares the actual squarefree
Fourier functionals \cite{waldspurger}. Local upper-unipotent projections
retain a nonzero initial functional; smoothness supplies lower-unipotent
invariance. The algebraic linear conditions choose a vector over a
number field. Galois conjugation followed by the matching finite
character twist shows that each normalized coefficient ratio is a
rational number times a two-power root of unity. This supplies the
valuation and nonvanishing comparison on the whole filter.
Only flat field extension of the algebraic model is used.
The comparison is proved at the original support $S_0$ before any extra
coprimality sieve $B$. The sieve is applied last, and every later good
recurrence prime avoids the full support $S_0\cup B$.

For both coefficient systems the integral operators have the required
omission and contraction rules. At a new prime $V_p$ raises the level
exponent only to one, with ramified quadratic character. The recurrence
makes the relevant differences divisible by the integer two before
residue reduction. Multiplying by ordinary theta, then taking a fresh
trace, removes the integral unary corrections. Two fresh traces vanish.
The matrix trace identity gives $t(gu^2)=t(g)$, hence elementary
quadratic labels. At a new odd level prime the weight-two local representation
has one unramified and one ramified line. The normalized-induction
and Whittaker calculation below identifies the Fourier $U_p$
eigenvalue with arithmetic Frobenius on the unramified line.
The diagonal-matrix calculation gives the two inertia differences.
These identities prove omission and primary isolation for the same
original form and depth, without choosing another minimum.

If the intrinsic odd limit symbol is a unit at the empty configuration,
choose an actual stage carrying that value with $M_\theta>\delta_o$.
The index $D=1$ lies in every moving support and then satisfies
$v_2a_{M_\theta}(1)=\delta_o<M_\theta$. The visibility implication
(O.8) gives $\ord_{s=1}L(E,s)=1$, which is exactly (RC).
No final companion or determinant bound is needed in this case.
For the rest of the final-family and missing-address argument, assume
that the empty configuration is not a unit. The bounded-weight witness
then supplies a least nonempty odd seed, as in
Appendix~\ref{app:endpoints}.

\paragraph{Fixing the final companion.}
Use a separate even symbol on negative discriminants square at $2N$,
excluding the fixed odd support, the ramification support of $E[8]$,
and the exceptional imaginary fields. Choose a unit $k$ of this
symbol. Then $L(E^{(k)},1)\ne0$. Fix $k$ and $K=\Q(\sqrt{k})$
before choosing the binary dimension. At least one prime ramified
in $K$ is outside all the old fixed label extensions and torsion
ramification. Choose every later moving theta prime $\ell_0$ away
from $k$. The intrinsic odd symbol continues to omit the primes of $k$.

At any such stage, the normal closure of the old odd label extension
is unramified outside its old support and $\ell_0$. The fresh ramified
prime of $K$ lies outside this set, so its intersection with $K$ is
$\Q$. Thus each old odd label still has a lift split in $K$.
The added stationary tests at the primes of $k$ may have relations;
include those relations in the common alphabet, rather than requiring
independent arbitrary prescriptions. Stabilize its separate projection
maps on the same ultrafilter. This preserves the old odd minimum.
Fresh ramification also makes $K$ independent of $\Q(E[8])$ and of
the old fields $\Q(i)$ and $\Q(\sqrt{\ell})$.

The forward theorem at this known nonzero central value gives
$s_2(E^{(k)}/\Q)=0$. The sum-and-difference isogeny over $\Q$ gives
\[
 s_2(E/K)=s_2(E/\Q)+s_2(E^{(k)}/\Q)=1.
\]
Indeed its kernel and cokernel on the relevant restriction and
corestriction maps are killed by a fixed power of two. Tensoring the
compact Selmer modules, equivalently the Pontryagin duals of the discrete
Selmer groups, with $\Q_2$ gives the direct sum and hence this identity
of coranks.
No integral direct-sum assertion is used.

\begin{lemma}[The nonempty seed for the present curve]\label{ex:sec7-ksplit-seed}
For $E=E_\ell$ and the $k,K$ just chosen, the positive even source
for $E^{(k)}$, normalized at index one, has a nonempty unit seed
consisting of $K$-split primes.
\end{lemma}
\begin{proof}\leavevmode
This is the specialization to the roots $0,\ell,-3\ell$ of the
coefficient calculation in Proposition~\ref{ex:app-nonempty-seeds}.
If no such nonempty unit existed, its integral $p^3$ coefficient
and a fresh trace would force
$a_p(E^{(k)})\equiv p+1\pmod8$ at the required good $K$-split primes.
The cited calculation uses positive Fourier indices and the positive
twist character convention; it derives this congruence inside the
assumed absence of nonempty units.

Choose such a prime with $p\equiv1\pmod4$, $(\ell/p)=-1$ and
$(k/p)=1$, avoiding the finite support. The fresh ramification of $K$
permits these simultaneous conditions. Because $k$ is a square modulo
$p$, the reductions of $E^{(k)}$ and $E_\ell$ are isomorphic over
$\F_p$. None of the three nonzero rational points of order two can
be doubled: their necessary square differences include, respectively,
$-\ell,\ell,-4\ell$. Each is a nonsquare under the chosen conditions.
The two-primary reduction group therefore has order exactly four,
contradicting the asserted congruence. This proof uses the particular
curve $E_\ell$; no nonempty-seed theorem for an arbitrary elliptic
curve is being asserted.
\end{proof}

We record the finite combinatorial realization used here. Each label
space is a finite elementary two-group; the common alphabet projects
onto each separately. Product surjectivity is unnecessary. Take
sufficient repeated copies of every label, with total label zero.
Keeping the seed anchors, contract the unused copies into the last
anchor. Minimal positive seed size makes the internal bits irrelevant,
so each of the two terminal polynomials is nonzero. Grounded Laplacian
minors express the constant and pair-contraction operators, including
singular matrices. A cyclic singular network isolates its target
radical and a path network isolates its selected pair coefficient.
In the product of the two top forms, taken in the ordinary polynomial
ring before Boolean reduction, a variable of exponent two uses one
singular program in each factor; exponent one uses one varying program.
The total degree kills the other contractions and leaves coefficient
one. Thus one choice works for both forms. CRT and Dirichlet realize
the prescribed forward residue bits by distinct primes; quadratic
reciprocity supplies the reverse bits.

\begin{proposition}[One family for the two coefficient tests]\label{ex:sec7-final-family}
Keep the fixed $E,h_*$, normalized odd symbol and final $k,K$ of the
preceding paragraphs. Assume, as in the remaining branch above,
that the intrinsic odd empty configuration is not a unit.
In this remaining branch, for every sufficiently large $b$ there is one
finite set $Q$ of distinct good odd primes split in $K$ and outside the
fixed support, and nonzero linear forms $\lambda_q:\F_2^b\to\F_2$,
such that the family
\[
 h_x=\prod_{q\in Q}(q^*)^{\lambda_q(x)}
 \quad(x\in\F_2^b),\qquad q^*=(-1)^{(q-1)/2}q,
\]
has $h_0=1$, all fixed local types unchanged, and, for $x\ne0$,
\[
 \ord_{s=1}L(E^{(h_x)},s)=1,\quad L_*(h_xk)\ne0,\quad
 j_{\rm gen}(h_x)\le\omega(h_x)+C_o,\quad
 v_2L_*(h_xk)\le2\omega(h_x)+C_e.                                \tag{3.9}
\]
The constants are independent of $b$. By (3.4), now with fixed $k$,
$2j_{\rm rc}(h_x,k)\le4\omega(h_x)+C_J$.
\end{proposition}
\begin{proof}\leavevmode
The packet argument just given supplies separate nonzero terminal
polynomials in one common alphabet. Apply
Lemma~\ref{ex:sec7-two-symbol-address}. The full proof in
Proposition~\ref{ex:app-simultaneous-address} and
Proposition~\ref{ex:app-final-family} preserves the projected total
labels, all stationary tests and the original two normalizations.

The order is fixed data and $D_w$, then the genus lower bound and odd
minimum, then the final $K=\Q(\sqrt k)$ and common alphabet, then
determinant precision $M$, then $b$ and its finite operator budget,
then a sufficiently advanced actual theta stage $M_\theta$ and its
prime realization. In particular $M_\theta$ is not $M$.
The finite graph construction, its local conditions and the coefficient
tests all occur at that same final stage. Appendix~\ref{app:endpoints}
spells out the absolute five-prime completion, normalization constants,
and the final use of the one corrected coordinate.
\end{proof}

% End 03-restricted-converse.tex


% Begin 03a-arithmetic.tex
\section{Actual classes, local conditions and central comparison}
\label{sec:arithmetic}
Let $K=\Q(\sqrt k)$ be imaginary, split at $2N$, and different from
$\Q(i)$ and $\Q(\sqrt{-3})$. Active primes are good odd primes split
in $K$. Derivative primes are good inert primes outside all current
supports. The actual finite coefficients are
\[
 A_i=(\Z/2^i\Z)[t]/t^i,\quad
 M_i=(T_2E\otimes A_i)(\chi_h\Psi_i).
\]
For a fixed binary group $G$, use $A_i[G]$ and the universal quadratic
character. The character $\Psi_i$ is finite, trivial at $2N$, ramified
at a moving conductor prime $r_i$, and has prescribed odd scalar
exponents at the finite active list. In a binary family let $Q$ be
the union of its prime supports. Apply
Proposition~\ref{ex:arith-moving-character} to an orientation above
every $q\in Q$, prescribing $\xi_i(\operatorname{Fr}_q)=1$ at each
finite stage. This prescription includes the primes unused by an
individual character, including every graph prime at augmentation.
The same $\Psi_i$ is used for all characters of the family. The finite
set $Q$ is fixed before choosing the sequence of moving primes $r_i$;
the trace bound $c_E$ does not depend on its size. The prime $r_i$
varies with precision. There is no assertion
that one fixed odd prime supports an infinitely ramified $\Z_2$ character.

\begin{definition}[Coefficient and local-condition conventions]
\label{ex:arith-data}
In Sections~\ref{sec:arithmetic} and~\ref{sec:limits}, retain the fixed
non-CM curve $E$, parametrization $\pi$, Tate lattice, and old local
types of Section~\ref{sec:rc}. Put $\Lambda=\Z_2[[t]]$ and
$\mathscr O=\widehat{\Lambda_{(2)}}$. A binary family uses one finite
elementary abelian two-group $G$, one real sign and common local
squareclasses at $2N$. Its universal varying quadratic character
is therefore trivial at the real place and at $2N$.
The active primes are good odd primes split in the indicated imaginary
quadratic field $K$; the moving character is the finite character
specified above, trivial at $2N$. The primitive and relaxed complexes
use the actual derived Kummer, unramified, full and orthogonal local
conditions specified below, with their maps and homotopies.
At derivative primes the four local conditions are those of
Lemma~\ref{ex:limits-derivative-switch}. Every determinant coordinate
uses the retained integral determinant basis.
\end{definition}
\begin{lemma}[Conjugate coefficients and local duality]\label{ex:arith-conjugate-duality}
For the finite coefficient module $M_i$ and coefficient ring
$R_i=A_i$ or $A_i[G]$ above, let $\tau$ be an actual complex
conjugation. Pullback of the Galois action gives an $R_i$-linear
identification $\tau^*M_i\simeq M_i^*(1)$ through
$(v,w)\mapsto e_2(v,\rho_T(\tau)w)$; it fixes $t$ and the group-ring
basis. With the conjugate places, actual orthogonal local conditions
and their isotropy homotopies, finite Selmer duality gives
\[
 C_U(M_i)\simeq
 R\operatorname{Hom}_{R_i}(C_{U^\perp}(M_i^*(1)),R_i)[-3].
\]
The local Kummer and derivative comparisons used here are proved
below and in Lemma~\ref{ex:limits-derivative-switch}.
\end{lemma}
\begin{proof}\leavevmode
\paragraph{The dual coefficient system.}
Fix an actual complex conjugation $\tau\in G_{\mathbf Q}$, so that
$\tau^2=1$. The moving character is constructed on the ring-class group:
\[
 \xi_i:\operatorname{Gal}(H_{r_i}/K)\longrightarrow
       \mathbf Z/2^{n_i}\mathbf Z,
 \qquad \Psi_i=(1+t)^{\xi_i},
\]
where $n_i$ is large enough that $(1+t)^{2^{n_i}}=1$ in $A_i$.
Conjugation acts on this ring-class group by inversion. Consequently
\[
 \Psi_i(\tau g\tau^{-1})=\Psi_i(g)^{-1}\qquad(g\in G_K).
\]
Let $R_i=A_i$ in the scalar construction and $R_i=A_i[G]$ in the
binary construction. Write $\alpha_i=\chi_h\Psi_i$, respectively
$\alpha_i=\chi_0\varepsilon\Psi_i$, where $\varepsilon:G_K\to G$
is the restriction of the universal rational quadratic character.
Then $\varepsilon^\tau=\varepsilon=\varepsilon^{-1}$ and
$\alpha_i^\tau=\alpha_i^{-1}$. Here $\tau^*M_i$ denotes pullback
of the Galois action; it does not act on $R_i$.

The Weil pairing gives the following $R_i$-linear perfect pairing:
\[
 M_i\otimes_{R_i}\tau^*M_i\longrightarrow R_i(1),\qquad
 (v,w)\longmapsto e_2(v,\rho_T(\tau)w).
\]
Indeed $\rho_T(\tau)^2=1$, the two scalar characters multiply to
one, and $e_2$ transforms by the cyclotomic character. Thus
$\tau^*M_i\simeq M_i^*(1)$. We do not identify $M_i$ with its
ordinary dual when $\Psi_i$ is nonquadratic. In this convention $t$
and every element of the group-ring basis are fixed. Inverting $1+t$
is already accounted for by pullback of the character.

At a split place $v$, the dual local condition is transported from
the conjugate place $\tau v$. At a fixed place it is transported
by the actual local representative of conjugation. Kummer conditions
at $2N$ are preserved by this transport. A full condition is paired
with the zero condition; the full local complex is replaced by zero
only at the localizations where its displayed contraction exists.
At a derivative prime the finite and singular planes are checked
with this transported pairing. In the chosen coordinates this is
the pairing $\beta(x,y)=\epsilon\langle x,Jy\rangle$.
The $J$ is therefore part of the duality convention, not an extra
identification of an arbitrary twist with its dual.

Use the coefficient-valued pairing in the finite Selmer diagrams,
with the conjugate local places and their isotropy homotopies. The
resulting duality is
\[
 C_U(M_i)\simeq
 R\operatorname{Hom}_{R_i}
       (C_{U^\perp}(M_i^*(1)),R_i)[-3].
\]
These are the maps and homotopies retained in both limits.
\end{proof}
% Source reconstruction of the finite arithmetic input I08/I09.
% This is a proof fragment, not an assertion about completion of RC.
\subsection{Finite Heegner descent and the derivative identities at two}

Fix the non-CM elliptic curve $E/\mathbf Q$, its conductor $N$, its
Tate module $T=T_2E$, and a modular parametrization
$\pi:X_0(N)\longrightarrow E$. Normalize $\pi(\infty)=0$.
Let $K$ be imaginary quadratic, split at every prime dividing $2N$,
and different from $\mathbf Q(i)$ and $\mathbf Q(\sqrt{-3})$.
The conductor $f$ below is prime to $2ND_K$; it may contain all the
active primes in a finite binary family and the moving prime.
Write $H_f$ for the ring-class field of the order
$\mathcal O_f=\mathbf Z+f\mathcal O_K$.

Let $I$ be any finite set of good inert primes outside $6fND_K$,
and put $n_I=\prod_{\ell\in I}\ell$ and $M_I=H_{fn_I}$.
At the end, a finite character $\theta$ of
$\Delta_f=\operatorname{Gal}(H_f/K)$ will take values in
$A^\times$, where $A=(\mathbf Z/2^m\mathbf Z)[t]/t^b$, or in
$A[G]^\times$ for a finite binary group $G$.
It has the form $\chi_h\Psi$, or the corresponding universal quadratic
character times $\Psi$. At $2N$, $\Psi$ is trivial and the universal
quadratic character has the fixed local type already specified in
the manuscript. These are hypotheses on the actual finite character.
Its construction is separate from the descent proved here.

\begin{proposition}[Finite Heegner descent with one multiplier]\label{ex:arith-finite-heegner}
With $E,\pi,K,f,I,M_I,A$ and the finite character $\theta$ as just
specified, there are integers $e\geq0$ and $L>0$ depending only
on the fixed curve, parametrization and allowed old local types,
and on any other prescribed fixed multiplier requirements, with
the following properties. Form the derivative over the complete
extension $H_{fn_I}/H_f$ and descend it before applying the character.
This gives finite classes $P_{I,m}$ with the old Kummer conditions,
the stated good unramified conditions and full active and moving
conditions. For $I=\varnothing$ the class is the Kummer class of
the $L$-multiplied unaveraged weighted point.

The unaveraged norms are (4.1). If $a\geq m+e$ and every new
derivative prime has $\rho_T(\operatorname{Fr}_\ell)\equiv
\rho_T(\tau)\pmod{2^{a+1}}$, the classes projected to precision
$m$ satisfy (4.2). These assertions hold over $A[G]$ before
character evaluation. They are realized by finite continuous
cocycles, local-condition representatives and boundary witnesses,
with the weighted coefficient map applied to all this data.
There is no multiplier depending on $K,f,I,m$, the truncation
length, the number of active primes or $|G|$.
\end{proposition}
\begin{proof}\leavevmode
\paragraph{1. One uniform torsion annihilator.}
By the non-CM open-image theorem \cite{serre}, choose $n\ge2$ such that
$1+2^nM_2(\mathbf Z_2)\subset\rho_T(G_{\mathbf Q})$, and choose
$g\in G_{\mathbf Q}$ with $\rho_T(g)=uI$, $u=1+2^n$.
For an actual complex conjugation $\tau$ set
\[
 z=g^2\tau g^2\tau^{-1},\qquad e=v_2(u^4-1)=n+2.       \tag{H1}
\]
The element $g^2$ fixes every quadratic field. Its image in the
abelian group $\operatorname{Gal}(H_d/K)$ is inverted by $\tau$,
for every $d$ and every allowed $K$. Thus $z$ fixes every such
$H_d$, while $\rho_T(z)=u^4I$. If $P\in E(H_d)[2^\infty]$,
then $(u^4-1)P=0$, so
\[
                       2^e E(H_d)[2^\infty]=0.          \tag{H2}
\]
This bounds the whole invariant torsion group, with no restriction
on the number of conductor primes. It also holds for a compositum
of the ring-class fields in a finite construction, since $z$ fixes
each of them.

For the matrix-span statement used elsewhere, square preimages of
$I+2^nE_{11}$, $I+2^nE_{12}$ and $I+2^nE_{21}$. The squares belong
to every quadratic subgroup, and their images, together with $I$,
are
\[
 I,\quad I+(2^{n+1}+2^{2n})E_{11},\quad
 I+2^{n+1}E_{12},\quad I+2^{n+1}E_{21}.
\]
Their $\mathbf Z_2$-span contains $2^{n+1}M_2(\mathbf Z_2)$:
subtract $I$, use that $1+2^{n-1}$ is odd, and obtain $E_{22}$
from $I-E_{11}$. This bound also depends only on $E$.

\paragraph{2. Point relations with the full old conductor.}
Choose the CM points $x_d\in X_0(N)(H_d)$ by compatible proper
ideal lattices and the fixed orientations at $N$, and put
$y_d=\pi(x_d)$. The split Heegner hypothesis identifies the curve
$X_K(N)$ in Cai--Chen--Liu, Theorem 2.15 \cite{ccl}, with $X_0(N)$:
$N_1=1$, and the condition $(*)$ reduces to $(d,N)=1$.
Take their old conductor to be any $d$ currently in use and the
new prime outside $6dND_K$. Their factor $u_m$ is one because
$\mathcal O_K^\times=\{\pm1\}$. Applying the elliptic quotient
to that theorem gives
\[
 \operatorname{Tr}_{H_{d\ell}/H_d}y_{d\ell}=a_\ell y_d
       \quad(\ell\text{ inert}),\qquad
 \operatorname{Tr}_{H_{dq}/H_d}y_{dq}
       =(a_q-\sigma_q-\sigma_q^{-1})y_d
       \quad(q\text{ split}).                            \tag{4.1}
\]
The inverse Frobenius pair at a split prime follows because their
product is the Artin symbol of the rational principal ideal $(q)$.
These are unaveraged sums. On degree-zero divisor classes,
$\varphi T_q=[a_q]\varphi$. Since $q\nmid N$,
$T_q[\infty]=(q+1)[\infty]$, so the displayed norm relations
for $y_d=\pi(x_d)$ have no cusp correction.
For other changes of CM convention, let $\mathcal C_N$ be the
subgroup of $J_0(N)(\overline\Q)$ generated by
$[c]-[\infty]$ for all cusps $c$. This is a finite group by
Manin--Drinfeld \cite[p.~59, Theorem]{elkik}.
Fix $C_\pi\geq1$ annihilating $\varphi(\mathcal C_N)$.
This single integer depends only on $N,\pi$.
For an Atkin--Lehner involution $W$ acting on $f$ by $w$, the exact
formula is $\pi(Wx)=[w]\pi(x)+\pi(W\infty)$; the final constant
lies in $\varphi(\mathcal C_N)$ and vanishes after multiplication
by $C_\pi$.

At an inert new prime $\ell$, a compatible place gives
\[
                      \overline y_{d\ell}
                         =F_\ell\overline y_d.           \tag{H4}
\]
Here $F_\ell$ is rational arithmetic Frobenius on the reduced
elliptic curve. This is the usual CM congruence, and it does not
require the old conductor to be an inert squarefree product.
Indeed $(d,\ell)=1$ makes the local CM order the maximal
unramified quadratic order. The CM elliptic curve has supersingular
reduction at this inert prime. Each of its $\ell+1$ cyclic raising
isogenies specializes to the quotient by the unique subgroup of
order $\ell$, namely the relative Frobenius kernel. The prime-to-$\ell$
cyclic $N$-subgroup specializes by the same Frobenius map.
The resulting equality of reduced modular points passes through
$\pi$, which extends over the good integral models at $\ell$.
This also proves (H4) after every old ideal-class translate, with
its place transported. Howard, Section 1.7, formula (13) \cite{howard2012}, records
this congruence in his conductor subfamily; his odd-prime descent
and residual-surjectivity theorem are not used here.

\paragraph{3. The relative group and invariant derivative class.}
The idelic description
$\operatorname{Pic}(\mathcal O_d)=
K^\times\backslash\mathbf A_{K,f}^\times/\widehat{\mathcal O}_d^\times$
identifies the kernel of
$\operatorname{Pic}(\mathcal O_{fn_I})\to
\operatorname{Pic}(\mathcal O_f)$ with
\[
 \prod_{\ell\in I}
 \frac{(\mathcal O_K/\ell\mathcal O_K)^\times}
      {(\mathbf Z/\ell\mathbf Z)^\times}
       \simeq\prod_{\ell\in I}C_{\ell+1}.               \tag{H5}
\]
There is no additional quotient by units: the only global units
are $\pm1$ and already belong to the denominator at every factor.
This proves (H5) without choosing a splitting of the full class
group. Let $\sigma_\ell$ be a generator of its $\ell$-factor,
with the local Artin orientation fixed. Set
\[
 D_\ell=\sum_{j=1}^{\ell}j\sigma_\ell^j,\qquad
 D_I=\prod_{\ell\in I}D_\ell,\qquad
 (\sigma_\ell-1)D_\ell=(\ell+1)-\operatorname{Tr}_\ell. \tag{H6}
\]
If $2^a$ divides $a_\ell$ and $\ell+1$ for all $\ell\in I$,
(4.1) and (H6) show that
\[
 c_{I,a}=\delta_a(C_\pi D_Iy_{fn_I})
 \in H^1(M_I,E[2^a])^{\operatorname{Gal}(M_I/H_f)}.     \tag{H7}
\]
For example its difference under $\sigma_\ell$ is the Kummer
class of
$C_\pi D_{I\setminus\{\ell\}}
 ((\ell+1)y_{fn_I}-a_\ell y_{fn_I/\ell})$, an actual
$2^a$-multiple of an $M_I$-point. The factors in (H5) commute.

\paragraph{4. Descend once, before twisting.}
Put $B=H_f$, $M=M_I$, $\Gamma=\operatorname{Gal}(M/B)$ and
$V_a=E[2^a]$. The five-term sequence for finite discrete coefficients
\cite[formula (4.1.4.1), p.~96]{nekovar} gives
\[
 0\longrightarrow H^1(\Gamma,V_a^{G_M})
 \longrightarrow H^1(B,V_a)
 \xrightarrow{\mathrm{res}} H^1(M,V_a)^\Gamma
 \longrightarrow H^2(\Gamma,V_a^{G_M}).                \tag{H8}
\]
By (H2), multiplication by $2^e$ is zero on the coefficient
module $V_a^{G_M}$ and hence on every cochain and cohomology
group of $\Gamma$ with those coefficients. For an invariant class
$c$, choose $b$ with $\mathrm{res}(b)=2^e c$ and define
\[
                         d_a(c)=2^e b.                  \tag{H9}
\]
Any two choices of $b$ differ by the first group in (H8), so (H9)
is independent of the choice. It is additive, natural for
coefficient reduction, and equivariant for transport of the whole
field diagram. Its restriction is exactly $2^{2e}c$.
This is one application to $M_I/H_f$, not one application for
each factor of (H5). The coefficient module in (H8) is untwisted
$E[2^a]$; the scalar and binary coefficient push comes only later.
Thus no invariant module containing a varying nonquadratic scalar
character has been silently substituted in this argument.
For $I=\varnothing$, the relative group is trivial and this same
definition gives $d_a(c)=2^{2e}c$. In particular it uses the same
normalizing multiplier for the primitive class and its derivatives.

\paragraph{5. Integral local descent at the old places.}
We record the local fact needed to put $d_a(c_{I,a})$ in the old
Kummer conditions. If $F'/F$ is a finite unramified extension of
nonarchimedean local fields and $B/F$ is an elliptic curve, then
\[
 H^1(\operatorname{Gal}(F'/F),B_0(F'))=0,\qquad
 c_B H^1(\operatorname{Gal}(F'/F),B(F'))=0,             \tag{H10}
\]
where $c_B=\#\Phi_B(\overline{k}_F)$ is the geometric component
group order. To prove the first assertion, reduce a cocycle to
the connected smooth special fiber. Lang's theorem \cite[Theorem~1 and its corollary]{lang1956} makes its
$H^1$ over the finite residue field zero; inflation--restriction
then makes its $H^1$ for the finite residue extension zero.
Lift the resulting coboundary through the smooth reduction map.
The remaining cocycle takes values in the formal group. For a
uniformizer of $F$, fixed by the unramified Galois group, each
successive formal quotient is an additive copy of the residue
field of $F'$ \cite[IV.3.2(a), p.~124]{silverman}. A normal basis makes this an induced module, so
its positive-degree cohomology vanishes. Correct the coboundary
successively in these quotients. The corrections converge in
$B_1(F')$, and the group is finite, so their limit solves all
the cocycle equations. The component exact sequence proves the
second assertion of (H10). No division by $[F':F]$ was used.

If a class $v\in H^1(F,B[2^a])$ becomes Kummer over $F'$, its
image in $H^1(F,B)[2^a]$ belongs to the kernel of restriction.
That kernel is the image of
$H^1(\operatorname{Gal}(F'/F),B(F'))$. Thus $c_Bv$ is Kummer
over $F$, by the Kummer exact sequence and (H10).

All ring-class extensions in use are unramified at $2N$, since
the conductors are prime to $2N$. The base fields there are
unramified extensions of $K_v=\mathbf Q_p$. Geometric component
groups are unchanged under unramified extension. Let $C_{\rm bad}$
be a common multiple of their orders for the finitely many
allowed local quadratic types at $2N$. The same integer works
for every conductor, every place above $2N$ and every $K$.
Choose a fixed positive integer $m_0$, allowing the other fixed
multiplier requirements in the argument to be included, and define
\[
 L=m_0C_\pi C_{\rm bad}2^{2e},\qquad
 b_{I,a}=m_0C_{\rm bad}d_a(c_{I,a})\in H^1(H_f,E[2^a]). \tag{H11}
\]
Thus any common fixed multiple of this descent multiplier and
the detection multiplier is permitted; every class is scaled by
the same $m_0$.
Then $b_{I,a}$ is Kummer at $2N$ and restricts to
$\delta_a(LD_Iy_{fn_I})$. At a good place away from $2fn_I$,
the same assertion has component order one. It gives the
unramified local condition, since at good odd places the finite
Kummer image equals the unramified image. At active and moving
primes the required condition is full. The integer $L$ is now
fixed, independently of $K,f,I,m,b$ and $G$.

\paragraph{6. Exact finite precision for one derivative prime.}
Let the retained torsion precision be $m$, choose $a\ge m+e$,
and require, for every retained derivative prime,
\[
 \rho_T(\operatorname{Fr}_\ell)\equiv J\pmod {2^{a+1}},
 \qquad J=\rho_T(\tau),\quad J^2=I.                    \tag{H12}
\]
Then $2^{a+1}$ divides $a_\ell$ and $\ell+1$, since
$\operatorname{tr}J=0$ and $\det J=-1$.
Choose $\ell\notin I$, and perform the construction for
$I\cup\{\ell\}$. Apply the other derivatives first, writing
\[
 X=D_Iy_{fn_I},\qquad Y=D_Iy_{fn_I\ell}.
\]
The $K$-prime $(\ell)$ splits completely in $H_{fn_I}$:
it is prime to the conductor and is a rational principal ideal
in the ring-class Artin map. Thus its residue field is
$\mathbf F_{\ell^2}$, and
\[
 F^2\overline X=\overline X,\qquad
 \operatorname{Tr}Y=a_\ell X,\qquad
 \overline Y=F\overline X,\qquad F^2-a_\ell F+\ell=0.  \tag{H13}
\]
The same splitting makes the old class $b_{I,a}$ at this place
exactly the Kummer class of $LX$, by its restriction identity
in (H11); there is no local descent ambiguity for this old class.
For $2^aQ=LX$ its finite coordinate is
\[
 f_{\rm old}=(F^2-1)\overline Q
             =(a_\ell F-(\ell+1))\overline Q.           \tag{H14}
\]
Although the two integer coefficients are highly divisible, $Q$
is a division point, so (H14) does not assert $f_{\rm old}=0$.

Choose a continuous cocycle representing $b_{I\cup\{\ell\},a}$
and choose $Z$ with $2^a Z=LD_\ell Y$.
Its restriction to $G_{H_{fn_I\ell}}$ is cohomologous to
$g\mapsto(g-1)Z$. Their difference is the coboundary of a single
element of $E[2^a]$; subtract the same global coboundary to make
the restriction exactly this Kummer cocycle. Denote the adjusted
representative by $b$. Restrict it to $G_{H_{fn_I}}$ and choose
a local inertia lift of $\sigma_\ell$. Put
\[
                         U=(\sigma_\ell-1)Z-b(\sigma_\ell).
\]
For $g\in G_{H_{fn_I\ell}}$, compare $b(g\sigma_\ell)$ with
$b(\sigma_\ell(\sigma_\ell^{-1}g\sigma_\ell))$.
Using the exact restricted Kummer representative gives $gU=U$.
Consequently $U$ is an actual point of $E(H_{fn_I\ell})$, and
\[
 2^aU=L((\ell+1)Y-a_\ell X).
\]
The point
\[
 U_0=L\frac{\ell+1}{2^a}Y-L\frac{a_\ell}{2^a}X        \tag{H15}
\]
also belongs to this field, with the same $2^a$-multiple. Its
coefficients are integers. Hence $U-U_0$ is upper-field
$2$-primary torsion and is killed by $2^e$, by (H2).
This comparison with an integer quotient point is necessary;
comparison with an arbitrary division point would give no bound.

Inertia acts trivially on the reduction of every algebraic point
at a good prime. Reduction is injective on $2$-primary torsion.
Thus the singular coordinate $b(\sigma_\ell)$ is
\[
 s_{\rm new}=(a_\ell-(\ell+1)F)\overline Q+\epsilon
             =(F-a_\ell)f_{\rm old}+\epsilon,
                    \qquad 2^e\epsilon=0.              \tag{H16}
\]
The second equality uses the exact endomorphism identity
\[
 (F-a_\ell)(a_\ell F-(\ell+1))
                         =a_\ell-(\ell+1)F.            \tag{H17}
\]
Reduction $T/2^aT\to T/2^mT$ is multiplication by
$2^{a-m}$ on geometric torsion points. It kills $\epsilon$.
Conditions (H12) then turn (H16) into $s_{\rm new}=Jf_{\rm old}$.

For the new finite coordinate, all inertia translates of $Y$
have the same reduction, so
\[
 \overline{LD_\ell Y}
       =L\frac{\ell(\ell+1)}2F\overline X.              \tag{H18}
\]
This is $2^a$ times an $F^2$-fixed reduced point, because
$2^{a+1}\mid\ell+1$. The local extension obtained by adjoining
the conductor $\ell$ is totally ramified, of degree $\ell+1$.
Its residue field is still $\mathbf F_{\ell^2}$. Lift the indicated
$2^a$-division point in the residue field; the discrepancy lies
in the formal group, where multiplication by $2^a$ has an
integral formal inverse because $\ell$ is odd
\cite[IV.2.3 and IV.3.2, pp.~122--124]{silverman}. This produces a
local $2^a$-division point of $LD_\ell Y$ in the upper local
field itself. The restricted Kummer class is therefore zero.
Choose a Frobenius lift in that upper local Galois group.
It acts trivially on $E[2^a]$ by (H12), so its cocycle value is
zero, independently of a coboundary adjustment. Hence
\[
                     f_{\rm new}=0,\qquad
                     s_{\rm new}=Jf_{\rm old}
                  \quad\hbox{at precision }2^m.         \tag{H19}
\]
The generator $\sigma_\ell$ fixes the tame orientation. A different
generator changes the corresponding singular coordinate by the
matching unit. Frobenius lifts and places in (H19) are part of
the retained local diagram. The extra bit in (H12) pays exactly
for the division by two in (H18). Repeat this calculation with
each prime of $I\cup\{\ell\}$ singled out: all other lower
conductors split at that prime by the same principal-ideal argument.
Thus the previously imposed singular conditions are retained.
The maximum of the required finite precisions suffices for the
whole finite set; neither (H11) nor its exponent is repeated.

\paragraph{7. The actual coefficient push and its local conditions.}
Reduce $b_{I,a}$ to $E[2^m]$, extend coefficients freely to $A$,
and only now use the character $\theta$ of $\Delta_f$.
Under Shapiro, followed by the equivariant coefficient map,
\[
 E[2^m]\otimes A[\Delta_f]\longrightarrow
 (E[2^m]\otimes A)(\theta),\qquad
                     v\otimes[\delta]\longmapsto\theta(\delta)v,
                                                               \tag{H20}
\]
obtain the class $P_{I,m}$ over $K$.
Here the induced module is written with the diagonal Tate action
and the left regular action on its class-group basis. Formula
(H20) is equivariant since left translation multiplies the value
by $\theta(g)$. It remains equivariant with values in $A[G]$,
before evaluation at any sign character.

For an explicit degree-one cochain formula, choose representatives
$r_\delta$ of $G_K/G_{H_f}$ using the normal quotient
$\Delta_f$, and write
$g r_\delta=r_{g\delta}h(g,\delta)$ with $h(g,\delta)\in G_{H_f}$.
The transferred cocycle is
\[
 (\mathcal T_\theta b)(g)=
 \sum_{\delta\in\Delta_f}
  \theta(r_{g\delta})\rho_T(r_{g\delta})b(h(g,\delta)). \tag{H21}
\]
The identity
$h(gk,\delta)=h(g,k\delta)h(k,\delta)$ proves the cocycle
identity with target action $\theta\rho_T$.
In degree zero the formula is the weighted sum
$\sum_\delta\theta(r_\delta)\rho_T(r_\delta)Q$; applying its
boundary gives (H21) applied to the boundary of $Q$.
More generally, put $R(g)=\theta(g)\rho_T(g)$ and, for
$g_1,\ldots,g_n$, set $\delta_i=g_{i+1}\cdots g_n\delta$.
The degree-$n$ map is
\[
 (\mathcal T_\theta c)(g_1,\ldots,g_n)
 =\sum_{\delta\in\Delta_f}R(r_{g_1\cdots g_n\delta})
       c\bigl(h(g_1,\delta_1),\ldots,h(g_n,\delta_n)\bigr).
\]
With the ordinary inhomogeneous differential, the relation for $h$
above gives $d\mathcal T_\theta=\mathcal T_\theta d$.
These are the integral Shapiro and transfer maps of
\cite[Sections 8.1.1--8.1.6, pp.~189--194]{nekovar}, written
with left cosets. Changing sections gives the integral homotopy
in Section 8.1.4; compatible sections give strictly transitive
transfer. No averaging factor occurs.

Localizing (H21) separates the cosets into decomposition-group
orbits: this is the local Mackey formula, a sum of the local
corestrictions from the transported places.
At $2N$, the deformation character $\Psi$ is trivial locally;
the fixed quadratic character need not be trivial. Under the
fixed quadratic-twist identification, the local map is the
ordinary integral norm on Kummer classes. One can see this
without a projector: the sum $\sum_\delta\chi(\delta)\delta Q$
is a point of the quadratic twist over the lower field, and
Kummer of this sum equals the cohomological transfer.
The extra $A$ coefficients take finite linear combinations of
such maps. Thus the old Kummer conditions proved in paragraph 5
survive (H20), also in the universal family with its common local
quadratic type. Local Kummer corestriction is natural even for
a ramified field extension; in the present old-place descent
the extensions used in (H10) are unramified.

At a derivative prime, $H_f/K$ splits completely and $\theta$
is trivial on the local Galois group. Each Mackey term uses a
transported place and a transported Frobenius. Transporting the
Tate basis transports $F$ as well, so (H19) is equivariant,
not an assertion that an arbitrary change-of-place matrix
commutes with a fixed untransported $J$. The same weights occur
on the upper and lower class. Summing gives, in the common target coordinates,
\[
 f(P_{I\cup\{\ell\},m})=0,\qquad
 s(P_{I\cup\{\ell\},m})=Jf(P_{I,m}).\tag{4.2}
\] Active and moving conditions are full; all other
good conditions are unramified and preserved by the same maps.

The split norm formula in (4.1) also survives this construction.
For a new conductor prime the class-field squares obtained by
retaining the derivative factors are Cartesian, by (H5) and the
same local-unit calculation at that prime. Restriction and
corestriction commute in these squares. If $b$ lifts $2^e c$
in (H9), its corestriction lifts $2^e$ times the corresponding
corestricted class; uniqueness after the second multiplication
by $2^e$ proves the required compatibility of $d_a$.
The weighted sums are transitive finite sums. Thus the two
horizontal Artin terms in (4.1) become $\theta(\sigma_q)$ and
$\theta(\sigma_q)^{-1}$, giving
$a_q-\theta(\sigma_q)-\theta(\sigma_q)^{-1}$ on the class.
For $I=\varnothing$, the class is precisely the Kummer image of
the $L$-multiplied weighted point, so the same statement gives
the central split-norm relation with that exact $L$.

\paragraph{The local diagrams underlying the coefficient push.}
At a retained derivative place put $q=\ell^2$ and choose tame
and unramified generators $\tau_v,\phi_v$ with
$\phi_v\tau_v\phi_v^{-1}=\tau_v^q$. The local coefficient
character is trivial. The two-generator tame resolution is
\[
 [M\xrightarrow{u\mapsto((\Phi_v-1)u,0)}M\oplus M
       \xrightarrow{(x,y)\mapsto(\Phi_v-q)y}M],
 \qquad \Phi_v=\rho(\phi_v).
\]
At the retained exponent $m\le a$, (H12) makes this the
zero-differential model $D^0\oplus(F_v\oplus S_v)[-1]
\oplus D^2[-2]$. Moreover $q\equiv1\pmod{2^{a+2}}$.
Represent the singular generator by a uniformizer adjusted by
a unit so its Kummer value on the chosen $\phi_v$ is zero.
Residue-unit Kummer characters allow this adjustment. The latter
congruence makes its self-cup zero: $-1$ is a $2^m$-th power,
so the adjusted uniformizer Kummer class has square zero. The finite self-cup is zero by unramified cohomological
dimension one. Retain the transported mixed pairing and the
$D^0,D^2$ pairing.

In the derivative diagram the old ring-class fields split
locally, and the new conductor extension is totally ramified.
For the latter choose the lower Frobenius in the upper subgroup
and take $\tau_{\rm upper}=\tau_{\rm lower}^e$; here $f=1$.
For split terms $e=f=1$. Transport these choices with the Tate
basis. With these compatible generators the local matrices, in the order
$D^0,F_v,S_v,D^2$, are
\[
 \operatorname{res}=\operatorname{diag}(1,f,e,ef),\qquad
 \operatorname{cor}=\operatorname{diag}(ef,e,f,1).
\]
The possible transfer shear is divisible by $(q-1)/2$ and
vanishes at retained precision. These matrices preserve each
of the four conditions retaining $D^0$, and satisfy the local
cup-product adjunctions. Split Mackey terms have $e=f=1$;
coefficient weights and reduction commute with these matrices.
An arbitrary new Frobenius lift is not used without transporting
the resulting finite/singular decomposition.

Compare the small free resolution with the continuous bar
resolution and choose augmentation-preserving comparison maps
and their composite homotopies. Continuous Hom into $M$ retains
these identities. Since the small model has zero differential,
the composite on it is the identity; normalize the other homotopy
to obtain a contraction. Here are its compatibility maps.
For this retained contraction $1-i_vp_v=dh_v+h_vd$ and a full
local cochain arrow $F$, take $F_L=p_wFi_v$ and restrict it to
$U_v$. The inclusion square has homotopy $h_wFi_v|_{U_v}$.
For composable arrows the transferred composition homotopy is
$p_zGh_wFi_v$, together with the transferred original bar
homotopy when present. Every such degree-minus-one map on
$U_v$ lands in $D^0$, which is retained in every condition;
thus these homotopies restrict too. For three composable maps,
write $K(G,F)=pGhFi$ and $L(H,G,F)=pHhGhFi$, with the
intermediate subscripts understood. Direct expansion gives
\[
 K(H,GF)-K(HG,F)+H_LK(G,F)-K(H,G)F_L=Ld-dL.
\]
The degree-minus-two map $L$ vanishes on every derivative
condition. Transfer the original bar comparison homotopies and
cubes as well; unrelated classwise boundaries are not substituted
for them. For the isotropy diagrams, $U_v\otimes U'_v$ is
free in degrees $0,1,2$, whereas the trace target is $A[-2]$.
Its negative Hom groups vanish, so the restricted zero pairing
has compatible null-homotopies and the requisite 2-homotopies.
This constructs the local data required in
\cite[8.6.2.1--8.6.2.4]{nekovar} before using 8.6.4.
At old Kummer places use the natural restriction/norm diagrams
of the multiplication-by-$2^m$ exact sequences, including
local torsion, and their derived scalar extensions; at good
unramified, full and zero places use the corresponding standard
subcomplexes. The witnesses in (H22) are chosen in these same
diagrams.

\paragraph{8. Actual finite Selmer-cone witnesses.}
Represent $P_{I,m}$ by a continuous global cocycle $p$.
At each constrained place, membership in the indicated local
condition gives a local cycle $u_v$ and a degree-zero cochain
$q_v$ with
\[
                   \operatorname{loc}_v p-\iota_vu_v=dq_v.
                                                               \tag{H22}
\]
These entries give a cycle in
$\operatorname{Cone}(C^\bullet(G_{K,S},M)\oplus\bigoplus U_v
\to\bigoplus C^\bullet(K_v,M))[-1]$.
For (H19), choose in the same way a degree-zero cochain whose
boundary is the difference of the two representatives being
compared. At the derivative primes use the already specified
$D^0$-retaining finite/singular local complexes, not full/zero
conditions. Their coordinate equality is the finite local
cohomology equality just proved. All coefficient groups and
coset sets at this stage are finite; the cocycles and the finite
list of witnesses factor through finite quotients.

Apply the coefficient push to these witnesses at the same time
as to the cocycle, including the natural Kummer comparison
homotopies. If $(i,p,h)$ is a retained contraction to a finite
target model, a relation $z-z'=dq$ becomes
$pz-pz'=d(pq)$. For a comparison $F$ to another model with
projection $p'$, the correction between $p'Fz$ and $p'Fi pz$
is the boundary of $p'Fhz$ (with the consistent contraction sign).
Thus the finite arithmetic provides actual cycles, boundaries
and comparison entries before compression. Adding a derivative
prime extends this same finite diagram. Existence and size
bounds for those target contractions, and the subsequent two
limits, are separate inputs; this proof makes no inference
from abstract isomorphism of cohomology groups to those bounds.

The resulting finite arithmetic uses the single multiplier (H11)
and the precision rule $a\ge m+e$, (H12). No factor depending on
$[H_f:K]$, $|G|$, $|I|$ or the number of active primes occurs.
\end{proof}

\subsection{The derived Kummer condition}
\begin{lemma}[The actual derived Kummer condition]\label{ex:arith-local-kummer}
For the actual local twist $B=E^{(h)}$ over $F=K_v$, the condition
$U_F=B(F)^\wedge_2[-1]\to R\Gamma(F,T_2B)$ includes the finite
two-primary torsion. Its reduction modulo $2^a$ has
$B(F)[2^a]$ in degree zero and $B(F)/2^aB(F)$ in degree one.
Its Kummer image is its exact local annihilator, and the comparison
has integral isotropy homotopies. At $v\mid2N$, triviality of the
moving character and the common old quadratic type identify the
scalar and binary versions by derived coefficient extension.
These are the local maps used in the Selmer-cone duality above.
\end{lemma}
\begin{proof}\leavevmode
For the actual local twist $B=E^{(h)}$ over $F=K_v$, put
\[
 U_F=B(F)^\wedge_2[-1]\longrightarrow R\Gamma(F,T_2B).
\]
The completion includes finite two-primary torsion. Derived reduction
modulo $2^a$ has $B(F)[2^a]$ in degree zero and
$B(F)/2^aB(F)$ in degree one. At $v\mid2N$, $\Psi_i$ is trivial,
so derived coefficient extension defines the scalar $A_i$ condition
for this actual twist. In the universal binary family every $h_x$
has the same local squareclasses at $2N$ as $h_0=1$. Hence every
sign character evaluates the universal character $\varepsilon$
trivially on $G_{K_v}$, and therefore $\varepsilon|_{G_{K_v}}=1$
in $G$, before reduction modulo two. The $A_i[G]$ Kummer condition
and its null-homotopy are consequently derived scalar extensions
of the fixed local data. Omitting the scalar twist, the universal
local-squareclass condition, or degree-zero torsion changes this argument.

The local finite dualities \cite[I.2.3, I.3.2--I.3.5]{milne} make the
Kummer map adjoint to $H^1(F,B[2^a])\to H^1(F,B)[2^a]$.
The latter has Kummer image as kernel, so the Kummer image is its exact
annihilator. To obtain the integral chain statement, write
$M=B(F)^\wedge_2$. The restricted cup map
$M[-1]\otimes^LM[-1]\to\Z_2[-2]$ has homotopy class determined by
$\operatorname{Hom}(M\otimes M,\Z_2)$. Finite isotropy makes every
value divisible by every power of two, hence zero.
A projective resolution therefore supplies a null-homotopy.
The possible ambiguity group in degree minus one vanishes.
Reduction modulo two, including the torsion in degree zero, gives the
perfect quotient-to-dual comparison; derived Nakayama lifts it.
The global Selmer duality theorem
\cite[Sections 5.1, 5.2.4, 5.4.3 and 6.3.4]{nekovar} is applied
to these actual local maps and homotopies.

For clarity the Selmer complex is the cone
\[
 C_U=\operatorname{Cone}\left(
 R\Gamma(G_{K,S},M_i)\oplus\bigoplus_{v\in S}U_v
 \longrightarrow\bigoplus_{v\in S}R\Gamma(K_v,M_i)\right)[-1].
                                                                    \tag{4.3}
\]
The compatible trace in (4.3) is coefficient-valued. A top-coefficient
functional used to test a finite Frobenius ring's perfection is not
used as a coefficient-reduction-compatible arithmetic trace.
\end{proof}

\subsection{Actual support comparisons}
At a tame split prime the local free model has
\[
 d^0(v)=(Av,Cv),\quad d^1(v,w)=-Cv+Dw,\quad
 A=\Phi-1,\ C=g-1,\ D=\Phi-N_q(g),
\]
\[
 N_q(g)=\frac{q+1}{2}+\frac{q-1}{2}g,\qquad DC=CA.                 \tag{4.4}
\]
At augmentation $g=1$, this is the direct sum of the unramified
complex $[\Phi-1]$ in degrees zero and one and the singular quotient
$[\Phi-q]$ in degrees one and two.

The support-change theorem \cite[Proposition 7.8.8(ii)]{nekovar}
retains arbitrary old local conditions and puts the unramified
condition at new places. Its finite-module proof applies at two.
Use it only when the actual coefficient representation is unramified
at the removed place. There are five uses here:
universal augmentation at $Q$; an unused scalar prime $q\nmid h$;
characteristic-two residual removal after $\chi_h=1$; the moving
prime after $t=0$; and a new good inert derivative prime outside
$h,r_i$ and the old support. In the last case the coefficient module
is unramified and the chosen finite local plane is the unramified
condition. It is not applied to an active ramified scalar prime
or to $r_i$ before central specialization. Replacing the unramified
condition by the full condition then gives the singular quotient
by a cone of Selmer cones. In particular the old complex is an actual
subcomplex of an upper triangular model, not just a complex with the
same Euler characteristic.

\begin{proposition}[One moving character for the complete active list]\label{ex:arith-moving-character}
Fix an allowed $K$ and any finite list of oriented active primes.
For every $n,B\geq1$ and finite avoidance set there are infinitely
many split rational primes $r\equiv1\pmod{2^B}$ and characters
$\xi:\operatorname{Gal}(H_r/K)\to\Z/2^n\Z$ that are surjective
on inertia at $r$, trivial on the local decomposition groups at
$2N$, and take value one at every prescribed active orientation.
They can be chosen with
\[
 v_2(2\pm a_r)\leq c_E,\qquad
 v_2\det(1\pm\operatorname{Fr}_r)\leq c_E,
\]
where $c_E$ depends only on the fixed Tate image, not on $K,n,B$
or the active list. One such character serves the entire binary
family before evaluation.
\end{proposition}
\begin{proof}\leavevmode
\paragraph{The simultaneous moving character.}
We give the construction for the non-CM curve used here. For every
$n,B\ge1$ and every finite avoidance set, there are infinitely many
split rational primes $r\equiv1\pmod{2^B}$ and characters
\[
 \xi:\operatorname{Gal}(H_r/K)\longrightarrow\Z/2^n\Z
\]
which are surjective on inertia at $r$, trivial at $2N$, and equal
to one at each prescribed oriented active prime. The trace bounds
below have one constant $c_E$, independent of $n,B,K$ and the active
list. The degrees of the finite fields and the sizes of the primes
chosen in the construction need not be uniformly bounded.

Choose prime ideals $\mathfrak p_1,\ldots,\mathfrak p_s$ in distinct
conjugate pairs, including the active orientations, one prime over
each prime dividing $2N$, and enough additional split prime ideals
to generate $\operatorname{Cl}(K)$. The latter exist outside any
finite set by Chebotarev in the Hilbert class field. Set $w_j=1$ at
the active primes and $w_j=0$ at the other chosen primes. Put
\[
 \mathcal R=\ker(\Z^s\longrightarrow\operatorname{Cl}(K)),
 \qquad (\beta_a)=\prod_j\mathfrak p_j^{m_j^{(a)}},
 \qquad \alpha_a=\beta_a/\overline\beta_a,
\]
where the $m^{(a)}$ form a $\Z$-basis of $\mathcal R$.
The generators $\beta_a$ may be fractional. Disjointness of the
conjugate pairs gives $v_{\mathfrak p_j}(\alpha_a)=m_j^{(a)}$.

Let $F_\infty=K(\mu_{2^\infty})$ and adjoin the $2^n$-th roots of
all the $\alpha_a$. Prescribe the multiplier
\[
 \zeta_{2^n}^{-\sum_jm_j^{(a)}w_j}
 \quad\hbox{on }\alpha_a^{1/2^n}.                              \tag{4.4a}
\]
These prescriptions respect every Kummer relation. Indeed, if
$\prod_a\alpha_a^{b_a}$ is a $2^n$-th power in $F_\infty$, its
root lies in a finite cyclotomic layer. At every nonzero-weight
prime that layer is unramified over $K$, since the prime is odd.
Consequently $2^n$ divides $\sum_ab_am_j^{(a)}$ at each such prime,
and hence divides $\sum_{a,j}b_am_j^{(a)}w_j$. The dyadic primes
have weight zero. This tests all multiplicative relations, including
those involving units; it does not assume independence of the
radicands. Kummer theory therefore gives the specified automorphism
of the finite abelian extension
\[
 L_\infty=F_\infty(\alpha_a^{1/2^n}:a)/F_\infty.
\]

Here is why these are precisely the extension conditions for $\xi$.
For a split $r$ avoiding the chosen primes, the ring-class sequence is
\[
 1\longrightarrow\mathbf F_r^\times
 \longrightarrow\operatorname{Pic}(\Z+r\mathcal O_K)
 \longrightarrow\operatorname{Cl}(K)\longrightarrow1.
\]
The kernel is the ratio of the two residue factors in
$(\mathcal O_K/r\mathcal O_K)^\times$ modulo diagonal rational
residues. There is no additional unit quotient because
$\mathcal O_K^\times=\{\pm1\}$. Write $q_j$ for the ideal-class
lift of $\mathfrak p_j$, and orient the connecting map $\delta$ so
that its principal relations are
\[
 \sum_jm_j^{(a)}q_j+\delta(\operatorname{res}\alpha_a)=0.
\]
The kernel and the $q_j$ generate the ring-class group. Any relation
between them projects to $\mathcal R$; subtracting its basis
relations leaves a relation within the kernel. Thus these are a
complete presentation, not merely necessary relations.
For $r\equiv1\pmod{2^n}$ the power-residue character
$\kappa_r:\mathbf F_r^\times\to\Z/2^n\Z$ is surjective.
The assignments $\xi(q_j)=w_j$ and $\xi\delta=\kappa_r$ extend
exactly when
\[
 \kappa_r(\operatorname{res}\alpha_a)
       =-\sum_jm_j^{(a)}w_j.                                  \tag{4.4b}
\]
Below we choose $r$ so that its radical Frobenius is (4.4a), which
gives (4.4b). Reversing the connecting-map convention reverses the
power-residue convention as well. No class number is divided out.

We next preserve all these conditions while controlling the Tate
trace. Let $J=\rho_T(G_{\Q(\mu_{2^\infty})})$. The determinant is
the cyclotomic character, so $J$ is the determinant-one part of the
fixed Tate image. Non-CM open image \cite{serre} gives a fixed
$d\ge2$ with
\[
 \Gamma_d:=\ker\bigl(\operatorname{SL}_2(\Z_2)
             \to\operatorname{SL}_2(\Z/2^d\Z)\bigr)\subset J.
\]
The subgroup $J_K=\rho_T(G_{F_\infty})$ has index at most two in
$J$, and therefore contains its squares. Write $U(a)$ and $V(b)$
for the upper and lower unipotent matrices. Then $J_K$ contains
$U(2^{d+1}\Z_2)$ and $V(2^{d+1}\Z_2)$. The identity
\[
 U(a)V(b)U\bigl(-a/(1+ab)\bigr)V\bigl(-b(1+ab)\bigr)
       =\operatorname{diag}(1+ab,(1+ab)^{-1})
\]
and Gaussian decomposition show that a subgroup containing
$U(2^k\Z_2),V(2^k\Z_2)$ contains $\Gamma_{2k}$.
Hence $\Gamma_m\subset J_K$ for $m=2d+2$.
Commuting $\operatorname{diag}(1+2^m,(1+2^m)^{-1})$ with the
upper and lower unipotents of level $m$ gives all upper and lower
unipotents of level $2m+1$, since
$v_2((1+2^m)^2-1)=m+1$. Applying the same decomposition gives
\[
 \Gamma_{4m+2}\subset\overline{[J_K,J_K]}.
\]
Every matrix commutator lifts to a commutator in $G_{F_\infty}$.
Such lifts fix $L_\infty$, whose Galois group over $F_\infty$ is
abelian. The image of the compact group $G_{L_\infty}$ is closed.
Thus, for the fixed integer $\nu=4m+2=8d+10$,
\[
                  \Gamma_\nu\subset\rho_T(G_{L_\infty}).       \tag{4.4c}
\]
This is a subgroup in the image of the actual radical-fixing
kernel. It requires no disjointness of the radical and Tate fields,
and is independent of the number of radicands and of $K$.

For the trace estimate one can take $c_E=2\nu+1=16d+21$.
We verify the elementary bound used here. Put $\delta=2^\nu$ and
$D(v)=\operatorname{diag}(v,v^{-1})$. For any
$A=\left(\begin{smallmatrix}a&b\\c&f\end{smallmatrix}\right)
\in\operatorname{SL}_2(\Z_2)$, one of
\[
       A,\quad AU(\delta),\quad AD(1+\delta),\quad AD(1+2\delta)
                                                                    \tag{4.4d}
\]
has both trace distances to $2$ and $-2$ of valuation at most $2\nu+1$.
If $A$ itself does not, its trace is within $2^{2\nu+2}\Z_2$ of
one of these two values, while its distance to the other has valuation
two. Its diagonal entries have the same parity. If both are even,
$c$ is odd by $af-bc=1$, and $AU(\delta)$ changes the trace by
$c\delta$, of valuation $\nu$. If both are odd, the two diagonal
multipliers change the trace by
\[
 \delta\frac{(a-f)+a\delta}{1+\delta},\qquad
 2\delta\frac{(a-f)+2a\delta}{1+2\delta}.
\]
If $v_2(a-f)<\nu$, the first change has valuation at most $2\nu-1$.
Otherwise write $a-f=\delta z$. For even $z$ the first change has
valuation $2\nu$; for odd $z$ the second has valuation $2\nu+1$.
All the chosen changes have valuation at least three. They therefore
leave the distance to the other forbidden value of valuation two,
and give the claimed bound for the nearer one.

Lift the prescribed Kummer automorphism to
$\sigma_0\in G_{F_\infty}$. By (4.4c), every right multiplier in
(4.4d) lifts to an element of $G_{L_\infty}$. Choose one giving the
trace bound and multiply $\sigma_0$ by that lift. The resulting
actual Galois element $\sigma$ fixes $F_\infty$, preserves every
prescribed radical multiplier, and satisfies
$v_2(2\pm\operatorname{tr}\rho_T(\sigma))\le c_E$.

Choose $B'\ge\max(B,n)$ with $B'>c_E$, and take a finite normal
extension of $\Q$ containing $K,\mu_{2^{B'}},E[2^{B'}]$ and all
the chosen radical roots. Restrict $\sigma$ to this extension.
Its restrictions on all intersections agree because it is already
an actual Galois element. Finite Chebotarev supplies infinitely many
rational primes in its Frobenius class. Choose a place realizing
this representative and avoid all the finitely many excluded primes.
Then $r$ splits in $K$ and $r\equiv1\pmod{2^{B'}}$. The Frobenius
formula on $\alpha_a^{1/2^n}$ gives (4.4b), and hence the character
$\xi$. Each local unit factor at $r$ maps onto the residue-ratio
kernel, so local reciprocity proves that $\xi$ is surjective on
inertia there. At $2N$ it is unramified with Frobenius value zero,
and is therefore trivial on the whole local decomposition group.
Conjugation negates its values because a prime ideal times its
conjugate is rational principal.

Finally $a_r\equiv\operatorname{tr}\rho_T(\sigma)\pmod{2^{B'}}$
and $\det(\operatorname{Fr}_r)=r$. Since $B'>c_E$, we obtain
\[
 v_2(2\pm a_r)\le c_E,\qquad
 v_2\det(1\pm\operatorname{Fr}_r)
       =v_2(1\pm a_r+r)\le c_E.
\]
The same bound controls $a_r-2\chi_h(r)$ for every constituent
unramified at $r$. Repeat with increasing $n,B$ and avoidance sets
to obtain distinct $r_i$. Taking $n_i\ge2i$ makes
$(1+t)^{2^{n_i}}=1$ in $A_i$: for $1\le j<i$,
$v_2\binom{2^{n_i}}j\ge n_i-v_2(j)\ge i$.
Thus $\Psi_i=(1+t)^{\xi_i}$ is the required finite character.
Its inertia generator can be oriented to have value $1+t$, and its
active oriented Frobenius values have odd exponent. One character
works for the entire fixed active list before any binary evaluation.
\end{proof}

\subsection{The same functional at every localization}
\begin{proposition}[The paired functional and its integral representatives]\label{ex:arith-paired-functional}
Use the retained inner-limit diagrams of
Lemma~\ref{ex:limits-inner-limit}, with the local conditions of
Definition~\ref{ex:arith-data}, and let $z$ be the actual retained
class. Transport $z$ by conjugation, lift it through the orthogonal
comparison at the generic point, and pair by the retained duality.
This defines one generic closed degree-two functional $\lambda_z$.
It has an integral representative after multiplication by $2^a$
over $\mathscr O[G]$, for a fixed old-support exponent $a$
independent of the active list and $G$.
For a scalar constituent with its primitive local conditions, its
paired functional has an integral representative at every height-one
prime $\mathfrak p\ne(2)$ of $\Lambda$, including $(t)$ and $(t+2)$.
These representatives have the same generic value for that primitive
constituent. The universal relaxed functional is asserted integral
after the fixed power of two over $\mathscr O[G]$; scalar regularity
of its corrected coordinate is obtained instead from the actual
primitive comparison of Lemma~\ref{ex:arith-common-coordinate}.
All these assertions concern the inner limit and the indicated
localizations, not inverses of $t$ in a finite Artin ring.
\end{proposition}
\begin{proof}\leavevmode
After the inner limit, put
$\mathscr O=\widehat{\Z_2[[t]]_{(2)}}$.
In (4.4), the residual active Frobenius is $(1+t)^aI$ with $a$ odd.
Thus $A,D$ are invertible over $\mathscr O[G]$.
The formulas
\[
 h^1(v,w)=A^{-1}v,\qquad h^2(u)=(0,D^{-1}u)
\]
contract the complex; the identities follow from $DC=CA$.
\paragraph{Which functional is paired with the class.}
Let $z$ denote the retained class and transport it by $\tau$ and the
preceding Weil identification to the conjugate coefficient system.
At a localization where the full local complexes at active and
moving primes are acyclic, the comparison from the orthogonal
condition to the transported condition is a quasi-isomorphism.
Apply its inverse in the derived category to the transported class,
and call the resulting class $z^\perp$. Define
\[
 \lambda_z(w)=\operatorname{PD}(w,z^\perp),\qquad
           w\in H^2(C_U).
\]
Thus $\lambda_z$ is a closed degree-two functional, with a cochain
representative obtained using the displayed local contractions and
the retained global duality map. The inverse is unique in the derived
category, so representatives built at distinct localizations define
the same generic functional. Integrality still requires the stated
integral contractions and fixed old-support correction; it does not
follow from abstract duality alone. The coordinate definition below uses this same $\lambda_z$.

Here the representatives must be constructed separately. Over
$\mathscr O[G]$, the active Frobenius operators $A,D$ are units,
and moving-prime inertia minus one is $t$ times a unit; the displayed
homotopies are integral. For the height-one assertion take the scalar
primitive constituent. Its active primes are precisely the primes
ramified in that constituent, so inertia minus one is $-2$, a unit
at $\mathfrak p\ne(2)$. At the moving prime inertia minus one is
invertible unless $\mathfrak p=(t)$; at $(t)$ the nonzero central
Frobenius determinants give integral adjugate contractions. In
particular $(t+2)$ is treated by the inertia contraction. Unused
good primes in the primitive constituent retain the unramified
condition and its orthogonal comparison; no inverse of a full local
Frobenius complex at such a prime is required. This proves the
scalar assertion only for the primitive functional. In the universal
relaxed family an unused prime can specialize to trivial inertia,
so the $-2$ contraction is not available there. Its correction is
the determinant identity of Lemma~\ref{ex:arith-common-coordinate}.
The common generic value alone would not prove any of these
integrality statements.

In the inertia contraction the maps are
\[
 h^1(v,w)=C^{-1}w,\qquad h^2(u)=(-C^{-1}u,0).
\]
\paragraph{The fixed old-support lattice and the exponent.}
Here is the integral content of the old-support comparison.  For each
$p\mid2N$ and each of the fixed allowed local quadratic types $\delta$,
choose once a bounded free $\mathbf Z_2$-model of the diagram
\[
 U_{p,\delta}\longrightarrow L_{p,\delta},\qquad
 U_{p,\delta}=(E^{(\delta)}(\mathbf Q_p)^\wedge_2)[-1],\qquad
 L_{p,\delta}=R\Gamma(\mathbf Q_p,T_2E^{(\delta)}),
\]
together with its Kummer map, cup pairing, isotropy homotopy, and
quotient-to-dual comparison.  Finite torsion is retained in $U_{p,\delta}$;
it is represented by a free resolution, not removed.  The local
duality and Kummer exact sequence identify that comparison integrally:
\[
 \operatorname{Cone}(U_{p,\delta}\to L_{p,\delta})
       \simeq R\operatorname{Hom}_{\mathbf Z_2}
                         (U_{p,\delta},\mathbf Z_2)[-2].
\]
This is the chain-level comparison of
Lemma~\ref{ex:arith-local-kummer}, with its specified isotropy
homotopy.  In particular it is not a rational isomorphism of free
parts of cohomology.  A quasi-isomorphism between these bounded free
models has an integral homotopy inverse: its bounded acyclic free
cone is contractible, by splitting from its last nonzero term
and proceeding downwards.  Consequently neither torsion lengths
nor inverse determinants of torsion differentials are denominators
in this comparison.

Every allowed $K$ splits at $2N$.  Its two completions above $p$ are
therefore $\mathbf Q_p$, with precisely two copies of this same
diagram, interchanged by conjugation.  The moving character and
the universal varying quadratic character are trivial on their
decomposition groups.  All old-place diagrams over $A_i[G]$ are thus
derived scalar extensions of this finite fixed list.  Tensoring the
chosen free models and their homotopy identities preserves all these
identities.  The remaining good unramified conditions use the
integral unramified comparison, and the derivative conditions use
their integral orthogonal planes.  Neither requires a new lattice
comparison when a prime is added.

For completeness, there is also no hidden group-order denominator in
duality.  The coefficient pairing is the $A_i[G]$-valued Weil pairing.
If one checks its perfection over $A_i$, the auxiliary pairing on the
group algebra is
\[
 (a,b)\longmapsto[1](ab),\qquad
 ([1](gh))_{g,h\in G}=(\mathbf1_{h=g^{-1}})_{g,h\in G}.
\]
Its matrix is a permutation matrix, with integral inverse.  It is not
the normalized trace $|G|^{-1}\operatorname{Tr}$.  The global
cup-product and trace maps are integral as well.  The hypotheses of
Selmer duality here are a perfect coefficient pairing and the actual
orthogonal local diagrams with their isotropy homotopies
\cite[Sections 6.2.5--6.2.8 and Theorem 6.3.4]{nekovar}; the old
Kummer diagrams just described supply those hypotheses.  Since $K$
is imaginary there is no real-place Tate correction in this argument.

Over $S=\mathscr O[G]$ the only remaining differences between the
transported conditions and their orthogonal conditions occur at the
active and moving places.  Their full local complexes have the
integral contractions displayed above.  Their direct sum, of any
finite size, is still an integral contraction.  The resulting
comparison $j:C_{U^\perp}\to\tau^*C_U$ is an integral homotopy
equivalence over $S$.  Write $j^{-1}$ for such a homotopy inverse and
$c$ for conjugate transport.  On the retained bounded free models,
\[
 z^\perp=j^{-1}c(z),\qquad
 \lambda_z=\operatorname{PD}(-,j^{-1}c(z))
\]
are integral cochains, and the second is closed.  Every map in this
formula is integral; no inverse on a torsion cohomology quotient is
taken.  The finite-model contractions and their retained limits
preserve these identities.
In detail, represent the comparison on the bounded free models and
contract its acyclic cone over $S$.  The cone contraction yields a
chain map $j^{-1}$ and homotopies $jj^{-1}\simeq1$ and
$j^{-1}j\simeq1$ over $S$.  Thus
$d(j^{-1}c(z))=j^{-1}c(dz)=0$ as a cochain identity, and its image
differs from $c(z)$ by an integral boundary.  Applying the integral
duality chain map gives a closed element of
$\operatorname{Hom}_S(C_U^2,S)$, with the claimed generic value.
This proves the existence of an actual integral closed lift, rather
than only a lift after rationalization.

With the conventions of this paper the actual class $z$ already
contains the fixed multiplier $L$.  Thus the preceding construction
permits $a=0$.  One may keep any larger fixed nonnegative $a$ in the
subsequent inequalities.  This does not change $z$, divide out $L$,
or change the determinant basis.  It proves the existence and
independence of the exponent from $G$, the active graph, and the
allowed varying $K$; a numerical computation of a new denominator
is unnecessary because no such denominator occurs in these integral
diagrams.  This conclusion concerns $S$.  The separate primitive
height-one argument above, and the corrected relaxed coordinate
comparison, retain exactly their stated scopes.
\end{proof}

\begin{definition}[The class-functional coordinate]
\label{ex:arith-coordinate}
For a constituent whose generic cohomology has dimensions $(1,1)$
in degrees one and two and vanishes otherwise, use the determinant
line
\[
 \mathcal D=\det\nolimits^{-1} C
       \simeq H^1(C)\otimes H^2(C)^*.
\]
The class-functional coordinate $u$ is the coordinate of
$[z]\otimes[\lambda]\in\mathcal D$ in the retained integral
determinant basis. Here $z$ is the actual class and $\lambda$ is
the closed functional just constructed. For a generically acyclic
constituent set $u=0$. The primitive and relaxed subscripts
specify the corresponding local conditions on $C$.
\end{definition}

\subsection{Central identity, torsion and the primitive coordinate}
\begin{lemma}[Central specialization preserves the same class]\label{ex:arith-central-specialization}
The finite classes of Proposition~\ref{ex:arith-finite-heegner}
satisfy the retained boundary (4.5), with the same fixed $L$.
For the functional of Proposition~\ref{ex:arith-paired-functional},
the actual central Kummer comparison $i$ satisfies
\[
 z_0=i(LmP),\qquad
 (\lambda_z)_0(i x)=\operatorname{PD}_K(x,c_K(LmP)).
\]
Here $P$ is the unaveraged ring-class point and $c_K$ is its
conjugate-dual transport. On a rank-one Mordell--Weil free lattice,
both free-line coefficients consequently have valuation
$j_{\rm rc}(h,k)+v_2(Lm)$. The statement identifies the same point
and functional; their integral torsion factors are still computed
in the integral Kummer complex and local quotient triangles.
\end{lemma}
\begin{proof}\leavevmode
The split norm relation gives an actual finite-stage boundary
\[
 z_i(0)-Lm_i\,j_i(b_h)=dv_i,\qquad
 m_i=a_{r_i}-2\chi_h(r_i).                                      \tag{4.5}
\]
The equality in global cohomology lifts to the Selmer cone because
the local degree-zero maps are onto and the degree-one maps inject.
Retain $v_i$. The uniform trace bound bounds $v_2(m_i)$ and the local
determinants by $c_E$; their adjugate identities multiplied by a fixed
power of two survive the limit.

% Insert after baseline line 1583, before the ordinary compact Kummer complex.
\paragraph{Specialization of the same paired functional.}
Put $A=\Lambda_{(t)}$, with residue field $\mathbf Q_2$.
The representative of $\lambda_z$ over $A$ is the one obtained from
the actual orthogonal comparison and the retained cup pairing.
At an active prime its contraction uses $C=-2\in A^\times$.
At the moving prime the central Frobenius determinants are nonzero;
their adjugate inverses therefore belong to $A$. Both contractions
commute with $A\to\mathbf Q_2$: their defining identities are matrix
identities with invertible denominators in $A$. In particular the
central full local complexes are rationally acyclic. The Kummer local
objects at these changed odd places are also rationally zero:
local points have finite two-primary completion, and the retained
central Frobenius matrices have nonzero determinants. Thus the
full/Kummer quotients are rationally acyclic as well. At the fixed old
places the Kummer maps and their duality homotopies are retained unchanged.

Let $j$ denote the resulting orthogonal comparison over $A$ and
let $c(z)$ be the conjugate-dual transport of the actual class.
The definition is $z^\perp=j^{-1}c(z)$ and
$\lambda_z=\operatorname{PD}(-,z^\perp)$.
Derived base change preserves the specified chain maps, cup products,
and homotopies. It also preserves $j^{-1}$, since the base change of
the two inverse identities is again the two inverse identities.
Consequently
\[
 (\lambda_z)_0
 =\operatorname{PD}_0\bigl(-,j_0^{-1}c(z_0)\bigr).
\]
To make the central identification explicit, write $R_0=C_U(M)_0$,
$R_0^c=C_{cU}(M^*(1))_0$, and $S_0=C_{U^\perp}(M^*(1))_0$.
Let $K_0$ be the ordinary rational Kummer complex and $K_0^c$ its
conjugate-dual transport. The actual comparison factors as
\[
 i:K_0\longrightarrow R_0,\qquad
 s:S_0\longrightarrow K_0^c,\qquad
 q:K_0^c\longrightarrow R_0^c,\qquad j_0=q s.
\]
All three maps are quasi-isomorphisms, by the preceding local
vanishing and the unchanged old conditions. The retained duality
homotopies give
\[
 \operatorname{PD}_{R}(i x,y)=\operatorname{PD}_{K}(x,s y),
 \qquad c(iP)=q c_K(P).
\]
The boundary (4.5) gives $z_0=i(LmP)$. Therefore
\[
 (\lambda_z)_0(i x)
 =\operatorname{PD}_{K}
       \bigl(x,s(qs)^{-1}q c_K(LmP)\bigr)
 =\operatorname{PD}_{K}\bigl(x,c_K(LmP)\bigr).
\]
Thus this is the Kummer-dual functional of the same transported point.
On the integral rank-one Mordell--Weil free lattice, the central
conjugation transport is an involutive automorphism (including the
quadratic sign), hence acts by $+1$ or $-1$. It preserves the point
index; it need not fix the point. The two free-line coefficients both have valuation
$j_{\rm rc}(h,k)+v_2(Lm)$. An orientation sign is a fixed unit;
there is no independently chosen scalar multiplying the functional.
The integral torsion correction is still computed from the ordinary
integral Kummer complex and the integral local quotient triangles,
not from their rationally acyclic reductions. This is the compatibility
needed to apply the arithmetic length calculation in (4.6).
\end{proof}
\begin{proposition}[Arithmetic calibration at known simple zeros]\label{ex:arith-central-calibration}
Let $K=\Q(\sqrt{k})$ and $h$ satisfy the ring-class hypotheses
above, with every prime dividing $h$ split in $K$. Suppose
$L(E^{(h)},s)$ has a simple zero at $1$ and
$L(E^{(hk)},1)\ne0$. Then the coordinate of
Definition~\ref{ex:arith-coordinate}, with the central class and
functional of Lemma~\ref{ex:arith-central-specialization}, satisfies
\[
 v_2u_{\rm prim}(0)
 =2j_{\rm rc}(h,k)-4\omega(h)-s(h)-s(hk)
       -\sum_{p\mid k}t_p+O_{E,\pi}(1).
\]
The error depends on the fixed old local types and fixed multiplier,
but not on $h$ or $k$. Here $0\leq t_p\leq2$, and $t_p=2$ in the
present full rational two-torsion setting. Only after a bound on
$\omega(k)$ is imposed may the companion sum be absorbed.
The finite-Sha assertion in this calibration is supplied by the
two stated analytic hypotheses.
\end{proposition}
\begin{proof}\leavevmode
For the ordinary compact Kummer complex $C_K$, universal coefficients
and the local degree-zero identifications give
\[
 H^1(C_K)_{\rm tors}=E^{(h)}(K)[2^\infty],\qquad
 0\to H^1(C_K)\otimes\Q_2/\Z_2\to\Sel_{2^\infty}
       \to H^2(C_K)_{\rm tors}\to0.
\]
The last finite group is Selmer modulo its maximal divisible subgroup.
Self-duality identifies the degree-three torsion with the degree-one
torsion and gives a unimodular pairing of the free degree-one and
degree-two parts. Thus their rational dimensions equal full Selmer
corank, including a possible divisible part of $\Sha$.

At known simple zeros, Lemma~\ref{lem:forward-rank}
gives rank one and finite $\Sha$.
For restriction of scalars and the product of the two twists, the
sum/difference isogenies are dual and compose to $[2]$.
Their free-lattice indices satisfy $IJ=2$, with regulator ratio
$2/I^2$. At odd primes the differential lattices agree; at two the
pullback determinants have bounded total valuation.
Milne's isogeny ratio \cite[proof of I.7.3, (7.3.1)]{milne} follows
from duality without assuming the BSD leading-term equality.
It applies here with finite $\Sha$, canonical heights and pullback
measures.

At every active $p\mid h$, $K$ splits, so the isogeny local factors
cancel exactly. At either split place the full/Kummer quotient
has length
$\dim_{\F_2}E[2]^{\operatorname{Fr}_p}=2$, since all of $E[2]$
is rational. The two places therefore contribute exactly $-4$. At a companion prime $p\mid k$, the rational component number
of the ramified twist is $\#E[2]^{\operatorname{Fr}_p}$, whereas
the restriction of scalars has connected special fiber.
Thus $t_p=\dim_{\F_2}E[2]^{\operatorname{Fr}_p}\le2$;
in the present full rational two-torsion case it is exactly two.
The result is
\[
 v_2u_{\rm prim}(0)
 =2j_{\rm rc}(h,k)-4\omega(h)-s(h)-s(hk)
                    -\sum_{p\mid k}t_p+O_{E,\pi}(1).             \tag{4.6}
\]
Only the companion sum can be absorbed when $\omega(k)$ is bounded.
No class number or unbounded degree-two torsion length is absorbed.

Here is the determinant and isogeny bookkeeping in more detail. This
calculation is used only at the known simple-zero constituents, where
the forward theorems give finite $\Sha$. Put
$\mathcal A=\operatorname{Res}_{K/\Q}(E^{(h)}_K)$ and
$\mathcal B=E^{(h)}\times E^{(hk)}$, and let $g_{\mathcal A}$ and
$s_{\mathcal A}$ be the lengths of their two-primary rational torsion
and Tate--Shafarevich group, respectively, with analogous notation
for $\mathcal B$. The ordinary Kummer complex has torsion lengths
$g_{\mathcal A},s_{\mathcal A},g_{\mathcal A}$ in degrees one, two
and three. A torsion disk with cokernel of length $e$ in degree $i$
contributes $(-1)^{i+1}e$ to the inverse determinant coordinate.
The central class $Lm_iP_{h,k}^{\rm rc}$ and its paired functional
each have free-coordinate valuation $j_{\rm rc}(h,k)+v_2(Lm_i)$.
Thus their ordinary Kummer coordinate has valuation
$2j_{\rm rc}+2v_2(Lm_i)+2g_{\mathcal A}-s_{\mathcal A}$.

At either place above the split moving prime $r_i$, the central
full/Kummer quotient has degree-two length
\[
 v_2(N_i),\qquad
 N_i=r_i+1-\chi_h(r_i)a_{r_i}
     =(r_i-1)-\chi_h(r_i)m_i.
\]
Use the moving-character construction with
$r_i\equiv1\pmod{2^{B_i}}$, $B_i>c_E$ tending to infinity.
Since $v_2(m_i)\le c_E$, both local lengths equal $v_2(m_i)$.
Their total negative contribution cancels the two occurrences of
$v_2(m_i)$ in the class-functional tensor. This equality persists
in the common limit. At each active split place the corresponding
quotient has only degree-two cohomology $E[2]$, of length two.
The two places above each active prime therefore contribute $-4$.

For clarity, orient the isogeny indices as
\[
 I=[\mathcal B(\Q)_{\rm free}:\phi\mathcal A(\Q)_{\rm free}],
 \qquad
 J=[\mathcal A(\Q)_{\rm free}:\psi\mathcal B(\Q)_{\rm free}].
\]
Then $IJ=2$ and
$\operatorname{Reg}_{\mathcal B}/\operatorname{Reg}_{\mathcal A}
=2/I^2$. Choose global generators $\omega_{\mathcal A}$ and
$\omega_{\mathcal B}$ of the top cotangent lattices of their
N\'eron models, and use their period measures over $\Q$ throughout.
The integral pullbacks
$\phi^*\omega_{\mathcal B}=a\omega_{\mathcal A}$ and
$\psi^*\omega_{\mathcal A}=b\omega_{\mathcal B}$ satisfy $ab=4$.
Thus $a,b\in\{\pm1,\pm2,\pm4\}$. Real change of variables gives
\[
 |a|\Omega_{\mathcal A}
 =\frac{\#\ker\phi_{\mathbf R}}{\#\operatorname{coker}\phi_{\mathbf R}}
   \Omega_{\mathcal B},
\]
where both finite groups have order at most four. Hence the period
ratio has uniformly bounded two-adic valuation, without a separate
discriminant normalization. With these differential choices the
finite measure factors are the products $C_{\mathcal A},C_{\mathcal B}$
of component numbers. The arithmetic isogeny identity gives
\[
 s_{\mathcal A}-s_{\mathcal B}
 =2g_{\mathcal A}-2g_{\mathcal B}
  +v_2\frac{\operatorname{Reg}_{\mathcal B}}{\operatorname{Reg}_{\mathcal A}}
  +v_2\frac{\Omega_{\mathcal B}}{\Omega_{\mathcal A}}
  +v_2\frac{C_{\mathcal B}}{C_{\mathcal A}}.
\]
The $2g_{\mathcal A}$ terms cancel on substitution into the
determinant formula. The remaining rational torsion is uniformly
bounded by the fixed local types at two. Component factors cancel
at every split prime. At $p\mid k$, put $E_0=E^{(h)}$; it has good reduction, and
$K_p/\Q_p$ is a tame ramified quadratic extension.
By \cite[Theorem~5.3 and Remark~5.4.1, pp.~299--300]{edixhoven},
the component group of the ramified twist is the fixed subgroup
of the good special fiber under its descent action $[-1]$:
\[
 \Phi(E_0^{(k)})\simeq\widetilde E_0[2].
\]
This is an isomorphism of finite \'etale residue-field group
schemes, so its rational component number is
$\#E[2]^{\operatorname{Fr}_p}$.
The Weil restriction of the good model has special fiber an
extension of $\widetilde E_0$ by a vector group, and is connected.
Thus its component number is one.
All other component factors cancel at split primes or equal one
at the remaining good primes. Consequently
\[
 v_2(C_{\mathcal B}/C_{\mathcal A})
   =\sum_{p\mid k}t_p,\qquad
 t_p=\dim_{\F_2}E[2]^{\operatorname{Fr}_p}.
\]
Substitution into the coordinate containing $-s_{\mathcal A}$
gives the negative sign in (4.6).
In the full rational two-torsion case each $t_p=2$.
\end{proof}
\begin{lemma}[Unused-prime comparison and one corrected coordinate]\label{ex:arith-common-coordinate}
At an unused good split prime, use the actual primitive-to-relaxed
maps and the retained class and functional above. Write $z$ for
the scalar Frobenius value. With $e_q=a_q-z-z^{-1}$ and $l_q$
the two-place quotient factor,
they satisfy (4.7) with $\varepsilon_q\in\Lambda^\times$.
For a finite binary family the universal correction $\eta$
lies in $\mathscr O[G]$, and its corrected coordinate
$\widetilde u=\eta u_{\rm rel}$ satisfies
$\widetilde u_x=\varepsilon_xu_{{\rm prim},x}$ at every character,
with $\varepsilon_x\in\Lambda^\times$.
At any character where the primitive coordinate is regular at
$t=0$, this identity preserves its central valuation and vanishing.
Only the product is specialized; $\eta_x$ itself may have a pole.
\end{lemma}
\begin{proof}\leavevmode
At an unused prime the exact point factor is $e_q=a_q-z-z^{-1}$.
The second factor comes from the same closed functional, as follows.
Let $i:C_P\to C_R$ be the primitive-to-relaxed map, and let
$j:C_S\to C_{P^\perp}$ be the paired strict-to-primitive map.
All old orthogonal conditions and their homotopies are retained.
At $q$ these are the maps from the unramified condition to the full
condition and from zero to the unramified condition. Naturality of
the cup pairing \cite[Sections 6.4.1--6.4.3]{nekovar} gives
\[
             \beta_R(v,iw)=\beta_P(jv,w).
\]
The actual norm identity is $z_R=e_q i(z_P)$, with its retained
boundary. The generic paired strict lift satisfies
$jz_S=e_q^c z_{P^\perp}$. Conjugation sends $z$ to $z^{-1}$, so
$e_q^c=e_q$. Consequently
\[
                 \lambda_R\circ i=e_q\lambda_P.
\]
Thus the class and its paired functional each contribute one
$e_q$; no integral self-duality of the primitive condition is assumed.
The two quotient determinants satisfy
\[
 \delta_q=\det(z\operatorname{Fr}_q-q)
              \det(z^{-1}\operatorname{Fr}_q-q)=q^4l_q.
\]
In the original integral determinant bases, the support triangle
therefore gives
\[
       l_q u_{\rm rel}=e_q^2\varepsilon_q u_{\rm prim},
             \qquad \varepsilon_q\in\Lambda^\times.             \tag{4.7}
\]
To see the unit assertion, retain the finite mapping-cone bases.
Their comparison is a sign; replacing them by the retained homotopy
equivalent free models contributes only determinants of unit disks.
These are $A_i$-units. Their limits and inverse limits are
$\Lambda$-units, and $q^{-4}$ is one also. An arbitrary
$\mathscr O$-unit would not suffice.

We check the integrality needed for the universal correction.
The graph-wide prescription at the start of this section applies
also to this unused prime. Modulo the maximal ideal
$(2,I_G)$ of $\mathscr O[G]$, every group element and fixed quadratic
sign becomes one, and the scalar $z$ becomes $1+t$. Since $E[2]$
is rational, $a_q$ is even at every good odd prime. Thus
\[
 \overline e_q=(1+t)+(1+t)^{-1}=\frac{t^2}{1+t}
       \ne0\quad\hbox{in }\mathbf F_2((t)).
\]
Consequently $e_q$ is a unit in $\mathscr O[G]$. This assertion uses
the prescribed exponent at every $q\in Q$, not merely at the primes
ramified in a particular constituent. With
\[
 l_q=\det(1-q^{-1}z\operatorname{Fr}_q)
       \det(1-q^{-1}z^{-1}\operatorname{Fr}_q),
\]
the trace and determinant of $\operatorname{Fr}_q$ are $a_q$ and $q$.
Expanding these two degree-two polynomials gives the exact identity
\[
 q^2l_q=e_q^2+(q-1)\bigl(q-1-(z+z^{-1})e_q\bigr).
\]
Because $q$ is odd and $e_q$ is a unit, this identity implies
$l_q/e_q^2-1\in2\mathscr O[G]$. Set
$v_q=(l_q/e_q^2-1)/2\in\mathscr O[G]$.
Then $1+v_q(1+g_q)$ selects $l_q/e_q^2$ at an inactive character
and one at an active character. Their finite product is therefore
one integral $\eta\in\mathscr O[G]$, without any division by the
group order. Put $\widetilde u=\eta u_{\rm rel}$.
It already is one common family; no interpolation of the individual
$\varepsilon_x$ is required. Formula (4.7) makes
$\widetilde u_x=\varepsilon_xu_{{\rm prim},x}$.
The central zero and valuations therefore transfer unchanged.
The multiplier $\eta_x$ may itself have a pole at $t=0$.
Only the product $\widetilde u_x$, identified with the regular
primitive coordinate times a $\Lambda$-unit, is specialized there.
\paragraph{Simultaneous prescriptions and local lifts.}
We spell out why the preceding calculation applies to the whole
graph. First choose its actual finite set $Q$ by
Proposition~\ref{ex:app-final-family}. Its primes are good, odd,
split in $K$, and outside $2ND_K$. The theta prime has already
been chosen and is an exclusion. Choose one prime of $K$ above
each $q\in Q$. In Proposition~\ref{ex:arith-moving-character},
``active list'' can be this entire list of orientations: the
construction uses only these splitting and avoidance properties,
not ramification in any particular quadratic constituent.

Here is the compatibility check in its class-group presentation.
Adjoin orientations above $2N$ and auxiliary split primes generating
$\operatorname{Cl}(K)$, all in distinct conjugate pairs. Assign weight
one to every member of $Q$ and weight zero to the other primes.
For the relation basis $m^{(a)}$ and
$\alpha_a=\beta_a/\overline\beta_a$ used there, a relation
$\prod_a\alpha_a^{b_a}\in
K(\mu_{2^\infty})^{\times 2^n}$ implies
\[
 2^n\mid\sum_a b_am_j^{(a)}\qquad(j\text{ corresponding to }Q).
\]
Indeed, a root belongs to a finite cyclotomic layer, which is
unramified at each of these odd primes, and its valuation is integral.
Multiplying by the weights and summing proves that the prescribed
Kummer multipliers respect every relation. At a resulting split
prime $r$, their power-residue values therefore satisfy every
relation in the presentation of $\operatorname{Pic}(\mathbf Z+r\mathcal
O_K)$. The resulting character has value one at every chosen
orientation in $Q$, value zero at $2N$, and is surjective on the
residue-ratio kernel, hence on inertia at $r$. The conjugate
orientation has the negative value because the product ideal is
rational principal. The radical-fixing subgroup (4.4c) preserves
all these prescriptions while imposing the two trace bounds.
Its level depends only on the Tate image. Enlarging $Q$ consequently
changes the finite Chebotarev field, but not $c_E$.

Perform this construction separately at each coefficient stage,
excluding $Q$, the theta prime, the old support and the previously
chosen moving primes. Taking $n_i\geq2i$ gives
$(1+t)^{2^{n_i}}=1$ in $A_i$, so these are actual finite characters.
They need not be restrictions of one character of $G_K$: the
retained inner limit is the limit of the finite diagrams of
Lemma~\ref{ex:limits-inner-limit}. At each fixed graph orientation
their values are already exactly $1+t$, so that identity survives
every coefficient reduction and that limit. The norm identities
(4.1), their coefficient maps, and their boundary witnesses hold
before character evaluation. This is the same construction for
all constituents. At $t=0$ the moving character is trivial, so
its added prescriptions do not alter the central quadratic family.

There is a local-lift issue only for the universal quadratic
character. Write $g_q\in G$ for its inertia generator at $q$;
thus $x(g_q)=(-1)^{\lambda_q(x)}$. Fix a lift of arithmetic
Frobenius at the chosen orientation and define $z_q$ using that
lift. The moving character is unramified there, and its value is
$1+t$. The fixed quadratic character is unramified there also.
Every other graph prime contributes an unramified quadratic
character at $q$. Changing the Frobenius lift can therefore only
replace $z_q$ by $g_qz_q$. For an inactive character $x$,
$x(g_q)=1$, and $x(z_q)$ is the actual unramified Frobenius scalar.
For an active character its value is immaterial to the correction,
which is one. Thus we never assign an unramified Frobenius
eigenvalue to a ramified constituent.

This independence also holds inside $R=\mathscr O[G]$, before
evaluation. Both choices of $z_q$ have reduction $1+t$ modulo
$(2,I_G)$, so both corresponding $e_q$ are units. Put
$r_q=l_q/e_q^2$ and $r'_q=l'_q/(e'_q)^2$. They agree in
$R/(g_q-1)$, hence $(1+g_q)(r'_q-r_q)=0$. Since $R$ is
two-torsion-free and $r_q=1+2v_q$, $r'_q=1+2v'_q$, this gives
\[
 (1+g_q)(v'_q-v_q)=0.
\]
The correction $1+(1+g_q)v_q$ is consequently independent of the
lift. Reversing the orientation replaces $z_q$ by $z_q^{-1}$
and leaves both $e_q$ and $l_q$ unchanged.

For clarity, write $s_q=z_q+z_q^{-1}$. Direct expansion gives
\[
 q^2l_q=(q-1)^2+a_q^2-a_q(q+1)s_q+qs_q^2,
 \qquad
 v_q=\frac{q-1}{2q^2e_q^2}
       \bigl(-(q+1)e_q^2+(q-1)-s_qe_q\bigr).
\]
Every term in the second expression is integral in $R$ because
$q$ is odd and $e_q$ is a unit. These inverses are taken only
after the inner limit and localization at $(2)$, never in the
finite ring $A_i[G]$. Reduction modulo $(2,I_G)$ gives
$\overline e_q=t^2/(1+t)$ and
$\overline l_q=t^4/(1+t)^2$. The latter identity also holds at
each sign evaluation modulo two. At augmentation the local
two-place determinant is $q^4l_q$, so its residual $t$-order is
exactly four. Multiplying the correction factors over the finite
set $Q$ proves the asserted universal integrality and all the
characterwise identities simultaneously.
\end{proof}

\begin{lemma}[The unknown base has zero coordinate if its derivative vanishes]\label{ex:arith-base-zero}
Use the base $h_0=1$ in the final fixed-$K$ construction, with
$L(E^{(k)},1)\ne0$ and full two-primary Selmer corank one for
$E/K$. If $L'(E/K,1)=0$, then
$u_{{\rm prim},0}(0)=\widetilde u_0(0)=0$ for the same coordinates
as above. The primitive coordinate is regular at $(t)$ by the local
argument below, and its corrected relaxed coordinate has the same
property by Lemma~\ref{ex:arith-common-coordinate}. No finite
$\Sha$ assertion at the base is required.
\end{lemma}
\begin{proof}\leavevmode
Put $A=\Lambda_{(t)}$, with residue field $\mathbf Q_2$ and
fraction field $F$. Use the actual finite free primitive complex
and its retained class and functional over $A$. The latter is
integral over $A$ by Proposition~\ref{ex:arith-paired-functional}.
The conductor-one height formula makes the specified point torsion
when $L'(E/K,1)=0$. Thus its rational Kummer class vanishes, and
the retained boundary (4.5) makes the specialization of the actual
cycle $z$ zero. No new normalization of $L$ is made here.

Full Selmer corank one gives central dimensions $(1,1)$ in degrees
one and two and zero in every other degree. This uses the full
Selmer group, including a possible divisible part of $\Sha$.
Over the DVR $A$, a bounded finite free complex decomposes, by
successive Smith reductions, into free zero-differential summands
and two-term disks $[A\xrightarrow{t^r}A]$, up to units. A disk
with $r>0$ contributes one additional residue-field cohomology
class in each of its two adjacent degrees. Semicontinuity and
Euler characteristic therefore leave only two generic possibilities:
the complex is acyclic, or its dimensions are $(1,1)$ in degrees
one and two.

In the acyclic case the specified coordinate is zero by definition.
In the second case the generic and central dimensions agree in
every degree. Hence no disk with $r>0$ occurs. All cancellable
disks have unit differential, so integral determinant base change
identifies the coordinate of $[z]\otimes[\lambda_z]$ with the
coordinate of its central tensor, in the retained determinant basis.
That tensor is zero because $[z]_0=0$. This proves regularity at
$(t)$ and $u_{{\rm prim},0}(0)=0$ directly. Finally
$\widetilde u_0=\varepsilon_0u_{{\rm prim},0}$ with
$\varepsilon_0\in\Lambda^\times$ by
Lemma~\ref{ex:arith-common-coordinate}, so the corrected coordinate
is also regular and zero at the center. This proof uses neither
the numerical lower bound of Proposition~\ref{ex:limits-uniform-lower-bound}
nor finite $\Sha$ at the unknown base.
\end{proof}

% End 03a-arithmetic.tex


% Begin 03b-limits.tex
\section{The two limits and determinant transfer}\label{sec:limits}
This section uses the actual diagrams of Section~\ref{sec:arithmetic}.
Their cohomology ranks alone would not suffice: maps to original
cochains, local comparisons, cup products, classes and boundary
witnesses are retained together. The same retained data include the
integral duality maps, their homotopy inverses, and the homotopies
expressing compatibility of duality with localization.

\subsection{Finite free models and the inner limit}
\begin{lemma}[Finite free models and retained inner-limit diagrams]\label{ex:limits-inner-limit}
Let $A$ be an Artin local ring with nilpotent maximal ideal, and
let a complex of flat $A$-modules have finite bounded residual
cohomology. This includes the actual continuous cochains with
finite free coefficient modules used here. It admits a bounded finite free model
with actual inclusion, projection and contraction homotopy.
Every fixed finite diagram of maps and homotopies, including the
duality and local-condition data used here, transports to these
models.

For one fixed active support and binary family, retain a tuple
whose model ranks are bounded over the finite coefficient stages.
One fixed nonprincipal ultrafilter gives a diagram over $\Lambda$
at all coefficient precisions, preserving its matrix identities
and retained homotopy inverses. This assertion imposes no rank
bound uniform over all active supports.
\end{lemma}
\begin{proof}\leavevmode
Let $A$ be an Artin local ring with nilpotent maximal ideal $\mathfrak m$.
A flat $A$-module is free, including in infinite rank: lift a basis
modulo $\mathfrak m$, use nilpotence for surjectivity, and flatness and
nilpotence for vanishing of the kernel. A continuous finite-coefficient
cochain module is a filtered union of modules of functions on finite
clopen partitions, hence is flat.

Suppose the residual cohomology of a complex of such modules is finite
and bounded. Split the residual complex into cohomology and disks,
and lift graded bases. Let $d_0$ be the lifted disk differential and
$d=d_0+\epsilon$, where $\epsilon$ has entries in $\mathfrak m$.
Use the graded decomposition into cohomology, disk sources and disk
targets. Let $i_0,p_0$ include and project the cohomology summand,
and let $h_0$ invert $d_0$ from disk targets to disk sources, zero
on the other summands. Then
\[
 p_0i_0=1,\quad h_0^2=h_0i_0=p_0h_0=0,\quad
 1-i_0p_0=d_0h_0+h_0d_0.
\]
With this normalized contraction, the finite geometric series give
\[
 \begin{aligned}
 d_C&=p_0\epsilon(1+h_0\epsilon)^{-1}i_0,&
 i&=(1+h_0\epsilon)^{-1}i_0,\\
 p&=p_0(1+\epsilon h_0)^{-1},&
 h&=h_0(1+\epsilon h_0)^{-1}.
 \end{aligned}
\]
Substitution using $d^2=0$ gives $pi=1$ and $1-ip=dh+hd$.
The resulting bounded finite free model carries actual maps and
homotopies to the original cochains. Transport any fixed finite
collection of arithmetic maps by $f\mapsto p'fi$. The identity
$1-ip=dh+hd$ transports their compatibility homotopies as well,
including the second homotopies in the cup-product diagrams.

Fix one nonprincipal ultrafilter on coefficient stages. Bounded ranks
can be fixed on a set in that ultrafilter. Each finite tuple of matrix
entries then has a compatible limit at every coefficient precision,
giving matrices over $\Lambda$. Every polynomial identity between
the matrices survives. Retained homotopy inverses remain homotopy
inverses. Mere stagewise quasi-isomorphisms without those maps would
not justify a nonflat coefficient change.
\end{proof}

\begin{lemma}[The reduced generic residual rank bound]\label{ex:limits-residual-bound}
Fix $E$, an integer $A\geq0$, and a number $j$ of derivative labels.
For the ring instances with $\omega(k)\leq A$, reduction to
characteristic two followed by the stated support comparison
and inversion of $t$ gives a generic residual model of rank at most
\[
 N(E,A,j)\leq56\omega(2N)+8A+56j+70.
\]
This bound is independent of the number and sizes of active primes
and of $v_2(r_i-1)$. It bounds the reduced generic model, not the
whole retained support-comparison tuple before $t$ is inverted.
\end{lemma}
\begin{proof}\leavevmode
All ranks in this proof are sums of the ranks in all degrees of one
complex. They are not the size of the diagram retaining all comparison
maps. Put $v=\omega(2N)$ and $m=v+1+j$. Let $S_0$ be the finite
places of $K$ above $2N$, the moving prime, and the derivative primes.
Let $T$ be the set of rational primes below $S_0$, together with the
primes ramified in $K/\Q$, and put $s=|T|$. Since $2\mid2N$,
\[
                    s\le v+A+1+j.                            \tag{5.1a}
\]
Primes added to $T$ solely because they ramify in $K$ are used for
the global estimate. They do not impose additional Selmer conditions.

We first give the global estimate at the residue field
$B_i/(t)=\F_2$, where $B_i=\F_2[t]/(t^i)$. The residual coefficient
module is the trivial module $W=E[2]=\F_2^2$. Write $S_T$ for all
finite places of $K$ above $T$, and let $r$ be the number of finite
rational primes ramified in $K$. Kummer theory gives the exact sequence
\[
0\longrightarrow\mathcal O_{K,S_T}^{\times}/
                 (\mathcal O_{K,S_T}^{\times})^2
\longrightarrow H^1(G_{K,S_T},\F_2)
\longrightarrow\operatorname{Cl}(\mathcal O_{K,S_T})[2]
\longrightarrow0.
\]
Here $2\in T$ and $\mu_2=\F_2$ as a Galois module. The unit term
has dimension $1+|S_T|$: the unit rank is $|S_T|$ and the roots of
unity contribute one dimension. Moreover
\[
 |S_T|\le2s-r,\qquad
 \dim_{\F_2}\operatorname{Cl}(\mathcal O_{K,S_T})[2]\le r-1.
                                                               \tag{5.1b}
\]
For completeness, the class-group inequality follows directly from
the imaginary quadratic genus argument. Conjugation acts by inversion
on the ideal class group. If $[\mathfrak a]$ has order dividing two,
write $\overline{\mathfrak a}=(\alpha)\mathfrak a$. Taking ideal
norms gives $N_{K/\Q}\alpha=1$, since this norm is positive. Hilbert
90 gives $\alpha=\beta/\bar\beta$ (take $\beta=1+\alpha$
unless $\alpha=-1$, and take $\beta=\sqrt{k}$ in that case), so
$\beta\mathfrak a$ is a
conjugation-invariant fractional ideal. Modulo rational principal
ideals such ideals are generated by the ramified prime ideals.
Their classes have exponent two, and the principal ideal of
$\sqrt{d}$, where $d$ is the squarefree radicand of $K$, supplies
a nonempty relation among them. Thus their
span has dimension at most $r-1$. The excluded field $\Q(i)$
would also satisfy the inequality, since its class group is trivial.
Finally an $S$-class group is a quotient of the ordinary class group,
and the two-rank of a quotient of a finite abelian group cannot
increase. This proves (5.1b), including the possible ramification
at two. The bound on $|S_T|$ counts one prime over each ramified
rational prime and at most two over each other prime.

Inflation from $G_{K,S_0}$ into $G_{K,S_T}$ is injective on $H^1$
with trivial coefficients. Consequently
\[
 h^0(G_{K,S_0},W)=2,\qquad
 h^1(G_{K,S_0},W)\le4s.
\]
The global Euler characteristic and cohomological-dimension statements
\cite[I.4.10 and I.5.1]{milne}, applied over the imaginary quadratic
field $K$, give
\[
 h^0-h^1+h^2=-2,\quad H^{>2}=0,\quad
 h^2=h^1-4.
\]
Thus the total residual cohomology dimension of the global complex
is at most $8s-2$, and in particular at most $8s+2$.
The latter weaker bound can also be obtained without genus theory:
the maximal pro-two group over $\Q$ unramified outside $T$ and
infinity has $s+1$ generators, by the squareclasses of $-1$ and
the primes in $T$; its subgroup for $K$ has at most $2s+1$
generators by the index-two Schreier bound. The maximal pro-two
group for $G_{K,S_0}$ is a quotient of this subgroup. Hence
$h^1(G_{K,S_0},W)\le4s+2$, giving total dimension at most $8s+2$.
In either computation the $H^2$ estimate is supplied by the global
Euler characteristic, not by a generator estimate alone.

We next bound the local terms in the Selmer cone (4.3). For a finite
place $w$, put $d_w=[K_w:\Q_2]$ if $w\mid2$ and $d_w=0$
otherwise. Local duality and the local Euler characteristic
\cite[I.2.3 and I.2.8]{milne} give, for the trivial module $W$,
\[
 (h^0,h^1,h^2)(K_w,W)=(2,4+2d_w,2).
\]
At an old Kummer place, the derived condition has residual dimensions
$(2,2+d_w)$ in degrees zero and one. Indeed
$E^{(h)}(K_w)^\wedge_2$ is the direct sum of a free
$\Z_2$-module of rank $d_w$ and finite two-primary torsion;
that torsion has two generators because $E^{(h)}[2]$ is rational.
Derived reduction retains its two-dimensional degree-zero torsion.
For the other conditions at the odd moving or derivative places,
the following table records the sum of the dimensions of the local
condition and of the full local complex. Cone shifts do not change
this sum.
\[
\begin{array}{c|c|c}
\text{condition }U_w&\sum h^a(U_w)&
       \sum h^a(U_w)+\sum h^a(K_w,W)\\ \hline
\text{old Kummer}&4+d_w&12+3d_w\\
\text{unramified, }w\nmid2&4&12\\
0,\ w\nmid2&0&8\\
\text{full, }w\nmid2&8&16\\
\text{derivative strict}&2&10\\
\text{derivative finite or singular}&4&12\\
\text{derivative relaxed}&6&14
\end{array}
\]
The last three rows use exactly the conditions retaining $D^0$
in Lemma~\ref{ex:limits-derivative-switch}; the relaxed condition
does not include $D^2$. There are at most two places over each
rational prime, and $\sum_{w\mid2}d_w=2$. The old prime two
therefore costs at most $30$, any other old rational prime at most
$24$, the moving prime at most $32$, and a derivative prime at
most $28$ (in fact it is inert and costs at most $14$).
Each is bounded by $48$. Thus all finite local terms cost at most
$48m$. No prime dividing $k$ alone is charged a local condition.

There are no real places of $K$. Its one complex place has full
cohomology $W$ in degree zero only. With the usual modified
archimedean convention its local complex is zero; even retaining
both an ordinary full term and a full local condition costs at
most four additional generators. Complex conjugation used to
transport duality is not an additional real local condition over
$K$. In particular no infinite two-periodic real resolution is
being included in this count.

At each finite coefficient stage, apply
Lemma~\ref{ex:limits-inner-limit} separately to these global and
local complexes and transport the actual local maps. The mapping
cone then has a finite free model whose total rank is bounded by
the sum just computed. Further cancellation can only decrease it.
The estimate applies before inverting $t$ to the smaller-support
model, uniformly in the moving prime and in the order of its
finite character: reduction modulo $t$ kills that character.
Allowing the weaker global estimate and the archimedean allowance
still gives
\[
 \begin{split}
 \operatorname{rank}_{\rm tot} C_{S_0}
 &\le(8s+2)+48m+4\\
 &\le8(s+1)+6+48(v+1+j)\\
 &\le56v+8A+56j+70.
 \end{split}\tag{5.1}
\]

It remains to transfer this estimate to the original support.
First consider a scalar ring instance. Reduction modulo two kills
$\chi_h$, so the coefficient system is unramified at every active
prime. The actual support comparison of
\cite[Proposition 7.8.8(ii)]{nekovar} therefore identifies the
smaller-support complex with the complex having unramified
conditions at those primes. At an oriented active place, reduction
of (4.4) has
\[
 C=0,\qquad A=D=\big((1+t)^e-1\big)I_2,
\]
where $e$ is odd; at the conjugate place the exponent is $-e$.
Both $((1+t)^e-1)/t$ and $((1+t)^{-e}-1)/t$ are units of
$\F_2[[t]]$. Thus the full, unramified and singular comparison
blocks all become contractible after inverting $t$. For each fixed
finite active support retain the comparison maps and homotopies,
take the inner limit, and only then invert $t$. Every active
comparison cone is then contractible, so the original complex has
the same generic residual cohomology as the bounded smaller-support
model. No count of active primes remains.

For clarity, a universal binary character is not trivial merely
after reduction modulo two. In the binary setting make the preceding
calculation after augmentation $\F_2[G]\to\F_2$. After the inner
limit and inversion of $t$, the ring $\F_2((t))[G]$ is local with
nilpotent augmentation ideal. A finite free model over this ring
has a minimal model obtained by cancelling unit entries; its ranks
are the dimensions of the cohomology after augmentation, and hence
obey (5.1). These are ranks over $\F_2((t))[G]$, not dimensions
over $\F_2((t))$. For scalar instances they are precisely the
generic residual ranks needed in the outer limit. The whole retained
comparison diagram can still grow with $G$ or with the active
support; no bound for that diagram is asserted. Finally the minimum
ranks over $\mathscr O$ are the dimensions after reduction to
$\F_2((t))$, so the same estimate bounds the scalar
$\mathscr O$-minimal models used below.
\end{proof}

\subsection{Evaluation and the outer valuation ring}
\begin{lemma}[Evaluation, coefficient change and the outer DVR]\label{ex:limits-outer-evaluation}
For the retained global models and for the primitive and derivative
Selmer models, evaluations at sequences of Galois elements satisfy
(5.2). A cycle defines an abstract crossed homomorphism of the
indicated product group. The induced map on $H^1$ is injective,
and the finite detecting-list comparison remains injective after
field extension. Direct evaluation of other Selmer models is asserted
only when their local maps induce isomorphisms on residual $H^0$
and injections on residual $H^1$. For orthogonal zero/full comparisons,
the assertion may instead be transported at a test ring where the
actual local comparison complexes are contractible.

For a sequence of ring instances with $\omega(k_j)\leq A$ and
a fixed number of derivative labels as in
Lemma~\ref{ex:limits-residual-bound}, the quotient
$\mathscr O_\infty$ in (5.3) is a DVR with uniformizer two.
The bounded minimal-model matrices, evaluations, duality maps and
their compatibility homotopies pass to it. Sequences whose
differentials tend to zero give actual crossed homomorphisms
there. No continuity assertion on a single original Galois
group, or completeness of $\mathscr O_\infty$, is required.
\end{lemma}
\begin{proof}\leavevmode
For fixed $K$ let $\Gamma_K=\prod_iG_K$ and include one simultaneous
complex conjugation. Evaluation of retained cochain maps at arbitrary
sequences $g=(g_i),h=(h_i)$ gives
\[
 e^1_gd=(\rho(g)-1)e^0,\qquad
 e^2_{g,h}d=\rho(g)e^1_h-e^1_{gh}+e^1_g.                          \tag{5.2}
\]
One fixed ultrafilter supplies these identities for all sequences.
A cycle therefore defines an abstract crossed homomorphism of the
product group. We do not assert that its restriction to one original
Galois group is a continuous pointwise limit.

Evaluation remains injective after field extension. At each finite
stage take representatives $V$ for $H^1$ and coefficient space $W$. Evaluate
$(c,w)$ by $b_c(g)-(\rho(g)-1)w$. The intersection of all kernels is
$(0,W^\Gamma)$. A finite list of at most $\dim V+\dim W$ evaluations
already has that intersection. Choose and pad these lists at the finite stages before taking their common limit. The comparison to $[W\to W^d]$ is an
isomorphism on $H^0$ and injective on $H^1$. Its cone has no residual
cohomology in degrees at most zero. Cancel unit disks there so that
the free cone starts in degree one; the same is true after every
coefficient extension. The long exact sequence gives the claimed
injection. For the primitive and derivative Selmer complexes the required
residual comparison with global cochains follows from their actual
local conditions. At an old Kummer place, derived reduction retains
$B(F)[2]$ in degree zero, isomorphically onto the local invariants,
and has the Kummer injection in degree one. The unramified condition
at a good place has the same degree-zero invariants as the full local
complex and injects in degree one. A full condition has the identity
map. Every derivative condition retains $D^0$ and includes its stated
degree-one subspace. The Selmer-cone exact sequence therefore gives
an isomorphism on global $H^0$ and an injection on global $H^1$.

This does not apply directly to the zero conditions in the retained
orthogonal complexes: at $(2,t)$ their local degree-zero invariants
need not vanish. Indeed, two zero local conditions with residual
invariant space $W$ can contribute
$(W\oplus W)/\operatorname{diag}(W)$ to Selmer $H^1$, and global
cochain evaluation kills this contribution. Such orthogonal complexes
are retained for duality, without asserting residual evaluation
injectivity. If an actual zero/full comparison is contractible at a
chosen test ring, its retained contraction identifies the two Selmer
models there and transports the evaluation injection. This use requires
that contraction at the chosen ring; it is not an assertion for arbitrary
unused relaxed primes at every scalar height-one localization.

Now allow $K_j,h_j$ to vary with $\omega(k_j)\le A$.
For copies $\mathscr O_j$ of $\mathscr O$, form
\[
 \mathscr O_\infty=
 \left(\prod_{\mathcal U}\mathscr O_j\right)/
 \{(x_j):v_2(x_j)\to_{\mathcal U}\infty\}.                         \tag{5.3}
\]
A nonzero element has valuation one fixed nonnegative integer on a
set in $\mathcal U$ and is that power of two times a unit.
Every nonzero ideal has an element of least valuation and is
principal. Thus (5.3) is a DVR with uniformizer two and characteristic-zero
fraction field. Completeness is unnecessary.

Integral minimal-model matrices and (5.2) pass to (5.3).
So do the retained duality maps, inverse homotopies and
duality-to-localization compatibility homotopies. After minimal
cancellation an integral homotopy equivalence is degreewise invertible:
its residual matrices are invertible, and Nakayama applies.
These inverse matrices and all compatibility identities pass to the
same outer quotient. Thus perfect duality and exact orthogonality
are preserved there; they are not inferred merely from isotropy.
Approximate cycles suffice: if $v_2(dw_j)\to\infty$, the degree-two
term in (5.2) tends to zero, so $w_\infty$ is an actual crossed
homomorphism. The finite detecting lists have bounded length by
(5.1); retain their cone contractions to preserve injection.
One must not replace approximate cycles by unexplained exact cycles.
\end{proof}

\subsection{One new prime detects two classes}
\begin{lemma}[The common kernel and finite realization of a rank-two test]\label{ex:limits-kernel-detection}
Use the exact kernel of sequences acting trivially on the full
Tate module, the relevant ring-class fields and quadratic
characters. On the characteristic-zero tests used below,
restriction of the evaluated cohomology to this kernel is injective.
When the moving scalar character is nonquadratic, two independent
classes have a rank-two localization at a single actual kernel
element. Finite Chebotarev realizes the required conjugation
coset and basis-cochain evaluations, with exact inner-limit
matrix identity (5.4) before any $\mathscr O$-minimal-model
denominators are applied.

At the residual test only scalar-inertia vanishing, bounded
cohomology and finite detecting lists are asserted. The quadratic
character exception at characteristic-zero height-one tests is
treated in Proposition~\ref{ex:limits-uniform-lower-bound}.
\end{lemma}
\begin{proof}\leavevmode
The exact kernel consists of sequences trivial on the full Tate
module, all relevant ring-class fields and quadratic characters.
The scalar element of Section~\ref{sec:arithmetic} is central modulo
this kernel, though not in the kernel itself.
For a cocycle zero on the kernel,
\[
 (u^4-1)b(g)=(g-1)b(z)
\]
makes it a coboundary over characteristic zero. The same fixed scalar
survives the outer quotient. To eliminate residual invariants, use moving-prime
inertia: its scalar difference is $t$, invertible in $\F_2((t))$.
These are two separate arguments.

On the characteristic-zero tests used for rank reduction, the fixed
finite matrix-span witness gives absolute irreducibility of each
rank-two summand. When the scalar character is not quadratic, moving
inertia has distinct characters on the two summands and simultaneous
conjugation exchanges them. The induced rank-four representation is
therefore absolutely irreducible with scalar endomorphisms. The outer
fraction field is one such characteristic-zero test; the quadratic
character exception at an inner test is treated separately below.

For two independent cohomology classes on these tests, the span of
their joint values on the exact kernel projects onto both copies of
this absolutely irreducible representation. A proper submodule with
both projections onto would be the graph of a scalar isomorphism,
contradicting independence. Thus the span is the full sum. If the
determinant of their two finite localizations vanished on every actual
kernel element, additivity and polarization would make it vanish on
the entire span, impossible. Hence one actual element supplies a
rank-two localization on these characteristic-zero tests.

The polarization identity itself is also valid in characteristic two,
but no residual irreducibility or residual rank-two localization is
asserted. Here $E[2]$ is trivial of rank two, so its scalar twist is
isotypic and the induced residual representation has multiplicity two.
At the residual test we use only scalar-inertia vanishing of $H^0$,
the bounded residual cohomology dimensions, and finite detecting lists
from intersections of linear evaluation kernels. Those lists need not
consist of one rank-two localization. The rank-two selection for the
uniform valuation bound occurs only after passage to the
characteristic-zero fraction field of the outer DVR.

Chebotarev is applied at finite stages, not to the product group.
Each of the finitely many original basis cochains in degrees zero
through two factors through a finite quotient. Include their quotient,
the required Tate precision, ring-class and quadratic fields and old
diagrams in one compositum. The element just found, multiplied by
conjugation, is a single actual automorphism of that compositum.
Its restrictions agree on intersections. Choose the new prime and
a place above it to realize that representative.

This also resolves the precision issue caused by localization.
Choose the new prime to match all original basis-cochain evaluations
at a common retained precision $n(i)\to\infty$. Taking the same inner
ultrafilter gives an exact equality of $\Lambda$-matrices
\[
                 L_j=E_{\sigma_j^2}.                            \tag{5.4}
\]
Only afterward compose with the inclusion $I_j$ from an
$\mathscr O$-minimal model:
$L_jI_j=E_{\sigma_j^2}I_j$.
The denominators in $I_j$ may depend on $j$; equality (5.4) makes this
irrelevant. Long $t$-jets are never treated as approximations in the
two-adic valuation of $\mathscr O$, where $t$ is a unit.
\end{proof}

\subsection{Switches and the varying-companion estimate}
\begin{lemma}[Derivative local conditions and determinant-preserving switch]\label{ex:limits-derivative-switch}
At a derivative prime use
$L_v=D^0\oplus(F_v\oplus S_v)[-1]\oplus D^2[-2]$ with the
transported pairing, and retain $D^0$ in all four conditions.
Take $U_{\rm str}=D^0$,
$U_{\rm rel}=D^0\oplus(F_v\oplus S_v)[-1]$,
$U_F=D^0\oplus F_v[-1]$, and $U_S=D^0\oplus S_v[-1]$.
Then strict and relaxed are orthogonal, the two plane conditions
are self-orthogonal, and their cone triangles are the ones below.

Over the stated fraction or test field, suppose the actual old
class spans generic $H^1$ and has nonzero finite localization $v$.
Its derivative from Proposition~\ref{ex:arith-finite-heegner}
has finite coordinate zero and singular coordinate $Jv$.
The strict generic $H^1$ is zero, the new generic $H^1$ is spanned
by this derivative, and the two class-functional determinant
coordinates have equal valuations, even when $v$ is not primitive
in the local integral lattice.
\end{lemma}
\begin{proof}\leavevmode
\paragraph{The strict and relaxed conditions at a derivative prime.}
The words strict and relaxed have a different local meaning at a
derivative prime from zero and full at an active or unused split
prime. At a derivative prime use the actual split local model
\[
 L_v=D^0\oplus(F_v\oplus S_v)[-1]\oplus D^2[-2]
\]
with zero differential at the retained precision, and retain its
pairing and isotropy homotopies. Here $D^0$ and $D^2$ are dual,
and $F_v,S_v$ are the finite and singular isotropic planes. Put
\[
 \begin{aligned}
 U_{\rm str}&=D^0,&
 U_{\rm rel}&=D^0\oplus(F_v\oplus S_v)[-1],\\
 U_F&=D^0\oplus F_v[-1],&
 U_S&=D^0\oplus S_v[-1].
 \end{aligned}
\]
Thus $U_{\rm str}^{\perp}=U_{\rm rel}$ and the two plane
conditions are self-orthogonal under the transported duality.
All four maps retain local degree-zero cohomology. In particular,
the relaxed condition here omits $D^2[-2]$ and is not the full
local complex. The actual cone triangles are consequently
\[
 C_{\rm str}\longrightarrow C_F\longrightarrow F_v[-1]
 \longrightarrow C_{\rm str}[1],\qquad
 C_{\rm str}\longrightarrow C_S\longrightarrow S_v[-1]
 \longrightarrow C_{\rm str}[1].
\]
These are the triangles used in the determinant switch. Zero/full
conditions at the other primes retain their previously stated
meaning. Exact annihilator and dimension assertions in the switch
argument are made after passage to the stated fraction or test
field, using these integral diagrams and their retained dualities.
At a new inert prime the finite and singular planes are free of rank
two, with cross pairing
$\beta(x,y)=\varepsilon\langle x,Jy\rangle$.
Suppose the old actual class $A$ spans generic $H^1$ and has
$f(A)=v\ne0$. Its derivative $B$ satisfies $f(B)=0$, $s(B)=Jv$.
The strict $H^1$ vanishes. The transported duality and localization
identities identify the relaxed localization image with the kernel
of its dual localization map, hence with its exact annihilator
\cite[Sections 6.3.3--6.3.4]{nekovar}. This image
contains $(v,0)$ and $(0,Jv)$, so these span
it. Hence the new generic $H^1$ is spanned by $B$.

The determinant comparison is exact even if $v$ is divisible.
In the fraction field choose $x$ with $v\wedge x$ of unit volume,
and similarly $y$ with $Jv\wedge y$ of unit volume.
The finite and singular determinant exact sequences have the same
strict factor, so their class-functional coordinates have equal
valuations. All strict torsion remains in that common factor.
Requiring $x$ integral would be wrong when $v$ is divisible.
\end{proof}

\begin{proposition}[Uniform ring lower bound and height-one regularity]\label{ex:limits-uniform-lower-bound}
Fix $E,\pi$, the old local types and a bound $A$.
For every ring instance satisfying the hypotheses of
Proposition~\ref{ex:arith-central-calibration} and
$\omega(k)\leq A$, the actual primitive coordinate satisfies
\[
 v_2u_{\rm prim}\geq-C(E,A),\qquad
 u_{\rm prim}\in\Lambda[1/2],
\]
where the first valuation is that of $\mathscr O$ and the
constant is independent of the active support and companion field.
Consequently its central specialization and (4.6) imply (3.5).

The regularity assertion also holds at the unknown full-corank-one
base used in Lemma~\ref{ex:arith-base-zero}. The uniform numerical
lower bound here is asserted only for the known simple-zero ring
instances; it is not asserted at that unknown base.
\end{proposition}
\begin{proof}\leavevmode
The starting generic kernel has rank one. Indeed, at a known simple
zero the ordinary class is nonzero. Equation (4.5), with its bounded
nonzero multiplier, makes the retained central class nonzero.
Over $\Lambda_{(t)}$ cancel the degree-zero and degree-three disks.
If the generic kernel of the resulting two-term free complex vanished,
the actual lifted cycle would vanish, contradicting its central value.
Semicontinuity bounds that kernel by the central dimension one.

Suppose no uniform lower bound exists. Choose instances with
$v_2u_j\to-\infty$, $\omega(k_j)\le A$, and write their minimal
completed models as $[\mathscr O_j^n\xrightarrow{D_j}\mathscr O_j^n]$,
fixing the bounded $n$ on an outer-large set. In each generic
one-dimensional kernel choose a primitive vector $v_j$.
Its outer image is nonzero, so the outer kernel dimension $s$ is
at least one. If $s>1$, one new derivative prime gives rank-two
localization including that vector. The relaxed localization image
is self-annihilating and contains the whole finite plane, hence equals
that plane. The strict kernel has dimension $s-2$ and the singular
switch adds no localized vector, so the new outer dimension is $s-2$.

The actual inner class is a nonzero scalar multiple of $v_j$.
Even if this scalar tends to zero, its generic localization is nonzero
on an outer-large set. Exact switch preservation therefore keeps
each inner determinant valuation. The new inner kernels again have
primitive vectors. Thus their outer dimension remains at least one;
a proposed drop to zero is already a contradiction.
There can be only finitely many drops. For each fixed number of them
(5.1) bounds the models; no bound uniform in an unbounded number of
switches is needed.

At terminal outer dimension one a maximal boundary minor $\Delta$
is nonzero. Its inner minors have one finite valuation $b$ on an
outer-large set. Determinant-one Schur operations in a square
two-term rank-one-kernel complex give
$u=\pm\alpha\beta/\Delta$, where $\alpha$ is an integral surviving
cycle coordinate and $\beta$ has valuation at least $-a$.
Consequently $v_2u_j\ge-a-b$, the required contradiction. We obtain
\[
                     v_2u_{\rm prim}\ge-C(E,A).                 \tag{5.5}
\]
Here (5.5) is asserted for the known simple-zero ring instances
used to obtain the varying-companion estimate. Their nonzero actual
central class supplied the rank-one start. It is not a uniform
valuation assertion at the unknown corank-one base.
Bounded matrix size alone would not prove it: $[2^j]$ is a
counterexample to that inference.

% Replace baseline lines 4112--4138.
\paragraph{Height-one regularity, including quadratic specializations.}
Fix a height-one prime $\mathfrak p\ne(2)$, put
$R=\Lambda_{\mathfrak p}$ and $k=\operatorname{Frac}(\Lambda/\mathfrak p)$,
and use the actual primitive scalar Selmer complex over $R$.
Here $R$ is a DVR and $k$ has characteristic zero. All subsequent
switches use the derivative conditions of
Lemma~\ref{ex:limits-derivative-switch}, retaining $D^0$.
The primitive old local conditions and these derivative conditions have
the evaluation injection of Lemma~\ref{ex:limits-outer-evaluation}.
Cancel the degree-zero and degree-three disks and then the unit disks.
Write $C=[R^s\xrightarrow{d}R^s]$ in degrees one and two, with
$d\otimes k=0$. Its generic kernel has dimension one. Choose a primitive
vector $v\in\ker d$. Such a vector exists because $\ker d$ is saturated
in $R^s$: divide a nonzero kernel vector by its largest common power
of a uniformizer. Thus $\bar v\ne0$ in
$V=H^1(C\otimes_R k)$ and $\dim_k V=s\ge1$.
The actual class $z$ is a nonzero fraction-field multiple of $v$.
We will only make switches whose finite localization is nonzero on
$\bar v$. Consequently their generic localization of $z$ is nonzero,
which is the hypothesis needed for the actual-class determinant switch.

If the scalar character specialized to $k$ is nonquadratic, the
rank-two detection of Lemma~\ref{ex:limits-kernel-detection} applies to
$\bar v$ and an independent vector whenever $s>1$. The finite
localization then has rank two. Exact duality makes the relaxed
localization image equal to the finite plane, so switching to the
singular plane lowers $s$ by two. We give the quadratic case in detail.

Let $\alpha$ be the specialized character on the index-two group.
Since the moving inertia generator has value $1+t$, quadraticity
implies $t(t+2)=0$ in $k$. The only possible primes are therefore
$(t)$ and $(t+2)$. The equality $\alpha^\tau=\alpha^{-1}=\alpha$
uses the actual conjugation $\tau$; it involves no automorphism of
$R$ or substitution for $t$. The action
\[
 c^\tau(g)=Jc(\tau^{-1}g\tau),\qquad J=\rho_T(\tau),
\]
on cocycles gives a $k$-linear involution on cohomology. Extend
$\alpha$ to the index-two overgroup by choosing its value on $\tau$
to be $+1$. The other extension differs by the character of $K$.
The induced rank-four representation is consequently
$W_+\oplus W_-$, with each $W_\pm$ of rank two. The matrix-span
witness in (H1)--(H2) and the paragraph following them makes their
restriction to the index-two group absolutely irreducible. Twisting
its matrices by nonzero scalar character values does not change their
linear span. The summands are nonisomorphic: an intertwiner restricts
to a scalar, whereas their actions of $\tau$ are $J$ and $-J$.
A scalar intertwining these two actions is zero since $2$ is invertible.
We write $V=V_+\oplus V_-$ for the resulting cohomology signs.

At a new inert derivative prime write $F,S$ for the finite and
singular planes. The action on $F$ is $J$, and that on $S$ is $-J$.
Indeed, the local rational Frobenius acts on the Tate module by $J$;
it fixes its square, the finite Frobenius generator, and conjugates
tame inertia by its $\ell$-th power. The retained congruences give
$\ell=-1$ in the limit, so pullback on the tame coordinate contributes
$\ell^{-1}=-1$. These assertions hold at every previously chosen
derivative prime as well. A common scalar sign would simply interchange
the names $+$ and $-$.
The transported cross pairing is exactly
\[
 \beta(x,y)=\varepsilon e_2(x,Jy),\qquad \varepsilon\in k^\times,
\]
from Lemma~\ref{ex:limits-derivative-switch}. The finite and singular
planes are isotropic. These are the retained local cup-product
identities, including the extra finite precision in (H12), not a
residual alternatingness assertion. Since $\det J=-1$ and $J^2=1$,
$e_2(x,Jy)=e_2(y,Jx)$. Choose nonzero $J$-eigenvectors $e_+,e_-$.
They form a basis and $e_2(e_+,e_-)\ne0$. The two local cohomology
signs are therefore
\[
 H_+=F(e_+)\oplus S(e_-),\qquad
 H_-=F(e_-)\oplus S(e_+).
\]
They are orthogonal nondegenerate symmetric hyperbolic planes.
In either plane the self-pairing of $aF(e_\pm)+bS(e_\mp)$ is
$2ab$ times a nonzero constant. Its only isotropic lines are thus
its finite and singular axes.

All old conditions are stable under this involution. Conjugate old
Kummer conditions are transported into each other; primitive ramified
full conditions occur at the conjugate pair; unused good conditions
are unramified; and both $F$ and $S$ at a previous derivative prime
are stable under the actions just computed. At the new prime let
$C_{\rm str}$ and $C_{\rm rel}$ have the actual truncated conditions
$D^0$ and $D^0\oplus(F\oplus S)[-1]$. Put
\[
 U=\operatorname{im}\bigl(H^1(C_{\rm rel}\otimes k)\longrightarrow
 F\oplus S\bigr),\qquad
 K_0=H^1(C_{\rm str}\otimes k).
\]
The retained localization-duality diagrams give $U=U^\perp$.
Specifically, localization's connecting map is its dual under the
strict/relaxed duality, so its kernel is both $U$ and $U^\perp$;
this is the exactness in \cite[6.3.3--6.3.4]{nekovar} used in
Lemma~\ref{ex:limits-derivative-switch}. Isotropy alone would not suffice.
The cone triangles, with $D^0$ retained, give exact sequences
\[
 0\longrightarrow K_0\longrightarrow V
   \longrightarrow U\cap F\longrightarrow0,
 \qquad
 0\longrightarrow K_0\longrightarrow V'
   \longrightarrow U\cap S\longrightarrow0,                 \tag{5.5a}
\]
where $V'$ is the switched fiber cohomology. There is no preceding
local degree-zero term, since the quotient conditions are
$F[-1]$ and $S[-1]$.
The space $U$ is stable under signs. Each $U_\pm\subset H_\pm$
is its own annihilator in $H_\pm$, hence is one of that plane's
two isotropic axes.

We also need one prime detecting prescribed classes of opposite
signs. Choose nonzero $c_+\in V_+$ and $c_-\in V_-$.
Their restrictions to the exact kernel $H$ of
Lemma~\ref{ex:limits-kernel-detection} are nonzero additive equivariant
maps. Their joint $k$-linear evaluation span in $W_+\oplus W_-$
surjects to each summand. Simplicity and nonisomorphism force this
span to be the full direct sum. At an element $h\in H$, evaluation
at $(\tau h)^2$ on a sign-$\delta$ cocycle is
\[
 c_\delta((\tau h)^2)=(1+\delta J)c_\delta(h)
 \qquad(\delta=+1,-1).                                      \tag{5.5b}
\]
This follows from additivity on $H$ and the sign relation; a
coboundary is zero on $H$. Thus the two finite-coordinate maps on
the joint span surject respectively onto $k e_+$ and $k e_-$.
There is an actual $h$ for which both coordinates are nonzero.
For completeness, if their product vanished at every actual $h$,
additivity and the identity
$Q(x+y)-Q(x)-Q(y)$ for that product would make its polar form
zero on all pairs of actual evaluations. It would then vanish on
their linear span, contradicting surjectivity onto the two axes.
Finite Chebotarev and the exact matrix identity (5.4) realize this
choice by a new inert prime, avoiding every previously used prime.
All character, Tate, cochain and old-diagram constraints are included
in the same finite quotient as in that lemma. No residual
irreducibility or infinite-level Chebotarev assertion is used.
The same argument for one nonzero sign class gives nonzero
localization in its finite axis.

Suppose first that both $V_+$ and $V_-$ are nonzero. Choose the
class in one sign to be a nonzero component of $\bar v$ and choose
any nonzero class in the other. Detection as above puts both finite
axes in $U$, so $U=F$ and $U\cap S=0$. Equation (5.5a) gives
$\dim V'=s-2$. The localization of $\bar v$ is nonzero: its two
sign components lie on different finite axes, so cannot cancel.

Suppose instead that only one sign occurs, say $V=V_+$, with
$s>1$. Detect $\bar v$ in its finite axis. Then
$U_+=F(e_+)$. Since $V_-=0$, (5.5a) gives
$U_-\cap F(e_-)=0$. The isotropic-axis calculation forces
$U_-=S(e_+)$. Consequently $K_0$ has sign dimensions $(s-1,0)$,
while $V'$ has sign dimensions $(s-1,1)$. This proves that the first
switch introduces the other sign and preserves total dimension.
Since $s-1\ge1$, the next switch uses the preceding mixed-sign
case and lowers the total dimension by two. After the first switch
choose the primitive vector in the new integral generic kernel;
include one of its nonzero sign components in the second detection.
Thus the actual-class nonvanishing requirement holds at both steps.
The case $V=V_-$ is identical with signs interchanged.

At each step the new class is the actual derivative
$P_{I\cup\{\ell\}}$ of Proposition~\ref{ex:arith-finite-heegner},
with the same fixed multiplier $L$. Equation (4.2) gives finite
coordinate zero and singular coordinate $Jf(P_I)$, which is nonzero
generically. The construction retains the old Kummer conditions and
all earlier singular conditions, by the final assertion following
(H19). Lemma~\ref{ex:limits-derivative-switch} therefore makes this
class span the new generic kernel and preserves the valuation of
its class-functional determinant coordinate. In particular we do
not replace the actual class by the chosen primitive vector.
That vector is used only to select the prime. The paired functional
is the one transported from the same actual class. The determinant
comparison uses the common strict determinant factor; fractional
complements to $f(P_I)$ and $Jf(P_I)$ have unit exterior volume.
Consequently a divisible localization creates no extra valuation,
and all strict torsion factors remain present on both sides.

Each new generic kernel has rank one and has a primitive integral
vector, so its fiber dimension is at least one. A proposed drop
from dimension two to zero is therefore a contradiction; that
configuration cannot occur. Whenever the dimension exceeds one,
one or two switches give a strict drop by two. Thus finitely many
switches leave dimension one. After cancellation the terminal
complex has one free generator in each of degrees one and two;
its differential is zero because its generic kernel has rank one.
The actual class is integral. To apply
Proposition~\ref{ex:arith-paired-functional} after these switches,
we check its orthogonal comparison on the added local conditions.
In the integral derivative model of
Lemma~\ref{ex:limits-derivative-switch}, $D^0$ pairs with $D^2$ and
$S$ pairs perfectly with $F$, while the pairing on $S$ is zero
with its retained isotropy homotopy. Hence
$(D^0\oplus S[-1])^\perp=D^0\oplus S[-1]$ as the specified local
condition, with the retained integral comparison maps. Conjugate
transport preserves this condition as well. At a derivative prime
the comparison from the orthogonal condition to the conjugate
condition is therefore the identity under those integral maps.
It requires no inverse of a localization vector or of a local
Frobenius operator. Using this identity in the Selmer-cone comparison,
together with the same old-place comparisons as before, constructs
the integral orthogonal lift of the new transported class. Composing
with the retained integral global duality gives a closed integral
degree-two functional. Generically this is exactly the functional
of Proposition~\ref{ex:arith-paired-functional}: both use the inverse
of the same orthogonal comparison, which is unique in the derived
category. Thus its primitive integrality proof extends to every
finite list of these derivative conditions. At the
active ramified primes this uses the unit $-2$; at the moving prime
it uses $t$ except at $(t)$, where the central Frobenius determinants
are nonzero and hence units in $R$. Unused good primes retain their
unramified condition. No inverse at an unused relaxed prime is used.
The terminal coordinate is the product of the two integral
coefficients, up to a determinant-basis unit, so belongs to $R$.
Valuation preservation transports this conclusion through every
switch. This proves regularity at every $\mathfrak p\ne(2)$,
including both quadratic specializations. Finally the normal-domain
intersection formula for $\Lambda[1/2]$ gives
$u_{\rm prim}\in\Lambda[1/2]$.

The same proof applies at the unknown full-corank-one base. If its
actual generic class is zero, its coordinate is zero. Otherwise its
nonzero generic cycle and the central dimension bound give the
rank-one start used above. This requires no uniform outer valuation
bound at that base. Combining (5.5), for the known simple-zero
instances, with (4.6) and $\sum_{p\mid k}t_p\le2A$ proves (3.5).
No fixed-field class number is absorbed into the varying-field bound.


\end{proof}

\subsection{Actual augmentation and the fixed-field count}
\begin{lemma}[Augmentation and the fixed-field rank budget]\label{ex:limits-fixed-field-budget}
Fix the final field $K$ and the family of
Section~\ref{sec:rc}, in the branch where the odd empty
configuration is not a unit. There are no derivative labels.
At augmentation all graph primes are unramified.
The actual old-support subcomplex gives (5.6). With
$d_K=\dim_{\F_2}\Sel_2(E/K)$,
$t_K=\operatorname{length}E(K)[2^\infty]$ and
$h_K=\operatorname{length}H^2(C_K)_{\rm tors}$, it has the rank,
Smith-disk and torsion bounds (5.7). These are fixed-$K$ bounds,
independent of the binary dimension and graph size.

At every sign character, the inverse new local blocks and the
resulting old-rank Schur complex have no Laurent terms below
$-4M$ modulo $2^M$, uniformly in the number of new blocks.
This is a bound on sign evaluations, not on coefficients in the
group basis. The finite number $h_K$ does not assume finiteness
of the whole $\Sha$ at the base.
\end{lemma}
\begin{proof}\leavevmode
We now use the final $K$ in the remaining branch of Section~\ref{sec:rc},
where the intrinsic odd empty configuration is not a unit.
There are no derivative labels in this
last binary family. At augmentation every graph prime is unramified,
and the actual support triangle gives
\[
 C_{\rm old}\longrightarrow(C_b)_0
       \longrightarrow\bigoplus_{q\in Q}Q_q\longrightarrow C_{\rm old}[1],
                                                                    \tag{5.6}
\]
with $C_{\rm old}$ an actual subcomplex of an upper triangular model.
The graph-wide prescription gives scalar $1+t$ at each chosen
orientation, even though the quadratic character is trivial at
augmentation. The Euler-factor calculation in
Lemma~\ref{ex:arith-common-coordinate} therefore gives
\[
 \overline{l_{q,0}}=\overline e_{q,0}^{\,2}
                  =\frac{t^4}{(1+t)^2}.
\]
The two-place local determinant differs only by the unit $q^4$.
Its residual $t$-order is consequently four, uniformly in $q$ and
the number of graph primes.

Put $T=T_2E$ and let $C_K$ be the ordinary compact Kummer
Selmer complex over the fixed imaginary quadratic field $K$,
with the derived Kummer condition, including its degree-zero
torsion, at every finite place. Lengths in this paragraph are
lengths over $\mathbb Z_2$. Write
\[
 d_K=\dim_{\mathbb F_2}\operatorname{Sel}_2(E/K),\qquad
 t_K=\operatorname{length}E(K)[2^\infty],\qquad
 h_K=\operatorname{length}H^2(C_K)_{\rm tors}.
\]
The finite Selmer group defines $d_K$, and perfection of $C_K$
makes its cohomology finitely generated, so $h_K$ is finite.
This does not assert that the full two-primary Tate--Shafarevich
group is finite. The already fixed companion has
$L(E^{(k)},1)\ne0$, so the sum-and-difference isogeny and the
forward corank-zero theorem give full Selmer corank one for $E/K$.
This includes any divisible part of the Tate--Shafarevich group.

Here are the cohomology calculations underlying the budget.
At finite coefficients, the local degree-zero maps are
isomorphisms onto local invariants and the local degree-one maps
are the Kummer injections. The Selmer cone therefore has
$H^0=E(K)[2]$ and $H^1=\operatorname{Sel}_2(E/K)$.
The Weil pairing and the exact orthogonality of the Kummer
conditions give the finite self-duality with shift three
\cite[5.1.1--5.1.6 and 6.3.4]{nekovar}. Since $K$ has no
real places and $E[2]$ is rational of rank two, this gives
\[
 \dim_{\mathbb F_2}H^i(C_K\otimes^{\mathbf L}\mathbb F_2)
       =(2,d_K,d_K,2)\quad(i=0,1,2,3),                         \tag{5.7a}
\]
with zero cohomology in other degrees. In particular a minimal
finite free $\mathbb Z_2$ model of $C_K$ has total rank $2d_K+4$.

For the integral torsion calculation, $H^0(C_K)=T^{G_K}=0$:
a compatible system of rational points of unbounded two-power
order cannot lie in the finite group $E(K)_{\rm tors}$.
The exact sequence for $0\to T\to T\otimes\mathbb Q_2
\to E[2^\infty]\to0$, and the local Kummer conditions, identify
\[
 H^1(C_K)_{\rm tors}=E(K)[2^\infty].
\]
Indeed this finite group injects as torsion global cohomology,
all of its elements satisfy the Kummer conditions, and a torsion
class maps to zero with rational coefficients, so is in this
image. Integral Kummer self-duality
$C_K\simeq R\operatorname{Hom}_{\mathbb Z_2}
(C_K,\mathbb Z_2)[-3]$ now gives
\[
 H^3(C_K)\simeq
 \operatorname{Ext}^1_{\mathbb Z_2}(H^1(C_K),\mathbb Z_2),
 \qquad \operatorname{length}H^3(C_K)=t_K.                     \tag{5.7b}
\]
It also identifies the free ranks in degrees one and two.
They both equal one: universal coefficients with
$\mathbb Q_2/\mathbb Z_2$ give
\[
0\longrightarrow H^1(C_K)\otimes\mathbb Q_2/\mathbb Z_2
 \longrightarrow\operatorname{Sel}_{2^\infty}(E/K)
 \longrightarrow H^2(C_K)_{\rm tors}\longrightarrow0,
\]
and the middle group has full corank one. Thus the total integral
torsion length of $C_K$ is $h_K+2t_K$.

At $t=0$ the moving character is trivial. At either place over
the good split moving prime $r_i$, the unramified condition is
the ordinary Kummer condition. Replacing it by the full condition
adds the singular complex
\[
 Q_{r_i,w}=[\,\mathbb Z_2^2
       \xrightarrow{\Phi_{r_i}-r_i}\mathbb Z_2^2\,]
       \quad\text{in degrees one and two}.
\]
Here $\Phi_{r_i}$ is arithmetic Frobenius on $T$, with trace
$a_{r_i}$ and determinant $r_i$. In particular
\[
 \det(\Phi_{r_i}-r_i)
   =r_i(r_i+1-a_{r_i})
   =r_i\det(1-\Phi_{r_i}),
\]
so this determinant is nonzero and has valuation at most $c_E$.
The same bound holds at the conjugate place. The singular
complexes have zero degree-one cohomology and finite degree-two
cokernels of total length at most $2c_E$.

This comparison is retained before the inner limit. The field
$K$ and $C_K$ are fixed, while the two moving Frobenius matrices
and the comparison maps may vary with $i$ and with the graph.
Pass to their limit at every coefficient precision using the
same retained finite diagram. A matrix whose determinant has
valuation at most $c_E$ cannot acquire a zero determinant in
this limit: the property is determined modulo $2^{c_E+1}$.
Consequently, if
$D_c=C_{\rm old}\otimes^{\mathbf L}_\Lambda\mathbb Z_2$,
the resulting triangle is
\[
 C_K\longrightarrow D_c\longrightarrow Q_r
      \longrightarrow C_K[1],                               \tag{5.7c}
\]
where $Q_r$ is the sum of two rank-two singular disks and
$H^2(Q_r)=J_r$ has length at most $2c_E$. The notation $r$
in this limiting complex does not designate a fixed rational
prime supporting characters of unbounded order.

Reducing (5.7c) modulo two, the two singular disks contribute at
most eight to total cohomology dimension. Equation (5.7a) thus
gives a finite free model of $C_{\rm old}$ over $\Lambda$ of
total rank at most $2d_K+12$, by cancelling unit entries at its
maximal ideal $(2,t)$. The integral long exact sequence of
(5.7c) gives
\[
 H^1(D_c)\simeq H^1(C_K),\qquad
 0\to H^2(C_K)\to H^2(D_c)\to J_r
       \to H^3(C_K)\to H^3(D_c)\to0.                         \tag{5.7d}
\]
In particular $H^3(D_c)$ is a quotient of $H^3(C_K)$.
The image of the torsion of $H^2(D_c)$ in $J_r$ has length at
most $\operatorname{length}J_r$, and its kernel is
$H^2(C_K)_{\rm tors}$. Hence, defining the symbol used below by
\[
             h_1:=\sum_i\operatorname{length}H^i(D_c)_{\rm tors},
\]
we have $h_1\le h_K+2c_E+2t_K$.

Finally, a finite free complex over $\mathbb Z_2$ decomposes
into free cohomology generators with zero differential and
Smith disks $[\mathbb Z_2\xrightarrow{2^e u}\mathbb Z_2]$,
$u\in\mathbb Z_2^\times$, $e\ge0$. Its only free cohomology
generators here are one in degree one and one in degree two.
If $m$ counts all its Smith disks, including those with $e=0$,
then $\operatorname{rank}_{\rm total}C_{\rm old}=2+2m$.
Every disk contributes $e$ to $h_1$, so the sum of the valuations
of all its nonzero Smith entries equals $h_1$. We have proved
\[
 \operatorname{rank}_{\rm total}C_{\rm old}\le2d_K+12,\qquad
 m\le d_K+5,\qquad h_1\le h_K+2c_E+2t_K.                     \tag{5.7}
\]
The model and its matrices may vary with the moving characters,
but these bounds depend only on the final fixed $K$ and fixed
old data. They do not depend on the binary dimension or graph
size. The varying-field estimate (5.1) is not used with $j=b$,
and $h_K$ is not treated as a varying-field constant.

The lift is made over $R=\Lambda[G]$ before passing to
$\mathscr O[G]$. Let $C$ be the retained finite free family model.
Its augmentation $C_0$ is the actual scalar cochain model, by the
retained coefficient-change maps. After augmentation the representation
is unramified at $Q$, so the support-change triangle (5.6) gives a
finite free model $D$ of $C_0$ whose underlying graded module is
$C_{\rm old}\oplus\bigoplus_q Q_q$ and whose differential is upper
triangular. This follows by representing the connecting map between
these bounded free complexes by a chain map and forming its cone.

Cancel unit disks in $C_0$. The same pivots and elementary operations
lift to $C$ over $R$, since the augmentation ideal is contained in
the Jacobson radical of $R$. Let $C'$ be the resulting family model,
so that $C'_0$ is minimal. Decompose $D$ over $\Lambda$ into a minimal
complex $D_{\min}$ and unit disks $E$. The actual comparison gives a
homotopy equivalence $C'_0\simeq D_{\min}$. A homotopy equivalence
between minimal bounded free complexes over a local ring is degreewise
invertible, by reduction modulo its maximal ideal and Nakayama.
Add to $C'$ the same unit disks $E_R$ and lift the resulting graded
isomorphism $C'_0\oplus E\simeq D$ to $R$. Its matrices are invertible
because their augmented determinants are units. Conjugating the
family differential by these lifted matrices makes its augmentation
exactly the differential of $D$. All matrices and their inverses at
this step lie in $\Lambda[G]$; thus their sign evaluations are power
series. The added disks are precisely those needed for $D$, so this
procedure does not enlarge its stated old-support rank.
The block containing the local quotients has
\[
 B_0=A=\operatorname{diag}(A_q),\qquad T_0=0
\]
in a differential $\left(\begin{smallmatrix}P&Q\\T&B\end{smallmatrix}\right)$.
Over $\mathscr O[G]$, cancel $B$. The surviving Schur differential
$P-QB^{-1}T$ augments to $P_0$.

Write $\det A_q=t^{d_q}v_q+2w_q$, $d_q\le4$, with $v_q$ a power-series
unit, and put
$C_q=\operatorname{adj}(A_q)/(t^{d_q}v_q)$.
For a character $x$, $B_xC=1+2F$ with both $C$ and $F$ having no
Laurent terms below degree $-4$. Therefore
\[
             B_x^{-1}=C\sum_{j=0}^{M-1}(-2F)^j\pmod {2^M}        \tag{5.8}
\]
has pole at most $4M$, independent of the number of blocks.
For completeness the adjacent-differential assertion is an exact
block identity. In the two degrees containing the new local block,
write the differential as
$d=\left(\begin{smallmatrix}P&Q\\T&B\end{smallmatrix}\right)$,
with $B$ invertible. Then
\[
 \begin{pmatrix}1&-QB^{-1}\\0&1\end{pmatrix}
 d\begin{pmatrix}1&0\\-B^{-1}T&1\end{pmatrix}
  =\begin{pmatrix}P-QB^{-1}T&0\\0&B\end{pmatrix}.
\]
If the preceding differential has components $(X,Y)^{\mathsf T}$,
then $d^2=0$ gives $TX+BY=0$. In the new source coordinates its
components are $(X,B^{-1}TX+Y)^{\mathsf T}=(X,0)^{\mathsf T}$.
If the following differential has components $(U,V)$, then
$UQ+VB=0$, and in the new target coordinates they are
$(U,UQB^{-1}+V)=(U,0)$. Thus the surviving entries of both adjacent
differentials are literally unchanged; no second inverse is
introduced there.

Here $A_q$ is the direct sum of the two rank-two singular
differentials at the conjugate graph places. Modulo two their
scalar entries are $t$ and $-t/(1+t)$, since their Frobenius
scalars are $1+t$ and $(1+t)^{-1}$ and $E[2]$ is trivial.
Consequently
\[
                  \det\overline A_q=\frac{t^4}{(1+t)^2}.
\]
This independently gives the residual order four in (5.6).
The factor $q^4$ in the Euler-factor convention is a unit and
does not affect it. Moreover every sign evaluation of a matrix
over $\Lambda[G]$ has the same reduction modulo two as its
augmentation, because each sign is congruent to one modulo two.
Thus $B_x=A+2V_x$ with $V_x$ a power-series matrix. For the
block-diagonal matrix $C$ in (5.8), the identities
\[
 A_q C_q=1+\frac{2w_q}{t^{d_q}v_q},\qquad
 B_xC=1+2F_x
\]
show that every entry of $C$ and $F_x$ has Laurent order at
least $-4$. Formula (5.8) therefore has Laurent order at least
$-4M$, because its $j$th summand has order at least $-4(j+1)$.
The estimate counts the largest pole of an entry, not the number
of summands or blocks. The remaining Schur differential has the
same bound since $P_x,Q_x,T_x$ are power-series matrices, and
the adjacent surviving entries have no poles at all. These
arguments take place after localization where $t$ is invertible;
they do not invert $t$ in a finite Artin coefficient ring.
\end{proof}

\subsection{Determinant clearing and the missing-address contradiction}
\begin{proposition}[A common denominator for the corrected coordinate]\label{ex:limits-fixed-field-clearing}
Under Lemma~\ref{ex:limits-fixed-field-budget}, use the paired
functional exponent $a$ of
Proposition~\ref{ex:arith-paired-functional} and the universal
correction $\eta$ of Lemma~\ref{ex:arith-common-coordinate}.
There are common $\Delta,Q\in\mathscr O[G]$ such that, wherever
$\Delta_x\ne0$,
$\widetilde u_x=2^{-a}\eta_xQ_x/\Delta_x^3$.
Putting $f=2^a\Delta^4$ and $d=\eta\Delta Q$ gives
$f_x\widetilde u_x=d_x$ at every character, including
$\Delta_x=0$, and $f_0\in\Lambda$. Moreover
\[
 v_2f_0(0)\leq a+4(h_K+2c_E+2t_K)=c_{\det},\qquad
 L_T(M)=16(d_K+5)M
\]
bounds the negative Laurent support of every $f_x$ modulo $2^M$,
uniformly in the sign character $x$. The elements $f,d$ themselves
belong to $\mathscr O[G]$ in the original group basis.
These bounds concern the same corrected coordinate throughout;
no pole bound for $d_x$ or for the group coefficients of $f$ is required.
\end{proposition}
\begin{proof}\leavevmode
Put $S=\mathscr O[G]$ and first perform the simultaneous new-local
cancellation of Lemma~\ref{ex:limits-fixed-field-budget}.  Its
determinant comparison factor $\kappa$ is an $S$-unit: the canceled
block is invertible over $S$, so its determinant and inverse are
integral over $S$.  Retain this factor, with the determinant-line
sign prescribed by the original basis.  Transport the actual cycle
$z$ and the closed functional $\mu=2^a\lambda_z$ to the remaining
complex $D$.  Both transported cochains are integral over $S$.
The augmentation of $D$ is the old complex over $\Lambda$, and
each sign-evaluated differential entry has pole at most $4M$
modulo $2^M$.  No bound on the entries of the transported $z$ or
$\mu$ is needed below.

Write $D_c=D_0\otimes_\Lambda\mathbf Z_2$.  Its integral Smith
decomposition gives graded decompositions
\[
 D^i=B_i\oplus H_i\oplus E_i,
\]
lifted by constant $\mathbf Z_2$ basis changes, where centrally the
only nonzero differentials are the Smith disks $E_i\to B_{i+1}$.
Here $H_1,H_2$ have rank one and all other $H_i$ vanish.  Let
$m_i=\operatorname{rank}E_i$, $m=\sum_i m_i$, and let $A_i$ be the
actual differential submatrix from $E_i$ to $B_{i+1}$.  Set
\[
 \delta_i=\det A_i,\qquad \Delta=\prod_i\delta_i.
\]
An empty determinant is one.  The entries in this definition are
the original entries of $D$, before any of these old-disk
cancellations.  Thus $\Delta$ has degree $m$ in those entries,
$\Delta_0\in\Lambda$, and
\[
 v_2\Delta_0(0)=\sum_{\text{central Smith disks}}v_2(\text{disk})
       \le h_1\le h_K+2c_E+2t_K.
\]

Here is the full denominator calculation.  Work temporarily in
$S[1/\Delta]$.  Group all sources and targets for a given degree,
and cancel the whole block $A_i$ once.  For that differential the
two elementary matrices are
\[
 \begin{pmatrix}1&0\\-CA_i^{-1}&1\end{pmatrix}
 \begin{pmatrix}A_i&B\\C&D'\end{pmatrix}
 \begin{pmatrix}1&-A_i^{-1}B\\0&1\end{pmatrix}
 =\begin{pmatrix}A_i&0\\0&D'-CA_i^{-1}B\end{pmatrix}.
\]
Both elementary matrices have determinant one.  In the adjacent
differentials, the equation $d^2=0$ makes the components into and
out of this canceled pair zero after these changes; the other
entries of those adjacent differentials are unchanged.  Hence a
cancellation does not modify any other surviving $A_j$.  Every
inverse used is $\operatorname{adj}(A_i)/\delta_i$, and each grouped
block is used only once.

More explicitly, in degree $k$ the projection associated to the
incoming cancellation is
\[
 (b,v)\longmapsto v-C A_{k-1}^{-1}b,
\]
and the inclusion associated to the outgoing cancellation is
\[
 v\longmapsto(-A_k^{-1}Bv,v).
\]
The other map in each of these two degrees is the corresponding
coordinate inclusion or projection.  Compose these maps in the
chosen cancellation order.  At each step each entry is linear in
the one inverse $A_i^{-1}$ at that step, and no $A_i$ is encountered
again.  Products of denominators therefore divide $\prod_i\delta_i$,
not a power depending on the number of individual Smith disks.
In particular, if $p:D\to H$ and $\iota:H\to D$ are the final
projection and inclusion, the two scalars
\[
 \alpha=(pz)^1\in H_1,
 \qquad \beta=(\mu\iota)^2\in H_2^*
\]
satisfy $Z:=\Delta\alpha\in S$ and $W:=\Delta\beta\in S$.
The final complex $H$ has just the map $H_1\to H_2$.

For the inverse determinant line, a disk in degrees $i,i+1$
contributes $\delta_i^{(-1)^i}$ to the coordinate of a cohomology
tensor.  This can be checked on the based disk
$e\mapsto\delta_i b$: replace the boundary basis $b$ by
$\delta_i b$ in $\bigotimes_k\det(D^k)^{(-1)^{k+1}}$.
Thus the remaining comparison factor is
\[
 T=\epsilon\prod_i\delta_i^{(-1)^i},\qquad
 P:=\Delta T=\epsilon\prod_{i\text{ even}}\delta_i^2\in S,
\]
where $\epsilon$ is the integral unit from the fixed basis ordering
and the constant Smith basis changes.  Including the previously
retained new-local factor, put
\[
 Q=\kappa ZWP\in S.
\]
This definition takes place over $S$, before sign evaluation.  It
uses no interpolation and no division by $|G|$.

If $\Delta_x\ne0$ and the constituent has the rank-one cohomology
pattern, its coordinate in the original determinant basis is exactly
\[
 \widetilde u_x=\eta_x\kappa_x\alpha_x(2^{-a}\beta_x)T_x
              =2^{-a}\eta_xQ_x/\Delta_x^3.
\]
This is the same $z$, functional, basis and multiplier $\eta$ as in
the statement.  If the constituent is generically acyclic, the
remaining map between the two lines is nonzero.  Its cycle
$\alpha_x$ is then zero over the fraction field, so $Q_x=0$ and the
formula agrees with the defined coordinate zero.  These are the
only generic patterns in the present family, by central
specialization and the central corank-one bound.

Now define $f=2^a\Delta^4$ and $d=\eta\Delta Q$ in $S$.
The preceding formula proves $f_x\widetilde u_x=d_x$ whenever
$\Delta_x\ne0$.  When $\Delta_x=0$, both $f_x$ and $d_x$ are zero;
the already defined coordinate $\widetilde u_x$ is finite in its
fraction field, and the same equality follows directly.  No rational
expression with a zero denominator is specialized.  In particular
$f_0\in\Lambda$ and
\[
 c=v_2f_0(0)\le c_{\det}:=a+4(h_K+2c_E+2t_K),\qquad
 L_T(M)=16(d_K+5)M.                                      \tag{5.9}
\]
Indeed $f$ is a polynomial of degree $4m$ in the differential
entries of $D$, with constant integral coefficients.  Each such
entry, at a sign character, has pole at most $4M$ modulo $2^M$.
A product of $4m$ entries has pole at most $16mM$, and sums cannot
increase this bound.  Finally $m\le d_K+5$.

Only the last pole estimate is characterwise.  The actual elements
$f,d$ belong to the original group ring $S$.  For $i<-L_T(M)$ all
sign values of $[t^i]f$ are divisible by $2^M$; the group-ring
coefficient itself need not be divisible by $2^M$.  No Laurent
support bound for $d$ or for the group coefficients of $f$ has been
used.  The constants $d_K,t_K,h_K$ are fixed only after the final
$K$ is fixed; $a$ and $c_E$ precede that choice.  None depends on
the graph or binary dimension.  A varying-$K$ estimate cannot be
deduced by treating $h_K$ as an absolute constant.
\end{proof}

\begin{proposition}[Finite coefficient transfer and the final contradiction]\label{ex:limits-coefficient-transfer}
Use the fixed field, common family, regular coordinates and
integral identity of Proposition~\ref{ex:limits-fixed-field-clearing}.
Let $B$ be the fixed upper bound on
$v_2\widetilde u_x(0)$ at every nonzero address supplied by
(3.9), (3.4), (4.6) and Lemma~\ref{ex:arith-common-coordinate}.
Choose $M,b$ as in (5.10), before the finite graph and its actual
theta stage are realized. For some nonzero address $x$,
\[
 v_2\bigl(\widetilde u_x(0)-\widetilde u_0(0)\bigr)
 \geq M-2c>B,\qquad c=v_2f_0(0)\leq c_{\det}.
\]
The convolution includes every negative Laurent coefficient
of $f_x$. If $L'(E/K,1)=0$, Lemma~\ref{ex:arith-base-zero}
makes $\widetilde u_0(0)=0$, contradicting the upper bound.
Since the fixed complementary value $L(E^{(k)},1)$ is nonzero,
this proves $\operatorname{ord}_{s=1}L(E,s)=1$.
\end{proposition}
\begin{proof}\leavevmode
For a coefficient of any element of $\mathscr O[G]$, its sign-value
function on $\F_2^b$ expands as
\[
 A(x)=a_\varnothing+
        \sum_{I\ne\varnothing}2^{|I|}a_I\prod_{i\in I}x_i .
\]
Its low bit is constant and its bit of index $j\ge1$ has Boolean
degree at most $2^{j-1}$. One can see this by extracting the
$Y^{2^j}$ coefficient of $(1+Y)^{A(x)}$ modulo two:
the first possible contribution of $x_I$ has degree at least
$2^{|I|}$ in $Y$, and $|I|\le2^{|I|-1}$.
Thus $q$ coefficient equalities modulo $2^M$ have sum of Boolean
degrees at most $q(2^{M-1}-1)$.

If that sum is less than $b$, multiply the complementary indicators
of the equations. Its degree is less than $b$, so its sum over the
Boolean cube is zero: every monomial misses a variable.
The support of this indicator is the simultaneous solution set.
It has even cardinality and contains zero, hence another address.

Apply this to the coefficients of $f$ from $-L_T(M)$ through zero
and the single constant coefficient of $d$. It suffices to choose
\[
 M>\max(c_{\det},B+2c_{\det}),\qquad
 b>(L_T(M)+2)(2^M-1),                                           \tag{5.10}
\]
where $B$ bounds the nonzero-address central valuations.
These choices precede the graph and its theta precision.
At the resulting nonzero address, negative coefficients of $f_x$
vanish modulo $2^M$, and its constant coefficient agrees with
$f_0(0)$ and has valuation $c$.
Positive coefficients may have smaller valuation. Nevertheless
$v_{\mathscr O}(f_x)\le c$, so integrality of
$d_x=f_x\widetilde u_x$ gives
$v_{\mathscr O}(\widetilde u_x)\ge-c$.
This estimate comes from the displayed product at the selected
nonzero address; it does not apply (5.5) to the unknown base.
Both $\widetilde u_x$ and $\widetilde u_0$ lie in $\Lambda[1/2]$
by the height-one argument above, so neither has negative powers
of $t$. In particular, since $f_0\in\Lambda$,
$[t^0]d_0=f_0(0)\widetilde u_0(0)$.
Put $F_0=[t^0]f_0$ and $F_x=[t^0]f_x$. Since $d_0$ is integral
and $[t^0]d_0=F_0\widetilde u_0(0)$, we also have
$v_2\widetilde u_0(0)\ge-c$, without using (5.5) at the base.
The convergent convolution has the form
\[
 [t^0]d_x=F_x\widetilde u_x(0)+E_x,\qquad
 E_x=\sum_{i<0}[t^i]f_x[t^{-i}]\widetilde u_x,
 \qquad v_2(E_x)\ge M-c.
\]
The untested coefficients below the Laurent cutoff have the same
divisibility by the pole bound, so this sum includes every negative
coefficient. Subtracting the two constant terms gives
\[
 F_x(\widetilde u_x(0)-\widetilde u_0(0))
 =([t^0]d_x-[t^0]d_0)-(F_x-F_0)\widetilde u_0(0)-E_x.
\]
Every term on the right has valuation at least $M-c$ and
$v_2(F_x)=c$. Consequently
\[
 v_2(\widetilde u_x(0)-\widetilde u_0(0))\ge M-2c>B.\tag{5.11}
\]

By (3.9), (3.4), (4.6) and the unit comparison (4.7), every nonzero address has
$v_2\widetilde u_x(0)\le B$ for one $B$ fixed before $b$; the finite
$\Sha$ terms have negative sign and may be discarded for this upper
bound. If the base derivative vanished, Section~\ref{sec:arithmetic}
would give $\widetilde u_0(0)=0$, contradicting (5.11).
Hence $L'(E/K,1)\ne0$. The complementary value $L(E^{(k)},1)$ is
nonzero, and
$L(E/K,s)=L(E,s)L(E^{(k)},s)$, so $E$ has analytic rank one.
This proves (RC), and hence Theorem~\ref{hyp:RC}, using the classical
applications and finite constructions specified above.
\end{proof}

% End 03b-limits.tex


% Begin 04-computation-value.tex
\section{Exact precision computation and polynomial compilation}\label{sec:compiler}
Encode a signed integer by
\[
 \operatorname{enc}(n)=\begin{cases}2n,&n\ge0,\\-2n-1,&n<0.\end{cases}
\]
The five natural inputs decode in the fixed order $(n,\kappa,a,l,j)$.
The decision predicate rejects unless
$n\notin\Cset$, $\kappa\ne0$, $a,l\ge1$, and $j=na$.
For nonzero multiples of $P_0$, $z=-x/y$ is well-defined. Set
\[
 z([l]P_0)=u_l+v_ls,\qquad\epsilon_l=|u_l|+|v_l|,
 \qquad r(a,j)=\tfrac12\Tr_{K_0/\Q}\frac{z([j]P_0)}{z([a]P_0)}.
\]
With the fixed constant $A_{\rm ht}^{(0)}=H_6(S_{\rm ht})$, its final
tests are
\[
 |r(a,j)-n|<\epsilon_l,\qquad
 4\epsilon_l<4^{-A_{\rm ht}^{(0)}|\kappa F(n)|^{24}}.
\]
These are strict tests on exactly computable rationals and integers.

The reference program S21 uses exact quadratic-field arithmetic and capped
integer operations. The ideal register intermediate representation S22
computes the same predicate on every five-natural tuple. Its independent
all-input argument covers normalized signed arithmetic, origin and
zero-ordinate cases, equal-abscissa addition, and every early rejection.
Nine exact polynomial identities verify the nonexceptional projective
formulas and preservation of $Y^2=X^3+AXZ^4$. Scalar loops decrease by
halving; the precision loop divides by four; factorial iteration is bounded.
This is a proof about the ideal intermediate representation, rather than
an assertion that a budgeted simulator always finishes.

For example, if $\epsilon=U/V>0$ with natural $U,V$ and integer $B\ge0$,
\[
 4\epsilon<4^{-B}\quad\Longleftrightarrow\quad
 4^B\le\left\lfloor\frac{V-1}{4U}\right\rfloor.
\]
This permits a finite cap. For positive $L$, $\sigma(x)=\min(x,L+1)$
commutes with saturated addition and multiplication. Powers and factorials
can be computed with this saturation. Capping an exponent is safe for
base at least two because an exponent above $L$ already makes the value
exceed $L$. Base zero, base one and exponent zero are separate cases.
This proves equality with the fixed-height expression without requiring
its full evaluation when the available precision is too small.

\subsection{The natural polynomial and the power interface}
The saved precision polynomial has $1\,305\,811$ natural variables,
the first five being inputs. Its SHA256 is
\begin{center}\small\ttfamily
f3e163f5c60b972c98b48647274f956dd7d78428e27ddc0543574a5f18120b8c.
\end{center}
Source S23 specializes a fixed universal natural halting polynomial at
the code of a wrapper. The wrapper inverts four Cantor pairs and then
runs the unchanged precision machine. Its pairing equations are
\[
 2t_i=(c_i+t_{i+1})(c_i+t_{i+1}+1)+2t_{i+1}\quad(0\le i<4),
 \qquad t_4=c_4.
\]
Their triangular intervals prove unique natural unpairing. The new
polynomial sums the square of the old halting polynomial and the four
squared differences above.

The universal compiler packs a finite run of a fixed $m$-state
increment/decrement/branch program into base-$B$ integers. Acceptance
means reaching the unique final halt with all registers zero; the fixed
interpreter includes a finite cleanup stage. Choose
\[
 c>0,\quad K=2^c,\quad B=4K>m,\quad T=B^q,\quad
 F_q=1+B+\cdots+B^q,\quad D=(K-1)F_q,
\]
and require every input to be below $K$. For natural integers, say
$X$ is masked by $Y$ when every set binary bit of $X$ occurs in $Y$.
Require each register history $R_j$ to be masked by $D$ and satisfy
$R_j<T$. Each state history $S_i$ is masked by $F_q$, and
\[
 \sum_{i=0}^{m-1}S_i=F_q,\qquad S_{\rm halt}=T.
\]
Thus register digits are below $K$, terminal register digits are zero,
and state digits are zero or one. Since $B>m$, the state sum has no
carries and selects exactly one state at each time.

The nonzero mask $N_j$ is determined by
\[
 K N_j=(R_j+D)\mathbin{\mathrm{AND}}(K F_q).
\]
Indeed each digit $r+K-1<2K<B$ has its $K$-bit set exactly when
$r>0$. For a decrement/branch instruction at state $i$ testing
register $j$, impose
\[
 V_i=S_i\mathbin{\mathrm{AND}}N_j,\qquad S_i=V_i+W_i.
\]
These are the nonzero and zero branch histories. Form the packed
increment indicator $A_j$ from states incrementing $j$, the decrement
indicator $C_j$ from the appropriate $V_i$, and each incoming-state
indicator $I_i$ from the states or branches directed to $i$.
They have zero-one digits, since there is just one active instruction
and branch at each time. The transition equations are
\[
 R_j+BC_j=\operatorname{initial}_j+BR_j+BA_j,\qquad
 S_i=\mathbf1_{i=0}+BI_i.
\]
Every register coefficient on either side is at most $K<B$;
uniqueness of base-$B$ digits gives the initial conditions and each
exact transition. The possible extra final slot vanishes because
halt has no outgoing instruction and the terminal registers are zero.
For $q=0$ the same equations require the initial state already to be
halt and the inputs to vanish. Conversely, any finite accepting run
fits after choosing $c$ sufficiently large, and supplies all these
masks and equations.

\subsection{The literal alpha and power formulas}
The alpha system below is \texttt{Exp\_Matrices.alpha\_equations},
with equivalence \texttt{alpha\_equivalence}; the power system is
\texttt{Exponential\_Relation.exp\_repr} and its polynomial expansion
\texttt{exp\_dioph} in the pinned AFP source \cite{afp}.
Their variables range over natural numbers. The alpha recurrence itself
is integer-valued; the bounds below justify its natural-valued uses.
Define $U_n=\alpha_b(n)$ by
\[
 U_0=0,\qquad U_1=1,\qquad U_{n+2}=bU_{n+1}-U_n\quad(n\ge0).
\]
For $b\ge3$ this recurrence gives
$0\le U_n<U_{n+1}$ and $U_n\ge n$. For $n\ge1$, the matrix
\[
 M=\begin{pmatrix}b&-1\\1&0\end{pmatrix},\qquad
 M^n=\begin{pmatrix}U_{n+1}&-U_n\\U_n&-U_{n-1}\end{pmatrix}
\]
has determinant one. Thus adjacent terms are coprime and
$U_n^2-bU_nU_{n-1}+U_{n-1}^2=1$.
Parameter congruence follows directly from the recurrence:
$b\equiv b'\pmod d$ implies
$\alpha_b(n)\equiv\alpha_{b'}(n)\pmod d$.
In particular $\alpha_2(n)=n$. This is parameter congruence,
not merely congruence between indices.

Every natural solution of $x^2+y^2=1+bxy$, $b>2$, is an unordered
adjacent pair. To prove this, order it as $X>Y$. If $Y=0$, $X=1$.
Otherwise $Z=bY-X$ satisfies $XZ=Y^2-1$, hence $0\le Z<Y$,
and $(Y,Z)$ satisfies the same equation. Descent reaches $(1,0)$;
reversing the steps generates exactly the adjacent pairs.
No ordering assumption on the original pair is needed.

For $k>0$, reduction of $M^k$ modulo $U_k$ is the invertible scalar
$U_{k+1}I$. Dividing an index by $k$ and using strict growth proves
$U_k\mid U_m$ if and only if $k\mid m$. More precisely, for $m>0$,
\[
                    U_k^2\mid U_m\quad\Longleftrightarrow\quad
                         kU_k\mid m.                           \tag{C.1}
\]
For $k=0$ both sides are false. For $k>0$ write $m=kq$ and expand
$M^{kq}=(U_kM-U_{k-1}I)^q$ modulo $U_k^2$. Its lower-left entry is
$qU_k(-U_{k-1})^{q-1}$, so coprimality proves (C.1), including $U_k=1$.

For $m>0$ put $v=U_{m+1}-U_{m-1}=\operatorname{tr}(M^m)$.
Cayley--Hamilton gives $M^{2m}\equiv-I\pmod v$.
Extend the recurrence by $U_{-j}=-U_j$. Every $n\ge0$ has
$n=2lm+\varepsilon j$ with $l\ge0$, $0\le j\le m$ and
$\varepsilon\in\{1,-1\}$; then
$U_n\equiv(-1)^l\varepsilon U_j\pmod v$.
The value sign need not equal the index sign. If $b\ge4$,
$v=bU_m-2U_{m-1}>2U_m$.

The actual alpha formula has natural witnesses $r,s,t,u,v,w,x,y$
and the following fifteen conditions:
\[
\begin{gathered}
 b>3,\quad u^2+t^2=1+but,\quad s^2+r^2=1+bsr,\quad r<s,\quad
 u^2\mid s,\quad v+2r=bs,\\
 w\equiv b\pmod v,\quad w\equiv2\pmod u,\quad w>2,\quad
 x^2+y^2=1+wxy,\\
 2X<u,\quad2X<v,\quad x\equiv X\pmod v,\quad
 2N<u,\quad x\equiv N\pmod u.
\end{gathered}                                                   \tag{C.2}
\]
A strict inequality $a<b$ is $b=a+1+d$; divisibility is $b=ad$;
congruence is $A+Mu=B+Mv$, with fresh natural witnesses.
At zero modulus the last equation says equality, so there is no
exception hidden in this replacement.

For soundness, the unordered Pell result gives $u=U_k$, $k>0$,
and $x=\alpha_w(n)$. The ordered pair gives $s=U_m,r=U_{m-1}$,
$m>0$. By (C.1), $u\mid m$. Parameter congruence gives
$x\equiv U_n\pmod v$ and $x\equiv n\pmod u$.
Reflect $n$ to $j\in[0,m]$. Since $2X<v$ and $2U_j<v$,
the congruence $X\equiv\pm U_j\pmod v$ forces $X=U_j$,
with the negative case forcing both zero.
Then $j\le X$, $2j<u$ and $2N<u$ force $N=j$ from
$N\equiv\pm j\pmod u$. Thus $X=U_N$.

Conversely choose an odd $u=U_k>2U_N$. Such terms are unbounded,
since adjacent terms are coprime and the sequence grows.
Set $t=U_{k+1}$, $m=ku$, $s=U_m$, $r=U_{m-1}$ and $v=bs-2r$.
Equation (C.1) gives $u^2\mid s$, and the required strict bounds hold.
The moduli $u,v$ are coprime: a common divisor is odd, divides $s$
and $2r$, hence divides both adjacent terms. CRT gives
$w\equiv2\pmod u$, $w\equiv b\pmod v$ with $w>2$.
Take $x=\alpha_w(N)$, $y=\alpha_w(N+1)$.
Parameter congruence verifies every remaining condition.
This includes $N=X=0$.

The literal power formula separates $Q=0,R=0,P=1$ and
$Q=0,R>0,P=0$. For $Q\ge1$ it imposes
\[
 a_0=\alpha_{Q+4}(R+1),\quad b=a_0+Q^2+2,\quad
 m+Q^2+1=bQ,\quad P<m,
\]
\[
 a_1=\alpha_b(R),\quad a_2=\alpha_b(R+1),\quad
 P+ba_1\equiv Qa_1+a_2\pmod m.                                  \tag{C.3}
\]
Induction gives $a_0\ge(Q+3)^R$, so $m>Q^R$ and $m>0$.
Let $E_R=Q U_R-bU_R+U_{R+1}$. The recurrence and
$bQ\equiv Q^2+1\pmod m$ give $E_0=1$ and
$E_{R+1}\equiv QE_R\pmod m$. Thus (C.3) says $P\equiv Q^R\pmod m$.
Both values lie in $[0,m)$, forcing equality. Conversely equality
supplies every witness. This includes $R=0$, $Q=1$ and $0^0=1$.

These formulas make the binomial and bit operations polynomial.
For $B=2^{n+1}$, the equation
\[
 (1+B)^n=B^k(Bq+C)+r,\quad C<B,\quad r<B^k
\]
selects exactly $\binom nk$, since all its base-$B$ coefficients
are below $B$, including zero for $k>n$.
The identity
$(1+t)^n=\prod_{i:\,\text{bit}_i(n)=1}(1+t^{2^i})$ over $\F_2$
makes odd binomial coefficients equivalent to bit inclusion.
Bit disjointness of $a,b$ is equivalent to inclusion of $a$ in $a+b$:
at the least shared bit there is no incoming carry and that bit
is cleared. Finally $x=z+a,y=z+b$ with $z,a,b$ pairwise bit-disjoint
says exactly $z=x\mathbin{\rm AND}y$.

\subsection{Interpreter, lowering and five-input specialization}
The universal program numbering is $e+1=2^w(2D+1)$.
At $w=0$ it halts. Otherwise the four base-$2^w$ digits at positions
$4p,\ldots,4p+3$ give the opcode, register and two targets at state $p$.
Opcode one increments, opcode two decrements or takes its zero branch;
every other opcode halts. Missing digits are zero, so decoding is total.

The interpreter stores a finite tape as
$0$ or $\operatorname{cons}(h,t)=2^h(2t+1)$.
Repeated halving decodes a nonempty head and tail; the tail is smaller
than the code. To update a register, save the finite reversed prefix,
replace the entry and restore the prefix. Absent entries are zero.
The reserved jump register remains zero. Copies restore their source;
aliased addition snapshots its operand before overwriting it.
Division maintains $x=qy+r$ and subtracts only when $r\ge y$;
the interpreter's divisors are at least two.
These invariants prove that one interpreter loop executes exactly one
decoded instruction, with finite work. A simulated halt enters cleanup
and clears all registers. An infinite simulation never enters cleanup.
Thus its all-zero halt is equivalent to virtual halting, without
requiring the virtual final tape itself to be zero.

The precision IR uses normalized signed pairs for integer arithmetic
and pairs over the basis $1,\sqrt2$ for field arithmetic.
Norm zero implies the field element is zero: after clearing
denominators, a primitive solution of $a^2=2b^2$ would have both
entries even. Every division is guarded by that exact test.
The Jacobian point formulas satisfy the curve equation and the affine
chord or tangent formulas after multiplying denominators.
Equal abscissas split into doubling and opposite-point branches;
zero ordinate doubles to the identity. The fixed projective point
is the same affine $P_0$. These are polynomial identities, independent
of the generator theorem.

The scalar loop halves its remaining count. The precision loop divides
its remainder by four. The capped arithmetic computes exactly
$\min(x,L+1)$; zero bases, unit bases and zero exponents are separate.
For bases at least two an exponent exceeding $L$ already exceeds the
cap. Thus the ideal IR terminates and agrees with the stated rational
predicate on all five natural inputs, including every early rejection.

Each primitive lowering uses fresh auxiliary registers, snapshots
overlapping sources, and preserves the live-register map at the end
of each finite block. Saturating decrement gives natural subtraction,
and repeated addition gives multiplication. All variable divisions
are guarded. Induction on completed loop iterations proves the same
claim for loops. Acceptance enters cleanup; rejection enters a
permanent self-loop. Hence the primitive machine halts exactly on
acceptance, although the ideal decision program is total.

For $\operatorname{pair}(a,b)=(a+b)(a+b+1)/2+b$, unpairing maintains
$\text{remainder}=\text{input}-s(s+1)/2$ and subtracts $s+1$
while $s<\text{remainder}$. At exit
$0\le\text{remainder}\le s$, so $b=\text{remainder}$, $a=s-b$.
Four repetitions put the five codes into registers zero through four.
Fresh scratch lies above the old register range and is cleared;
old work registers remain zero. Relocating every target by the prefix
length therefore preserves the precision machine.

The polynomial specialization fixes the universal program input to
this wrapper code, maps the old packed input to variable five, and
maps every old variable $i\ge2$ and private-block start by $i\mapsto i+7$.
The four pairing variables occupy five through eight, disjoint from
the five inputs. The appended four squared equations force their
unique Cantor values. Thus, for the exact saved polynomial $C$,
\[
 \forall\mathbf c\in\N^5,\qquad
 \exists\mathbf z\in\N^{1305806}\ C(\mathbf c,\mathbf z)=0
 \quad\Longleftrightarrow\quad
        \operatorname{accepts\_naturals}(\mathbf c)=\mathrm{True}. \tag{C.4}
\]
The natural compiler proof has no height, restricted-converse or
rational-model premise. Its use of the height expression concerns
only computation of that fixed integer expression.

\subsection{Translation to point codes}
The typing order is independent of the natural predicate. First,
\texttt{member} forces a canonical code in $D$, and raw addition on two
canonical elliptic inputs forces their canonical sum, including the
origin and two-torsion cases. The fixed scalar graph gives $[m]Q$ on
$D$. In \texttt{multiply}, the three root membership calls are forced
conjuncts inside \texttt{height\_certificate}. Thus the four scalar shifts
are in $D$ before they are passed, unchanged, to the 27 slope calls and
the six test-group calls within each slope. This supplies the hypotheses
of every raw addition in those groups. Only after globally typed
multiplication is established do the four square products and three
successive additions justify \texttt{natural}. Their outputs are canonical
by the raw-addition postcondition, so no separate output membership
assumption is needed. The signed-encoding and exceptional branches start
their addition chains at this natural code, the literal origin, or the
fixed signed generator. An inactive product branch is not assumed true.
The fixed generator equations are conjuncts outside the internal products
in all three templates that use them. These observations give the base
cases for the source-DAG induction below.
Apply (C.4) using the exact signed index model. Induction on the actual
source DAG gives each node the code of its signed integer value.
Natural guards constrain source variables, not intermediate sums or
products. Raw addition receives canonical codes; globally typed
multiplication justifies the natural guards. Fixed generator equations
and scalar templates supply constants.

The power-value recipe has three signed code inputs, seventy-nine
private natural witnesses, and an output code for the value of the
power-template polynomial. That value is not independently required
to vanish. Its saved root equality is $[17751579,15]$.
The precision parent has $54056$ explicit variables and $15845$
disjoint private blocks of seventy-nine variables. No parent variable
node accesses a private coordinate. Replacing each block by its
private power-value call therefore preserves the full existential
assignment. Only the final parent root is equated to the identity.

At each primitive formula, equality is a difference, conjunction is
a sum of squares, and disjunction is a product. A topological induction,
including the globally conjoined fixed equations, proves equality of
polynomial values under substitution, not merely equisatisfiability
of an unrelated formula. This is the required link to the saved AST.

\section{A variable-shift relation for ordinary values}\label{sec:bridge}
We prove the ordinary-value deduction from the exact index model,
the compiled precision relation, the height theorem and $\mathsf{PC}$.
These are supplied by Sections~\ref{sec:model}, \ref{sec:compiler},
\ref{sec:height} and Theorem~\ref{hyp:PC}, respectively. Capital letters below denote five-coordinate point codes;
their lower-case indices are ordinary integers. The external variable
$t$ is any rational number.

\subsection{Approximation for every fixed shift}
Define $\Close(N,\mathcal K,t)$ by choosing $A,L,J$, linking all five
signed codes to precision inputs in the order $(n,\kappa,a,l,j)$, and
imposing $|t-r(a,j)|<\epsilon_l$ on those same coordinates. Put
\[
 B(n,\kappa)=H_6(S_{\rm ht})|\kappa F(n)|^{24}.
\]
The precision contract gives
\begin{align}
 \Close(\iota(n),\iota(\kappa),t)&\Longrightarrow
                  |t-n|<e^{-B(n,\kappa)},\label{eq:close-sound}\\
 \Close(\iota(n),\iota(\kappa),n)&\quad\text{for every }
                  n\notin\Cset,\ \kappa\ne0.\label{eq:close-complete}
\end{align}
The first follows from $|t-n|<2\epsilon_l<4^{-B}/2<e^{-B}$.
For the second, positive multiples of $P_0$ approach the identity
simultaneously in the compact group $E(\mathbf R)\times E(\mathbf R)$.
This follows, for instance, by partitioning the compact closure of the
cyclic subgroup into arbitrarily small neighborhoods and subtracting
two sufficiently close multiples. If the resulting positive differences
were bounded, a fixed nonzero multiple would be the identity; infinite
order excludes this. Choose $l$ first to make the positive $\epsilon_l$
small enough. Since the derivative of $[n]$ at the identity is $n$
\cite[IV.2.3(a), p.~122]{silverman},
$r(a,na)\to n$ along positive multiples $a$ approaching it. Choose such
an $a$ and put $j=na$. This works for every already fixed nonzero shift
$\kappa$, regardless of its size.

\subsection{The six local tests}
For $N=\iota(n)$, $\mathcal K=\iota(\kappa)$, form the code for
$\kappa n$, choose $H=\iota(h)$, $U=\iota(s_0)$, and put
$V_0=H-\mathcal K$, $V_1=U-\iota(\kappa n)$.
Require $\kappa\ne0$, $H,V_0\ne O$. Define $w(Q)=1/x(Q)$ for
$Q\ne O$, $w(O)=0$. For $t\notin\Cset$ use
\begin{align*}
 D_1&=w(H),&D_2&=w(U),&D_3&=\sigma w(V_0),&D_4&=\sigma w(V_1),\\
 D_5&=z(U)/z(H)-t,&D_6&=\sigma(z(V_1)/z(V_0))-t.
\end{align*}
Here $\sigma$ is conjugation in $K_0$. Choose $\eta,q\in\Q$ and impose
\begin{align*}
 \eta^2&=1,&R_0&=\eta2730^6F(t),&a_0&=1-R_0q^2\ne0,\\
 R&=R_0a_0^2,&C_j&=D_j/F(t),&W&=\prod_{j=1}^6(1+RC_j^2)\ne0,\\
 \Theta&=\Tr_{K_0/\Q}(W^{-1})\ne0,&B_{\rm pole}&=R/\Theta^2,
 &\Phi(B_{\rm pole})&.
\end{align*}
Call this relation $\Test(N,\mathcal K,t)$. All quotients are guarded
polynomially. There is one common denominator $F(t)$, one product $W$,
one trace and one pole instance.

\begin{lemma}[Divisibility]\label{lem:test}
If $\Test(\iota(n),\iota(\kappa),t)$ holds, then
$M_{S_{\rm ht}}(t)\mid\kappa F(n)$.
\end{lemma}
\begin{proof}
Let $p\mid M_{S_{\rm ht}}(t)$. Exactly one $c\in\Cset$ has
$\delta=v_p(t-c)>0$ odd. The contact differences are units, so
$v_p(F(t))=\delta$; the field is unramified at $p$. Put $b=v_p(a_0)$.
If $q=0$, then $b=0$. Otherwise $\delta+2v_p(q)$ is odd and
nonzero, so $b=\min(0,\delta+2v_p(q))\le0$. In either case
$\alpha=v_p(R)=\delta+2b$ is odd. At each place above $p$,
\[
 v(1+RC_j^2)=\min(0,\alpha+2(v(D_j)-\delta)).
\]
Let $\lambda=\max_{v\mid p}v(W)\le0$. Then
$v_p(\Theta)\ge-\lambda\ge0$. The pole input has odd valuation
$\alpha-2v_p(\Theta)$, which is positive by Lemma~\ref{lem:pole}.
Therefore $\alpha>0$ and $\lambda> -\alpha/2$.

At a place attaining $\lambda$, each nonpositive factor valuation is
at least $\lambda$. Consequently
\[
 v(D_j)\ge\delta-3\alpha/4=\delta/4-3b/2>0
 \quad(1\le j\le6).
\]
Integral valuations are at least one, at the same place for all tests.
The first four put $H,U$ in one formal group and $V_0,V_1$ in the
conjugate formal group. At the identity we set both $z(O)=0$ and
$\log(O)=0$; thus the cases $U=O$ and $V_1=O$ obey the same ratio
identities below. Their denominators are nonzero by the guards on
$H$ and $V_0$. For a nonzero point in either formal group,
\[
 \log(Q)=z(Q)(1+e_Q),\qquad v(e_Q)\ge1.
\]
Here $p>324$ and $v(p)=1$. By
\cite[IV.5.5 and IV.6.4, pp.~129--132]{silverman},
the logarithm is an injective isometry on $v(z)\ge1$.
The integral invariant differential gives in relative degree $r-1$
the bound $(r-1)v(z)-v_p(r)\ge1$ for every $r\ge2$.
Thus the guarded nonzero points $H,V_0$ have nonzero logarithms.
The fifth test first makes $z(U)/z(H)$ integral and congruent to
$t$; multiplication by the units $1+e_Q$ preserves that congruence.
The relation $[s_0]H=[h]U$ then gives
$\log(U)/\log(H)=s_0/h$. Likewise the sixth gives
$(s_0-\kappa n)/(h-\kappa)\equiv t$.
Multiplying by the integer denominators and subtracting yields
$v_p(\kappa(n-t))\ge1$. Since $t\equiv c$, we get
$p\mid\kappa(n-c)$, hence $p\mid\kappa F(n)$. The squarefree
product divides as asserted. No unit assumption on $\kappa$ was used.
\end{proof}

\begin{lemma}[One suitable shift]\label{lem:test-complete}
For every $n\in\Z\setminus\Cset$ there is
$\kappa\ne0$ with $\Test(\iota(n),\iota(\kappa),n)$.
\end{lemma}
\begin{proof}
Choose the sign so $R_0=\eta2730^6F(n)>0$ and let $S'$ be its finite
prime support. Each sufficiently small local formal neighborhood has
finite index in the compact local group
\cite[VII.6.3, p.~200]{silverman}. A positive common multiple
of the quotient orders gives $\kappa$ with $[\kappa]P_0$ in every
chosen neighborhood. Take
\[
 H=[2\kappa]P_0,\quad U=[2\kappa n]P_0,\quad
 V_0=[\kappa]P_0,\quad V_1=[\kappa n]P_0.
\]
The formal parameter and the derivative of $[n]$ make each $C_j$
integral above $S'$ after sufficiently shrinking the neighborhoods.
Infinite order preserves the nonzero guards.

Choose rational $q$ integral at $S'$ and arbitrarily close to
$1/\sqrt{R_0}$ in $\mathbf R$, without equality. Denominators
$R_0m+1$ with nearest integer numerators give such choices, adjusting
by one if equality occurs. At $S'$, $a_0$ is a unit and
$W\equiv1$, $\Theta\equiv2$ modulo $p^{v_p(R_0)}$.
Integral trace justifies this also at the ramified place two.
The pole input has valuation at least four at two and at least six
at the other primes of $S_0$. Outside $S'$, its valuation is
$2v_p(a_0)-2v_p(\Theta)$ and is even, so its reduced denominator is
a square. In both real embeddings, $R$ tends to zero positively,
$W\to1$ and $\Theta\to2$. All guards hold and
$0<B_{\rm pole}<1/2$ for a sufficiently close choice.
Theorem~\ref{hyp:PC} supplies the pole witnesses.
\end{proof}

\subsection{Exact ordinary value}
Define $\Val(N,t)$ by membership of $N$ and the disjunction
\[
 \bigvee_{c\in\Cset}(N=\iota(c)\wedge t=c)
 \quad\vee\quad
 \exists\mathcal K\,
 (\Close(N,\mathcal K,t)\wedge\Test(N,\mathcal K,t)).
\]
The same five coordinates $\mathcal K$ are used in both calls.
In an accepting main branch $n\notin\Cset$, $\kappa\ne0$, and
Lemma~\ref{lem:test} and the height estimate imply
\[
 h(t)\le H_6(S_{\rm ht})|\kappa F(n)|^{24}=B(n,\kappa).
\]
For $t=a/b\ne n$ in reduced form, $|a-nb|\ge1$, hence
\[
 |t-n|\ge1/b\ge e^{-h(t)}\ge e^{-B(n,\kappa)},
\]
contradicting~\eqref{eq:close-sound}. Thus $t=n$.
Conversely choose $\kappa$ by Lemma~\ref{lem:test-complete}, then use
\eqref{eq:close-complete} for that same shift. Private witness blocks
are disjoint, so the two calls can be satisfied simultaneously.
The six remaining integers have exactly their displayed branches.

Therefore $\Val(\iota(n),t)$ is exactly the graph $t=n$.
Existentially quantifying $N$ gives ordinary $\Z$ in $\Q$.
S01 and S16 multiply the exception and main-branch polynomials and
conjoin membership by a sum of squares. The literal substitutions in
Section~\ref{sec:syntax} and Appendix~\ref{app:literal} identify this
polynomial with the saved root $P_*$. This proves the zero-set part of
Lemma~\ref{lem:assembly}.

% End 04-computation-value.tex


% Begin 05-bounds-appendices.tex
\section{The cutoff and polynomial size}\label{sec:size}
The literal constants in S03 define the catalogue product $H_{\rm cat}$,
its maximum $H_{\max}$, $N_s=P_sH_{\rm cat}$ with
$P_s=\prod_{p\le324}p$, the least formal multiple $M_0$, and $c_*=4M_0$.
They take
\[
 R=\max(D,H_{\max},c_*,10^{12}),\quad
 A_{\rm ht}=2\cdot10^{13}\max(1,C_{\rm cm},Q_R^{17}),\quad
 T=(48A_{\rm ht})^{13}.
\]
The pole clearing constant is $\ell_0=12P_s^{2916}H_{\rm cat}^{732}$.
The symbol $F_0=324!$ is distinct from the six-contact polynomial $F(u)$.

\begin{proposition}[Literal cutoff]\label{prop:cutoff}
The fixed prescription satisfies
\[
 R=D,\qquad T=(96\cdot10^{13})^{13}D^{2210G_0^2D^3}
            <2^{\,2^{\,2^{9013}}}.
\]
This is a statement about the prescribed integer, independently of its
semantic role in excluding false multiplication triples.
\end{proposition}
\begin{proof}
We give the bounds of S25. Exact small calculations give
$10^5<W<2^{512}$ and $2^{2000}<F_0<2^{2250}$.
For $q\ge10^5$, the unconditional progression estimate
\cite[Theorem 1.2, (1.10)--(1.11), Corollary 6.17 with $A=1$]{bmor}
applies for $x\ge\exp(0.03\sqrt q(\log q)^3)$, including the possible
exceptional-zero case. At $X=2^{W^3}$ its error is at most
$X/(160\log X)$ and the hypothesis follows from
$0.03\sqrt W(\log W)^3<(\log2)W^3$.
Hence $\vartheta(X;W,a)>X/(2W)$, giving more than
$X/(2W^4)>22$ primes in each unit class. No prime at most 324
belongs to such a class. Thus
\[
 H_{\max}<X,\quad H_{\rm cat}\le X^{22W},\quad
 \log_2N_s<2^{2054}.
\]
For every odd prime dividing $N_s$, take the actual reduction order of
$P_0$ at each place above it, and let their product be $m_0$.
All preceding positive multiples stay in the dyadic formal subgroup,
since $P_0$ starts there and the subgroup is closed under addition.
A nonidentity point there has negative abscissa valuation, so its
reciprocal is defined; the constructor accepts the identity directly.
Thus every earlier search step is defined without a nontorsion premise.
At the bad place $s\equiv3\pmod7$, the point reduces to $(2,6)$
on $y^2=x^3$ and is nonsingular because $2y\not\equiv0$.
The other place $s\equiv4\pmod7$ has good reduction.
Take the actual order of the reduced point in the nonsingular
reduction group, using
\cite[VII.2.1, p.~188]{silverman}.
At a good place the residue degree is at most two, so Hasse's
bound \cite[V.1.1, p.~138]{silverman} gives order at most
$(p+1)^2<p^3$ for odd $p$. The singular smooth group over
$\F_7$ has order seven. There are at most two places over each
rational prime, so $m_0\le N_s^6$.
If $[N_s^4m_0]P_0$ is the identity, the search has already succeeded.
Otherwise its reciprocal coordinate is defined and the following
estimates apply. At each odd place, $[m_0]P_0$ is formal, and multiplication by
$N_s^4$ gives
\[
 v(1/x([N_s^4m_0]P_0))\ge2+8v_p(N_s).
\]
At two normalize $v_s(s)=1$, so $v_s(2)=2$. Initially
$v_s(x)=-6$, $v_s(y)=-9$, and $v_s(z)=3>2$.
The log isometry therefore gives
$v_s(z([2]Q))=v_s(z(Q))+2$ on this subgroup, and hence
\[
 v_s(1/x([N_s^4m_0]P_0))\ge6+16v_2(N_s).
\]
Writing this reciprocal as $a+bs$, extraction of $a$ and $b$
uses division by $2$ and $2s$ and loses at most three units of
$v_s$. Thus both rational coefficients lie in
$N_s^4\Z_{(N_s)}$, where denominators in $\Z_{(N_s)}$ are
prime to $N_s$. This proves the required divisibility and guard. Therefore the literal search
succeeds by $N_s^4m_0$ and $M_0\le N_s^{10}$.
This uses actual finite orders, not numerical upper bounds incorrectly
treated as multiples of those orders.

Consequently $c_*<2^{2^{2059}}$, whereas
$2^{2^{8008}}<D<2^{2^{9010}}$, so $R=D$.
Next $C_{\rm cm}<2^{50000D^2}$ and
\[
 \log_2Q_D^{17}=170G_0^2D^3\log_2D>50000D^2.
\]
The maximum is therefore $Q_D^{17}$, giving the identity for $T$.
Finally $G_0<2^{37}$, $E_0<D^3$, $\log_2Q_D<D^5$,
$\log_2A_{\rm ht}<D^7$, and $\log_2T<D^8<2^{2^{9013}}$.
\end{proof}

\subsection{Quantitative syntax bounds}
Use coefficient $\ell^1$ norm for integer polynomials and the weighted
pair norm $\|a+bs\|=\|a\|_1+2\|b\|_1$. They are submultiplicative.
The division-polynomial recurrences give
\[
 \deg f_n\le(n+1)^3,\qquad\|f_n\|\le2^{100(n+1)^3}.
\]
For the even recurrence the exponent slack is
$4m^3-24m-27\ge9$ for $m\ge3$; for the odd recurrence it is at least
ten for $m\ge2$. This absorbs the fixed factors and addition.

The following facts have been checked for the actual saved expression:
the absolute values of all literal fixed constants and all occurring
division-polynomial indices are at most $T$; the specified syntax-height
traversal has root height at most 256; each model recipe has at most
$2^{20}$ constraints and each composition recipe has at most $2^{20}$
terms. The check accounts for all 23 templates, seven recipes, private
witness blocks, 173 scalar constant nodes and 50 finite constructors. Its exact
artifact identity and independent replay are recorded in the source map.
Together with the constructor bounds below, this discharges the
quantitative part of $\mathsf B$. The all-input arithmetic interpretation of the same constituents is the
separate substitution induction of Section~\ref{sec:compiler}.

Put $B_*=T+1$. Coordinates of $\iota(n)$ for
$|n|\le T$ have absolute numerator and positive denominator at most $2^{K_*}$, with
$K_*=2^{30}(T+1)^3$. Evaluate the division polynomials at $P_0$ and use
the fixed coordinate formulas; at most eight quadratic-field operations
are needed, each taking an available height bound $H$ to at most $3H^4$.
There are at most $8(T+1)^3$ guard factors. The resulting bounds for
each finite constructor are
\[
 \deg\le B_*^4,\qquad\log_2\max(1,\|\cdot\|_1)\le B_*^8.
\]

Give variables syntax height zero, constants height one, and add one
for each arithmetic operation or finite constructor. A template call
adds one plus the callee's intrinsic height and the maximum argument
height. The substitution inequalities
\[
 \deg P(Q)\le\deg P\max(1,\max_i\deg Q_i),\qquad
 \|P(Q)\|_1\le\|P\|_1\max(1,\max_i\|Q_i\|_1)^{\deg P}
\]
give degree and logarithmic norm at most $B_*^{8h}$ at positive height
$h$, with the bounded sums in the actual recipes treated as in S26.
This general call rule applies to polynomial arguments. In each outer
composition recipe the arguments are coordinate variables, so the call
is only a renaming, even when coordinates repeat. For any such map
$\sigma$,
\[
 \deg P(Y_{\sigma(i)})\le\deg P,\qquad
 \|P(Y_{\sigma(i)})\|_1
 =\sum_\beta\left|\sum_{\alpha\mapsto\beta}c_\alpha\right|
 \le\sum_\alpha|c_\alpha|=\|P\|_1.
\]
Thus these calls retain the callee's majorant height before charging
the sum-of-squares or product connective. The seven recipe heights are
$228,233,235,251,252,254,256$. A rule charging an extra unit for each
of these mere renamings would describe a looser syntax-depth bound,
not this proved degree-and-norm majorant.
For the model wrappers use the safe height bound four for the
generator equations and root differences; the callee contribution
dominates this bound. The literal height-two claim for those equations
in S26 is not used.
This height is a majorant for the compact constructor syntax, not for
an expanded monomial tree. The height-256 bound yields
\[
 \deg P_*\le(T+1)^{2048},\qquad
 |[X^\alpha]P_*|\le2^{(T+1)^{2048}}.
\]

Since $T$ is an integer, $T+1\le U=2^{2^{2^{9013}}}$. Hence
\[
 (T+1)^{2048}\le U^{2048}=2^{2^{2^{9013}+11}}
                     <2^{2^{2^{9014}}}.
\]
As $9014<65536=\tw4$, the degree is less than $\tw7$ and the
coefficient magnitude less than $\tw8$. The root recipe declares
$292\,890\,231\,904$ witnesses. Together with the preceding estimates,
these complete the bounds for the same $P_*$ in
Theorem~\ref{thm:target}.

\section{Consequences of the integer definition}
The proof of Theorem~\ref{thm:target} combines the following results.
\begin{center}\small
\begin{tabular}{p{.20\textwidth}p{.43\textwidth}p{.27\textwidth}}
\toprule
Result & Mathematical inputs & Proof location\\\midrule
Index group & Two-isogeny descent, reductions, finite height search &
Section~\ref{sec:model}\\
Six-contact bound & Quaternionic model, tame pullback, compatible-isogeny
height theorem, theta comparison & Section~\ref{sec:height}\\
Pole completeness & Simultaneous descent, restricted converse, forward
rank-one theorem & Section~\ref{sec:pole}\\
Multiplication & Index group, height bound, selected primes,
pole completeness, finite guard & Section~\ref{sec:mult}\\
Precision relation & Decision program, lowering, power and packed histories &
Section~\ref{sec:compiler}\\
Ordinary value & Shared-shift Close and Test, height and rational separation &
Section~\ref{sec:bridge}\\
Bounded $P_*$ & Literal substitutions and quantitative syntax &
Sections~\ref{sec:syntax}, \ref{sec:size}; Appendix~\ref{app:literal}\\\bottomrule
\end{tabular}
\end{center}

\subsection{Three corollaries}\label{sec:four-corollaries}
We apply the exact equality between the rational zero projection of
$P_*$ and ordinary $\Z$ proved in Theorem~\ref{thm:target}.

\begin{corollary}[Negation of Mazur's finite-component conjecture]
\label{cor:mazur}
Mazur's
finite-connected-components conjecture for the real-topological closure
of rational points is false.
\end{corollary}
\begin{proof}
Mazur's Conjecture~3 in \cite[p.~165]{mazur1995} asserts that the closure
of the rational points of any variety over $\Q$ in its real locus has
finitely many connected components. In that source a variety means a
reduced scheme of finite type; irreducibility is not required.
Put $r=292\,890\,231\,904$ and
\[
 V=\bigl(\{P_*(t,w_1,\ldots,w_r)=0\}\subset\mathbf A^{r+1}_{\Q}\bigr)_{\rm red},
 \qquad C=\overline{V(\Q)}\subset V(\mathbf R).
\]
Reduction changes neither the rational nor the real points. Let
$\pi:V(\mathbf R)\to\mathbf R$ be projection to $t$.
Theorem~\ref{thm:target} gives $\pi(V(\Q))=\Z$. Since $\Z$ is closed in $\mathbf R$
and $\pi$ is continuous, $C\subset\pi^{-1}(\Z)$, so $\pi(C)\subset\Z$.
Every integer already occurs on $V(\Q)$, hence $\pi(C)=\Z$.
The image of any connected component of $C$ is a connected subset of
$\Z$, therefore a singleton. Points projecting to distinct integers
belong to distinct components. Thus $C$ has infinitely many connected
components, contradicting the stated conjecture.
\end{proof}

\begin{corollary}[A universal rational Diophantine polynomial]
\label{cor:rational-universal}
There exist $m\ge1$ and a fixed polynomial
$U\in\Z[e,n,x_1,\ldots,x_m]$ such that, for every $e,n\in\N$,
\[
 \exists\mathbf x\in\Q^m\ U(e,n,\mathbf x)=0
 \quad\Longleftrightarrow\quad
 \text{machine $e$ halts on input $n$}.
\]
Here $e$ and $n$ are ordinary integer parameters, with $\N$ embedded
in $\Q$, and are not elliptic-curve coordinates. In particular, every
computably enumerable subset of $\N$ is represented by a suitable
fixed value of $e$ in this one polynomial.
\end{corollary}
\begin{proof}
Fix an effective enumeration of machines. The Davis--Putnam--Robinson--Matiyasevich theorem, \cite[Basic Theorem, p.~3]{matiyasevich1971},
applied to the universal halting relation gives a fixed
$D\in\Z[e,n,z_1,\ldots,z_s]$, with $s\ge1$, such that, for $e,n\in\N$,
\[
 \text{machine $e$ halts on $n$}
 \quad\Longleftrightarrow\quad
 \exists\mathbf z\in\Z^s\ D(e,n,\mathbf z)=0.
\]
To use integer witnesses in this statement, replace each nonnegative
witness by a sum of four integer squares. In the positive-integer
convention of the cited source, apply the theorem to the predicate
that machine $a-1$ halts on $b-1$, substitute $a=e+1$, $b=n+1$,
and use one plus that sum for each witness; see also
\cite[p.~2]{matiyasevich1971}. A dummy witness ensures $s\ge1$ if needed.
For each $z_i$ take a separate block $\mathbf w_i$ of $r$ rational
variables, where $r$ is the witness count of $P_*$. Define
\[
 U(e,n,\mathbf z,\mathbf w_1,\ldots,\mathbf w_s)
   =D(e,n,\mathbf z)^2+\sum_{i=1}^s P_*(z_i,\mathbf w_i)^2.
\]
A rational zero makes every summand zero. Theorem~\ref{thm:target}
then forces every $z_i$ to be an integer and yields the displayed
halting condition. Conversely, integer witnesses for $D$ have separate
rational witness blocks for $P_*$, giving a zero of $U$. All blocks
are private and the polynomial is fixed independently of $e,n$.
No variable, degree or coefficient bound proved for $P_*$ is asserted
for this new polynomial $U$.
\end{proof}

\begin{corollary}[The many-one degree of $\operatorname{HTP}(\Q)$]
\label{cor:htp-degree}
\[
 \operatorname{HTP}(\Q)\equiv_m 0',
 \qquad 0'=\{e\in\N:\text{machine $e$ halts on input $e$}\}.
\]
Here $\operatorname{HTP}(\Q)$ is the set of codes of integer-coefficient
polynomials, of arbitrary finite arity, having a rational zero, under a
fixed effective coding.
\end{corollary}
\begin{proof}
Given $e$, specialize the fixed polynomial of
Corollary~\ref{cor:rational-universal} at its ordinary integer parameters
to obtain $U(e,e,\mathbf x)$. Computing its coefficient list and code is
effective, and
\[
 e\in0'\quad\Longleftrightarrow\quad
 \ulcorner U(e,e,\mathbf x)\urcorner\in\operatorname{HTP}(\Q).
\]
This proves $0'\le_m\operatorname{HTP}(\Q)$.

Conversely, from the code of a polynomial $f$ construct a machine that
ignores its input and enumerates all rational tuples of the appropriate
arity. It evaluates $f$ using exact rational arithmetic and halts as soon
as it finds a zero. Its index $j(f)$ is computable from the finite code
of $f$, and
\[
 f\in\operatorname{HTP}(\Q)\quad\Longleftrightarrow\quad j(f)\in0',
\]
where the left side uses the code of $f$. Thus
$\operatorname{HTP}(\Q)\le_m0'$. Invalid polynomial codes, if present in
the coding convention, are sent to a fixed nonhalting index. Using
rational-coefficient polynomial codes gives the same result, since
clearing their finitely many denominators is effective and preserves
rational zeros. The two reductions prove the assertion.
\end{proof}

More generally, an integer Diophantine problem translates to a rational
one by adjoining $P_*$ for every integer witness, with private witness
blocks. Natural witnesses can first be expressed as sums of four integer
squares. Bounds for the resulting polynomial require their own composition
calculation; the bounds for $P_*$ do not automatically transfer.

For example, a machine that enumerates ZFC proofs and halts upon finding
a contradiction has a rational Diophantine halting condition by
Corollary~\ref{cor:rational-universal}. This equivalence does not assert
that a contradiction exists, and no additional explicit polynomial is
specified here.

% Begin degree-eight construction
\section{The base residual compiler}
\label{sec:degree8-original}

This section records the base residual compiler used as a construction
stage for the quartic in Section~\ref{sec:quartics}. Its degree-eight
sum-of-squares encoding fixes the arithmetic interfaces and conventions;
it is not a third final polynomial specification. The construction uses
a joint comparison of four elliptic indices, a Lucas equation, and the
five-contact height estimate proved below. The archived counts describe
this base compiler before its reciprocal revision and quadratic lift.

\begin{proposition}[Finite specification of the base compiler]\label{prop:o8-finite}
The prescription below specifies a polynomial
\[
 F_8\in\Z[X,V_1,\ldots,V_{15671}]
\]
of degree exactly eight, with exactly $129497$ nonzero monomials and
maximum absolute coefficient
\[
 685783949153203534233601<2^{80}.
\]
These statements are finite identities about the displayed construction
and its integer coefficient table.
\end{proposition}

\begin{theorem}[Arithmetic interpretation]\label{thm:o8-conditional}
For every
$x\in\Q$,
\[
 x\in\Z\quad\Longleftrightarrow\quad
 \exists(v_1,\ldots,v_{15671})\in\Q^{15671},\qquad
 F_8(x,v_1,\ldots,v_{15671})=0.
\]
\end{theorem}

The required height estimate is proved below; Theorem~\ref{thm:height}
uses a different contact set. The seven-variable pole lemma below is proved from
Theorem~\ref{hyp:RC}; it is not inferred from Theorem~\ref{hyp:PC},
which has different variables and a restricted completeness domain.
The finite coefficient calculation does not establish the height bound.
The supplied unpublished Freitag archive is the provenance of this
construction and of the original coefficient table. It is not cited here
as an arithmetic theorem. The published inputs already used in
Sections~\ref{sec:model} and \ref{sec:mult} remain
available, together with the quaternion trace theorem specified below.

\subsection{The height input and the exact pole predicate}

Put
\[
 S_* =\{p:p\le76176\},\qquad
 \mathcal D=\{0,1,2401,-343/32\},
\]
\[
 H_*=2(4608^{165888}+6^{100})
          \left(\prod_{p\le76176}p\right)^{24}+128.
\]
For $a\in\Q\setminus\mathcal D$, let $M_5(a)$ be the product of
the primes outside $S_*$ at which the point $(a:1)$ has positive odd
contact order with one of $\infty,0,1,2401,-343/32$. Explicitly, choose
local homogeneous coordinates $(u:v)$ and subtract
$\min(v_p(u),v_p(v))$ from the valuations of $v$ and of
$u-\delta v$ for $\delta\in\mathcal D$.

\begin{theorem}[The five-contact height estimate]\label{hyp:o8-height}
For every $a\in\Q\setminus\mathcal D$,
\[
 h(a)\le H_*M_5(a)^{24}.
\]
The displayed constant satisfies $\log_2H_*<4783227$.
\end{theorem}
\begin{proof}
Write $s=a$, $M(s)=M_5(a)$, $H=H_*$,
$B_*=\prod_{p\le76176}p$, and
$\mathcal C=\{\infty,0,1,2401,-343/32\}$.
We prove the geometric comparison, the non-CM estimate through the
modular level, and the CM estimate over the same fiber field.
All heights use the original stable Faltings normalization.
\paragraph{Recovering the contact line from published quotients}
Let $B/\mathbb Q$ be the indefinite quaternion algebra of discriminant $210$,
$\mathcal O$ a maximal order, $X$ its maximal-level Shimura curve, and $W$
its Atkin--Lehner group. The standard genus formula gives
\[
 g(X)=1+\frac{\varphi(210)}{12}=5,
\]
because the elliptic-point terms vanish: $5$ splits in $\mathbb Q(i)$ and
$7$ splits in $\mathbb Q(\sqrt{-3})$, so neither quadratic field embeds in
$B$. The group $W$ is $(\mathbb Z/2\mathbb Z)^4$ with involutions $w_m$
indexed by square classes of divisors of $210$.

The published models of \cite[Theorem 4.2 and Table 2, pp.~938--939]{GonzalezRotger2006} have the following coefficient rows
$(a,b,c,d,e)$ for $y^2=ax^4+bx^3+cx^2+dx+e$:
\[
\begin{array}{c|rrrrr}
 m&a&b&c&d&e\\\hline
30&-43&-686&-2915&-1372&-172\\
42&-43&600&-986&-600&-43\\
70&-43&256&-130&-768&-387\\
105&-43&2&505&12&-1548
\end{array}
\]
The cited theorem identifies each with $X/\langle w_m\rangle$.
The genus formula and the Atkin--Lehner group law used here also appear
in the same published paper, Proposition 5.2 and the paragraph following
Proposition 5.3, pp.~940--941.

Here is a deduction of all the required branch information, rather than
an additional assumption about it. Each of these four genus-one quotients
has eight fixed points upstairs, by Riemann--Hurwitz. A stabilizer for the
$W$ action on a characteristic-zero smooth curve is cyclic, and therefore
has order at most two. Each nonempty fixed set of an involution is a union
of $W$-orbits of size eight. Riemann--Hurwitz for $X\to X/W$ now gives
\[
 8=16(2g(X/W)-2)+\sum_{1\ne w\in W}\#\operatorname{Fix}(w).
\]
The four known fixed sets already contribute $32$. Thus $g(X/W)=0$,
the total is $40$, and there is exactly one further orbit of eight fixed
points. Inertia generators around the branch points have product one.
The four classes $30,42,70,105$ are independent in $W$, and their product
is the class $210$. Hence there are exactly five branch points, each of
index two, with distinct labels $30,42,70,105,210$. Since the involutions
are defined over $\mathbb Q$ and each label occurs only once, each branch
point is rational. Three of them trivialize the genus-zero quotient.

Normalize $b_{210}=\infty$, $b_{105}=0$, $b_{70}=1$, and write
$b_{42}=\lambda$, $b_{30}=\mu$. On a genus-one quotient the residual
group $(\mathbb Z/2)^3$ has translation subgroup of order four, necessarily
the full two-torsion group after choosing an origin. Its quotient has the
same Jacobian $j$-invariant: it is multiplication by two over the algebraic
closure. The remaining double cover of $X/W$ branches at the four points
other than $b_m$. Therefore the four quartic invariants determine the four
Legendre invariants of these four-point subsets.

For a quartic put
\[
 I=12ae-3bd+c^2,\quad
 J_q=72ace+9bcd-27ad^2-27b^2e-2c^3,\quad
 j=\frac{6912I^3}{4I^3-J_q^2}.
\]
The exact computation gives the Legendre parameters
$1/2401,32/375,1/225,5/1029$, in row order. Enumerating the six cross-ratio
values of each of the first two parameters and testing the last two
invariants leaves exactly
\[
 \lambda=2401,\qquad\mu=-343/32.
\]
The script \texttt{check\_height.py} reproduces this rational computation.
Every nonzero determinant of the five primitive contact forms has prime
support in $\{2,3,5,7\}$.

\paragraph{Integral geometry and the degree 576 field}
Choose $\mu_0\in\mathcal O$ with $\mu_0^2=-210$. The embedding criterion
allows this choice in a maximal order; maximal orders are conjugate here.
Use the maximal moduli stack $\mathscr X$ and the entire principal-level
three scheme $Y$, including its geometric components. Yuan's integral
moduli theorem \cite[Theorem 1.1 and Sections 2.1--2.2]{Yuan2024} makes $\mathscr X$ regular, proper and semistable over
$\mathbb Z$, with a universal abelian surface, and smooth away from $210$.
Its generic stabilizer is precisely $\{\pm1\}$. Principal level three is
fine and gives a universal full three-torsion marking.

Work first over $R_0=\mathbb Z[1/210]$. The effective level-change group
has order $36$, prime to every residue characteristic of $R_0$. Its generic
action is free. For a finite tame group acting on a smooth relative curve,
the fixed locus of an element is smooth locally along its fixed points;
an isolated special fixed point consequently lifts. Thus no new fixed
point occurs here, and $Y_{R_0}\to X_{R_0}$ is finite etale. The $W$ action
extends to the smooth stable model of $X$. Its stabilizers remain cyclic,
of order at most two. Distinct involution fixed loci cannot intersect;
each fixed locus is etale over the base. Consequently $X_{R_0}/W$ is a
smooth genus-zero curve with five disjoint branch sections of index two.
The three rational sections already chosen identify it with
$\mathbb P^1_{R_0}$. Their values are the five elements
of $\mathcal C$.

Set $K=\widehat{\mathcal O}^{\times}$, $U=U(3)$,
$Z=\mathbb A_f^\times$, and $N=N_{B(\mathbb A_f)^\times}
(\widehat{\mathcal O})$. Normalization preserves $3\widehat{\mathcal O}$,
so $U$ is normal in $N$. The canonical adelic action is defined over the
reflex field $\mathbb Q$ \cite[2.2.5 and 2.3.1]{Deligne1979}; this refers to the entire canonical model, not
an individually selected complex component. Every central finite idele
is a rational scalar times a unit idele. At three its scalar residue is
$1$ or $-1$, which a rational sign removes. Thus $Z\subset\mathbb Q^\times U$
acts trivially on $Y$.

At three the quaternion algebra is ramified. Its maximal order modulo
three has $81$ elements, radical of size $9$, and hence $72$ units.
It follows that
\[
 1\longrightarrow (\mathcal O/3\mathcal O)^\times/\{\pm1\}
 \longrightarrow\Gamma=N/(ZU)\longrightarrow W\longrightarrow1,
 \qquad |\Gamma|=36\cdot16=576.
\]
Above an unbranched point the actions of $W$ and of the effective
level-change kernel are free. The fiber therefore is a torsor of size
$576$ for this constant group. The Galois permutations commute with its
regular action, so belong to its regular centralizer, also of order
$576$. The normal splitting field $L$ consequently satisfies
$[L:\mathbb Q]\mid576$: the Galois group is a subgroup of that
regular centralizer. In particular $[L:\mathbb Q]\le576$.

All fiber surfaces and their three-torsion markings are defined over $L$.
For two such surfaces, a homomorphism annihilating the three-torsion
factors through $[3]$. Hence reduction embeds
\[
 \operatorname{Hom}(A,B)/3\operatorname{Hom}(A,B)
 \hookrightarrow\operatorname{Hom}(A[3],B[3]).
\]
The finite Galois image on the integral Hom lattice is trivial modulo
three. The principal congruence subgroup modulo three in an integral
general linear group is torsion-free, so this image is trivial. Thus
every geometric homomorphism between fiber surfaces is defined over $L$.

At a prime outside $210$, tame local coordinates for the quotient are
$w=z^2$ up to a unit at a branch section, and etale elsewhere. A base
point with even contact gives no ramification after an unramified
extension. The universal abelian scheme extends at the resulting
integral point of $Y_{R_0}$. Therefore
\[
 \operatorname{rad}(D_LN_{\rm bad})\mid B_*M(s).
\tag{R}\label{eq:h5-R}
\]
This assertion includes contact with infinity. It is not a statement
obtained merely from the numerical prime support of pairwise determinants.

The normalizer correspondences give isogenies between the conjugate
surfaces. In the non-CM case the geometric endomorphism algebra is $B$;
Skolem--Noether adjusts a conjugate isogeny to commute with its chosen
quadratic subfield. All such maps lie over $L$ by the Hom argument.
These facts establish the usual GL$_2$-type and compatibility hypotheses;
they do \emph{not} establish the separate bounded descent to a
GL$_2$-type variety over $\mathbb Q$.

\paragraph{The metric comparison, with its exact additive budget}
Put $t=32s$, $(a_1,a_2,a_3,a_4)=(0,32,76832,-343)$ and
$P(t)=\prod_i(t-a_i)$. The sections
\[
 q_0=(dt)^2/P(t),\qquad q_1=tq_0
\]
are holomorphic bicanonical sections on $X$, with no common zero. At a
finite branch point use $t-a_i=z^2$; at infinity use $t=z^{-2}$.
At infinity $q_0$ vanishes to order two and $q_1$ does not vanish.

For each prime let $\delta_p=\max_{i<j}v_p(a_i-a_j)$, so
\[
 J=\prod_p p^{\delta_p}=2^{10}3^2 5^3 7^4=2765952000.
\]
Both $J^2q_i$ extend across the maximal regular stack. Indeed, at the
generic point of a vertical component semistability makes $p$ a
uniformizer, and the relative differential satisfies $v(du)\ge v(u)$.
Writing $r_i=v(t-a_i)$ gives $v(dt)\ge\max r_i$. If $v(t)<0$, all $r_i$
equal $v(t)$ and both sections have nonnegative valuation. Otherwise all
$r_i\ge0$ and at most one exceeds $\delta_p$. In either case
$2\max r_i-\sum r_i\ge-2\delta_p$. This proves the required vertical
integrality; regularity handles codimension two.

For the archimedean estimate write $r_i=|t-a_i|$, $U_0=1+\sum_i r_i$ and
\[
 ds_\rho^2=\frac{U_0}{\prod_i r_i}|dt|^2,\qquad
 C_{\rm met}=16(12\cdot76832+9)^3=12540153541944538512.
\]
Its pullback is smooth and positive at every branch point and infinity.
The curvature formula, $\Delta U_0=\sum_i r_i^{-1}$ and
$|\nabla U_0|^2\le16$ give
\[
 -K_\rho\ge\frac{\sum_i\prod_{j\ne i}r_j}{2U_0^3}.
\]
When $|t|\le2\cdot76832+2$, at most one $r_i<1/2$, and
$U_0\le12\cdot76832+9$. The displayed bound is at least
$C_{\rm met}^{-1}$. Outside this disk, $|t|/2\le r_i\le3|t|/2$
and $U_0\le7|t|$, giving the stronger bound $1/1372$.
The maximum principle against the curvature $-1$ metric implies
$ds_\rho^2/C_{\rm met}\le ds_{\rm hyp}^2$. Yuan's Petersson norm is
$\|d\tau\|=2\operatorname{Im}\tau$. Consequently
\[
 \|q_0\|_{\rm Pet}\le4C_{\rm met}/U_0,\qquad
 \|q_1\|_{\rm Pet}\le4C_{\rm met}|t|/U_0\le4C_{\rm met}.
\]

Yuan's isometry \cite[Theorem 1.1 and metric definitions, pp.~2--3]{Yuan2024} identifies the square of the Hodge bundle with
$210\omega_{\mathscr X}^{\otimes2}$ with the induced metric. The Hodge
height in his normalization equals the original stable Faltings height
plus $\log\pi$ in dimension two. Thus
\[
 h_{\bar\omega^{\otimes2}}(y)=2h_F(A_y)+2\log\pi+\log210.
\]
The sign of $\log210$ is positive: the sublattice $210\omega^{\otimes2}$
has height $h_{\bar\omega^{\otimes2}}-\log210$. Evaluate on an extension
where the point extends to the proper maximal stack; its universal
abelian scheme computes stable height. The integral sections and the
supremum estimate yield
\[
 h(s)\le2h_F(A_y)+2\log\pi+\log210+2\log J
                       +\log(4C_{\rm met})+\log32.
\]
The integer comparison
$210J^2(4C_{\rm met})32\cdot16<2^{160}$, together with $\pi<4$ and
$\log2<3/4$, bounds the additive term by $120<128$. Therefore
\[
 h(s)\le2\max(0,h_F(A_y))+128.
\tag{G}\label{eq:h5-G}
\]


% Proposed repair of the original H5 arithmetic step. Independent review pending.
% Uses the existing baseline bibliography keys, plus GaudronRemond2023 and
% the published 2021 edition of Voight, chapter 43.
\paragraph{A level bound for the non-CM fibers}
\label{sec:h5-level-repair}

This argument allows the scalar character order and coefficient field to
vary. It does not assert bounded same-field descent to a rational
GL$_2$-type variety. Let $s$ be a rational point outside the five contacts,
let $L/\mathbb Q$ be the normal fiber field constructed above, and let
$A/L$ be a non-CM fiber surface. Thus $[L:\mathbb Q]\mid576$, all
geometric homomorphisms between conjugate fibers are defined over $L$,
and $\operatorname{End}^0(A_{\overline{\mathbb Q}})=B$, the quaternion
division algebra of discriminant $210$. Put
\[
 R=B_*M(s),\qquad T=\{p:p\mid R\}.
\]
The geometric argument gives that $L$ is unramified outside $T$ and
$A$ has good reduction outside the primes above $210$. All conjugate
QM structures lie in the same normalizer orbit.

\paragraph{Good reduction after full torsion.}
First a QM surface with the indicated division action has potentially
good reduction at every finite place. Indeed, pass to an extension with
semistable reduction, using \cite[Theorem 3.4, p.~406]{silverbergzarhin}.
Every endomorphism extends uniquely to the N\'eron model by its mapping
property. On the identity component of the geometric special fiber it
preserves the maximal torus $\mathcal T$: this torus is the maximal
connected affine subgroup, and every homomorphism from a torus to an
abelian variety is zero. If its rank $t$ is positive, restriction gives
a unital homomorphism
\[
 B\longrightarrow\operatorname{End}^0(\mathcal T)=M_t(\mathbb Q),
 \qquad 1\le t\le2.
\]
The kernel is a two-sided ideal of the simple algebra $B$, and the map
is unital, so it is injective. For $t=1$ dimensions prohibit this;
for $t=2$ it would be an isomorphism from a division algebra to
$M_2(\mathbb Q)$, also impossible. Hence $t=0$. The semistable identity
component is abelian, which is good reduction. These N\'eron-model and
reduction conventions are those of \cite[\S3, pp.~405--406]{silverbergzarhin}.

Set $E=L(A[11])$. We next justify that $E/\mathbb Q$ is Galois.
For each $\sigma$, the normalizer relation supplies a compatible
isogeny ${}^{\sigma}A\to A$ of degree prime to $11$. Here is the
prime-to-$11$ detail. In the analytic description
$A_z=\mathbb C^2/(\mathcal O(z,1)^t)$ of the published 2021 edition
of Voight \cite[equations (43.1.2)--(43.1.3)]{voight2021}, a positive normalizer element
$b=\left(\begin{smallmatrix}a&b'\\c&d\end{smallmatrix}\right)$
gives the scalar compatible map $(cz+d)I:A_{bz}\to A_z$.
When $\mathcal O b\subset\mathcal O$, its degree is
$[\mathcal O:\mathcal O b]=\operatorname{Nrd}(b)^2$.
At $11$ the algebra is split and the maximal-order normalizer is
$\mathbb Q_{11}^*\operatorname{GL}_2(\mathbb Z_{11})$.
Rescale $b$ rationally to make it an $11$-adic unit, then clear the
remaining denominators by an integer prime to $11$. This produces
the claimed isogeny. Changes of level marking use isomorphisms.
All these geometric homomorphisms are $L$-defined by construction.
Consequently the $11$-torsion representations of $A$ and its conjugates
over $L$ have the same kernel. Thus every conjugate of $E$ equals $E$.

The faithful action on $A[11]$ now gives
\[
 [E:\mathbb Q]\le576|\operatorname{GL}_4(\mathbb F_{11})|
 <2^{10}11^{16}<2^{66}.
\tag{LR1}
\]
For residue characteristic different from $11$, full rational
$11$-torsion gives semistable reduction over $E$ by
\cite[Theorem 3.5, p.~406]{silverbergzarhin}. The preceding torus
argument makes this good reduction. In particular it applies at
each of $2,3,5,7$.

\paragraph{The scalar lift with controlled support.}
Choose $B$-compatible isogenies $\mu_\sigma:{}^\sigma A\to A$ over $L$,
with $\mu_1=1$, and inflate this finite choice to $G_{\mathbb Q}$.
Their discrepancies $c(\sigma,\tau)$ commute with $B$, so lie in
$\mathbb Q^*$. Put $a_\sigma=(\deg\mu_\sigma)^{-1/4}>0$ and
$\nu_\sigma=a_\sigma\mu_\sigma$, with trivial Galois action on the
added algebraic coefficients. Since a scalar $c$ has degree $c^4$,
the discrepancies of the $\nu_\sigma$ form a sign cocycle
$\varepsilon\in Z^2(G_{\mathbb Q},\{\pm1\})$, inflated from
$\operatorname{Gal}(L/\mathbb Q)$.

The sign class is unramified, hence zero, locally outside $T$.
At each odd prime where its Brauer invariant is $1/2$, choose an odd
Dirichlet character modulo $p$; at $2$ use the odd character modulo
$4$ if required. Their product $\chi$ is tame at every odd prime,
has local conductor at most $4$ at $2$, and has order
\[
 m\le2\prod_{\substack{p\in T\\p\ne2}}(p-1)\le2R.
\]
For constant coefficient groups, the boundary in
$1\to\{\pm1\}\to\mu_\infty\xrightarrow{2}\mu_\infty\to1$
is the obstruction to a character square root. Local reciprocity
identifies this obstruction with the value at $-1$: the cyclic
torsion factor has precisely this obstruction, while the valuation
and principal-unit factors admit character square roots.
Thus the boundary of $\chi$ equals $[\varepsilon]$ at every finite
place. Brauer reciprocity and the global-to-local injection give
equality globally. A pointwise square root, corrected by a sign
cochain, gives
\[
 b:G_{\mathbb Q}\longrightarrow\mu_{2m},\qquad
 b^2=\chi,\qquad\delta b=\varepsilon.
\tag{LR2}
\]
Only this \emph{sign} cochain is asserted to take root-of-unity values;
the full scalar splitting also involves the $a_\sigma$.

On $G_L$, $b$ is a conjugation-invariant character. Outside $T$ its
inertia characters are quadratic. Only finitely many are nontrivial;
twist by the quadratic character of the product of those odd primes
to kill them. At each odd prime there is a unique nontrivial quadratic
inertia character. This changes no square or coboundary in (LR2),
and introduces possible ramification only at $2\in T$.
The kernel field $L'$ of $b|G_L$ is normal over $\mathbb Q$ and
\[
 [L':\mathbb Q]\le2m[L:\mathbb Q]\le2304R.
\tag{LR3}
\]

\paragraph{The compatible modular representations.}
Choose a number field containing all normalization scalars, the values
of $b$, and a splitting field of $B$. For each coefficient prime $\ell$
and an embedding into $\overline{\mathbb Q}_\ell$, set
\[
 \rho_\ell(\sigma)=b(\sigma)^{-1}a_\sigma\mu_\sigma\circ\sigma.
\]
The cocycle identity makes this a representation commuting with $B$;
after splitting $B$, matrix units give $\rho_\ell=r_\ell\oplus r_\ell$
with $\dim r_\ell=2$. On $G_{L'}$, $\rho_\ell$ is the Tate action of
$A$. Faltings' semisimplicity and endomorphism theorems and $\operatorname{End}^0(A)=B$
show that the multiplicity representation is absolutely irreducible.
Proper-smooth comparison and finite descent, as in
Lemma~\ref{lem:de-rham-descent}, make $r_\ell$ de Rham with Tate
Hodge--Tate weights $\{0,-1\}$; Berger's monodromy theorem makes it
potentially semistable.

It is odd. On Betti homology complex conjugation followed by the
holomorphic isogeny $\mu_c$ exchanges the two Hodge summands, each
a simple $B\otimes\mathbb C$-module. Its action on the multiplicity
space is therefore nonscalar. After the scalar correction its square
is one, so its eigenvalues are $1,-1$. The representation is unramified
outside $T\cup\{\ell\}$, by (LR2) and the reduction and field properties.
Pan's regular two-dimensional modularity theorem applies at $\ell=5$.
We use its published formulation in \cite[Section 7, p.~624]{calegari-icm2022},
with the definition of a geometric representation on pp.~611--612,
and the detailed author proof in \cite[Theorem 1.0.4]{pan2022}.
It requires no additional residual-image restriction at this prime and
identifies $r_5$ with a cuspidal modular representation up to a character
twist. The published formulation is used as stated; no byte or text
identity between Pan's author version and the journal version is asserted.

The twist can be taken finite and the weight is two, without invoking
a separate local-algebraicity theorem for characters. The twisting
character is a de Rham subrepresentation of the relevant Hom
representation. Comparing its integer weight with $\{0,-1\}$ and
the modular weights $\{0,1-k\}$ forces $k=2$ and twist weight zero.
On $G_{L'}$, the identity
$(\det r_5)^2=\det V_5(A)=\chi_{\rm cyc}^2$ shows that
$\det r_5/\chi_{\rm cyc}$ has finite image, since $G_{L'}$ has finite
index in $G_{\mathbb Q}$. The modular determinant
divided by $\chi_{\rm cyc}$ also has finite image. The square of the
twisting character consequently has finite image, so the character
itself has finite image. Absorb it in the modular form and take the
primitive weight-two newform $f$ of level $N$.

The conductor at $5$ is not omitted. For each prime
$p\notin T\cup\{3,5\}$, reduction of $\mu_\sigma$ composed with
relative $p$-Frobenius is a quasi-endomorphism of the good reduction
of $A$ at a chosen place. After clearing denominators its characteristic
polynomial is independent of the Tate prime. Half its trace, multiplied
by $b(\sigma)^{-1}a_\sigma$, is therefore one algebraic number giving
the trace of every $r_\ell$. Enlarge the coefficient field to contain
the coefficients of $f$. The equality with the trace of $f$ at five
is an equality of algebraic numbers; transporting it to an embedding
at three and applying Chebotarev and semisimplicity identifies $r_3$
with the three-adic realization of $f$. Thus use $r_3$ to compute the
conductor at $5$ and $r_5$ at every other prime. The modular level is
the conductor of these realizations by
\cite[\S\S0.3--0.7, Theorem A, pp.~409--411]{carayol}
and newform theory \cite[Theorems 1 and 4]{casselman}.
Here the totally real field is $\mathbb Q$, of odd degree one, so
Carayol's extra finite-place hypothesis for even degree does not apply.
His unramified normalization twist and contragredient preserve the
conductor; his theorem allows arbitrary nebentypus.
Li's Fricke relation and Shimura's
Jacobian construction, as in the baseline modularity argument, also
place $r_5$ in the Tate module of $J_1(N)$.

\paragraph{The primitive level equals the local conductor.}
Let $f\in S_2(\Gamma_0(N),\chi)$ be the primitive newform obtained
after absorbing the finite character twist, and let $\Pi_f$ be its
unitary cuspidal automorphic representation. The character $\chi$
is arbitrary; its conductor need only divide $N$.
For every prime $p$, write $a(\Pi_{f,p})$ for the local conductor
exponent. Then
\[
 a(\Pi_{f,p})=v_p(N).
\tag{CL1}
\]
Here is the exact published newform input. In
\cite[Theorem~1, p.~302]{casselman}, the local conductor ideal is
the largest ideal $\mathfrak c$ for which the appropriate central
character eigenspace under $\Gamma_0(\mathfrak c)$ is nonzero;
at that ideal the eigenspace is one-dimensional. The remark on
pp.~306--307 identifies this ideal with the exponent in the local
epsilon factor. Theorem~4 on p.~313 gives the global level spaces:
they vanish unless $\mathfrak c(\Pi)\mid M$, and at
$M=\mathfrak c(\Pi)\prod_p p^{b_p}$ their dimension is
$\prod_p(b_p+1)$. The paragraph immediately following that theorem
identifies all primitive newforms with the minimal level
$M=\mathfrak c(\Pi)$. This proves (CL1).

This statement includes nebentypus. Casselman's classical setup on
p.~309 allows every character $\epsilon$ modulo $M$ and uses
$f|_k\gamma=\epsilon(a)^{-1}f$ for
$\gamma=\left(\begin{smallmatrix}a&b\\c&d\end{smallmatrix}\right)
\in\Gamma_0(M)$. Since $ad\equiv1\pmod M$, this is exactly
$\epsilon(d)f$. Thus the usual nebentypus convention is already
covered. No squarefree-level hypothesis occurs.

For a coefficient characteristic $\ell\ne p$, Carayol's published
Theorem~A gives the full local Weil--Deligne compatibility, not only
an equality at unramified primes
\cite[\S\S0.3--0.7, pp.~409--411]{carayol}.
Its base field here is $\mathbb Q$, its real weight is two, and its
extra finite discrete-series-place condition applies only to even
degree totally real fields, so imposes no additional hypothesis
here. Theorem~A includes the dyadic cases; the weaker Theorem~B
is not the result being used.

Carayol normalizes reciprocity by sending geometric Frobenius to
a uniformizer. His Hecke correspondence in \S0.5 differs from
local Langlands by a contragredient and an unramified norm twist.
It is characterized there by equality of local $L$- and
epsilon-factors, including character twists. The conductor is
unchanged by either duality or an unramified twist. In particular,
the passage to the Tate-module convention used for $r_\ell$ changes
no conductor exponent. Equality of good Frobenius traces identifies
the global representation supplied by Carayol with the already
constructed irreducible modular representation, in this convention.
This uses Chebotarev and Brauer--Nesbitt; it does not use the
unpublished irreducibility observation in Carayol's remark after
Theorem~A. Since the matching representation is irreducible, its
matching semisimplification also forces the original representation
to be irreducible.

The epsilon-factor conductor in (CL1) therefore equals the Artin
conductor of $r_\ell|_{G_{\mathbb Q_p}}$:
\[
 a_p(r_\ell)=a(\Pi_{f,p})=v_p(N),\qquad p\ne\ell.
\tag{CL2}
\]
The Artin conductor here includes its tame term and Swan term;
equivalently one uses the full Weil--Deligne conductor, including
the monodromy contribution. One must not discard monodromy by
replacing the local representation with its semisimplification on
inertia. Carayol's compatibility supplies precisely the required
full local object. His epsilon-factor normalization and
Casselman's conductor-zero additive character give the same
exponent; changing an additive character only introduces the usual
common dimension-dependent shift, so no such change is needed.

Apply (CL2) with $\ell=5$ at every $p\ne5$ and with $\ell=3$ at
$p=5$. The companion proof identifies both $r_5$ and $r_3$ with
the realizations of this same $f$. Thus the asserted conductor
bounds control every prime divisor of the primitive level, even
when $p^2\mid N$. No compatibility theorem at $p=\ell$ is needed.

\paragraph{The conductor bound.}
Let $F=E\mathbb Q(\chi)$, where $\mathbb Q(\chi)$ means the kernel
field of the character, not its value field. On $G_F$, $b$ is a
conjugation-invariant quadratic character. Its kernel gives a Galois
extension $F'/\mathbb Q$ of degree at most two over $F$.
For $p=2,3,5,7$, let $e_E,e_{F'}$ denote local ramification indices
over $\mathbb Q_p$. From (LR1), $v_p(e_E)\le65$.
At odd $p$ both added extensions are tame, whereas at $2$ each
adds ramification index at most two. Consequently
\[
 v_p(e_{F'})\le65\ (p=3,5,7),\qquad v_2(e_{F'})\le67.
\tag{LR4}
\]
The possibly unbounded tame ramification index causes no problem.
For a finite Galois extension of $\mathbb Q_p$ with ramification
index $e$ and different exponent $\delta$, an Eisenstein polynomial
over the maximal unramified subfield gives
\[
 \delta\le e-1+ev_p(e).
\]
Indeed the derivative terms have distinct valuations modulo $e$,
and its leading derivative term has the stated valuation. In upper
numbering the different formula is
\[
 \delta/e=1-1/e+\int_0^\infty(1-|G^u|^{-1})\,du.
\]
Before the last positive break the integrand is at least $1/2$.
Thus (LR4) bounds this break for $F'/\mathbb Q_p$ by $136$.
These ramification and conductor conventions are those of
\cite[\S\S4.1--4.3, pp.~155--159]{serrelocal}.
Over $F'$ the scalar correction is trivial and $A$ has good reduction.
Above that break, therefore, the two-dimensional representation is
trivial on upper wild groups. Its Swan conductor is at most $272$,
and its tame conductor is at most two. Hence
\[
 a_p(r)\le274<300\qquad(p=2,3,5,7).
\]
For $p>7$, $L/\mathbb Q$ is tame because its degree divides $576$,
$A/L$ has good reduction, and $b$ is trivial on wild inertia:
on that pro-$p$ group its square is tame and its remaining image has
order at most two. Thus $a_p(r)\le2$ there; it is zero outside $T$.
Using the coefficient-prime comparison above proves
\[
 N\le210^{300}R^2.
\tag{LR5}
\]

\paragraph{The original constant is preserved.}
Apply the accompanying modular-Jacobian transfer calculation with
$C=210^{300}$ and (LR3)--(LR5). It gives
\[
 h_F(A)<2^{24}210^{3600}R^{24}
 <2^{28824}R^{24}<4608^{165888}R^{24}.
\tag{LR6}
\]
The last comparison uses $210<2^8$, $4608>2^{12}$, and
$28824<1990656=12\cdot165888$. Together with the separately proved
CM bound and geometric height comparison, (LR6) yields exactly
\[
 h(s)\le
 \bigl(2(4608^{165888}+6^{100})B_*^{24}+128\bigr)M(s)^{24}.
\]
No contact, coefficient constant, or radical exponent has been changed.
% UNAPPLIED CANDIDATE: global review and published Pan match are pending.

% Actual level-based repair: no uniform coefficient-field bound is used.
% The modular representation and level estimate are proved in the companion
% height-level-repair fragment. This file proves its height consequence.
\paragraph{Transfer from the level, allowing unbounded coefficient degree.}
Let $A/L$ be the non-CM quaternionic surface. Write $R\geq1$ for the
squarefree product of the prescribed rational bad primes, including $2$,
and assume $[L:\mathbb Q]\leq576$. The modularity construction provides
a weight-two newform of level $N$ and a field $L'$ over which the Tate
modules of $A$ and $J=J_1(N)$ share an irreducible constituent. In the
notation of the sign-splitting construction, the auxiliary Dirichlet
character has order
\[
 m\leq 2\prod_{\substack{p\mid R\\p\ne2}}(p-1)\leq2R.
\]
The sign-splitting cochain has values in $\mu_{2m}$. If one instead
uses the full scalar splitting map, it also includes the square root
of the degree map; that factor equals one on $G_L$. Thus in either
notation the restriction to $G_L$ is a character valued in
$\mu_{2m}$, and its kernel field satisfies
\[
 d'=[L':\mathbb Q]\leq2m[L:\mathbb Q]\leq2304R.
\tag{LT1}
\]
The quadratic twist used to remove extra ramification does not increase
this bound, since $2$ divides $2m$. There is no bound here on $m$
independent of $R$.

Faltings' Hom theorem, followed by scalar extension in the Hom space,
gives a nonzero $L'$-homomorphism $A\to J$. Absolute simplicity of $A$
and Poincar\'e reducibility therefore give an $L'$-isogeny
$A\times Z\to J$ for an abelian variety $Z/L'$. Put $g=\dim J$;
then $g\geq2$. The modular genus estimate
\cite[(5.9), p.~13]{vonKaenel2021} gives $g\leq N^2/24$.
The standard Belyi map has degree at most $N^2$ and
\cite[Theorem~1.1.1 and \S2.3]{javanpeykar} gives
\[
 h_F(J)\leq13000000N^{12}.
\tag{LT2}
\]
All heights in this argument use the original stable Faltings metric.

Suppose the companion local argument has proved $N\leq C R^2$,
where $C\geq1$. Then
\[
 g\leq C^2R^4/24,\qquad
 h_F(J)\leq13000000C^{12}R^{24}.
\]
Set $l=\log C$ and $r=\log R$. The published
Gaudron--R\'emond Theorem~1.9(1) \cite[pp.~7--8]{GaudronRemond2023}, combined with half-logarithmic isogeny-height variation,
gives
\[
 h_F(A)+h_F(Z)\leq h_F(J)+V,
 \quad
 V=2g^2\left(8g^2\log(7g)+\log d'
       +\log\max(1,\log d',h_F(J))\right).
\tag{LT3}
\]
Here the source variety for the isogeny estimate is $J_{L'}$.
The three logarithms in (LT3) satisfy respectively
\[
 \log(7g)\leq2l+4r,\quad
 \log d'<8+r,\quad
 \log\max(1,\log d',h_F(J))<17+12l+24r.
\]
For the last assertion, the bound in (LT2) also dominates $1$ and
$\log d'$; use $8+r\leq9R^{24}$.
It follows that
\[
 V\leq
 \frac{C^8R^{16}}{24^4}(32l+64r)
 +\frac{C^4R^8}{288}(25+12l+25r).
\tag{LT4}
\]
Using $l\leq C^4/4$, $r\leq R^8/8$, and also
$r\leq R^{16}/16$ in the second summand gives
\[
 \frac{V}{C^{12}R^{24}}
 \leq\frac{16}{24^4}+\frac{25+3+25/16}{288}<1.
\tag{LT5}
\]
The Bost lower bound in \cite[(2.2), p.~5]{vonKaenel2021}
gives $h_F(Z)\geq-\tfrac12\dim Z\log(2\pi^2)>-\tfrac32g$.
Consequently
\[
 h_F(A)<(13000000+1+1)C^{12}R^{24}
       <2^{24}C^{12}R^{24}.
\tag{LT6}
\]
This calculation needs no uniform bound on the degree of the newform's
coefficient field and no assertion that all its geometric factors are
isogenous to $A$.

For the companion conductor estimate one may take $C=210^{300}$.
Since $210<2^8$ and $4608>2^{12}$,
\[
 2^{24}C^{12}=2^{24}210^{3600}
 <2^{28824}<2^{1990656}<4608^{165888}.
\tag{LT7}
\]
Thus the non-CM height inequality is strictly within the original
coefficient budget, once the stated modularity and level estimate have
been established. In the five-contact application $R=B_*M(s)$, so
the radical exponent is exactly $24$, as required.

% Original 15671 architecture: independent CM branch, version 1.
% This lemma is conditional on the separately audited field/ramification lemma.
\paragraph{The CM height estimate over the same field}
Let $A/L$ be an abelian surface of CM type, let $B$ be a quaternion
division algebra over $\mathbb Q$ embedded in
$\operatorname{End}^0_{\overline L}(A)$, and suppose that all geometric
endomorphisms of $A$ are defined over $L$. Assume
\[
 d=[L:\mathbb Q]\le576,\qquad
 \operatorname{rad}(|D_L|)\le R,\qquad R\ge1.
\]
All heights below are stable Faltings heights in Faltings's original
normalization. We prove
\begin{equation}\label{eq:original-cm-bound}
 h_{\rm F}(A)<854R<6^{100}R^{24}.
\end{equation}
No assertion about the existence or degree of $L$ is included in this lemma.

First, $A_{\overline L}$ is isogenous to $E^2$ for a CM elliptic curve
$E$. Here is a dimension argument which makes the alternative cases
explicit. If $A$ were geometrically simple, its division endomorphism
algebra would act faithfully on the four-dimensional rational homology.
Its dimension therefore divides four. As it contains $B$, it would equal
$B$; but $B$ has no commutative semisimple subalgebra of dimension four,
contrary to CM type. If the two elliptic factors were not isogenous,
their endomorphism algebra would be a product of commutative fields and
could not contain $B$. Finally, if $E$ had no CM, the endomorphism algebra
would be $M_2(\mathbb Q)$, and the embedding of the four-dimensional
division algebra $B$ would be an impossible isomorphism.

The elliptic factors and the isogeny may be taken over $L$. Indeed,
the rational idempotents in $\operatorname{End}^0(A)$ are defined over
$L$, and the image of an integral multiple of such an idempotent is an
elliptic subvariety over $L$. Endomorphisms and homomorphisms of the
resulting factors extend to rational endomorphisms of $A$ by the
inclusions and rational projections; these are fixed by
$\operatorname{Gal}(\overline L/L)$. Thus the factor isogenies, and all
endomorphisms of $E$, are defined over $L$.

Put $F=\operatorname{End}^0(E)$, an imaginary quadratic field, and
$\mathcal O=\operatorname{End}(E)$. Differentiation gives an injective
unital map $F\longrightarrow\operatorname{End}_L(\operatorname{Lie}E)=L$:
it is nonzero, and every nonzero elliptic endomorphism in characteristic
zero is separable. Consequently $F\subset L$.

We next pass to the maximal order without extending $L$. Choose
$f\ge1$ with $f\mathcal O_F\subset\mathcal O$, choose an integral
basis $\alpha_1,\alpha_2$ of $\mathcal O_F$, and form the finite
subgroup scheme
\[
 K=\sum_{i=1}^2(f\alpha_i)(E[f])\subset E.
\]
Each $f\alpha_i$ is defined over $L$, so $K$ and the quotient
$E'=E/K$ are defined over $L$. Under a complex uniformization
$E(\mathbb C)=\mathbb C/\Lambda$, this subgroup is
$\mathcal O_F\Lambda/\Lambda$; thus $E'(\mathbb C)$ has lattice
$\mathcal O_F\Lambda$ and endomorphism ring $\mathcal O_F$.
The ring cannot be larger, since the rational endomorphism algebra
remains $F$. This is also a direct special case of the same-field
maximal-order theorem of R\'emond, but the construction just given is
the proof used here. It follows that $A'=(E')^2$ is $L$-isogenous to $A$.

Write $\Delta_F$ for the fundamental discriminant of $F$. The
discriminant tower formula implies that every rational prime dividing
$\Delta_F$ divides $D_L$. A quadratic fundamental discriminant has
odd prime exponents one and exponent at two at most three. Therefore
\[
 |\Delta_F|\le4\operatorname{rad}(|D_L|)\le4R.
\]
The multiplier four includes the possible factor $2^3/2$.

For completeness we give an elementary analytic bound on $h_{\rm F}(E')$.
CM elliptic curves have potentially good reduction. After a finite
extension with everywhere good reduction, the elliptic discriminant
formula of Silverman, reproduced as equation (6.2) in
\cite{cmheight}, gives
\[
 h_{\rm F}(E')=\frac1{[L':\mathbb Q]}\sum_{\sigma:L'\hookrightarrow\mathbb C}
 \left(-\log(2\pi)-\log|\eta(\tau_\sigma)|^2
                     -\tfrac12\log\operatorname{Im}\tau_\sigma\right).
\]
Here $L'$ is used only to evaluate the stable height; it is not the
field used in the isogeny estimate below. The discriminant term is zero
because reduction is good. We use the individual-embedding formula
(6.2), not the class-average assertion of Proposition 6.1, so no
hypothesis on Weierstrass coefficients lying in $\mathbb Q(j)$ is needed.

For each embedding choose a reduced quadratic parameter of discriminant
$\Delta_F$. If $y=\operatorname{Im}\tau$, then
\[
 \frac{\sqrt3}{2}\le y=\frac{\sqrt{|\Delta_F|}}{2a}
                   \le\frac{\sqrt{|\Delta_F|}}2,
 \qquad a\in\mathbb Z_{\ge1}.
\]
Put $r=|e^{2\pi i\tau}|<1/16$. The product defining $\eta$ gives
\[
 -\log|\eta(\tau)|^2
 \le\frac{\pi y}{6}+2\sum_{n\ge1}\frac{r^n}{1-r^n}
 \le\frac{\pi y}{6}+\frac{2r}{(1-r)^2}
 <\frac{\pi\sqrt{|\Delta_F|}}{12}+\frac{32}{225}.
\]
Also $-\tfrac12\log y\le\tfrac14\log(4/3)<1/12$.
Discarding the negative term $-\log(2\pi)$ and using $\pi<4$ shows
\[
 h_{\rm F}(E')<\sqrt{|\Delta_F|}+1\le2\sqrt R+1\le3R.
\]
The height is additive on products, hence $h_{\rm F}(A')<6R$.

Apply \cite[Theorem 1.9(1)]{GaudronRemond2023} over $L$, with source
$A'$ and target $A$ and with no extension of the base field.
For a surface, its quantity $\Xi(A')$ is
\[
 \Xi(A')=\left(14^{32}d\max(1,\log d,h_{\rm F}(A'))\right)^8.
\]
The general assertion supplies an isogeny $A'\to A$ of degree at most
$\Xi(A')^2$; no identification of CM type with that paper's
separate ``MM'' condition is made or needed. The isogeny-height
inequality then gives
\[
 h_{\rm F}(A)\le6R+
 8\log\left(14^{32}\,576\max(1,\log576,6R)\right).
\]
One can see the last inequality directly from invariant differentials:
after semistable base change, pullback on determinant differentials has
integral cokernel, while its archimedean squared norm is multiplied by
the isogeny degree. Thus the height difference is half the logarithm
of that degree minus a nonnegative normalized cokernel length.

Finally, $\log14<3$, $\log576<7$, and
$\max(1,\log576,6R)\le7R$. Since $\log7<2$ and $\log R\le R$,
\[
 h_{\rm F}(A)<6R+8(96+7+2+\log R)
             \le854R.
\]
As $6^4=1296>854$ and $R\ge1$, this proves
\eqref{eq:original-cm-bound}. In particular the original
$6^{100}R^{24}$ allowance absorbs this CM estimate with the field
degree still exactly bounded by 576.

\paragraph{The unconditional numerical certificate and scalar detector}
The specified integer $H$ satisfies
\[
 \log_2H<4783227.
\]
Indeed the exact sieve enumerates $7499$ primes in $S_*$, with largest
prime $76163$, and their product has bit length $109445$. Hence
$B_*<2^{109445}$. With $A_0=4608^{165888}$, we have
$6^{100}<A_0$ and $128<A_0$, whence $H<8A_0B_*^{24}$. Since $4608<2^{13}$,
\[
 \log_2H<3+165888\cdot13+24\cdot109445=4783227.
\]
This is a check of the number H, not a proof that H is a height constant.

For completeness, the scalar product used by the original polynomial is
\[
 r(\beta)=\beta(\beta-1)(\beta-2401)(32\beta+343)(2\beta-1).
\]
At a prime $p>76176$, if $v_p(\beta)<0$, its valuation is
$5v_p(\beta)$ and so detects odd contact with infinity. If $\beta$ is
$p$-integral and has odd contact with a finite member of $\mathcal C$,
only that factor has positive valuation; every other factor is a unit.
The extra root $1/2$ introduces only exceptional collision primes
(notably factors of $4801$ and $718$). Thus every prime of $M(\beta)$ has
odd valuation in $r(\beta)$. If a separate comparison proves
$M(\beta)^{48}\mid s\ne0$, the five-contact estimate just proved gives
$h(\beta)\le H\sqrt{|s|}$.

Combining the non-CM and CM estimates with \eqref{eq:h5-R} and
\eqref{eq:h5-G} gives
\[
 h(s)\le 2(4608^{165888}+6^{100})(B_*M(s))^{24}+128
       \le H_*M(s)^{24},
\]
because $M(s)\ge1$. The numerical certificate above proves the
remaining assertion about $H_*$.
\end{proof}


Define $\Psi(b)$ by seven rational witnesses $t,U,e,p,q,s,h$ and four
equations, with $B=17b$:
\begin{align}
 U(1+bt^2)&=b(1-bt^2),&
 p^2&=B(1-e)(1+3e),\label{eq:o8-pole}\\
 h&=e(B-U^2),&
 Bq^2&=(s^2-Bh)(s^2+3Bh).\notag
\end{align}

\begin{lemma}[Exact seven-variable pole statement]\label{hyp:o8-pole}
If $b\ne0$ and $\Psi(b)$, then $v_p(b)$ is not a negative odd integer
for every $p\notin\{2,3,17\}$. Conversely, $\Psi(b)$ holds whenever
$b>0$, the denominator of $b$ in lowest terms is a square, and $b$ is
a square in each of $\Q_2,\Q_3,\Q_{17}$.
\end{lemma}
\begin{proof}
These are Lemmas~\ref{lem:original-seven-parity-soundness} and
\ref{lem:original-seven-parity-completeness}. Their proof gives the
simultaneous two-twist construction explicitly, and applies the existing
Theorem~\ref{hyp:RC} and Lemma~\ref{lem:corank}. It introduces no new
simultaneous-realization assumption.
\end{proof}

The lemma includes the degenerate values of the seven witnesses
permitted by \eqref{eq:o8-pole}. No inverse condition on $e$, $h$, $s$,
or $U$ may be inserted into this particular polynomial. For $b\ne0$,
the first equation itself implies $1+bt^2\ne0$. The predicate is soluble
at $b=0$; the norm argument below accounts for that case where necessary.

% Additive proof for the original 15671-variable polynomial. No circuit changes.
\subsubsection{The seven-witness parity predicate and its finite correction}
\label{subsec:original-seven-parity}
Put $B=17b$. The predicate used here is exactly the existence of
$t,U,e,p,q,s,h\in\mathbb Q$ satisfying
\begin{align}
 U(1+bt^2)&=b(1-bt^2),\label{eq:seven-param}\\
 p^2&=B(1-e)(1+3e),\label{eq:seven-first}\\
 h&=e(B-U^2),\label{eq:seven-h}\\
 Bq^2&=(s^2-Bh)(s^2+3Bh).\label{eq:seven-second}
\end{align}
Denote it by $\Psi(b)$. No nonzero conditions are added to this formula.

\begin{lemma}[The exact one-sided assertion]
\label{lem:original-seven-parity-soundness}
If $b\ne0$ and $\Psi(b)$, then $v_\ell(b)$ is not a negative odd
integer for any prime $\ell\notin\{2,3,17\}$.
\end{lemma}
\begin{proof}
Suppose $k=v_\ell(b)<0$ is odd. Equation~\eqref{eq:seven-param}
forces $1+bt^2\ne0$: if it vanished, its right side would be $2b\ne0$.
If $t=0$, then $U=b$. Otherwise $v_\ell(bt^2)$ is odd and nonzero,
so $(1-bt^2)/(1+bt^2)$ is a unit. Thus $v_\ell(U)=k$ in all cases,
and $v_\ell(B-U^2)=2k$.

If $e=0$, equation~\eqref{eq:seven-first} gives $p^2=B$, impossible.
Write $j=v_\ell(e)$. If $j<0$, the right side of
\eqref{eq:seven-first} has valuation $k+2j$, which is odd.
If $j>0$, it has valuation $k$, again odd. Consequently $j=0$.
This includes the possibilities $p=0$, $e=1$, and $e=-1/3$.
Now $h\ne0$ and $v_\ell(Bh)=3k$ is odd. If $s\ne0$, its square
has even valuation, different from $3k$, so the two factors on the
right of~\eqref{eq:seven-second} have the same finite valuation.
Their product has even valuation and is nonzero. The same conclusion
holds when $s=0$, since the product is $-3(Bh)^2$.
The left side is either zero or has odd valuation, a contradiction.
\end{proof}

The omitted guards therefore cause no loss of the asserted soundness.
At $b=0$, the equations require $U=p=h=s=0$ and permit arbitrary
$t,e,q$. For $b\ne0$, the six-witness source formula is recovered
exactly by $e_{\rm old}=Be$, $p_{\rm old}=Bp$, $q_{\rm old}=Bq$,
with $s,t,U$ unchanged. Its first curve residual is $B^2$ times
\eqref{eq:seven-first}'s residual and its second is $B$ times
\eqref{eq:seven-second}'s residual after~\eqref{eq:seven-h}.
Division by $B$ reverses this transformation.

\begin{lemma}[Simultaneous realization from the existing restricted converse]
\label{lem:original-seven-simultaneous-realization}
Use Theorem~\ref{hyp:RC}, the restricted converse $\mathsf{RC}$ for
$E_l:y^2=x(x-l)(x+3l)$ with $l>0$ squarefree and $l\equiv11\pmod{12}$.
Let $\beta,\theta>0$ be squarefree, with disjoint supports avoiding
$2,3$, such that $\beta\equiv17\pmod{24}$,
$\theta\equiv1\pmod{24}$, and each is a local square at every prime
dividing the other. Then there is a positive squarefree $n$, prime
to $6\theta$, for which both $E_n$ and $E_{n\theta}$ have rational
points with nonzero abscissa in $\beta\mathbb Q^{\times2}$.
The case $\theta=1$ is allowed.
\end{lemma}
\begin{proof}
Here are the descent steps, so that no separate realization theorem
is being assumed. Write $B_0=\operatorname{supp}(\beta)$,
$J=\operatorname{supp}(\theta)$, and $d=|B_0|\ge1$.
For a support prime $q_i$, put $s_i=[-1/q_i]$, $\lambda_i=[3/q_i]$,
and define over $\mathbb F_2$
\[
 (Kv)_i=\sum_{j\ne i}[q_j/q_i](v_j+v_i).
\]
Add $d-1$ prime slots $I$ of type $(s,\lambda)=(0,0)$ and one slot
$q_0$ of type $(1,1)$. Let the $B_0$-by-$I$ symbol matrix have columns
$e_i+e_d$ for $1\le i<d$. Set all other symbols incident with $I$
to zero, prescribe $[q_0/q_b]$ as the corresponding row sum, and
obtain reverse symbols by reciprocity. The $q_0$--$J$ symbols remain
free. Mutual squareness gives zero row sums for both directed
$B_0$--$J$ blocks; hence $K\mathbf1_{B_0}=0$ on both supports.

Choose $a$ with ones at one fixed member of $B_0$ and at $q_0$.
The matrix $K^\sharp=K+\lambda a^t$ kills $\mathbf1$ on both sides.
After deleting $q_0$, its smaller-support matrix has the form
\[
 M_0=\begin{pmatrix}D&C\\C^t&0\end{pmatrix},
 \qquad \mathbf1^tD\mathbf1=1.
\]
The columns of $C$ span the even-sum hyperplane. If $M_0(x,y)^t=0$,
then $x=c\mathbf1$; multiplication of the first equation by
$\mathbf1^t$ gives $c=0$, and injectivity of $C$ gives $y=0$.
This also proves the case $d=1$, where $M_0=(1)$.
Adding $J$ preserves $M_0$. Each free $q_0$--$J$ bit toggles just
one new diagonal entry. The enlarged determinant is multilinear in
these bits, with full-product coefficient $\det M_0=1$, so some
Boolean choice makes it nonzero. CRT and Dirichlet's theorem realize
the prescribed symbols by distinct fresh primes, with types imposed
modulo $12$. Thus $n=\prod_{B_0\cup I\cup\{q_0\}}q$ and $n\theta$
are squarefree and congruent to $11$ modulo $12$.

The projection $T=1-\mathbf1_{B_0}a^t$ satisfies $K=K^\sharp T$.
Nonsingularity of the deleted matrix gives
$\ker K^\sharp=\langle\mathbf1\rangle$, and therefore
\[
 \ker K=\langle\mathbf1,\mathbf1_{B_0}\rangle,
 \qquad \lambda\notin\operatorname{im}K
\]
on both supports. For the latter assertion, $Kx=\lambda$ would
make $Tx-\mathbf1_{B_0}$ constant, contradicted by applying $a^t$.

For clarity, the local descent needed here is the manuscript's full
two-descent in coordinates $(x+3l,x-l)$. The two torsion classes are
$(l,1)$ and $(3,-l)$. At an odd good prime, the local image consists
of the four pairs with even valuations; at a prime dividing $l$,
the two torsion classes span it by their independent valuation
vectors. At $3$ it is the set of all first-coordinate classes with
square second coordinate. At $2$ both coordinate valuations are
even: the only problematic case $x-l=2u$, $u$ odd, would give
$(y/2)^2\equiv l\pmod8$, impossible for $l\equiv3\pmod4$.
When both differences are divisible by $4$, their difference $4l$
and the cubic equation give even valuations for both. At infinity
the first coordinate is positive, and the positive-second-coordinate
subgroup has index two. Its possible representatives consequently are
$(3^\delta\prod q_i^{u_i},\prod q_i^{v_i})$ with necessary equations
\[
 Ku+\operatorname{diag}(\lambda)v=\delta\lambda,
 \quad (K+\operatorname{diag}(s))v=0,
 \quad (\lambda+s)^tv=0.
\]
Reciprocity gives $K+K^t=ss^t+\operatorname{diag}(s)$.
Hence $\ker(K+\operatorname{diag}(s))=\ker K^t\cap s^\perp$
is a line. The functional $\lambda^t$ is nonzero on this line,
because $\lambda\notin\operatorname{im}K$ and
$\lambda^t\mathbf1=0$, whereas $s^t\mathbf1=1$.
The displayed equations force $v=0$, $\delta=0$, and $u\in\ker K$.
There are at most eight Selmer classes.
The class $(\beta,1)$ is actually local everywhere: it is the identity
at $2$ and infinity, the first torsion class at $3$, satisfies the
torsion comparison at support primes by $K\mathbf1_{B_0}=0$, and
is unramified at the other odd primes. It is independent of the two
torsion classes since $B_0$ is nonempty and $q_0\notin B_0$.
Thus each two-Selmer group has dimension exactly three.

Lemma~\ref{lem:corank} gives full two-primary Selmer
corank one. Explicitly, the Kummer sequence gives
$r+\dim\Sha[2]=1$; writing
$\Sha[2^\infty]=(\mathbb Q_2/\mathbb Z_2)^t\oplus F$, the
alternating Cassels--Tate pairing on the finite quotient makes
$\dim F[2]$ even, so $F=0$ and $r+t=1$.
Apply $\mathsf{RC}$ to $n$ and $n\theta$, and the manuscript's
published forward analytic-rank-one theorem. Both curves have rank
one and finite $\Sha$, hence $\Sha[2]=0$. Kummer exactness now
represents $(\beta,1)$ by actual rational points. These are not
two-torsion, so their abscissas are nonzero of class $\beta$.
When $\theta=1$ the same point can be used twice.
\end{proof}

\begin{lemma}[The completeness scope actually used]
\label{lem:original-seven-parity-completeness}
Assume only the existing $\mathsf{RC}$ and the published descent and
forward-rank facts used above. If $b>0$ has square reduced denominator
and is a square in $\mathbb Q_2,\mathbb Q_3,\mathbb Q_{17}$, then
$\Psi(b)$ holds.
\end{lemma}
\begin{proof}
Choose $t\in\mathbb Q$ close to a solution of $bt^2=1$ at infinity
and at $2,3,17$. In~\eqref{eq:seven-param} this makes $U$ arbitrarily
small at those places. We can and do require
$V=B(B-U^2)>0$ and $V$ locally square at all three primes.
Explicit sufficient conditions are $U^2<B$ in $\mathbb R$,
$v_2(U^2/B)\ge3$, and $v_3(U^2/B),v_{17}(U^2/B)\ge1$.
The denominator tends locally to $2\ne0$ and is positive globally.

Let $\beta,\theta$ be the positive squarefree representatives of
$B,V$. The number $B=17b$ is square at $2$, a nonsquare unit class
at $3$, and has odd valuation at $17$. Thus
$\beta\equiv17\pmod{24}$ and $\beta>1$;
$\theta\equiv1\pmod{24}$ and $17\nmid\theta$.
At any prime $\ell\mid\beta$ outside $2,3,17$, the valuation of
$b$ is odd and therefore positive, since its reduced denominator is
square. The same quotient calculation as in soundness gives
$v_\ell(U)=v_\ell(b)$; hence $1-U^2/B$ is a square unit.
This proves that $V$, and hence $\theta$, is square at every prime
of $\beta$, including $17$ by construction. The supports are disjoint.
Conversely, if $\ell\mid\theta$, then $\ell\notin\{2,3,17\}$
and $v_\ell(B)$ is even. Odd valuation of $B(B-U^2)$ forces
$v_\ell(B)=2v_\ell(U)$ and strict cancellation, so
$B/U^2\in1+\ell\mathbb Z_\ell$ is square. This proves the other
mutual-square condition. If $U=0$, then $\theta=1$ and this last
argument is unnecessary.

Apply the preceding lemma. Write $V=\theta c^2$ and rescale the
second point by $(x,y)\mapsto(c^2x,c^3y)$ to put it on $E_{nV}$.
If the first point has abscissa $Ba^2$, divide both curve parameters
by $a^2$ and the point coordinates by $a^2,a^3$. The resulting
points have abscissas $B,Bs^2$ with $s\ne0$, on $E_{e_0}$ and
$E_{e_0V}$ respectively. Write their ordinates as $y_1,y_2$.
Set
\[
 e=e_0/B,\qquad p=y_1/B,\qquad
 q=y_2/(B^2s),\qquad h=e(B-U^2).
\]
Direct substitution verifies all four equations. Every division
made in choosing these witnesses is by a nonzero number; none is
an additional equation or guard in the original polynomial.
\end{proof}

\paragraph{Finite correction and the unguarded norm root.}
At $2$ use valuation parity followed by unit coordinates $[-1],[5]$;
at $3$ and $17$ use valuation parity followed by the nonsquare-unit bit.
The columns for $2,3,5,7,11,13,17$ are exactly
\[
\begin{pmatrix}
1&0&0&0&0&0&0\\0&1&0&1&1&0&0\\0&1&1&0&1&1&0\\
0&1&0&0&0&0&0\\1&0&1&0&1&0&1\\0&0&0&0&0&0&1\\
0&1&1&1&1&0&0
\end{pmatrix}.
\]
The accompanying exact computation verifies rank seven and all 128
distinct subset signatures. Rational solutions of $e^2-e=0$ are
exactly $0,1$. Therefore the original eight bits and two products
\[
 A=(1+e_2)(1+2e_3)(1+12e_{13})(1+16e_{17}),\quad
 D=(1+4e_5)(1+6e_7)(1+10e_{11})(1-2\eta)
\]
give exactly $AD=(-1)^\eta\prod p^{e_p}$, always nonzero.
For nonzero integral $r$, choose $\delta=\operatorname{sign}(r)$
and the unique subset product $\mu$ having the same local signature
as $\delta r$. Taking $AD=\delta\mu$ makes $ADr$ positive and
locally square at $2,3,17$.

Choose $J_0$ sufficiently deep at all places over primes dividing
$\mu r$, then $S$ in a sufficiently divisible subgroup satisfying
the anchor. Since $f=z^2$ times a formal unit, the coefficients of
$f(J_0),f(S),z(S)/z(J_0)$ can attain any prescribed finite lower
valuation bounds. This remains true at ramified and bad primes by
working in sufficiently deep formal subgroups. Thus at $\ell\mid r$
all normalized leaves are integral, and at $\ell\mid\mu$ with
$\ell\nmid r$ all nonconstant leaves can be made divisible by $\ell$.
The norm tree $N_r(a,b)=a^2+rb^2$ has its constant leaf $1$ on
the all-left path; its root $W$ is a unit at these primes.
In the actual application $r=\beta(\beta-1)(\beta-2401)
(32\beta+343)(2\beta-1)$ with $\beta=X^2+2$ and $X\in\mathbb Z$.
It is nonzero and $|r|>1$ since $X^2\ne2399$ and $\beta\ge2$.
Consequently $W\ne0$, and $b=ADrW^2$ is positive, locally square
at $2,3,17$, and has square reduced denominator. Completeness applies.

For soundness no global guard $W\ne0$ is needed. At any prime where
$v_\ell(r)$ is odd, the two nonzero summands of $a^2+rb^2$ have
different valuation parity. Such a norm vanishes only if both inputs
vanish. Recursing to the constant leaf $1$ proves $W\ne0$ at every
prime relevant to an odd-contact argument. Outside the fixed
exceptional set, $AD$ is a unit and
$v_\ell(ADrW^2)=v_\ell(r)+2v_\ell(W)$.
The exact original norm argument therefore applies unchanged. If
$r$ has no relevant odd valuation, there is no local conclusion to
prove. No guard or witness has been inserted, and the original
polynomial and coefficient bound are unchanged.

\subsection{Elliptic charts and the anchor}

Use the normalized curve and point
\[
 K=\Q(\sqrt2),\quad d=3-\sqrt2,\quad D=d^2=11-6\sqrt2,
 \quad E:y^2=x^3-Dx,
\]
\[
 P_*=(57/8-7\sqrt2/2,-165/8+417\sqrt2/32),
 \qquad G=2E(K)=\Z P_*.
\]
The identification of $G$ follows from Proposition~\ref{prop:generator}
by the unit change of variables $\lambda=3+2\sqrt2$.
Write $f=1/x$ and $z=-x/y$. A nonzero point of $G$ is represented by
eight field variables $u,v,w,z,f,h,x,y$ satisfying
\begin{equation}\label{eq:o8-chart}
 f=u^2,\quad v^2+df=1,\quad w^2+v^2=2,\quad h=vw,
 \quad zh+u=0,\quad xf=1,\quad yz+x=0.
\end{equation}
Each field equality means its two rational coefficient equalities.
\begin{lemma}[The sparse point chart]\label{lem:o8-chart}
The equations \eqref{eq:o8-chart} parametrize precisely the nonzero points
of $G$, and their reciprocal functions are defined.
\end{lemma}
\begin{proof}
It gives $x=u^{-2}$, $y=vw/u^3$ and the three square root differences
$x=(1/u)^2$, $x-d=(v/u)^2$, $x+d=(w/u)^2$.
Here is a direct verification that these squares give precisely the
required doubles. Write $e_1=0,e_2=d,e_3=-d$ and choose
$r_i^2=x-e_i$ with $r_1r_2r_3=-y$. Put
\[
 B=r_1r_2+r_1r_3+r_2r_3,\qquad
 X'=x+B,\qquad
 Y'=-(r_1+r_2)(r_1+r_3)(r_2+r_3).
\]
Since $X'-e_i=(r_i+r_j)(r_i+r_k)$, the point $(X',Y')$ lies on
$E$. None of the pair sums is zero, because the $e_i$ are distinct.
Its tangent slope is $-(r_1+r_2+r_3)$, and the duplication equations
give $[2](X',Y')=(x,y)$.
Conversely, if $(x,y)=[2](a,b)$ is nonzero, then $b\ne0$ and direct
duplication gives
\[
 x-e_i=\left(
 \frac{(a-e_i)^2-(e_i-e_j)(e_i-e_k)}{2b}\right)^2.
\]
Choose the signs of these square roots to have product $y$, and recover
$u,v,w,h,z,f$ from \eqref{eq:o8-chart}. The group $G$ is torsion-free, so its nonzero points have
$x,y\ne0$. Thus the sparse chart and the membership equations in
Section~\ref{sec:model} have the same nonzero image.
The inverse relation $xf=1$ excludes zero denominators. In particular
$u,h,z,y$ are nonzero.
\end{proof}

Subtraction uses the equations
\[
 \ell x-y=j,\quad \ell a+b=j,\quad
 x'=\ell^2-x-a,\quad \ell x'+y'=j,\quad
 f'x'=1,\quad z'y'+x'=0
\]
for $(x,y)-(a,b)$, with an inverse for the rational norm of $x-a$.
Doubling uses the usual tangent formula with its nonzero denominators.

The baseline duplication estimates give
\[
 4h(x(Q))-\log392\le h(x(2Q))
       \le4h(x(Q))+\log676.
\]
Since $h(x(P_*))=\log41$, passage to the canonical height gives
$2\widehat h(P_*)\ge\log41-\frac13\log392>1$, and hence
\begin{equation}\label{eq:o8-indexheight}
 h(x(nP_*))\ge n^2-6\qquad(n\ne0).
\end{equation}
The deliberately weaker constant six suffices throughout this section.

Let $\mathcal T(a)$ be the seven-variable formula
\begin{equation}\label{eq:o8-T}
 u_0^2-2u_1^2-3u_2^2+6u_3^2=1,\qquad
 (a-2u_0)^2-8v_1^2-12v_2^2+24v_3^2=4.
\end{equation}
Then $\mathcal T(\Q)=\Z_{(2)}\cap\Z_{(3)}$. Indeed the quaternion
algebra $(2,3)$ is ramified exactly at $2,3$ and split at infinity.
Koenigsmann's trace-sum theorem, Proposition~6 of
\cite{koenigsmann-defining}, identifies the sum of two norm-one trace
sets with the intersection of these two local rings. Eliminating the
scalar coordinate of the second norm-one element gives exactly
\eqref{eq:o8-T}, with no translation of $a$.

Introduce twenty-five rational variables
\begin{equation}\label{eq:o8-chain}
 c_0=3,\qquad c_{i+1}=c_i^2\quad(0\le i<24),\qquad C=c_{24}.
\end{equation}
Thus $C=3^{2^{24}}$, but this large integer is not a coefficient in
$F_8$. For $S=sP_*\ne O$ and $f(S)=f_0+f_1\sqrt2$, impose
$Ca_i=f_i$ and $\mathcal T(a_i)$, for $i=0,1$.
At the prime above $3$, $P_*$ reduces to the nonidentity point
$(\sqrt2,0)$ of order two, and $v_3(z(2P_*))=1$.
The formal group therefore gives
$v_3(f(sP_*))=2+2v_3(s)$ whenever this valuation is positive.
The anchor forces
\begin{equation}\label{eq:o8-anchor-bound}
 |s|\ge3^{2^{23}-1}>(20H_*)^2.
\end{equation}
For the last inequality use Theorem~\ref{hyp:o8-height},
$\log_2 3>3/2$, and
$3(2^{23}-1)=25165821>4(5+4783227)=19132928$.
The allowed indices contain a nonzero subgroup: sufficiently divisible
indices make $f(sP_*)/C$ integral at both $2$ and $3$. The anchor can
therefore be imposed together with any finite list of formal divisibilities.

% Original 15,671-witness anchor. No polynomial or constant is changed.
\subsubsection{The literal seven-witness semilocal predicate}\label{subsec:o8-anchor-proof}

For $t\in\mathbb Q$, let $\mathcal T(t)$ mean that there are
$u_0,u_1,u_2,u_3,v_1,v_2,v_3\in\mathbb Q$ with
\[
 u_0^2-2u_1^2-3u_2^2+6u_3^2=1,
 \qquad
 (t-2u_0)^2-8v_1^2-12v_2^2+24v_3^2=4.
\]
Then
\[
 \mathcal T(t)\quad\Longleftrightarrow\quad
 t\in\mathbb Z_{(2)}\cap\mathbb Z_{(3)}.
\]
Indeed, use the quaternion convention $i^2=2$, $j^2=3$ and $ij=-ji$.
Its reduced norm is $x_0^2-2x_1^2-3x_2^2+6x_3^2$ and its reduced trace
is $2x_0$. Put $v_0=(t-2u_0)/2$. The two displayed equations say
exactly that $t$ is a sum of two reduced traces of norm-one elements.
No eighth variable is needed, since $v_0$ is recovered by this rational
expression with fixed nonzero denominator~$2$.

In the notation of \cite[Definition~4 and Proposition~6]{koenigsmann-defining},
this is $T_{2,3}=S_{2,3}+S_{2,3}$. That proposition identifies
$T_{a,b}$ with the intersection of the local integer rings at the
ramified places, with a separate real interval only when infinity is
ramified. Here the Hilbert symbols are $-1$ at $2$ and $3$, and $1$
elsewhere. At $3$ this follows from $(2/3)=-1$; at every other odd
prime both parameters are units. At $2$, writing $2=2^1\cdot1$ and
$3=2^0\cdot3$, the usual exponent is
$0+1\cdot(3^2-1)/8+0=1$. The algebra splits over $\mathbb R$ because
both parameters are positive. Thus the ramification set is precisely
$\{2,3\}$, which proves the claimed equivalence in both directions.
This uses Koenigsmann's sum of two trace sets, not a predicate with
an additional finite offset.

\subsubsection{The original anchor and its exact index subgroup}\label{subsec:o8-anchor-subgroup}

Write $K=\mathbb Q(\rho)$, $\rho^2=2$, and use the literal normalized
curve and generator
\[
 E:y^2=x^3-(11-6\rho)x,
 \qquad
 P=\left(\frac{57}{8}-\frac72\rho,
          -\frac{165}{8}+\frac{417}{32}\rho\right).
\]
As in the original construction, use Proposition~\ref{prop:generator}, transported by the unit change
of variables, to obtain $2E(K)=\mathbb ZP$. For a nonzero $S=nP$, put
$f(S)=1/x(S)=f_0+f_1\rho$. Retain the exact constant chain
\[
 c_0=3,\qquad c_{j+1}=c_j^2\quad(0\le j<24),\qquad C=c_{24}.
\]
It forces $C=3^{2^{24}}$, with $25$ shared rational variables and
$25$ quadratic-or-linear equations. The anchor introduces two
variables $a_i$, imposes $Ca_i=f_i$, and imposes $\mathcal T(a_i)$
for each $i\in\{0,1\}$. It has $16$ private witnesses and six
quadratic equations, exactly as in the saved module.

Set $N=2^{24}$ and $m=N/2=2^{23}$. We claim that, for nonzero $n$,
\[
 \text{the anchor holds on }nP
 \quad\Longleftrightarrow\quad
 n\in 2\cdot3^{m-1}\mathbb Z.
\]
First consider the place above~$3$. It is unramified quadratic, with
residue field $\mathbb F_9$ and normalized valuation $v(3)=1$.
For rational $a,b$ one has
$v(a+b\rho)=\min(v_3(a),v_3(b))$: after dividing by the smaller
power of $3$, linear independence of $1,\bar\rho$ over $\mathbb F_3$
precludes cancellation.
The curve has good reduction there. Direct calculation gives
\[
 \bar P=(\bar\rho,0),\qquad
 z(2P)=
 \frac{3396356159088-4135204317756\rho}{3287811625151},
 \qquad z=-x/y.
\]
The reduction of $P$ has order two. The denominator in this expression
is a $3$-unit and both numerator coefficients have valuation one.
Consequently $nP$ is in the formal kernel exactly when $n$ is even,
and $v(z(2P))=1$.

For completeness, the valuation of every such multiple can be computed
without an unproved uniformity assertion. The integral formal group
in the parameter $z$ has inverse $-z$, so its multiplication series
are odd. Their linear coefficient is the corresponding integer.
Thus, for a $3$-unit integer $u$, $[u]z=uz+O(z^3)$ preserves positive
valuation, whereas $[3]z=3z+O(z^3)$ increases valuation by exactly one
at depth at least one. The higher terms have valuation at least
$3v(z)>v(z)+1$. Apply this repeatedly to $[n/2](2P)$, allowing
negative integers by negation. It gives
\[
 v(z(nP))=1+v_3(n),\qquad
 v(f(nP))=2+2v_3(n)\quad(2\mid n).
\]
The second equality follows from
$f=z^2(1-(11-6\rho)f^2)$ on the formal branch.
The integral formal-group construction and its linear multiplication
coefficient are given in \cite[IV\S1, Example~IV.2.2.3 and
Proposition~IV.2.3(a)]{silverman}. No minimal-model assertion at
the bad prime~$2$ is used below.

At the place above~$2$, normalize $v(2)=1$, so $v(\rho)=1/2$.
The two summands of a rational pair have valuations of opposite
integer parity after multiplication by~$2$, and hence cannot cancel.
The displayed coordinates give
\[
 v(x(P))=-3,\qquad v(y(P))=-9/2,\qquad v(z(P))=3/2.
\]
Thus $P$ lies on the formal branch at infinity, even though reduction
at~$2$ is bad. Every integral multiple lies in this formal subgroup.
One can see the integral branch directly by setting $w=-1/y$:
the equation becomes $w=z^3-(11-6\rho)zw^2$, which has a unique
solution $w\in z^3\mathcal O[[z]]$; its group law is the integral
formal group just cited. In particular $f(nP)$ is integral at~$2$.
Since
\[
 v(a+b\rho)=\min(v_2(a),v_2(b)+1/2),
\]
its rational coefficients are both $2$-integral. As $C$ is a $2$-unit,
the two anchor conditions at~$2$ impose no further restriction on $n$.

At~$3$, the two anchor conditions are equivalent to
$v(f(nP))\ge N$. Positive valuation of $f$ implies negative valuation
of $x$, hence membership in the formal kernel, so $n$ is even.
The preceding exact formula makes the remaining condition
$v_3(n)\ge m-1$. Conversely these two conditions give both required
coordinate divisibilities at~$3$, and the condition at~$2$ already
holds. The proved characterization of $\mathcal T$ supplies its
rational witnesses, proving the subgroup claim.

In particular, anchor indices contain a nonzero finite-index subgroup
and have positive representatives. They can be combined with any
finite collection of requirements that indices be divisible by fixed
integers: take a common multiple with $2\cdot3^{m-1}$. More generally,
any finitely many formal-neighborhood conditions whose preimages are
finite-index subgroups can be imposed in the same way. This statement
does not claim compatibility with arbitrary nonzero residue conditions.

The old and new reciprocal coordinates are equivalent for this anchor.
With $\lambda=3+2\rho$, one has $f_{\rm old}=\lambda^2f$ and
\[
 \begin{pmatrix}f_{{\rm old},0}\\f_{{\rm old},1}\end{pmatrix}
 =\begin{pmatrix}17&24\\12&17\end{pmatrix}
 \begin{pmatrix}f_0\\f_1\end{pmatrix}.
\]
This matrix has determinant one and integral inverse, so it preserves
coordinatewise membership in
$C(\mathbb Z_{(2)}\cap\mathbb Z_{(3)})^2$ in both directions.

\subsubsection{The unchanged numerical cutoff}\label{subsec:o8-anchor-cutoff}
In this calculation $H=H_*$ is the constant in Theorem~\ref{hyp:o8-height}.

The anchor gives $|n|\ge3^{2^{23}-1}$. Conditional only on the separately
audited original height bound $\log_2H<4783227$, it follows that
$|n|>(20H)^2$. Indeed $\log_2 3>3/2$ and $\log_2 20<5$, and
\[
 3(2^{23}-1)=25165821
 >19132928=4(5+4783227).
\]
Thus the original $24$ squarings suffice whenever that exact height
input is established. This calculation does not replace the required
five-contact height theorem by a six-contact theorem or by a larger
constant. No part of this proof changes the saved polynomial.

\subsection{The exact norm tests}

For a nonzero rational contact $r$ and rational tests $X_i$ of positive
orders $L_i$, apply the binary operation
$N_r(a,b)=a^2+rb^2$ to the ordered leaves
$1,X_1/r^{L_1},\ldots,X_k/r^{L_k}$ in the following fixed order.
First pair $1$ with $X_1/r$, where $L_1=1$. Group the remaining
tests by increasing order, retaining their original order within each
group. Pair successive tests within each group and carry any last
unpaired test. Then pair successive resulting nodes at each level and
carry any last node until only a root $W$ remains. The polynomial
equations for the bottom pairs are
\[
 rw-r-X_1^2=0,\qquad
 r^{2L}w-X_i^2-rX_j^2=0;
\]
an unpaired normalized leaf satisfies $r^Lw-X_i=0$.
Higher nodes satisfy $w-a^2-rb^2=0$. Powers of $r$ are computed
by binary squaring. Add an inverse for $r$, a variable for $W^2$,
and the predicate $\Psi(rW^2)$.

At a prime with $R=v_p(r)$ odd, the two summands in every norm node
have distinct valuation parity. There is no cancellation. The constant
leaf is on the all-left path, so $v_p(W)\le0$ if $R<0$.
Consequently $v_p(rW^2)>0$ first implies $R>0$.
It also gives $v_p(W)>-R/2$. Descending any norm node, a bound
$v_p(N_r(a,b))>-cR$ with $c<1$ gives
$v_p(a)>-cR/2$ and $v_p(b)>-(c+1)R/2$.
Every normalized leaf therefore has valuation greater than $-R$.
Thus $v_p(X_i)>(L_i-1)R\ge L_i-1$, or $v_p(X_i)\ge L_i$.
The root cannot vanish at such a prime. Indeed $N_r(a,b)=0$ forces
$a=b=0$ over $\Q_p$, and descent would make its constant leaf zero.
Thus Lemma~\ref{hyp:o8-pole} applies even though the polynomial
has no separately imposed inverse for $W$.

For completeness at a positive prime $r=p$, tests of orders $L_i$
can be made so deep that every nonconstant normalized leaf is
$p$-integral and the root $W$ is a $p$-adic unit. Then $rW^2$ has
square reduced denominator. If $r$ is a local square at $2,3,17$,
the completeness half of Lemma~\ref{hyp:o8-pole} applies.

\subsection{Forty-one contacts and one common comparison point}

Set
\[
 R=47027=31\cdot37\cdot41,\quad W_0=25878772920,
 \quad c=18484837801,\quad V_0=32W_0.
\]
For nine rational targets $a_i$, put $A_i=Ra_i+3$ and
$w_i=A_iv-u$. Use the forty-one nonzero contacts
\begin{equation}\label{eq:o8-contacts}
 v/V_0,\qquad u-\delta v\ (\delta\in\mathcal D),\qquad
 -(w_i-\delta v)\ (\delta\in\mathcal D,\ 1\le i\le9).
\end{equation}
For integer $a_i$, the prime-pattern argument of
Section~\ref{sec:mult} applies on
$v=32W_0(c+W_0y)$, $u=c+W_0x$ and gives positive ordinary prime
values for every distinct contact, all congruent to $c$ modulo $W_0$.
Repeated targets use the same prime values. These are local squares
at $2,3,17$. The precise application uses at most forty-one distinct
affine forms and their successor exclusions. At $31,37,41$ there are
at most eight slopes since every $A_i$ reduces to $3$; at $p\ge43$
there are at most forty. Choosing two slopes outside that set proves
local admissibility. The small-prime conditions are built into $c$.
The full uniform fiber sieve and same-prime order argument are given
in Subsection~\ref{subsec:o8-pattern-proof} below. No conditions relating orders at
different selected primes are needed.

An odd normalized contact of $u/v$ or $w_i/v$ outside $S_*$ forces
some raw contact in \eqref{eq:o8-contacts} to have odd valuation.
To see this, let $m$ be the minimum of the two coordinate valuations.
For a finite contact of odd depth, if $m$ is even its distinguished
linear form has odd valuation; if $m$ is odd another finite contact
form has valuation $m$. At an odd contact with infinity the same
argument uses the denominator if $m$ is even and a finite form if
$m$ is odd. All differences of distinct finite contact points are units
outside $S_*$. The signs and $V_0$ do not affect parity there.

For nonzero points $Q_i=n_iP_*$, $1\le i\le4$, take $Q_0=P_*$,
one common $J=jP_*\ne O$, and four rational proxies $b_i$.
At each raw contact introduce $X_0=hP_*$ and $X_i=m_iP_*$ and
the subtraction charts $Y_i=X_i-4Q_i$, $0\le i\le4$.
The points $4Q_i$ are computed by two doublings. Put
\[
 L(Z)=\sum_{j=0}^{15}
       \frac{(-1)^j\binom{2j}{j}D^j}{4j+1}Z^{4j+1}.
\]
This polynomial has degree sixty-one; the first omitted logarithm term
has degree sixty-five. Test order one for
\[
 f(X_i)\ (0\le i\le4),\quad f(J^\sigma),\quad
 f(Y_i^\sigma)\ (0\le i\le4),
\]
and order sixty-four for
\[
 L(z(X_i))/L(z(X_0))-b_i\ (1\le i\le4),\qquad
 z(Y_i^\sigma)/z(J^\sigma)\ (0\le i\le4).
\]
The inverses in these formulas are imposed explicitly. For each field
test $a+b\sqrt2$ of order $L_i$, use the rational tests
$(\theta^2-2)b$ and $a+\theta b$ of that order, together with
$\theta^2-2$ of order one. There is one common rational $\theta$
for all forty-one comparisons. This gives twenty-three rational tests
of order one and eighteen of order sixty-four.

At an odd raw contact the norm tests select a place above $p$, and
the conjugate parameter ratios imply
$h\equiv4$ and $m_i\equiv4n_i\pmod {p^{64}}$.
In particular $h$ is a unit. The logarithm expansion and its degree
sixty-five remainder give
\begin{equation}\label{eq:o8-proxies}
 b_i\in\Z_p,\qquad b_i\equiv n_i\pmod {p^{64}}.
\end{equation}
This uses the same formal-group estimates as Section~\ref{sec:mult}.

For completeness choose a consistent orientation at each distinct
selected prime. Let $d_{1,p},d_{2,p}$ be the two reduction orders.
The ordinary-prime conditions give $\gcd(d_{1,p},d_{2,p})=4$ and
$p\nmid d_{1,p}d_{2,p}$. Choose
$j=\operatorname{lcm}_p d_{2,p}$. For each occurrence at $p$ choose
\[
 h\equiv0\pmod {d_{1,p}},\qquad
 h\equiv4\pmod {d_{2,p}},\qquad
 h\equiv4\pmod {p^{64+v_p(j)}}.
\]
These congruences are compatible. Avoid $h=0,4,-4$, and set
$X_i=hn_iP_*$. They give all the required formal conditions and
parameter ratios relative to the single point $J$. Increase the last
precision if necessary and choose $\theta$ last by the Chinese
remainder theorem. Repeated occurrences of a prime use the same
orientation. Every subtraction denominator is nonzero, and the
formal-logarithm unit calculation at $p$ proves $L(z(X_0))\ne0$.

% Independent arithmetic justification for the unchanged original 15671-witness circuit.
\subsubsection{The original forty-one-prime pattern}\label{subsec:o8-pattern-proof}
Let
\[
 \mathcal D=\{0,1,2401,-343/32\},\quad R=47027=31\cdot37\cdot41,
 \quad W=25878772920=4\prod_{q\le29}q,
\]
\[
 c=18484837801,\qquad V_0=32W=828120733440.
\]
For integers $a_1,\ldots,a_k$, $k\le9$, put $A_i=Ra_i+3$,
$w_i=A_iv-u$. There are $5+4k$ displayed positions
\begin{equation}\label{eq:old41-list}
 v/V_0,\qquad u-\delta v,\qquad
 u-(A_i-\delta)v\quad(\delta\in\mathcal D,\ 1\le i\le k).
\end{equation}
For every lower bound these positions can simultaneously be positive
primes beyond that bound, congruent to $c$ modulo $W$, with the
following property. For the curve
$E:y^2=x^3-(3-\sqrt2)^2x$, both reductions above each prime are good,
and their cardinalities $d_1,d_2$ satisfy
\[
 \gcd(d_1,d_2)=4,\qquad p\nmid d_1d_2.
\]
This is a same-prime statement and asserts no cross-prime gcd condition.

\begin{proof}
Use $u=c+Wx$, $v=V_0(c+Wy)$. Put
\[
 \mathcal S=\mathcal D\cup\bigcup_i(A_i-\mathcal D),\quad
 K_s=V_0s\in W\mathbb Z,
 \quad\Psi_0=c+Wy,\quad\Psi_s=c+Wx-K_s(c+Wy).
\]
Remove duplicate targets. The remaining slopes are all distinct.
Indeed, the nonzero integral differences in $\mathcal D-\mathcal D$
are $\pm1,\pm2400,\pm2401$, none divisible by $31$.
The integral sums in $\mathcal D+\mathcal D$ are
$0,1,2,2401,2402,4802$, with residues $0,1,2,14,15,28$ modulo $31$,
none equal to $A_i\equiv3$. These facts exclude intersections between
secondary families and between primary and secondary families.
The $m=1+|\mathcal S|\le41$ forms have pairwise independent linear
parts. Repeated targets may be restored afterwards with repeated values.

At $q\mid W$, the forms equal $c$; direct reduction gives
$c\equiv1\pmod8$, $c\equiv3\pmod7$, and $c\equiv1\pmod q$ at
all other odd primes $q\le29$. Thus the forms and their successors
are units at every odd $q\mid W$. For $q>29$, the affine map from
$(x,y)$ to $(u,v)$ is invertible. At $q=31,37,41$, all $A_i=3$
modulo $q$, so $|\mathcal S\bmod q|\le8$. For $q\ge43$ this
cardinality is at most $40$. In either case choose distinct residues
$b,e$ outside the slopes. Put $\gamma=V_0^{-1}$ modulo $q$ and,
reversing $b,e$ if necessary, require $e-b\ne-\gamma$. Then
\[
 v=(e-b)^{-1},\quad u=bv
\]
satisfies
\[
 u-sv=v(b-s)\ne0,\quad u-sv+1=v(e-s)\ne0,
 \quad\gamma v\ne0,\quad\gamma v+1\ne0.
\]
This proves all the value and successor local conditions. No condition
on a successor is imposed at $2$.

Add an unused integer coordinate $z$, and choose a fixed rectangular
box $\Omega$ inside the cone $v>0$, $u/v>\max(0,|A_1|,\ldots,|A_k|)+2413$ in the
homogeneous leading coordinates. All forms on $X\Omega$ lie between
$\alpha X$ and $\beta X$ for positive fixed constants. Fix $Q>82$
containing every prime divisor of $W$ and every nonzero $K_s-K_t$.
Impose one good residue at every prime up to $Q$. For a cutoff $L>Q$,
exclude $\Psi_\nu=0,-1$ at every additional prime up to $L$.
The surviving proportion is
\[
 \gg_{\mathbf a}\prod_{Q<q\le L}(1-2m/q)
 \gg_{\mathbf a}(\log L)^{-2m},
\]
by \cite[Theorem~1, modulus~2]{williams-mertens}.

For fixed $L$, apply \cite[Main Theorem, p.~1763]{greentao} to every
surviving lattice residue in the convex box. The singleton partition
in \cite[Definition~1.5 and Lemma~1.6]{greentao} bounds complexity by
$m-2\le39$. The inputs $\mathrm{MN}(39)$ and $\mathrm{GI}(39)$ follow
from \cite[Theorem~1.1]{moebius} and \cite[Theorem~1.3]{gtz}.
At imposed primes the local factor is $(1-1/q)^{-m}$; at the other
primes it is $1+O(m^2/q^2)$. The singular product therefore has a
positive lower bound independent of $L$. All coefficient and modulus
bounds are fixed before $X$ tends to infinity. Summing the finitely
many asymptotics and removing proper prime powers gives
\begin{equation}\label{eq:old41-lower}
 \#\{\text{prime patterns}\}\ge
 c_{\mathbf a}\frac{X^3}{(\log L)^{2m}(\log X)^m}.
\end{equation}
The constant $c_{\mathbf a}>0$ is independent of $L$; the threshold
for $X$ may depend on it. Proper powers contribute at most
$O_{\mathbf a}(X^{5/2}(\log X)^{m+1})$ before removing von Mangoldt
weights, since each fixed value has $O_{\mathbf a}(X^2)$ fibers.

We give the uniform upper estimates needed to impose the order condition.
For forbidden sets of cardinality $\rho(q)$ in a two-dimensional box
of side $N$, put $g(q)=\rho(q)/q^2$, $h(q)=g(q)/(1-g(q))$, and extend
multiplicatively on squarefree support. CRT gives
\[
 A_d=N^2g(d)+R_d,\qquad |R_d|\ll Nd+d^2.
\]
With $G(Z)=\sum_{d\le Z}h(d)$, take
\[
 \lambda_d=\frac{\mu(d)}{G(Z)\prod_{q\mid d}(1-g(q))}
 \sum_{\substack{e\le Z/d\\(e,d)=1}}h(e)\quad(d\le Z).
\]
Then $\lambda_1=1$, $|\lambda_d|\le d$, and the Selberg quadratic
form equals $G(Z)^{-1}$. To verify the last assertion, use
$y_r=\sum_{r\mid d}\lambda_dg(d)=\mu(r)h(r)/G(Z)$ and
$1/g(n)=\sum_{r\mid n}1/h(r)$. The number of sifted points is thus
\begin{equation}\label{eq:old41-upper}
 \ll N^2/G(Z)+NZ^6+Z^8.
\end{equation}
For $\rho(q)=2q-1$, supported on $q\equiv1\pmod4$ with at most one
prime omitted, $G(Z)\gg_Q\log Z$, uniformly in that omission.
For $\rho(q)=(m-1)q$ at all $q>Q$,
$G(Z)\gg_{Q,m}(\log Z)^{m-1}$. These follow from
\cite[Theorem~1, moduli~4 and~2]{williams-mertens}. Truncation is
justified by using $z=Z^{1/512}$ and the probability distribution
$h(d)/\prod_{q\le z}(1+h(q))$: its expected $\log d$ is at most
$40\sum_{q\le z}(\log q)/q<\tfrac12\log Z$ for large $Z$.
Thus at least half the mass lies below $Z$. Removing one norm-sieve
prime changes its product by a factor at least $1/2$.

For a selected prime, $(7/p)=-1$ and $p\equiv1\pmod8$. The two
reductions of $3-\sqrt2$ have product $7$ and opposite square classes,
so their traces are opposite. Their discriminants $64d^6$ are units.
The automorphism $(x,y)\mapsto(-x,iy)$ is rational over $\mathbb F_p$.
By the Tate-module injection \cite[III.7.4]{silverman}, its rational
centralizer is $\mathbb Q(i)$: at the auxiliary prime $3$ the
centralizer has dimension two because $T^2+1$ is irreducible.
Frobenius is integral by \cite[V.2.3.1(b)]{silverman}; hence it is
$a+bi\in\mathbb Z[i]$. The trace and determinant formulas
\cite[III.8.6]{silverman} give
\[
 a^2+b^2=p,\qquad d_1=p+1-2a,\qquad d_2=p+1+2a.
\]
Full rational two-torsion gives $4\mid d_1,d_2$, and
$v_2(d_1+d_2)=2$, so their gcd has $2$-part exactly $4$.
Also $0<d_i<2p$ and $d_i\ne p$, so $p\nmid d_1d_2$.
An odd common divisor $\ell$ divides $a$ and $p+1$, and forces
$b^2\equiv-1\pmod\ell$. The successor conditions exclude $\ell\le L$.

The number of bad prime values of any fixed form is
\begin{equation}\label{eq:old41-bad}
 \ll_{\mathbf a}X/(L\log X)+X^{7/8}+\sqrt X\log X.
\end{equation}
Indeed, for $L<\ell\le X^{1/8}$, write
$a=\ell r$, $b=b_0+\ell s$, with at most two $b_0$ modulo $\ell$.
The box has side $N\asymp\sqrt X/\ell\gg X^{3/8}$.
Use the norm case of \eqref{eq:old41-upper} with $Z=X^{1/100}$;
its error is uniformly $o(N^2/\log X)$. Sum
$O(X/(\ell^2\log X))$. For $\ell>X^{1/8}$, $a\ne0$ gives
$O(\sqrt X/\ell)$ choices with no additive constant, and the
$b$ count is $O(\sqrt X/\ell+1)$. Summation proves the remaining terms.

On a fiber $\Psi_0=p$, use $(x,z)$; the other forms are
$c+Wx-K_s p$. On a fiber $\Psi_{s_0}=p$, use $(y,z)$ and
\[
 x=K_{s_0}y+\frac{p-c(1-K_{s_0})}{W},\qquad
 \Psi_s=p+(K_{s_0}-K_s)\Psi_0.
\]
The quotient is integral for an actual fiber. For
$Q<q\le X^{1/100}<p$, the remaining zero lines are distinct and
parallel, with total cardinality $(m-1)q$. The second case of
\eqref{eq:old41-upper} bounds each fiber by
$O_{\mathbf a}(X^2/(\log X)^{m-1})$, independently of $p,L$.
Multiply by \eqref{eq:old41-bad} and sum over $m$ positions.
The bad patterns are at most
\[
 C_{\mathbf a}\frac{X^3}{L(\log X)^m}
 +O_{\mathbf a}\left(
 X^{23/8}/(\log X)^{m-1}+X^{5/2}/(\log X)^{m-2}\right).
\]
Choose $L$ so large that $C_{\mathbf a}/L$ is less than half
$c_{\mathbf a}/(\log L)^{2m}$, then choose $X$ large.
A good pattern survives, proving the assertion.
\end{proof}

\subsubsection{A shared auxiliary point with private shifts}\label{subsec:o8-common-J-proof}
Let $\mathcal P$ be the finite set of distinct selected primes, all
larger than $76176$. At each $p$ fix an orientation of its two places,
and let $d_{p,1},d_{p,2}$ be the corresponding group cardinalities.
Set
\[
 j=\operatorname{lcm}_{p\in\mathcal P}d_{p,2}>0,
 \qquad J=[j]P_*.
\]
For each $p$, independently choose an integer $h_p$ satisfying
\begin{equation}\label{eq:old41-private-crt}
 h_p\equiv0\pmod{d_{p,1}},\qquad
 h_p\equiv4\pmod{d_{p,2}},\qquad
 h_p\equiv4\pmod{p^{64+v_p(j)}}.
\end{equation}
The first two congruences are compatible because their modulus gcd
is $4$. The last is compatible because $p\nmid d_{p,1}d_{p,2}$.
The resulting arithmetic progression contains an $h_p$ outside
$\{0,4,-4\}$. Put
\[
 X_0=[h_p]P_*,\qquad X_i=[h_p n_i]P_*,\qquad
 Y_i=[(h_p-4)n_i]P_*,\quad n_0=1.
\]
These are private for each contact occurrence, whereas $J$ is common.
Use the same orientation at repeated occurrences of $p$.

At the primary place all $X_i$ are formal; at the opposite place
$J$ and all $Y_i$ are formal. The formal logarithm is a homomorphism
and an isometry on $p\mathbb Z_p$ by
\cite[IV.6.4 and its proof; VII.2.1--2.2]{silverman}. The hypotheses hold
because $p$ is odd, split, and of good reduction. For any integers
$k,j$ whose multiples of $P_*$ lie in the formal group,
\[
 \frac{\log([k]P_*)}{\log([j]P_*)}=\frac{k}{j}.
\]
To see this without defining a logarithm of $P_*$ itself, apply the
homomorphism law to $[j]([k]P_*)=[k]([j]P_*)$. Nonzero formal
multiples have nonzero logarithm. Therefore
\[
 v_p\left(\frac{z(Y_i)}{z(J)}\right)
 =v_p\left(\frac{(h_p-4)n_i}{j}\right)\ge64,
 \qquad
 \frac{\log(X_i)}{\log(X_0)}=n_i.
\]
This proves the required local comparison conditions. Although other
$d_{q,2}$ may be divisible by $p$, their effect is exactly the term
$v_p(j)$ in \eqref{eq:old41-private-crt}; no cross-prime gcd assertion
is needed.

For the original truncation
\[
 L(z)=z\sum_{r=0}^{15}\frac{(-1)^r\binom{2r}{r}}{4r+1}(Dz^4)^r,
\]
the full logarithm has the same series extended over all $r\ge0$.
This follows by integrating
$(1+4Dz^4)^{-1/2}$, the invariant differential in the parameter
$z=-x/y$. If $v_p(z)\ge1$ and $p>76176$, every omitted term has
valuation at least $65v_p(z)$: its integer numerator can only increase
valuation, and for $n=4r+1\ge65$,
$n v_p(z)-v_p(n)\ge65v_p(z)$.
For $n<p$ this is immediate; for $n\ge p$ use $n-65\ge v_p(n)$.
Thus $L(z)/\log(z)\in1+p^{64}\mathbb Z_p$. The required primary
truncated-log ratios are congruent to the integers $n_i$ modulo $p^{64}$.

Conversely, in an accepted comparison at any $p>76176$, its conjugate
parameter-ratio tests imply
\[
 m_i-4n_i\equiv0\pmod{p^{64}},\qquad h-4\equiv0\pmod{p^{64}},
\]
since $j$ is a nonzero integer. Hence $h$ is a $p$-unit. The same
logarithm identity and truncation estimate give
$L(X_i)/L(X_0)\equiv m_i/h\equiv n_i\pmod{p^{64}}$ and show
that this ratio is integral. The rational proxy test then gives the
claimed integrality and congruence for every $b_i$.

\subsubsection{The common root approximation and parity boundary}\label{subsec:o8-root-proof}
For each distinct selected $p$, retain the chosen square root
$\rho_p$ of $2$. After constructing the private points, each field
test $A+B\sqrt2$ of order $k$ has $v_p(A+B\rho_p)\ge k$.
Choose a common integer $\theta$ sufficiently close to $\rho_p$ at
every selected prime that
$v_p(\theta-\rho_p)\ge k-v_p(B)$ for every nonzero coefficient $B$,
and at least one. There are finitely many requirements; CRT supplies
such an integer. Since $\theta+\rho_p$ is a unit,
\[
 v_p((\theta^2-2)B)\ge k,\qquad v_p(A+\theta B)\ge k.
\]
Repeated primes require the same root and impose compatible precision
bounds. Conversely, the root test and these two rational tests imply
the field test at the selected root. This verifies the circuit's one
common $\theta$, not a separate root variable per contact.

Every raw contact parameter in the complete assignment is a positive
prime congruent to $c$ modulo $W$. It is therefore a square in
$\mathbb Q_2,\mathbb Q_3,\mathbb Q_{17}$. The binary norm batch has
unit root at its contact prime, because normalized leaves are integral
and its leftmost constant leaf is $1$. Its root is positive and nonzero
in the real embedding. The argument $r\mathcal W^2$ consequently is
positive, has square reduced denominator, and is locally square at
$2,3,17$. These are the exact hypotheses of the original specialized
parity-completeness input. This verifies its consumption but does not
constitute a new proof of that input, which is supplied here by
Lemma~\ref{lem:original-seven-parity-completeness}. In particular no integrality of
the comparison points at these three small primes is required here.

\subsection{Scalar height and the defining relation}

Put $\beta=X^2+2$ and
\[
 r(\beta)=\beta(\beta-1)(\beta-2401)(32\beta+343)(2\beta-1).
\]
This is nonzero for every rational $X$, since $X^2=2399$ has no
rational solution. Introduce $J_0\ne O$ and apply a scalar norm
batch to the two coefficients of each of
$f(J_0),f(S),z(S)/z(J_0)$, of orders $1,1,1,1,48,48$.
In this batch replace the pole argument by $\delta\mu rW^2$, where
$\delta\in\{1,-1\}$ and $\mu$ is a squarefree product of
$2,3,5,7,11,13,17$. Precisely, eight Boolean variables produce
\[
 A=(1+e_2)(1+2e_3)(1+12e_{13})(1+16e_{17}),\quad
 B=(1+4e_5)(1+6e_7)(1+10e_{11})(1-2\eta),
\]
and the pole argument is $ABrW^2$. These factors have no valuation
outside $S_*$. Every odd five-contact of $\beta$ gives an odd
valuation of $r(\beta)$ there. At infinity its valuation is five
times that of $\beta$; at a finite contact all other factors are units.
Thus the scalar tests force $M_5(\beta)^{48}\mid s$, and
Theorem~\ref{hyp:o8-height} gives
\begin{equation}\label{eq:o8-scalarheight}
 h(\beta)\le H_*\sqrt{|s|}.
\end{equation}

For integral $X$, the scalar $r$ is a nonzero integer. The square
classes of $2,3,5,7,11,13,17$ form a basis of
$\Q_2^*/\Q_2^{*2}\times\Q_3^*/\Q_3^{*2}
\times\Q_{17}^*/\Q_{17}^{*2}$.
Choose $\delta$ and $\mu$ so that $\delta\mu r$ is positive and
a square at these three places. Choose $J_0$ in the required formal
subgroups and then choose $S$ sufficiently divisible relative to $J_0$
at all primes dividing $r\mu$, together with the anchor conditions.
All normalized nonconstant leaves are then integral at primes dividing
$r$, and vanish at primes dividing $\mu$ but not $r$. The root is
a unit at these primes. Consequently $\delta\mu rW^2$ has square
reduced denominator, and Lemma~\ref{hyp:o8-pole} supplies the
scalar pole witnesses.

Take $A=nP_*,S=sP_*,T=tP_*,Z=zP_*\ne O$, and use
$(\beta,\xi,\tau,\zeta)$ as the four proxies. Impose
\begin{equation}\label{eq:o8-Lucas}
 \tau^2-\xi(\xi-2)\tau\zeta+(\xi-2)^2\zeta^2=1.
\end{equation}
The nine rational targets in \eqref{eq:o8-contacts} are a nonzero
common denominator $e$ and the eight coefficients of
$ex(A),ex(S),ex(T),ex(Z)$. Require its forty-one comparison batches,
the scalar batch, and the anchor just described.

If every associated five-contact radical has its $K$th power dividing
a nonzero integer $N$, Theorem~\ref{hyp:o8-height} and
$Ra_i+3=u/v+w_i/v$ give
\[
 h(a_i)<2H_*|N|^{24/K}+14,
 \qquad h(x(Q))<7H_*|N|^{24/K}
\]
for each of the four points. The second bound follows by applying the
first to $e$ and both coefficients of $ex(Q)$ and absorbing the
constant smaller than fifty into $H_*$.

\begin{lemma}[Lucas growth]\label{lem:o8-lucas}
If $s,t,z\in\Z$, $|s|=a\ge7$, $z\ne0$, and
\[
 t^2-s(s-2)tz+(s-2)^2z^2=1,
\]
then $|t|>\exp(a-3)$. Conversely each such $s$ admits positive
integers $t,z$ satisfying this equation.
\end{lemma}
\begin{proof}
Put $w=\operatorname{sgn}(s)(s-2)z$ and
$\rho=(a+\sqrt{a^2-4})/2$. The equation becomes
$t^2-atw+w^2=1$. Every positive norm-one unit of $\Z[\rho]$ is a
power of $\rho$: multiply it by a power into $[1,\rho)$, write the
result as $v=A+B\rho$, and use
$B=(v-v^{-1})/(\rho-\rho^{-1})\in[0,1)\cap\Z$.
Thus $v=1$. If $U_0=0$, $U_1=1$ and
$U_{k+1}=aU_k-U_{k-1}$, the identities for $\rho^{\pm k}$ give
$|w|=U_k$ and $|t|\in\{U_{k-1},U_{k+1}\}$ for some $k\ge1$.
Modulo $a-2$ one has $U_k\equiv k$, and modulo $a+2$ one has
$U_k\equiv(-1)^{k-1}k$. The divisibility of $w$ by $|s-2|$ therefore
gives $k\ge a-2$ when $s>0$ and $k\ge a+2$ when $s<0$.
Induction gives $U_k\ge(a-1)^{k-1}$, whence
$|t|\ge(a-1)^{a-4}>e^{a-3}$ for $a\ge7$.
For existence take $k=a-2$, $z=U_k/(a-2)$ if $s>0$, and
$k=a+2$, $z=U_k/(a+2)$ if $s<0$, in both cases with $t=U_{k+1}$.
\end{proof}

\begin{proof}[Proof of Theorem~\ref{thm:o8-conditional}]
Suppose the relation holds, and let
$B=\max(|n|,|s|,|t|,|z|)$ and
$\Delta=t^2-s(s-2)tz+(s-2)^2z^2-1$.
If $\Delta\ne0$, the congruences \eqref{eq:o8-proxies} imply that
every coordinate contact radical to the sixty-fourth power divides
$\Delta$. Since $|\Delta|<5B^4$, the index-height and shared-height
bounds give
\[
 B^2-6<7H_*|\Delta|^{3/8}<14H_*B^{3/2}.
\]
Hence $B<(14H_*+6)^2\le(20H_*)^2$, contradicting the anchor.
Thus $\Delta=0$. Lemma~\ref{lem:o8-lucas} gives
$|t|>\exp(|s|-3)$ because $z\ne0$.

If $\beta\ne n$, put $N=|\operatorname{num}(\beta-n)|$ in lowest
terms. The same congruences give
\[
 B^2<7H_*N^{3/8}+6,\qquad N\le2e^{h(\beta)}B,
\]
whence $\log B<\frac8{13}\log(16H_*)+\frac3{13}h(\beta)$.
Writing $a=|s|>(20H_*)^2$, \eqref{eq:o8-scalarheight} implies
$h(\beta)<a/20$ and $\log(16H_*)<a/20$. Therefore
$\log B<11a/260<a-3$, another contradiction.
We obtain $\beta=n\in\Z$, and $X\in\Q$, $X^2=\beta-2\in\Z$
imply $X\in\Z$ by reducing its denominator.

Conversely let $X\in\Z$ and set $A=\beta P_*$. Choose the scalar
height data and then an anchored $S$ with all required divisibilities.
Lemma~\ref{lem:o8-lucas} gives positive integers $t,z$ satisfying the
index equation; set $T=tP_*$ and $Z=zP_*$. Choose an integer $e$
clearing their abscissa coefficients. The ordinary prime pattern
provides \eqref{eq:o8-contacts}. Set the proxies equal to their integer
indices, choose $J$ and the private points as above, and finally choose
the common $\theta$. The norm and pole completeness statements supply
all remaining witnesses. No subsequent choice changes the target points.
\end{proof}

\subsection{The base compiler prescription and finite verification}

Every equation above is implemented with named arithmetic intermediates
so that each rational residual has degree at most four. A field product
uses $(a+b\sqrt2)(c+d\sqrt2)=(ac+2bd)+(ad+bc)\sqrt2$.
Scalar inverses are equations $uv=1$. Field inverses use the two
coefficient equations of the field product being one; subtraction has
the separate scalar inverse of its abscissa-difference norm.
The logarithm uses Horner evaluation in $(Dz^4)^2$ with block length
two. The two doublings of the fixed point $P_*$ also use arithmetic
intermediates; inserting its large precomputed fourth multiple as a
coefficient would give a different polynomial. The chart is exactly
\eqref{eq:o8-chart}, including the intermediate $h=vw$.

For each nonzero rational residual, multiply by the least common multiple
of its coefficient denominators, then divide all coefficients by their
positive greatest common divisor. Denote the resulting integral rows
by $R_1,\ldots,R_{15107}$. In the variable order specified by
\texttt{degree8-original/build.py} and its residual table, define
\begin{equation}\label{eq:o8-polynomial}
 F_8=\sum_{i=1}^{15107}R_i^2.
\end{equation}
This equation is an exact finite prescription for every coefficient.
The residual table, its full collected expansion, and the construction
certificate accompany the manuscript. No huge constant is treated as
an additional free parameter. In particular, all twenty-five variables
in \eqref{eq:o8-chain} belong to the stated witness count.

The recursive module counts are
\[
\begin{array}{c|rr}
 \text{module}&\text{new variables}&\text{residuals}\\\hline
 \text{pole}&7&4\\
 \text{anchor}&16&6\\
 \text{scalar batch, including pole}&34&31\\
 \text{forty-one-test batch, including pole}&58&55\\
 \text{one joint comparison, including batch}&373&360\\
 \text{root, including all its calls}&15646&15082\\
 \text{shared anchor chain}&25&25
\end{array}
\]
The root has $303$ direct variables and $285$ direct residuals; its
calls are one anchor, one scalar batch, and forty-one comparisons.
Thus $303+16+34+41\cdot373=15646$ and
$285+6+31+41\cdot360=15082$. Expansion of \eqref{eq:o8-polynomial}
gives the finite assertions of Proposition~\ref{prop:o8-finite}.
Over $\Q$, a sum of squares vanishes exactly when every summand
vanishes, so the residual interpretation proves the theorem.
This computation checks the polynomial and its compilation. The proof
of the height estimate is the arithmetic argument given above.

% End degree-eight construction
% Begin explicit quartic constructions
\section[An explicit quartic definition]{An explicit quartic definition of the integers}
\label{sec:quartics}

We now give the second final polynomial of this manuscript. Its arithmetic
source uses the base residual compiler of Section~\ref{sec:degree8-original},
with the reciprocal logarithms and mixed precisions proved below. A fresh
quadratic lift and explicit rational solution maps then give the following
fully expanded sum of squares.

\begin{theorem}[An explicit quartic definition]\label{thm:quartic}
With the arithmetic inputs of Section~\ref{sec:degree8-original}, the
coefficient file specified below defines an integer polynomial $\Phi$ of
total degree $4$, with $82\,476$ nonzero monomials in $9\,625$ variables.
For every $q\in\Q$,
\[
 q\in\Z\quad\Longleftrightarrow\quad
 \exists(w_1,\ldots,w_{9624})\in\Q^{9624},\qquad
 \Phi(q,w_1,\ldots,w_{9624})=0.
\]
Thus there are $9\,624$ existential witnesses in addition to the input.
Every listed variable occurs in the expanded polynomial, and
\[
 \|\Phi\|_\infty=685783949153203534233601<2^{80}.
\]
\end{theorem}

\subsection{Literal specification and coefficient file}

The companion repository directory \nolinkurl{quartics/quartic-9625/}
contains the complete coefficient table \nolinkurl{polynomial.expanded.txt}.
Its gzip copy has identical uncompressed bytes. After one comment line,
each line contains a nonzero signed integer coefficient followed by a
nondecreasing list of variable indices. An index repeated $e$ times denotes
exponent $e$; a coefficient with no indices is the constant term. Thus
\texttt{-6 0 2 2} denotes $-6X_0X_2^2$. The input is always $X_0$.
Like terms are collected and the lines are ordered lexicographically by
their index lists. The witness indices are assigned by occurrence frequency;
the file \nolinkurl{variable-map.json} records the exact relabeling.

The uncompressed coefficient file has $1\,558\,552$ bytes, including its
comment and line endings. Its SHA-256 digest is
\begin{center}\footnotesize\ttfamily
6d67d45f3d615b7fd2242e3a9404dd9f5da570fc711a29ebe77f613e80ebbf99.
\end{center}
This is a literal finite specification of every coefficient. No unspecified
auxiliary constant or choice of witnesses enters the coefficient table.

% Complete replacement of the released reciprocal-source subsection.
% Reviewed source: Version 2, SHA256 38d458ba4178083766c5ae56a894da595fe27170c1ef07f6b49750fc5edf7c37.
% Selected source: H19 / high40 / beta12 / anchor25 / scalar32 / 29 full + 8 beta.
% No endpoint count, source pin, certificate hash, or coefficient claim is supplied here.
\subsection{The selective reciprocal comparison source}
\label{subsec:quartic-source}
We use the curve, cyclic group of integer indices, pole predicate and
arithmetic inputs of Section~\ref{sec:degree8-original}. The present
source uses a sharper consequence of the same height construction,
thirty-seven direct coordinate contacts, and two kinds of reciprocal
comparison. These arithmetic changes are proved here before the
quadratic lift. They are not justified by a change of source digest.
Throughout,
\[
 K=\Q(\sqrt2),\quad d=3-\sqrt2,\quad D=d^2=11-6\sqrt2,
 \qquad E:y^2=x^3-Dx,\qquad G=2E(K)=\Z P_*.
\]
An element of $K$ is represented by its two rational coefficients.
Multiplication and inversion have their literal meanings
\[
 (a_0,a_1)(b_0,b_1)=(a_0b_0+2a_1b_1,a_0b_1+a_1b_0),\qquad
 a^{-1}=\frac{(a_0,-a_1)}{a_0^2-2a_1^2}\quad(a\ne0).
\]
The denominator is nonzero by the irrationality of $\sqrt2$.
All affine subtraction charts and their nonzero-abscissa guards remain.

% BLOCK HEIGHT19.
\begin{theorem}[A sharper five-contact height estimate]\label{thm:quartic-height19}
Keep $S_*$, $\mathcal D$ and $M_5$ from Theorem~\ref{hyp:o8-height}, and put
\[
 B_*=\prod_{p\le76176}p,\qquad
 A_{19}=2^{109}210^{1800},\qquad H_{19}=2A_{19}B_*^{19}+128.
\]
For every $a\in\mathbb Q\setminus\mathcal D$,
\[
 h(a)\le H_{19}M_5(a)^{19},
 \qquad \log_2H_{19}<2093966.
\]
Here $h$ is the absolute logarithmic Weil height.
\end{theorem}
\begin{proof}
We use the geometric construction and compatible representations in the
proof of Theorem~\ref{hyp:o8-height}, not merely its final exponent-$24$
inequality. Write $R=B_*M_5(a)$ and $C=210^{300}$.
The geometric comparison there is
\[
 h(a)\le2\max(0,h_F(A))+128.
\]
All $h_F$ in this proof use the original stable Faltings metric.
The CM estimate in that proof is $h_F(A)<854R$.

In the non-CM case the same construction gives an absolutely simple
quaternionic surface $A/L$, a normal field $L'/\mathbb Q$, and the final
primitive weight-two newform $f$, after the finite twist has been
absorbed, with
\[
 [L:\mathbb Q]\le576,\quad d'=[L':\mathbb Q]\le2304R,\quad
 N\le CR^2,\quad \operatorname{rad}(N)\mid R.
\]
Its two-dimensional compatible representation $r_\ell$ satisfies
\[
 V_\ell(A)|_{G_{L'}}\simeq
       (r_\ell\oplus r_\ell)|_{G_{L'}},\qquad
 (\det r_\ell)^2|_{G_{L'}}=\chi_{\rm cyc}^2|_{G_{L'}}.
\]
These are (LR3)--(LR5) and the determinant calculation in the existing
modularity argument. There is no bounded-degree hypothesis on the
coefficient field of $f$.

Let $\epsilon$ be the nebentypus of this final form. In the Tate-module
convention, with arithmetic Frobenius, $\det r_\ell=\epsilon\chi_{\rm cyc}$;
see \cite[\S12.5, (12.5.1)--(12.5.2) and Remark 12.5.1,
pp.~119--120]{diamond-im}. Equivalently, at a good prime its
characteristic polynomial is $X^2-a_pX+p\epsilon(p)$.
A consistent contragredient convention replaces the finite character by
its inverse, with the same kernel and order. The character $\epsilon^2$
is trivial on $G_{L'}$, so its image order divides $d'$. Consequently
\[
 t=\operatorname{ord}(\epsilon)\le2d'\le4608R.
\]
This bounds the character order, not the coefficient-field degree.

Put $H=\ker\epsilon\subset(\mathbb Z/N\mathbb Z)^*$ and
\[
 \Gamma_H=\left\{
 \begin{pmatrix}a&b\\c&d\end{pmatrix}\in\Gamma_0(N):
 d\bmod N\in H\right\}.
\]
Since $\epsilon(-1)=1$, this group contains $-I$.
Let $X_H=X_1(N)/H$ for the diamond action and $J_H=\operatorname{Jac}(X_H)$.
The curve, diamond action and quotient are defined over $\mathbb Q$.
The rational model here parametrizes a point of exact order $N$, up to
$H$, and is the model specified by \cite[\S2.1, (2.1.8)]{mascot-moduli};
it introduces no cyclotomic field-degree factor.

Write $q:J_1(N)\to A_f$ for the newform quotient. Its pulled-back
differentials are spanned by all Galois conjugates of $f$; see
\cite[Theorem~1(i),(iv), p.~526]{shimura-factors} and
\cite[\S3.1]{baker-et-al-finiteness}. Every $h\in H$ fixes those differentials,
since its eigenvalue on each conjugate is a conjugate of $\epsilon(h)=1$;
the diamond convention is \cite[\S3.2.1]{baker-et-al-finiteness}.
Homomorphisms of abelian varieties in characteristic zero are determined
by their differentials. Hence $q\circ\langle h\rangle=q$.
For
\[
 e_H=\frac1{|H|}\sum_{h\in H}\langle h\rangle
       \in\operatorname{End}_{\mathbb Q}^0(J_1(N))
\]
we therefore have $qe_H=q$.

Let $H_{\rm eff}$ be the image of $H$ in
$\operatorname{Aut}(X_1(N))$, and let $\pi:X_1(N)\to X_H$ be the quotient.
On Jacobians,
\[
 \pi^*\pi_*=\sum_{h\in H_{\rm eff}}h,\qquad
 \pi_*\pi^*=[|H_{\rm eff}|].
\]
The averages over $H$ and $H_{\rm eff}$ agree. In particular the trivial
action of $-1$ is not counted twice in the first identity.
The image of $e_H$ is thus $\mathbb Q$-isogenous to $J_H$, and its
restriction of $q$ is surjective. Poincare reducibility makes $A_f$
a $\mathbb Q$-isogeny factor of $J_H$.

The accepted realization of $r_\ell$ in $A_f$ puts it in $V_\ell(J_H)$.
Over $L'$, Faltings' Hom theorem and extension of scalars give a nonzero
homomorphism $A\to J_H$. Absolute simplicity and Poincare reducibility
give an $L'$-isogeny $A\times Z\to J_H$. In particular
$g=\dim J_H\ge2$. The Jacobian has its canonical principal polarization;
no bound on a polarization degree of $A$ or $Z$ is being assumed.

The map $X_H\to X(1)$ has degree
\[
 d_H=[\operatorname{SL}_2(\mathbb Z):\Gamma_H]
     =Nt\prod_{p\mid N}(1+1/p).
\]
There is no factor of two here, because $-I\in\Gamma_H$.
It is branched only over the cusp and the two elliptic points, so
$\deg_B(X_H)\le d_H$. The congruence-curve genus bound gives
$g\le1+d_H/12\le d_H$, using $g\ge2$.
These degree and genus conventions, and the resulting height inequality,
are recorded in \cite[\S4, (4.1)--(4.2), p.~8]{vonKaenel2021}.
For every prime $p\ge5$, $(1+1/p)^6\le p$, while the possible factors
at two and three have product at most two. It follows that
\[
 d_H\le9216CR^{19/6}<T,\qquad
 T=2^{14}CR^{19/6},\qquad g<T.
\]
By \cite[Theorem~1.1.1 and \S2.3]{javanpeykar},
\[
 h_F(J_H)\le13000000\,d_H^5g<2^{24}T^6.
\]
This is the generic Belyi-curve estimate; no squarefree-level restriction
or trivial-nebentypus assumption enters it.

Apply \cite[Theorem~1.9(1), pp.~7--8]{GaudronRemond2023} with source
$(J_H)_{L'}$ and target $A\times Z$. Combining its degree bound with
half-logarithmic isogeny-height variation gives
\[
 h_F(A)+h_F(Z)\le h_F(J_H)+V,\quad
 V=2g^2\left(8g^2\log(7g)+\log d'
     +\log\max(1,\log d',h_F(J_H))\right).
\]
The Bost lower bound in \cite[(2.2), p.~5]{vonKaenel2021} gives
$h_F(Z)>-3g/2$. Put $u=\log T$. Then
\[
 V\le16T^4(2+u)+2T^2(25+7u)\le T^6.
\]
Indeed $\log d'<8+\log R\le8+u$, and the last logarithm in $V$ is
less than $17+6u$. For $T\ge16$ the inequalities
$2+\log T\le T^2/32$ and $25+7\log T\le T^4/4$ follow at $16$
and remain true by differentiation. Thus
\[
 h_F(A)<(2^{24}+1)T^6+\tfrac32T
       <2^{25}T^6=A_{19}R^{19}.
\]
The earlier CM estimate is also less than this bound.
The geometric comparison and $M_5(a)\ge1$ prove the assertion.
Finally $B_*<2^{109445}$, $210<2^8$ and $A_{19}B_*^{19}>64$ give
\[
 H_{19}<4A_{19}B_*^{19}
       <2^{\,2+109+1800\cdot8+19\cdot109445}
       =2^{2093966}.
\]
\end{proof}

% BLOCK ROOT37: no denominator-clearing variable is added to this source.
\begin{lemma}[Thirty-seven contacts for rational coefficients]
\label{lem:quartic-rational37}
Keep $K=\Q(\sqrt2)$, $d=3-\sqrt2$, $D=d^2$, and
$G=2E(K)=\Z P_*$ from Section~\ref{sec:degree8-original}.
Let $A=[n]P_*,S=[s]P_*,T=[t]P_*,Z=[z]P_*$ be nonidentity points.
Order the eight rational coefficients of
$x(A),x(S),x(T),x(Z)$ as $a_1,\ldots,a_8$. Put
\[
 R_{\rm c}=47027,\quad W_0=25878772920,\quad c=18484837801,
 \quad V_0=32W_0,
\]
\[
 A_i=R_{\rm c}a_i+3,\quad w_i=A_iv-u.
\]
Use the following thirty-seven nonzero raw contacts:
\begin{equation}\label{eq:quartic-contacts37}
 r_0=v/V_0,\qquad u-\delta v\quad(\delta\in\mathcal D),
 \qquad -(w_i-\delta v)\quad(1\le i\le8,\ \delta\in\mathcal D).
\end{equation}
Here and below the order within $\mathcal D$ is
$(0,1,2401,-343/32)$. The common five positions are numbered $0,\ldots,4$;
the other positions are numbered $5,\ldots,36$ in the displayed order.

Every normalized odd contact of $u/v$ or $w_i/v$ outside $S_*$ is
detected by an odd valuation at a raw position belonging to that slope.
If the radicals for both slopes $u/v,w_i/v$ have their $e$th powers
dividing an integer $N\ne0$, then, with $H=H_{19}$,
\[
 h(a_i)<2H|N|^{19/e}+14.
\]
Consequently the two coefficients of any one of the four abscissas give
\begin{equation}\label{eq:quartic-direct-height}
 h(x(Q))<4H|N|^{19/e}+30<7H|N|^{19/e}
 \quad(Q\in\{A,S,T,Z\})
\end{equation}
whenever all its required positions have those divisibilities.

Conversely, for any fixed eight rational coefficients, the inherited
ordinary-prime construction supplies these raw contacts with every finite
contact a positive selected prime and $r_0$ a positive rational square
times a selected prime. Repeated positions may have equal values.
The same pole-completeness construction applies to these values.
\end{lemma}
\begin{proof}
Nonzero raw contacts imply $v\ne0$ and exclude the four finite branch
values for every slope. At a prime outside $S_*$ all differences of
distinct members of $\mathcal D$ are units. Write $m$ for the minimum
of the valuations of the two coordinates of one slope. If an odd finite
normalized contact has $m$ even, its distinguished raw linear form has
odd valuation. If $m$ is odd, another finite form has valuation $m$.
For an odd contact at infinity use the denominator in the first case
and a finite form in the second. The unit $V_0$ and the signs do not
change parity. This proof uses rational coordinates and also covers
negative valuations; it does not presuppose that a contact is prime.

Apply Theorem~\ref{thm:quartic-height19} to the two slopes. The identity
$R_{\rm c}a_i+3=u/v+w_i/v$ and the elementary logarithmic height inequalities
give the first bound, since $\log(6R_{\rm c})<14$. For
$x(Q)=a+b\sqrt2$ use
$h(x(Q))\le h(a)+h(b)+h(\sqrt2)+\log2$.
This proves \eqref{eq:quartic-direct-height}, since $H>128$ and $|N|\ge1$.

For the converse choose a positive integer $D_0$ with $D_0a_i\in\Z$
for every $i$, after fixing the eight coefficients. This integer is a
choice in the existence proof, not a variable of the source. Set
\[
 T_0=V_0D_0^2,\qquad p_0=c+W_0y,\qquad
 u=c+W_0x,\qquad v=T_0p_0.
\]
For each distinct slope
$b\in\mathcal D\cup\bigcup_i(A_i-\mathcal D)$ put
$K_b=T_0b\in W_0\Z$ and $\Psi_b=u-K_bp_0$.
Together with $p_0$, these are at most thirty-seven integer affine forms,
each congruent to $c$ modulo $W_0$. Distinct slopes give independent
linear parts. Exact duplicate slopes are removed only while applying
the prime theorem and restored at all their original positions afterward.

The value and successor admissibility in
Section~\ref{subsec:o8-pattern-proof} remains valid. At primes dividing
$W_0$ use the same residue $c$. If a prime $\ell>29$ divides $D_0$,
all $K_b$ vanish modulo $\ell$, and $u=p_0=1$ avoids zero and minus one.
If $\ell>29$ does not divide $D_0$, the change from $(u,p_0)$ to $(u,v)$
is invertible. For $\ell=31,37,41$ there are at most eight slope residues,
since $R_{\rm c}a_i\equiv0$; for $\ell\ge43$ there are at most thirty-six.
Choose distinct missing residues $b_0,e_0$ and order them so that
$e_0-b_0\ne-T_0^{-1}$. Then
\[
 v=(e_0-b_0)^{-1},\qquad u=b_0v
\]
make every $\Psi_b$ different from zero and minus one, and make
$p_0=v/T_0$ different from zero and minus one. Interchanging the two
missing residues permits the required ordering, because $\ell$ is odd.

Choose a fixed positive cone
$u/v>\max_i|A_i|+2413$, $v>0$. Enlarge the fixed initial cutoff past
$82$, every prime dividing $D_0$, and the nonzero determinants of the
distinct forms. The inherited prime-pattern proof then applies with its
same unused third coordinate and with complexity at most $35$.
For its fiber upper estimates, fixing $\Psi_{b_0}=p$ gives the same
one-variable forms
$\Psi_b=p+(K_{b_0}-K_b)p_0$.
All coefficients, the cone, the denominator $D_0$ and the cutoff are
fixed before the box grows. Thus the proof uses precisely the inherited
fixed-system prime and successor estimates, not a uniform estimate in a
growing denominator. It gives the positive ordinary primes and the
same-prime elliptic-order conditions used there.

Now $r_0=D_0^2p_0$. At a prime $\ell\mid D_0$ outside $S_*$, both $u$
and $p_0$ are units, $v_\ell(v)=2v_\ell(D_0)$, and
$v_\ell(A_i)\ge-v_\ell(D_0)$. Hence $w_i=A_iv-u$ is a unit.
Both slopes have negative valuation $-2v_\ell(D_0)$, and their normalized
infinite depths are even. No new normalized odd contact is created at
such a denominator prime. At $p_0$, the odd valuation of $r_0$ is one.

For each occurrence choose its private comparison witnesses sufficiently
deep that its norm root is a unit at its assigned prime. The remaining
denominator valuations in the pole argument occur in the square of that
root. The raw contacts are positive, their reduced denominators are
squares, and they are squares in $\Q_2,\Q_3,\Q_{17}$ by the retained
residue $c$ and the square factor $D_0^2$. Lemma~\ref{hyp:o8-pole}
therefore supplies the pole witnesses. Equal prime occurrences use the
same orientation and may use separate private points. No condition on
reduction orders at different primes has been added.
\end{proof}

% BLOCK RECIPROCAL: the released complete proofs with the stated uniform range.
\begin{lemma}[Reciprocal logarithm and its error]
\label{lem:quartic-reciprocal-log}
\label{lem:quartic-uniform-reciprocal}
Put $f=1/x$, $z=-x/y$, and define
\[
 a_0=1,\qquad a_1=-\frac25,\qquad
 a_j=\frac{4j-5}{4j+1}a_{j-1}\quad(j\ge2),\qquad
 G(w)=\sum_{j\ge0}a_jw^j.
\]
On the formal branch at the identity,
\[
 \log_E(z)=zG(Df^2).
\]
For every $J\ge0$ put $G_J(w)=\sum_{j=0}^{J}a_jw^j$.
For every integer $10\le e\le64$, set $m=\lceil e/4\rceil$ and
\[
 G_{m-1}(w)=\sum_{j=0}^{m-1}a_jw^j,\qquad
 L_e(z,f)=zG_{m-1}(Df^2).
\]
Let $p>76176$ be split and choose either place above it. If the point is
nonidentity and lies in the formal group at that place, so
$r=v_p(z)\ge1$ and $v_p(f)>0$, then
\[
 v_p\bigl(\log_E(z)-L_e(z,f)\bigr)\ge(4m+1)r,
 \qquad \frac{L_e(z,f)}{\log_E(z)}\in1+p^e\Z_p.
\]
\end{lemma}
\begin{proof}
The curve equation gives
\[
 f=z^2(1-Df^2),\qquad w=Df^2,\qquad Dz^4=\frac{w}{(1-w)^2}.
\]
Coefficient comparison in the defining recurrence gives
\[
 4w(1-w)G'(w)+(1+w)G(w)=1-w.
\]
Indeed the coefficient of $w$ is $5a_1+a_0=-1$, and each later
coefficient gives $(4j+1)a_j=(4j-5)a_{j-1}$. Differentiating the
curve identity yields $z\,dw/dz=4w(1-w)/(1+w)$. Consequently
\[
 \frac{d}{dz}\{zG(w)\}=\frac{1-w}{1+w}
                   =(1+4Dz^4)^{-1/2}.
\]
The square root has constant term one. This is the normalized invariant
differential for $z$, and $zG(w)$ has zero constant term, which proves
the formal identity.

At the stated place, $D$ is a unit and the formal branch gives
$v_p(f)=2r$, hence $v_p(w)=4r$. The denominator of $a_j$, for $j\ge1$,
divides $5\prod_{k=2}^j(4k+1)$ and therefore divides $(4j+1)!$.
Thus
\[
 v_p(a_j)\ge-v_p((4j+1)!)\ge-\frac{4j+1}{p-1}.
\]
For $j\ge m$ with $4j+1<p$, the coefficient is $p$-integral and
$v_p(a_jw^j)\ge4mr$. Otherwise $j\ge19044$, so $3j\ge4m$ and
$(4j+1)/(p-1)\le j$. Therefore
\[
 v_p(a_jw^j)\ge4jr-j\ge4mr.
\]
These lower bounds tend to infinity. They control the entire infinite
tail, including terms whose denominators are divisible by $p$.
Multiplication by $z$ proves the first estimate. Formal logarithm isometry
gives $v_p(\log_E(z))=r$, and proves the relative estimate.
\end{proof}

\begin{lemma}[Global nonvanishing]
\label{lem:quartic-reciprocal-nonzero}
For every $f\in K$ and every $J\ge2$,
\[
                 G_J(Df^2)\ne0.
\]
In particular $L_e(z,f)\ne0$ whenever $z\ne0$ and $10\le e\le64$.
\end{lemma}
\begin{proof}
Let $\nu_2$ be the rational exponent-of-two valuation and put
\[
 v(a+b\sqrt2)=\tfrac12\nu_2(a^2-2b^2)
             =\min\{\nu_2(a),\nu_2(b)+\tfrac12\}.
\]
The two terms in the norm have valuations of opposite parity when finite.
Norm multiplicativity and the displayed minimum prove the valuation
identities. Its finite values lie in $\tfrac12\Z$, and $v(D)=0$.
Every $a_j$ with $j\ge1$ has valuation one, because $a_1=-2/5$ and
all subsequent recurrence factors have odd numerator and denominator.
If $f=0$, the assertion is immediate. If $t=v(f)\ge0$, the constant
term of $G_J(Df^2)$ is uniquely least in valuation. If $t<0$, then
$t\le-1/2$ and the nonconstant valuations $1+2jt$ strictly decrease;
the last is at most $1-J<0$, below the constant term as well. In either
case a uniquely least term proves nonvanishing. No good-reduction
assumption at two is used.
\end{proof}

% BLOCK SELECTIVE_SPECIFICATION: the selected beta-only source.
\subsubsection{Selective comparison systems}
For this source put
\begin{equation}\label{eq:quartic-selected-parameters}
 H=H_{19},\qquad \kappa=40,\qquad q=12,
 \qquad E_{\rm sc}=32,\qquad A_{\rm an}=25.
\end{equation}
Here $\kappa$ is a precision, while $K$ continues to denote the field.
Retain $X^2=\beta-2$, the proxy equation \eqref{eq:o8-Lucas}, all four
nonidentity root points $A,S,T,Z$, the two common auxiliary points,
and the base pole and norm predicates. In the scalar batch use orders
$(1,1,1,1,32,32)$ on the two coefficients each of
$f(J_0),f(S),z(S)/z(J_0)$. In the anchor use
$c_0=3$, $c_{i+1}=c_i^2$ for $0\le i<25$ and $C=c_{25}$.
Thus $C$ is computed by witnesses; it is not a new coefficient of the
polynomial. Replace the denominator-cleared height targets by precisely
the eight coefficients in Lemma~\ref{lem:quartic-rational37}.

Here is a uniform specification of the two comparator types. Write
$n_1=n,n_2=s,n_3=t,n_4=z$ and
$(b_1,b_2,b_3,b_4)=(\beta,\xi,\tau,\zeta)$; put
$Q_0=P_*$ and $Q_i=[n_i]P_*$. For a chosen subset
$I\subset\{1,2,3,4\}$ introduce private nonidentity primary points
$P_0=[h]P_*$ and $P_i=[m_i]P_*$ for $i\in I$. Use the base affine
subtraction charts for
\[
 R_0=P_0-[4]Q_0,\qquad R_i=P_i-[4]Q_i\quad(i\in I).
\]
Every old subtraction guard is retained: in particular the two input
abscissas are different and these differences are nonidentity. The
shared $J=[j]P_*\ne O$ and its reciprocal parameters are also retained.
No identity-allowing difference model is used in this specification.

For precisions $e_0$ and $e_i\le e_0$ define
\[
 C_i=L_{e_i}(P_i)/L_{e_0}(P_0).
\]
The ordered field tests are first
\[
 f(P_0),\ (f(P_i))_{i\in I},\ f(J)^\sigma,
 \ \phi_0,\ (\phi_i)_{i\in I},
 \qquad \phi_i=f(R_i)^\sigma,
\]
all of order one, followed by $(C_i-b_i)_{i\in I}$ of orders $e_i$,
and then $z(R_0)^\sigma/z(J)^\sigma$ of order $e_0$ and
$(z(R_i)^\sigma/z(J)^\sigma)_{i\in I}$ of orders $e_i$.
Here $\sigma(\sqrt2)=-\sqrt2$, and every list indexed by $I$ is in
increasing order. The already-conjugated common inputs are used once,
not conjugated a second time.

Before these paired field tests, retain the single rational test
$\varepsilon=\theta^2-2$ of order one. A field expression
$a+b\sqrt2$ of order $e$ gives the two rational tests
$\varepsilon b,a+b\theta$, both of order $e$.
Each resulting list has its own constant-left norm tree and the same
pole predicate as the base compiler. The two choices are
\[
\begin{array}{c|c|c|c}
 \text{type}&I&e_0&(e_i)_{i\in I}\\ \hline
 \text{full}&\{1,2,3,4\}&40&(12,40,40,40)\\
 \text{beta only}&\{1\}&12&(12).
\end{array}
\]
Their exact rational order lists are, respectively,
\begin{align*}
 &\bigl(1^{[23]},12^{[2]},40^{[8]},12^{[2]},40^{[6]}\bigr),\\
 &\bigl(1^{[11]},12^{[6]}\bigr).
\end{align*}
where $e^{[r]}$ means $r$ consecutive copies of $e$.
These have lengths $41,17$; the lists specify expressions as well as
orders. In particular the beta-only denominator has precision $12$.

Assign these types to the thirty-seven positions as follows:
\begin{equation}\label{eq:quartic-mode-partition}
\begin{array}{c|c|c}
 \text{positions}&\text{role}&\text{comparison}\\ \hline
 0\text{--}4&\text{common}&\text{full}\\
 5\text{--}12&x(A)&\text{beta only}\\
 13\text{--}20&x(S)&\text{full}\\
 21\text{--}28&x(T)&\text{full}\\
 29\text{--}36&x(Z)&\text{full}.
\end{array}
\end{equation}
There are $29$ full and $8$ beta-only positions. These are
source call counts, not polynomial variable counts. No raw position is
removed or merged when two slopes or primes coincide.

\begin{lemma}[The local conclusions of the two comparators]
\label{lem:quartic-selective-local}
At every $p>76176$ where a raw contact has odd valuation, a full
comparator gives
\[
 \beta\in\Z_p,\quad \beta\equiv n\pmod {p^{12}},\qquad
 (\xi,\tau,\zeta)\in\Z_p^3,\quad
 (\xi,\tau,\zeta)\equiv(s,t,z)\pmod {p^{40}}.
\]
A beta-only comparator gives the displayed beta conclusion, with no
Lucas conclusion asserted at that position.
No ordinary-prime hypothesis is needed in this direction.
\end{lemma}
\begin{proof}
Let $R=v_p(r)$ be the odd valuation of the raw contact. The all-left
constant path of the norm tree and Lemma~\ref{hyp:o8-pole} exclude
$R<0$. For $R>0$, descent through that tree gives, for a rational test
$Y$ assigned order $e$,
\[
 v_p(Y)>(e-1)R\ge e-1,\qquad v_p(Y)\ge e.
\]
This argument is unchanged for each of the two tree sizes and for
every positive integer order. It does not assume a prime raw value.

The root-error test makes $\theta$ a $p$-adic unit and selects a root
$\rho$ of $2$ with $\theta-\rho\in p\Z_p$ and $\theta+\rho$ a unit.
Thus $p$ splits. Since
$\theta-\rho=\varepsilon/(\theta+\rho)$, the paired inequalities
$v_p(\varepsilon b),v_p(a+b\theta)\ge e$ imply
$v_p(a+b\rho)\ge e$, even when $b$ is not integral.
The order-one $f$ tests put all primary points in the formal group at
this place, and $J,R_0,(R_i)$ in the formal group at its conjugate.

Formal-logarithm homomorphism and isometry
\cite[IV.6.4 and VII.2.1--2.2]{silverman} give
\[
 v_p((h-4)/j)\ge e_0,\qquad
 v_p((m_i-4n_i)/j)\ge e_i.
\]
Here $j$ and all indices are integers and $j\ne0$. Therefore
$h\equiv4\pmod {p^{e_0}}$, $m_i\equiv4n_i\pmod {p^{e_i}}$,
and $h$ is a unit. The same homomorphism at the primary place gives
$\log_E(P_i)/\log_E(P_0)=m_i/h\in\Z_p$.
Lemma~\ref{lem:quartic-uniform-reciprocal} and $e_i\le e_0$ imply
$C_i\equiv m_i/h\pmod {p^{e_i}}$. The remaining primary tests yield
$b_i\in\Z_p$ and $b_i\equiv n_i\pmod {p^{e_i}}$, as claimed.
\end{proof}

\begin{lemma}[Scalar height and the new anchor]
\label{lem:quartic-height19-anchor}
Every solution of the specified scalar batch and anchor has, with
$a=|s|$,
\[
 s\in2\cdot3^{2^{24}-1}\Z\setminus\{0\},\qquad
 h(\beta)\le H a^{19/32}<a/40,
\]
and
\begin{equation}\label{eq:quartic-cutoffs19}
 a>(24H)^{10},\qquad a>(40H)^{32/13},\qquad
 a>40,\qquad \log(40H)<a/40.
\end{equation}
Conversely this anchor can be imposed together with any fixed finite
list of formal divisibilities on $s$.
\end{lemma}
\begin{proof}
The exact anchor subgroup calculation in
Section~\ref{subsec:o8-anchor-subgroup}, with $25$ squarings, gives the
stated subgroup. Its two local computations are unchanged; the exponent
of the final power of three is now $2^{25}$.
In particular
\[
 \log_2 a>\tfrac32(2^{24}-1)=25165822.5.
\]
The binary height bound and $\log_2 24<5$, $\log_2 40<6$ give
\[
 10\log_2(24H)<20939710<25165822.5,
\]
\[
 \tfrac{32}{13}\log_2(40H)
 <\tfrac{32}{13}\,2093972<5154400<25165822.5.
\]
These prove the first two cutoffs and $a>40$.

Apply the same constant-left norm argument to the scalar batch. At each
odd contact prime of $\beta$ outside $S_*$ it excludes negative odd raw
valuation and gives the six assigned coefficient orders. The first four
make $J_0=[j_0]P_*$ and $S$ formal at every place over that prime. The
last two imply $v_p(s/j_0)\ge32$, hence $v_p(s)\ge32$. This reasoning
also applies at inert primes: both rational coefficients are tested and
the extension is unramified outside two. The retained scalar raw
polynomial and its small-prime correction detect every prime of
$M_5(\beta)$, including an odd infinite contact. Therefore
\[
 M_5(\beta)^{32}\mid s,
 \qquad h(\beta)\le H M_5(\beta)^{19}\le Ha^{19/32}.
\]
The domain hypothesis holds because $\beta=X^2+2$ is in none of the four
branch values: the only positive candidate $2401$ would require the
nonsquare integer $2399$ to be a rational square. The second cutoff
now gives $h(\beta)<a/40$.
Since $32/13>2$, that cutoff also gives $a>(40H)^2$.
For $H\ge1$, $40\log(40H)<(40H)^2$ follows at $H=1$ and then by
differentiating the difference. This proves the last bound.
Finally the displayed anchor subgroup contains arbitrarily divisible
nonzero integers, exactly as required in the base completeness argument.
\end{proof}

% BLOCK FIRST_CONTRADICTION.
\begin{lemma}[The Lucas comparison needs twenty-nine positions]
\label{lem:quartic-selective-lucas-growth}
The assignment in \eqref{eq:quartic-mode-partition}, together with
the proxy equation, implies
\[
 t^2-s(s-2)tz+(s-2)^2z^2=1,
 \qquad |t|>\exp(|s|-3).
\]
\end{lemma}
\begin{proof}
Put $a=|s|$, $B_L=\max(|s|,|t|,|z|)$ and
\[
 \Delta=t^2-s(s-2)tz+(s-2)^2z^2-1.
\]
Every finite position of the six coefficients of $x(S),x(T),x(Z)$,
as well as all five common positions, retains all three Lucas
comparisons. Lemma~\ref{lem:quartic-selective-local}
and the proxy equation therefore give $p^{40}\mid\Delta$ at every
odd raw-contact prime needed by these six coefficients. By the parity
argument in Lemma~\ref{lem:quartic-rational37}, each of their normalized
contact radicals has its fortieth power dividing $|\Delta|$ if
$\Delta\ne0$. Choose the point among $S,T,Z$ whose index has size $B_L$.
The index-height estimate \eqref{eq:o8-indexheight} and
\eqref{eq:quartic-direct-height} give
\[
 B_L^2-6<7H|\Delta|^{19/40}
       <7\cdot5^{19/40}HB_L^{19/10}<16HB_L^{19/10},
\]
where $|\Delta|<5B_L^4$. As $B_L,H\ge1$, this implies
$B_L^{1/10}<22H<24H$, hence $B_L<(24H)^{10}$.
It contradicts $B_L\ge a$ and \eqref{eq:quartic-cutoffs19}.
Thus $\Delta=0$. Since $z\ne0$ and $a>40$, Lemma~\ref{lem:o8-lucas}
gives the asserted bound on $|t|$ itself. The point $A$ and index $n$
have not been used in this first contradiction.
\end{proof}

% BLOCK BETA_ONLY_VARIANT: retain this block for the 29/8 endpoint.
\begin{lemma}[Soundness of the beta-only variant]
\label{lem:quartic-beta-only}
The beta-only assignment in \eqref{eq:quartic-mode-partition} implies
$\beta=n$ and $X\in\Z$.
\end{lemma}
\begin{proof}
Suppose $\beta\ne n$, write its difference in lowest terms, and put
$N=|\operatorname{num}(\beta-n)|\ge1$.
All thirty-seven positions retain a beta comparison. At every relevant
prime $p$, the local conclusion gives $v_p(\beta-n)\ge12$; reducedness
then shows $p\nmid\operatorname{den}(\beta-n)$ and $p^{12}\mid N$.
Every normalized radical for all eight coefficients is covered.
Set $B=\max(|n|,|s|,|t|,|z|)$. The upper and lower heights of the point
realizing $B$ give
\[
 B^2<7HN^{19/12}+6,
 \qquad N\le2\exp(h(\beta))B.
\]
The second estimate follows by writing $\beta=b/c$ in lowest terms:
$|b|,|c|\le\exp(h(\beta))$ and $|n|\le B$.
Since $7\cdot2^{19/12}+6<40$ and $H,B\ge1$,
\[
 B^{5/12}<40H\exp(19h(\beta)/12),
 \qquad
 \log B<\tfrac{12}{5}\log(40H)+\tfrac{19}{5}h(\beta)
         <\tfrac{31}{200}a<a-3.
\]
But $B\ge|t|$, contradicting
Lemma~\ref{lem:quartic-selective-lucas-growth}. Hence $\beta=n\in\Z$.
The equation $X^2=\beta-2$ forces a rational $X$ to have denominator
one, so $X\in\Z$. This second step uses the larger maximum $B$;
there is no unproved bound on $n$ in terms of $B_L$.
\end{proof}

% BLOCK SOURCE_THEOREM: the selected 29-full / 8-beta source.
\begin{theorem}[Arithmetic interpretation of the selective source]
\label{thm:quartic-source-arithmetic}
With the arithmetic inputs and source-version conventions of
Section~\ref{sec:degree8-original}, the selective system specified
above has rational witnesses at input $X\in\Q$ if and only if $X\in\Z$.
The assertion also holds for its compressed primary-chart presentation
described below. The original input coordinate is unchanged.
\end{theorem}
\begin{proof}
Soundness is Lemma~\ref{lem:quartic-beta-only}.
For completeness let $X\in\Z$, set $\beta=n=X^2+2$ and $A=[n]P_*$.
Choose the scalar sign and small-prime squareclass correction as in the
base construction. Choose $J_0$ sufficiently deep at its finite set of
scalar places and then choose a positive nonzero $s$ in
$2\cdot3^{2^{24}-1}\Z$ with all additionally required divisibilities.
Arbitrarily deep choices give the scalar order $32$ conditions, make the
normalized scalar leaves integral and the norm root a unit at the
assigned primes, and preserve the local squareclasses at $2,3,17$.
The pole theorem supplies the scalar witnesses.
Lemma~\ref{lem:o8-lucas} supplies positive $t,z$ with the required
index equation. Put $S=[s]P_*,T=[t]P_*,Z=[z]P_*$ and set
$(\xi,\tau,\zeta)=(s,t,z)$.

Now fix the eight rational abscissa coefficients and apply
Lemma~\ref{lem:quartic-rational37}. Let $\mathcal P$ be the set of
distinct selected ordinary primes, with the same orientation and
reduction orders $d_{p,1},d_{p,2}$ as in the base proof. Put
\[
 j=\operatorname{lcm}_{p\in\mathcal P}d_{p,2},\qquad J=[j]P_*.
\]
At a position assigned to $p$ choose a private integer $h_p$ satisfying
\[
 h_p\equiv0\pmod {d_{p,1}},\qquad
 h_p\equiv4\pmod {d_{p,2}},\qquad
 h_p\equiv4\pmod {p^{40+v_p(j)}}.
\]
The inherited same-prime conditions
$\gcd(d_{p,1},d_{p,2})=4$ and $p\nmid d_{p,1}d_{p,2}$ prove compatibility.
The factor $v_p(j)$ includes orders from every other selected prime;
no cross-prime coprimality is assumed. Avoid $h_p=0,4,-4$.
For the type at this position take only its required points
\[
 P_0=[h_p]P_*,\qquad P_i=[h_pn_i]P_*\quad(i\in I).
\]
All are nonidentity. If a primary and its target had equal abscissa,
they would be equal or opposite, forcing $h_p=4$ or $-4$ because
$G$ is torsion-free and the corresponding $n_i$ is nonzero. Thus every
retained affine subtraction guard holds. Their differences have indices
$(h_p-4)n_i$, with $n_0=1$.

The primary points are formal at the first place, and $J$ and all
required differences at the second. The displayed congruences and
Lemma~\ref{lem:quartic-uniform-reciprocal} give every retained test order;
in beta-only positions the common precision $40$ choice is stronger
than the required precision $12$. Increase the private congruence
precision when necessary to make the norm leaves as deep as required.
Repeated occurrences of a prime have the same orientation and otherwise
independent private witnesses. Choose the single rational $\theta$ last,
to meet all actual coefficientwise requirements across all positions
and comparator types. Finite Chinese remaindering supplies it.

Each norm tree retains its constant leaf on the all-left path. At the
assigned prime its root is a unit; the positive raw value times its
square has square reduced denominator and the three required local
squareclasses. The last part of Lemma~\ref{lem:quartic-rational37}
therefore supplies every remaining pole witness. This constructs a
solution with the given $X$. The compressed presentation is equivalent
by the following explicit maps.
\end{proof}

% BLOCK COMPRESSED: replaces the old five-point-only lemma and G15 circuit.
\begin{lemma}[Compressed charts and exact evaluation for each type]
\label{lem:quartic-compressed-primary}
For each required primary point use $u_i,v_i,f_i,x_i,p_i\in K$ with
\[
 f_i=u_i^2,\qquad v_i^2+df_i=1,\qquad
 p_i^2=x_i-Df_i,\qquad x_if_i=1.
\]
Put $G_i=G_{\lceil e_i/4\rceil-1}(Df_i^2)$, with the same convention
for $G_0$, and replace the primary ratios by
\[
 \Lambda p_0=G_0,\qquad B_i=C_ip_i,\qquad B_i\Lambda=G_i
 \quad(i\in I).
\]
This presentation has exactly the same rational solution projection as
the full charts and ratios for the specified comparator. Every external
point coordinate and the rational input $X$ are fixed.
\end{lemma}
\begin{proof}
From a full chart set $p_i=y_if_i$, $\Lambda=-L_{e_0}(P_0)=G_0/p_0$
and $B_i=C_ip_i$. The old equations give the four displayed chart
equations and $z_i=-1/p_i$, so the ratio equations have exactly this sign.
Conversely global nonvanishing of every $G_i$ makes $\Lambda,p_i$
nonzero. Also $x_if_i=1$ gives $f_i\ne0$, and
\[
 f_ip_i^2=1-Df_i^2=v_i^2(1+df_i)
\]
makes $v_i\ne0$. Restore
\[
 y_i=x_ip_i,\qquad z_i=-1/p_i,\qquad
 h_i=u_ip_i,\qquad w_i=u_ip_i/v_i,
 \qquad L_i=z_iG_i.
\]
These satisfy every equation of \eqref{eq:o8-chart}; the old common
logarithm inverse is $1/L_0$ and $C_i=L_i/L_0$.
In the subtraction equations replace each $y_i$ by $x_ip_i$, keeping
all their guards and all difference parameters. No other coordinate
has been removed. These forward and reverse maps are rational on the
specified nonzero loci and preserve each remaining equation.
\end{proof}

The finite evaluations used here can themselves be written quadratically.
For precision $40$ bind four field pairs
\[
 w=Df^2,\qquad w_2=w^2,\qquad w_3=w_2w,\qquad U_6,
\]
with the fourth pair defined by the quadratic equation
\[
 U_6=\frac{a_3+a_4w+a_5w_2+
                  w_3(a_6+a_7w+a_8w_2+a_9w_3)}{a_6}.
\]
The quadratic output expression is
\[
 G_9(w)=a_0+a_1w+a_2w_2+a_6w_3U_6,
 \qquad a_6=-\frac{418}{16575}\ne0.
\]
For precision $12$ bind only $w=Df^2$ and use
$G_2(w)=1+a_1w+a_2w^2$. Substitute these expressions directly into
$\Lambda p_0=G_0$ and $B_i\Lambda=G_i$. Every product is quadratic in
the bound rational coordinates. Expansion gives each required $a_j$
coefficient exactly. The scale is one; no finite-error estimate is used
as a substitute for these polynomial identities.

% BLOCK SOURCE_BOUNDARY: replace old source counts/pins only from final artifact.
The source defined by these equations is a new arithmetic specification.
In particular neither deletion of unused comparisons nor the changes to
the denominator logarithm and norm tree assert a coordinate-forgetting
map between arbitrary old and new witness tuples. The preceding theorem
proves both implications for the rational input projection. The explicit
maps in the compressed-chart lemma concern only those chart changes.
The base compiler remains construction machinery; the polynomial
presented below is the single selected quartic obtained from this source.
Its independent source check must verify the complete modules, the
ordered expression lists and position assignment, all nonzero guards,
the original affine height targets, scalar orders, anchor terminal
coordinate, and the full flattening. The subsequent quadratic lift and
certificate replay preserve the same input coordinate. Those finite
checks establish source and polynomial identities; the arithmetic
projection statement rests on the theorem above with the inherited
arithmetic inputs and their existing source-version qualifications.


The resulting source has $9\,318$ rational coordinates, including input
zero, and $8\,854$ residuals. The source file
\nolinkurl{source-residuals.jsonl.gz} in the quartic package has SHA-256
\begin{center}\footnotesize\ttfamily
97b7bc300084afaefe9980f1feea760be8b4dedaf6eb788fcc4c3af440262579.
\end{center}
Its uncompressed bytes have SHA-256
\begin{center}\footnotesize\ttfamily
d68d86cdf995a0e80c38953f1ad03262733f81145b949a0501b143eca6ac0a9c.
\end{center}
The entry \nolinkurl{source-entry.json} gives this source the distinct name
\begin{center}\footnotesize\ttfamily
selective-beta-k40-q12-a25-h19-s32-d68d86cdf995.
\end{center}
The \nolinkurl{source/} directory contains its eight modules, shared
chain, flattened constraints and compiler copies. The independent source
audit reconstructs the two comparison types and their exact ordered
arguments, verifies every retained baseline module, and checks the
complete flattening. These finite checks are separate from
Theorem~\ref{thm:quartic-source-arithmetic}.

\subsection{Quadratic reductions preserving the rational input projection}
\label{subsec:quartic-reductions}

We describe the algebraic reductions applied to the new, fixed source.
All scalar coordinates range over $\Q$, and coordinate zero is the original
input $X$. Each field equation denotes its two rational coefficient
equations. These reductions prove equality of rational input projections
to that source; its arithmetic interpretation is
Theorem~\ref{thm:quartic-source-arithmetic}.

\paragraph{Definitions, row operations, and interfaces.}
An equation $cv-F=0$, with $c\in\Q^*$ and $F$ independent of $v$,
permits substitution $v=F/c$ in every other equation and deletion of $v$.
Restriction is the forward solution map, and adjoining $v=F/c$ is its
inverse. Each use below retains degree at most two after substitution.
The initial quadratic lift uses the reverse operation: introduce targets
for products of two available coordinates until the source residuals are
quadratic. Evaluating these product definitions and forgetting their
targets give the two solution maps. A previously defined product can be
reused only with its retained defining equation.

An operation replacing $R_i$ by $aR_i+bR_j$, while retaining $R_j$,
has inverse $R_i=(R_i'-bR_j)/a$ when $a\ne0$ and $i\ne j$.
Thus every recorded linear row operation preserves the simultaneous zero
set. A shear replaces an old witness by $v=v'+cw$, with fixed rational
$c$ and another witness $w$ retained. Its forward and reverse coordinate
maps are $v'=v-cw$ and $v=v'+cw$. This is an invertible rational coordinate
change and preserves quadratic degree.

Every removed internal coordinate must have no occurrence outside the
replaced rows, unless the stated substitution is made in \emph{every}
such row. Retained coordinates keep their identities and values. All new
coordinates are fresh, and input coordinate zero is never a target of
deletion or redefinition. These conditions apply to the current system
after all preceding transformations.

\paragraph{Norm inverses.}
For rational affine expressions $A,B$, replace the constraint
$u(A^2-2B^2)=1$ and its private lifting definitions by
\[
 As+2Bt=1,\qquad At+Bs=0.
\]
The forward assignment is $s=Au$, $t=-Bu$; the reverse assignment is
$u=s^2-2t^2$. The identity
\[
 (A^2-2B^2)(s^2-2t^2)=(As+2Bt)^2-2(At+Bs)^2
\]
proves the reverse implication. All other private lifting coordinates
are reconstructed through their old definitions. No additional
nonvanishing premise or square-root choice is used.

\paragraph{Private rational inverse guards.}
Suppose the retained equations include
\[
 \varepsilon=\theta^2-2,\qquad R(F-1)=\varepsilon^2.
\]
Rational nonsquareness of $2$ gives $\varepsilon\ne0$ and hence $R\ne0$.
If a rational coordinate $u$ occurs only in $uR=1$, remove it and that
equation. The reverse map $u=1/R$ is defined on every remaining solution.

\paragraph{Two reductions of the point chart.}
Each retained full root chart has field coordinates $u,v,w,z,f,h,x,y$ and equations
\[
 f=u^2,\quad v^2+df=1,\quad w^2+v^2=2,\quad h=vw,
 \quad zh=-u,\quad xf=1,\quad yz=-x.
\]
Replace the private coordinates $w,h$ by $p$ and use
\[
 f=u^2,\quad v^2+df=1,\quad p=yf,\quad p^2=x-Df,
 \quad xf=1,\quad yz=-x.
\]
Forward, put $p=yf$. The inverse and parameter equations give
$x,f,u,y,z\ne0$. From $zh=-u$ one obtains $xh=yu$ and $h=up$.
Therefore
\[
 fp^2=h^2=v^2w^2=(1-df)(1+df)=1-Df^2,
\]
which proves $p^2=x-Df$. Conversely, the new equations give $p\ne0$
and $fp^2=v^2(1+df)$. Thus $v\ne0$. Define
\[
                         h=up,\qquad w=up/v.
\]
Then $h=vw$, $w^2=1+df=2-v^2$, and
$zh=zu yf=-u$. These formulas reconstruct the old coordinates exactly.
Both systems are quadratic in rational components; for example the
two new equations $p^2=x-Df$ are
\[
 p_0^2+2p_1^2=x_0-11f_0+12f_1,\qquad
 2p_0p_1=x_1+6f_0-11f_1.
\]

In a second step, retain $xf=1$, replace $p=yf$ and $yz=-x$ by
$pz=-1$, and substitute $y=xp$ in every outside equation.
Indeed the old equations imply both new relations, while adjoining
$y=xp$ gives $yf=p$ and $yz=-x$. This step is used only where all
substituted equations remain quadratic. For example a subtraction line
$\lambda x-y=\eta$ becomes $x(\lambda-p)=\eta$. No outside occurrence
of either rational component of $y$ is omitted.

\paragraph{A quadratic presentation of the parity block.}
This replacement concerns only its rational equations, with unchanged
input $b$. Its old private variables $t,U,e,p,q,s,h$ satisfy
\begin{align*}
 U(1+bt^2)&=b(1-bt^2),&
 p^2&=17b(1+2e-3e^2),\\
 h&=e(17b-U^2),&
 17bq^2&=(s^2-17bh)(s^2+51bh).
\end{align*}
Use instead twelve private rational variables
$u,v,a,z,e,p,c,w,s,q,j,k$ and nine quadratic equations
\begin{gather*}
 z=u^2,\quad a=bv,\quad z+va=b^2,\quad c=be,\\
 p^2-17b-34c+51ec=0,\qquad w-17bc+cz=0,\\
 j=s^2,\qquad k=bq,\qquad
 j^2+34jw-867w^2-17qk=0.
\end{gather*}
For the forward map put $\delta=1+bt^2$. It is nonzero: otherwise
$bt^2=-1$ and the first old equation gives $0=2b$, a contradiction.
Keep $e,p,s,q$ and set
\[
 u=U,\quad v=2bt/\delta,\quad a=bv,\quad z=U^2,\quad
 c=be,\quad w=bh,\quad j=s^2,\quad k=bq.
\]
The identity
$[b(1-bt^2)]^2+b(2bt)^2=b^2(1+bt^2)^2$ proves the first conic;
substitution proves the other equations.

Conversely $u^2+bv^2=b^2$. Set $h=e(17b-z)$ and keep $e,p,s,q$.
If $b=0$, take $U=t=0$. If $b\ne0$ and $u+b\ne0$, take
$U=u$ and $t=v/(u+b)$. In this case the first old residual, multiplied
by $(u+b)^2$, is
\[
 (u-b)(u+b)^2+b(u+b)v^2
       =(u+b)(u^2+bv^2-b^2)=0.
\]
If $b\ne0$ and $u=-b$, then $v=0$; take $U=b$ and $t=0$.
In every case $U^2=z$, and the remaining old equations follow from
$c=be$, $w=bh$, $j=s^2$, and $k=bq$. The possible sign change of $U$
is harmless because all old witnesses and deleted lifting coordinates
are private. These maps prove equality of existential projections;
they are not asserted to be inverse on the private witness tuples.

\paragraph{Absorbing the root error in a norm pair.}
Keep $\varepsilon=\theta^2-2$ and $R(F-1)=\varepsilon^2$, so $R\ne0$.
A bottom norm row, after expanding its private definitions, has the form
\[
 PN-\varepsilon^2B^2-RT^2=0,
\]
where $P$ is the retained even contact power, $B$ is affine, and $N$
is the retained bottom norm coordinate. Introduce $O,Q$ and replace it by
\[
 RO=P,\qquad Q=B^2,\qquad ON-FQ+Q-T^2=0.
\]
All three equations are quadratic. Multiplication of the last equation
by $R$ gives the old one using the first two equations and
$R(F-1)=\varepsilon^2$. Conversely, divide by the already nonzero $R$.
More explicitly,
\begin{align*}
 PN-\varepsilon^2B^2-RT^2
 ={}&-N(RO-P)+Q(RF-R-\varepsilon^2)\\
    &+\varepsilon^2(Q-B^2)+R(ON-FQ+Q-T^2).
\end{align*}
The forward assignments are $O=P/R$ and $Q=B^2$. For a retained
identity $P=R^2$, simply use $O=R$. A square coordinate can be shared
only when its proposed defining squares are identical polynomials.

Reconstruct deleted aliases and lifting targets in their original
acyclic order for the reverse map. Any definition with an outside use,
and every dependency it needs, is retained. The norm coordinate $N$,
all upper norm nodes, the final norm, and the parity argument keep their
values. Hence this step requires no invariance assertion about the
parity predicate under a change of square class or leaf order.

\subsection{The integer quartic and exact finite verification}

After the reductions, clear rational denominators in each nonzero quadratic
residual and divide its integer coefficients by their positive greatest
common divisor. These are nonzero rational rescalings and preserve each
equation. Relabel the active witness coordinates consecutively, leaving
input coordinate zero fixed. The file \nolinkurl{residuals.json.gz}
specifies the resulting $9\,161$ nonzero integer residuals. Every weight
in \nolinkurl{weights.json} is one, and the specified polynomial is
\[
 \Phi(X_0,X_1,\ldots,X_{9624})=\sum_{i=1}^{9161}R_i^2.
\]
Each $R_i$ has degree at most two. At rational values every square is
nonnegative, so this sum vanishes exactly when every residual vanishes.

The file \nolinkurl{certificate.jsonl.gz} records the fresh quadratic
lift from the reciprocal source and $1\,246$ subsequent transformations:
\[
\begin{array}{lr}
\text{scalar norm inverse replacements}&161\\
\text{private contact inverse removals}&37\\
\text{full root chart replacements}&6\\
\text{ordinate eliminations}&2\\
\text{parity block replacements}&38\\
\text{combined bottom norm absorption}&1\\
\text{invertible linear row operations}&312\\
\text{rational definition eliminations}&685\\
\text{invertible rational shears}&4
\end{array}
\]
The compressed primary charts and their reciprocal evaluations belong to
the source itself. They are not old logarithm-replacement certificate events.
The absorption record contains all its affected bottom norm pairs.
The certificate's SHA-256 digest is
\begin{center}\footnotesize\ttfamily
8aaa2249eda6972c441d58b212ec64a53e1202d5e9f74ea4cbc53e7e3f71a0bd.
\end{center}

The independent verifier
\nolinkurl{quartics/quartic-9625/verify_certificate.py}
checks the named source entry, every replacement identity, the defining
order of reconstructed coordinates, and every occurrence of a removed
variable. It verifies fresh coordinates and the current-row hypotheses of
each algebraic rule, rejects changes to input zero, and rejects a final
residual of degree above two. Its checks use the current system after all
preceding transformations. The final variable map is a checked bijection
fixing the input. The verifier's SHA-256 digest is
\begin{center}\footnotesize\ttfamily
cf587ff9c28af5593e346eb7484d1713bd3f1ff77d28aa5ada06b2b92909d6f4.
\end{center}

The source construction audit recorded in \nolinkurl{source-audit.json}
is a separate check, with checker SHA-256
\begin{center}\footnotesize\ttfamily
41a3d780636c1f28852762ee0ce7a252e207ec74e7367da275676200a6068ffa.
\end{center}
It checks the actual reciprocal polynomial identities, compressed-chart
equations, ordered tests and full module expansion. Its scope is source
construction. The local approximation estimates, global nonvanishing,
and the two arithmetic contradictions are proved above.

The algebraic verifier independently expands the final residual squares and
compares the complete collected coefficient table with the distributed
polynomial. This establishes the degree, monomial count, occurrence of every
variable, and exact coefficient maximum. A bound for a local replacement
alone would be insufficient: the full expansion includes all instances and
all cross terms. The version and file digests identify the exact finite
objects, rather than an unchecked change to an earlier source pin.

\begin{proof}[Proof of Theorem~\ref{thm:quartic}]
Theorem~\ref{thm:quartic-source-arithmetic} and
Lemma~\ref{lem:quartic-compressed-primary} identify the rational input
projection of the new source with $\Z$, using precisely the stated inherited
arithmetic inputs. For each fixed rational input, the initial product
definitions and the displayed solution maps preserve the existence of
rational witnesses through every recorded transformation. The simultaneous
vanishing of the final residuals is therefore equivalent to that of the
source. Their unweighted sum of squares has the same rational input
projection. The complete integer expansion gives the asserted counts and
coefficient maximum; it contains a nonzero degree-four term, so its degree
is exactly four.
\end{proof}

The finite source and algebraic checks do not replace the height estimate,
elliptic-group calculation, prime-pattern theorem, or norm and pole arguments
inherited from the base construction.
% End explicit quartic constructions

\appendix
\section{Dependency order in the arithmetic argument}\label{app:obligations}
The even coefficient bound uses stationary-base selection as well as
integrality. The odd minimum uses a fixed nonzero witness and the prior
genus-index lower bound. Support removal is applied only to the actual
unramified specialization. The functional bound and determinant clearing
are proved separately, and the same corrected coordinate is used at every
address. These constructions are given in Sections~\ref{sec:rc}--\ref{sec:limits}
and Appendices~\ref{app:even-source}--\ref{app:endpoints}.

The conversion to ordinary integers uses the canonical model and total
multiplication, followed by the exact computation relation. Only the
shared-shift separation argument in Section~\ref{sec:bridge} identifies
the external rational variable with the ordinary value of its integer
index. The literal substitution and size calculation then apply to the
same expression $P_*$.

\section{Source identities and scope}\label{app:sources}
The companion \texttt{SOURCE\_MAP.json} records exact paths, hashes and
locators, distinguishing mathematical and construction sources, artifacts,
and version records. The base collection contains 83 sources; parity, construction
and uniformity supplements are in separate addenda. The principal local
sources are listed below; full paths
are in the source map.
\begin{longtable}{p{.07\textwidth}p{.58\textwidth}p{.25\textwidth}}
\toprule ID & Source basename & Use\\\midrule\endhead
S01 & \texttt{ordinary-integer-polynomial.js.txt} & Outer construction\\
S02 & \texttt{rational-circuit-translator.js.txt} & Substitution semantics\\
S03 & \texttt{finite-polynomial-expressions.js.txt} & Fixed constructors\\
S04 & \texttt{variable-shift-polynomial-library.js.txt} & Close and Test\\
S05 & \texttt{elliptic-generator-proof.txt} & Generator certificate\\
S06 & \texttt{explicit-elliptic-data.txt} & Rank, torsion, codes\\
S07 & \texttt{verify-elliptic-generator.js.txt} & Exact finite algorithm\\
S08 & \texttt{six-contact-height-proof.txt} & Contacts and constants\\
S09 & \texttt{z-in-q-height-independent-audit.txt} & Height proof\\
S10 & \texttt{explicit-pole-predicate.txt} & Pole polynomial\\
S11 & \texttt{z-in-q-pole-completeness-audit.txt} & Simultaneous descent\\
S12 & \texttt{parity-frozen-4eaee875.tex} & Raw local matrices\\
S13 & \nolinkurl{multiplication-certificate-correctness-proof.txt} & Multiplication proof\\
S14 & \texttt{twenty-two-prime-pattern-proof.txt} & Selected primes\\
S15 & \texttt{z-in-q-variable-shift-bridge.txt} & Ordinary-value graph\\
S16 & \texttt{z-in-q-ordinary-integer-assembly-proof.txt} & Outer composition\\
S17 & \texttt{z-in-q-restricted-converse-assembly.txt} & RC assembly inputs\\
S18 & \texttt{ring-limits-frozen-4eaee875.tex} & Finite models and limits\\
S19 & \texttt{paired-functional-and-uniformity-proof.txt} & Paired functional\\
S20 & \texttt{determinant-transfer-audit.txt} & Clearing and transfer\\
S21 & \texttt{variable-shift-precision.py.txt} & Exact decision program\\
S22 & \texttt{precision-register.js.txt} & Ideal IR and lowering\\
S23 & \texttt{specialize-precision.js.txt} & Wrapper specialization\\
S24 & \texttt{integer-halting-correctness-proof.txt} & Natural compiler\\
S25 & \texttt{z-in-q-explicit-cutoff-bound.txt} & Literal cutoff\\
S26 & \texttt{z-in-q-polynomial-size-bound-proof.txt} & Syntax/norm argument\\
S27 & \texttt{z-in-q-numerical-polynomial-bounds.txt} & Tower comparison\\
S28 & \texttt{determinant-switch-audit.txt} & Switches and outer rank\\
S29 & \texttt{ring-determinants-frozen-4eaee875.tex} & Ring data and derivatives\\
S30 & \texttt{preliminaries-frozen-4eaee875.tex} & Selmer coranks and lengths\\
S31 & \texttt{fixed-certificate-constants.js.txt} & Exact least-multiple search\\
S32 & \texttt{pointwise-frozen-4eaee875.tex} & Distinct point indices\\
\bottomrule
\end{longtable}
The supplement manifest also records the revised fixed-degree CM proof,
the complete natural compiler proof, actual arithmetic and analytic
producer arguments, finite-Artin and determinant-count clarifications,
the fixed BFH witness application, exact substitution checks, and the
new fixed-form, weighted-theta, operator and normalization derivations.
For the height estimate, the derivation in Section~\ref{sec:height}
supersedes the historical arithmetic-height arguments in S08--S09 and
the earlier CM supplements. Those historical alternatives are not
additional premises. Primary theorem applications remain explicit in
the mathematical text.
The version records identify the source files used in the construction.

% End 05-bounds-appendices.tex


% Begin 06-literal-construction.tex
\section{Literal finite-constructor specification}\label{app:literal}
The supplementary file \nolinkurl{ordinary-integer-definition.json}
is part of this specification. It is a finite directed acyclic description,
not an expanded monomial list. The following rules determine one polynomial
from its bytes. Diagnostic labels and intended-theorem strings add no equations.

The root is \texttt{integers}. Variable zero is the external input and
variables one through $292890231904$ are its rational witnesses.
A nonnegative point reference $r$ means $X_r,\ldots,X_{r+4}$.
References $-1,-2,-3$ respectively mean
\[
 (1,0,0,0,0),\quad(0,g_0,-145,g_1,g_2),\quad
                      (0,g_0,-145,-g_1,-g_2).
\]
Every model recipe conjoins
$8g_0-1641=0$, $4g_1+15465=0$, $32g_2-87483=0$.

\subsection{Scalar nodes and fixed expressions}
Read nodes in their stored topological order.
A \texttt{const} is its decimal integer; \texttt{var} is its local variable;
\texttt{add} and \texttt{mul} are ordinary polynomial operations.
For a \texttt{call} to a template with $p$ parameters, substitute the
specified node polynomials for variables $0,\ldots,p-1$, and $X_{w+j}$
for variable $p+j$, where $w$ is the witness start. The value is the
callee's entire root polynomial.

A \texttt{fixed\_integer} evaluates a finite expression using literals,
references to preceding definitions, addition, multiplication,
nonnegative powers, maximum, factorial, the catalogue product or maximum,
and the least formal multiple. The catalogue uses the first 22 primes
in increasing order in each unit class. The searches are deterministic
and terminate by Section~\ref{sec:size}. They produce constants, not
quantified witnesses.

A \texttt{psi\_component} uses the division-polynomial recurrence in
Section~\ref{sec:syntax}, pair multiplication
$(a,b)(c,d)=(ac+2bd,ad+bc)$ and $A=(-11243,7950)$.
It returns component zero or one; every index is fixed.
The positive fixed point codes use
\[
 x_n=x(P_0)-\frac{\psi_{n+1}(P_0)\psi_{n-1}(P_0)}{\psi_n(P_0)^2},
 \qquad y_n=\frac{\psi_{2n}(P_0)}{2\psi_n(P_0)^4}.
\]
Nontorsion makes these fixed denominators nonzero.
Negative indices negate both ordinate components.

A \texttt{finite\_guard\_product} ranges over
$(a,b,c)\in[-T,T]^3\cap\Z^3$ with $ab\ne c$, with fixed $T$.
Write the fifteen coordinates of the three codes as reduced
$A_j/B_j$, $B_j>0$, and put $D=\prod_jB_j$. Its factor is
\[
                 \sum_{j=1}^{15}(DX_j-(D/B_j)A_j)^2.
\]
Multiply these factors, with empty product one. This is an integer
polynomial with precisely the finite false-triple zero set used above.

\subsection{Recipe constraints and witness scope}
A model-evaluation recipe sums squares of the generator equations,
listed natural calls, scalar-constant calls, add/multiply/power-value calls,
and the five coordinate differences of \texttt{rootEquality}.
A source-variable row with a non-null natural witness start contributes
exactly that natural call. Operation kinds zero, one and two select
addition, multiplication and power value. Source-node and provenance
fields add no equations.

A sum-of-squares composition sums squares of its callee polynomials;
a product composition multiplies them. Parameter \texttt{[var,i]} means
$X_i$. Each private block uses its recorded start and length.
The checked ranges are disjoint from parameters and other private blocks.
These rules include globally conjoined fixed equations in active callees.

The following parameter and witness counts are extracted from the exact
specification, without expanding its polynomial.
\begin{longtable}{p{.47\textwidth}rrr}
\toprule Template & Parameters & Witnesses & Root node\\\midrule\endhead
\nolinkurl{member} & 5 & 7 & 149\\
\nolinkurl{raw_add} & 15 & 6 & 308\\
\nolinkurl{raw_z} & 7 & 1 & 68\\
\nolinkurl{raw_inverse_x} & 7 & 0 & 54\\
\nolinkurl{log25} & 4 & 0 & 229\\
\nolinkurl{pole} & 1 & 11 & 106\\
\nolinkurl{test_group} & 44 & 162 & 515\\
\nolinkurl{slope} & 22 & 1,298 & 2681\\
\nolinkurl{height_certificate} & 36 & 35,106 & 260\\
\nolinkurl{multiply.scalar} & 10 & 1 & 136\\
\nolinkurl{multiply.guard} & 15 & 1 & 22\\
\nolinkurl{multiply} & 15 & 35,134 & 83\\
\nolinkurl{natural} & 5 & 140,604 & 79\\
\nolinkurl{value.signed_encode} & 10 & 140,654 & 138\\
\nolinkurl{value.abs} & 2 & 4 & 23\\
\nolinkurl{value.close_coordinates} & 16 & 53 & 271\\
\nolinkurl{value.test} & 11 & 35,282 & 775\\
\nolinkurl{value.exceptions} & 6 & 406 & 487\\
\nolinkurl{scalar.integer_constant_0} & 10 & 1 & 136\\
\nolinkurl{scalar.integer_constant_1} & 10 & 1 & 136\\
\nolinkurl{scalar.integer_constant_2} & 10 & 1 & 136\\
\nolinkurl{scalar.integer_constant_3} & 10 & 1 & 136\\
\nolinkurl{scalar.integer_constant_4} & 10 & 1 & 136\\
\bottomrule\end{longtable}

\begin{longtable}{lrrr}
\toprule Recipe & Parameters & Witnesses & Total variables\\\midrule\endhead
\nolinkurl{power_value} & 20 & 17,786,698 & 17,786,718\\
\nolinkurl{precision} & 25 & 292,889,492,836 & 292,889,492,861\\
\nolinkurl{close} & 11 & 292,890,196,199 & 292,890,196,210\\
\nolinkurl{main} & 6 & 292,890,231,486 & 292,890,231,492\\
\nolinkurl{value_union} & 6 & 292,890,231,892 & 292,890,231,898\\
\nolinkurl{value} & 6 & 292,890,231,899 & 292,890,231,905\\
\nolinkurl{integers} & 1 & 292,890,231,904 & 292,890,231,905\\
\bottomrule\end{longtable}


The five scalar-constant templates use the values in
\texttt{library.fixedConstants}. Large program literals are retained
byte-for-byte in the supplementary specification. Replacing them by
unspecified sufficiently large constants would change this definition.

\section{Primary theorem applications}
\begin{longtable}{p{.26\textwidth}p{.65\textwidth}}
\toprule Input & Exact application boundary\\\midrule\endhead
Selmer parity \cite{dd} & Rational two-isogeny, $p=2$, full Selmer corank;
no finite-$\Sha$ hypothesis.\\
Kato \cite{kato} & Fixed weight-two form and symbol, integral transfer and
fixed-valuation smoothing; rational divisibility is Theorem 13.4(2),
with oddness, rational irreducibility and the unipotent condition.\\
BFH and Hoffstein--Luo \cite{bfh,hl} & Actual signed-prime twist with
trivial character and arbitrary finite splitting set. Only the
nonvanishing clause has the four-prime bound before the extra prime.\\
Waldspurger and Ramsey \cite{waldspurger,ramsey} & Full good-Hecke
eigenspace, same dyadic squareclass, and algebraic vectors by flat
field extension. No arbitrary-ring base change in weight $3/2$.\\
CM norms and heights \cite{ccl,cst,cornutjetchev} & Split level,
conductors coprime to level and the new prime, no extra units,
unaveraged sums and one fixed parametrization.\\
Duality \cite{milne,nekovar} & Actual derived Kummer conditions and their
homotopies. Support change only after the coefficient representation
is unramified at the removed place.\\
Forward finiteness \cite{gz,kolyvagin} & Only at known analytic
nonvanishing/simple zeros or a known nontorsion Heegner trace.\\
Isogeny ratio \cite[proof of I.7.3]{milne} & Finite $\Sha$, dual
sum/difference maps, canonical heights and pullback measures;
no BSD leading-term assumption.\\
Height \cite{pan2022,javanpeykar,gr2014,fowler,cmheight,pazuki} &
Odd irreducible de Rham representation at five; explicit conductor
and modular-Jacobian transfer. CM order class number bounded by
the degree of the original Hom field.\\
Primes \cite{greentao,moebius,gtz,bmor} & Fixed finite-complexity system
after fixing sieve cutoff, and the unconditional explicit progression
estimate.\\\bottomrule
\end{longtable}
The source package includes the specification, natural precision polynomial
and construction sources in \nolinkurl{SUPPLEMENT_MANIFEST.json}.
The larger proof corpus is identified by hashes and stable locators in
\nolinkurl{SOURCE_MAP.json}. A source identity is not a proof verdict.

% End 06-literal-construction.tex


% Begin 07-even-source.tex
\section{Fixed-form arithmetic and the even coefficient source}\label{app:even-source}
This appendix supplies the constructions used in (3.1)--(3.3).
Throughout, $E=E_\ell$ is fixed, with the lattice, parametrization and
period convention of Section~\ref{sec:rc}. Constants may depend on these
data and on the prescribed finite list of local types. They do not depend
on the number or sizes of the later good conductor primes.

\subsection{One retained class, including discriminant \texorpdfstring{$4h$}{4h}}
All numbered primary references in this subsection are to Kato~\cite{kato}.

Fix a non-CM elliptic curve $E/\mathbf Q$ of conductor $N$, its weight-two newform $f$,
$T=T_2E$, and $V=T\otimes\mathbf Q_2$. In the application
$E=E_\ell$ is fixed first. Put
$\Lambda=\mathbf Z_2[[\Gamma]]=\mathbf Z_2[[t]]$, where $\Gamma$
is the real cyclotomic Galois group. Let
$\kappa:G_{\mathbf Q}\to\Gamma\subset\Lambda^\times$ be the quotient map,
and set $A=\Lambda[1/2]$.
We use arithmetic Frobenius throughout. The action on the cyclotomic
group-ring coefficient is the inverse deck action.

\paragraph{The precise output.}
\begin{proposition}[One retained Kato system and its two comparisons]
\label{ex:app-kato-system}
Use the fixed curve, lattice and cyclotomic convention just specified.
Fix $e\ge\max\{3,v_2(N)\}$ large enough for the fixed initial local data;
the smoothing choice is (K5).
Let $G$ be a finite elementary abelian two-group, let
$\psi:G_{\mathbf Q}\to G$ be a continuous homomorphism, and let
$\chi_x=\chi_0(x\circ\psi)$, for $x\in\widehat G$, be a finite
family of quadratic characters with one fixed real sign and fixed
local squareclasses at the primes dividing $2N$.
Since the characters of $G$ separate
its elements, $\psi(c)=1$ for complex conjugation $c$, and
$\psi|_{G_{\mathbf Q_v}}=1$ at every $v\mid2N$. All varying conductor primes
are good odd primes for $E$. Let $F$ be the squarefree least common
multiple of the odd parts of the primitive conductors. There is a
single retained-smoothed class
\[
 Z_F\in 2^{-c_E}H^1\bigl(G_{\mathbf Q,S},
        T\otimes\Lambda[G](\chi_0\psi\kappa^{-1})\bigr),
 \qquad S=\operatorname{prime}(2NF),                         \tag{K1}
\]
where $c_E$ is fixed before $F$ and $G$.
There is also one $F_E\in\Lambda$ with $F_E(0)\ne0$, fixed before
$F,G$, such that at every height-one prime $\mathfrak p$ of $A$,
\[
 \ell_{A_{\mathfrak p}}H^2(C_{S,x})_{\mathfrak p}
 \leq
 \ell_{A_{\mathfrak p}}
 \left(H^1(C_{S,x})/A F_E(Z_F)_x\right)_{\mathfrak p}.       \tag{K2}
\]
Here $C_{S,x}$ is the full rational global Iwasawa complex of the
constituent. For an odd squarefree signed parameter $h$ giving
$\chi_x=\chi_h$, central reciprocity for this same class gives
\[
 v_2\exp^*\bigl(F_E(Z_F)_x(0)\bigr)
 =v_2L_*(h)+\sum_{q\mid F,\ q\nmid h}
        v_2\#E^h(\mathbf F_q)+2e+v_2F_E(0)+\epsilon_h,       \tag{K3}
\]
when the value is nonzero. The numbers $\epsilon_h$ belong to a
bounded set depending only on the fixed curve, period conventions
and local types at two. A zero central $L$-value gives zero dual
exponential. Formula (K3) does not assert that its remaining unit is
one. No part of (K1)--(K3) compares integral lattices of varying newforms.
\end{proposition}
\begin{proof}\leavevmode
\paragraph{The retained route and its comparison.}
All classes in this subsection retain the smoothing integers of (K5).
The comparison used below is the exact identity (E.1), obtained by
integral weighted corestriction and the norm relation of
\cite[Proposition 8.12]{kato}. We do not divide by the smoothing
operators, or compare a chosen generator with a separately normalized
zeta class for the varying newform. In particular the central factor
is $s_F$ of (K6), whose valuation is $2e$ for every $F$.
The use of power-series zero finiteness in paragraph~7 establishes
nonvanishing of the supplied Euler system in both generic sign
components; it is not a comparison of two Euler systems.

\paragraph{1. One original lattice and one symbol combination.}
Kato (8.1.3), p.~180, constructs retained-smoothed classes
\[
 {}_{c,d_s}z_m(f,1,1,\xi,S_m)
 \in H^1\bigl(\mathbf Z[\zeta_m,1/2],L_f(1)\bigr),
 \qquad S_m=\operatorname{prime}(2mN),                      \tag{K4}
\]
where $L_f=V_{\mathcal O_\lambda}(f)$ is the image lattice of the
fixed form. For $\xi\in\mathrm{SL}_2(\mathbf Z)$, section~8.9,
p.~185, explicitly constructs the trace homomorphism into the
integral induced lattice
\[
 \mathbf Z_2[\operatorname{Gal}(\mathbf Q(\zeta_m)/\mathbf Q)]
       \otimes V_{2,\mathbf Z_2}(Y_1(N)).
\]
Section~8.11 projects this lattice onto $L_f$. Auxiliary modular
curve levels used in the trace do not change its target lattice.
Choose a rational identification $L_f(1)\otimes\mathbf Q_2=V$.
One fixed $2^{c_E}$ sends its lattice into $T$.

By Kato~13.6--13.7, p.~227, the symbol span has a nonzero projection
to each Betti sign line. The union of the two rational kernels does
not contain the whole span. Choose an integral finite linear
combination $\xi$ with both projections nonzero, clearing its fixed
denominators once. Norm relations and reciprocity are linear, so
the same combination works at every $m$. Although (K4) is written
with one symbol, it henceforth denotes this fixed combination.

\paragraph{2. Retain one smoothing choice for the whole family.}
Fix $e\geq\max\{3,v_2(N)\}$, enlarged only for fixed initial data,
and take
\[
                 c=d_s=1+2^e\,3N_{\rm odd}F.              \tag{K5}
\]
These integers exceed one, are prime to $6NF$, are one modulo $N$
and $F$, and are one modulo every dyadic quadratic conductor under
consideration. Distinctness or primality of $c,d_s$ is not a
hypothesis of Kato's construction. The exponents $(u,v)$ in (4.2.4),
p.~143, are $(1,1)$ when $k=2,r=r'=1$. The retained-smoothed
version of Theorem~6.6, p.~163, therefore has central scalar
\[
 s_F=(c^2-c)(d_s^2-d_s)=c(c-1)d_s(d_s-1),
                   \qquad v_2(s_F)=2e.                   \tag{K6}
\]
No smoothing factor is divided out. The prime factors of $cd_s$
exclude some Euler-system indices in Kato~13.1; they are not added
to the inverted arithmetic support $S$, as is explicit in (8.1.3)
and in the definition of $H^q$ in section~8.2. In particular their
number creates no local $H^2$ correction.

\paragraph{3. The integral $j_*$ condition survives ramified transfer.}
For a number field $K$, let $X_K=\operatorname{Spec}\mathcal O_K[1/2]$
and $j_K:\operatorname{Spec}K\to X_K$. At finite coefficients
$M=T/2^aT$, the low-degree Leray sequence gives an injection
\[
 H^1(X_K,j_{K,*}M)\longrightarrow H^1(K,M)
\]
whose image is the kernel of all inertia restrictions outside two.
Indeed the higher-direct-image stalk formula and Leray sequence
\cite[Tags 03Q9 and 03QA]{stackscohom} identify the stalks of
$R^1j_{K,*}M$ at a finite place with $H^1(I_v,M)$;
the edge map is restriction to these stalks. The generic-point map
$j_K$ is affine, hence quasi-compact and quasi-separated, as required
by the stalk formula. Kato's section~8.2,
pp.~180--181, defines the integral groups by the inverse limit of
these finite-coefficient groups, so this assertion applies to the
classes (K4) before rationalization.

Let $L/K$ be finite, with absolute Galois groups $H\subset G_K$.
For a class $z\in H^1(H,M)$ and an inertia subgroup $I\subset G_K$,
the Mackey formula is
\[
 \operatorname{res}_{I}^{G_K}\operatorname{cor}_{H}^{G_K}z
 =\sum_{g\in I\backslash G_K/H}
   \operatorname{cor}_{I\cap gHg^{-1}}^{I}
   \operatorname{res}_{I\cap gHg^{-1}}^{gHg^{-1}}(g z).    \tag{K7}
\]
The intersection in each summand is the upper inertia subgroup
for the corresponding place of $L$. If all upper inertia
restrictions vanish, every summand is zero. This does not require
$L/K$ to be unramified. It is a finite double-coset identity at each
finite coefficient precision. Uniqueness under the Leray injection
makes the extensions across $X_K$ compatible in $a$. Thus the
integral $j_*$ condition survives corestriction, including at primes
ramified in a quadratic twisting field.

\paragraph{4. Primitive twisting is an actual Euler system.}
Fix a constituent $\chi=\chi_h$ of primitive conductor
$D=|\Delta_h|=2^a d$, $d$ odd, and a basis $e_\chi$ of
$\mathbf Z_2(\chi)$. Put
$\Sigma=\operatorname{prime}(2Ndcd_s)$.
For $m$ with $\operatorname{prime}(m)\cap\Sigma=\{2\}$ set
\[
 K_m=\mathbf Q(\zeta_m),\quad
 L_m=\mathbf Q(\zeta_{Dm}),\quad
 y_m=\operatorname{cor}_{L_m/K_m}
             (z_{Dm}\otimes e_\chi).                      \tag{K8}
\]
The character is trivial over $L_m$, so tensoring is an integral
map there. Formula (K7) places (K8) in the precise $j_*$ group
required by Kato~13.1. Each class uses the original $f,\xi,c,d_s$.

For $m\mid m'$, transitivity of corestriction reduces the norm
of (K8) to Proposition~8.12, p.~186. For a deck automorphism $g$,
equivariance gives
\[
 \operatorname{cor}((gz)\otimes e_\chi)
       =\chi(g)^{-1}g\operatorname{cor}(z\otimes e_\chi).
\]
Consequently the factor at a newly introduced good prime $q$ is
\[
 E_q(\chi)=1-\frac{a_q\chi(q)}q\sigma_q^{-1}
                         +\frac1q\sigma_q^{-2}.           \tag{K9}
\]
In Kato~13.1, pp.~224--225, the polynomial is
$P_q(X)=\det(1-\operatorname{Fr}_qX\mid T\chi)$ and the norm
factor is $P_q(q^{-1}\sigma_q^{-1})$, exactly (K9).
Repeated prime powers introduce no factor because Proposition~8.12
has a product over newly added support primes. Thus $(y_m)_m$ is
an Euler system of the stated lattice, not a class inferred from
proportionality of rank-one Iwasawa modules.

For $m=2^n$ with $n\geq a$, the field $L_m$ can be replaced by
$\mathbf Q(\zeta_{d2^n})$: apply repeated-two norm compatibility
from exponent $n+a$ to $n$. This introduces no factor $2^a$.
The construction therefore covers discriminant $4h$ on the same
cofinal tail. At the coefficient level it retains the full character
$\chi_2\chi_{\rm odd}$. A dyadic character nontrivial on $\Gamma$
changes its action by the integral automorphism
$1+t\mapsto-(1+t)$, or $t\mapsto-2-t$. A dyadic sign character
instead changes the finite real-sign component. Neither operation
is a division by a sign projector.

\paragraph{5. One coefficient map produces every constituent.}
For $n$ large enough to contain the dyadic characters, put
$A_{F,n}=\operatorname{Gal}(\mathbf Q(\zeta_{F2^n})/\mathbf Q)$.
Apply Shapiro to $z_{F2^n}$, then the map on group bases
\[
 \Phi_n:\mathbf Z_2[A_{F,n}]\longrightarrow
                  \mathbf Z_2[\Gamma_n\times G],\qquad
 [s]\longmapsto\chi_0(s)[\bar s][\psi(s)].                 \tag{K10}
\]
With inverse deck actions, both sides transform under $g\in G_{\mathbf Q}$
by the inverse of its group image. On the target the extra scalar
is $\chi_0(g)$; since $\chi_0^{-1}=\chi_0$, (K10) is equivariant.
This is an integral coefficient homomorphism. Shapiro restriction
to an inertia subgroup is the direct sum of the upper restrictions;
the same zero restrictions used in (K7) therefore remain zero after
(K10). In particular the map acts on actual integral classes.
The map in (K10) is defined on this finite group for the following
reason. The constituent indexed by the trivial character of $G$ is
$\chi_0$, and
$x\circ\psi=\chi_x\chi_0$ for every $x\in\widehat G$.
All these quadratic Dirichlet characters have conductor dividing
$F2^n$ once $n$ is large enough. If an element of $G_{\mathbf Q}$
acts trivially on $\mathbf Q(\zeta_{F2^n})$, every character of $G$
is therefore one on its image under $\psi$. Character separation
makes that image the identity. Thus both $\chi_0$ and $\psi$
factor through $A_{F,n}$. Explicitly, for $g\in G_{\mathbf Q}$,
\[
 \Phi_n([g^{-1}s])
 =\chi_0(g)\,[\bar g]^{-1}[\psi(g)]^{-1}\Phi_n([s]).
\]
This verifies equivariance on each integral group basis before
evaluation at any character. The quotient to the real tower sends
each basis element to its image; its matrix has entries zero and
one. The transition maps and character evaluations are also maps
of integral coefficient modules. None of these operations is the
idempotent $|G|^{-1}\sum_g x(g)[g]$.


The maps $\Phi_n$ commute with the norm transition maps.
Proposition~8.12 has no new prime when $n$ increases, so their
images define the single class $Z_F$ in (K1).
The finite-coefficient degree-zero groups are finite and satisfy
Mittag--Leffler. The inverse-limit exact sequence
\cite[Corollary 4.1.3, p.~96]{nekovar} therefore has zero
$\varprojlim^1 H^0$ term, so it gives this degree-one class
in completed cochains.
The quotient from the full cyclotomic tower to $\Gamma$ is the
integral quotient of group bases, equivalently corestriction to
the real subfield. No factor $(1+\tau)/2$ occurs.

After evaluation at $x$, (K10) is the coefficient form of the
unnormalized corestriction weighted by $\chi_x$. Factor that
corestriction through the primitive conductor field in (K8).
Transitivity and Proposition~8.12 give the exact identity
\[
              (Z_F)_x=
       \prod_{q\mid F/d}E_q(\chi_x)\,y_{2^\infty}.        \tag{E.1}
\]
It is an identity of the given classes. The same smoothing choice
appears on both sides. There is no group-order denominator and no
comparison with a canonical zeta class of $f\otimes\chi_x$.

For clarity, at a constituent write $\chi_h=\chi_2\chi_{\rm odd}$ and
$|\Delta_h|=2^a d$. Before taking the cofinal limit, the corresponding
scalar coefficient map is exactly
\[
 [s]\longmapsto\chi_2(s)\chi_{\rm odd}(s)[s|_{\Gamma_n}],
                    \qquad n\ge\max\{a,2\}.\tag{E.2}
\]
Thus the dyadic character is part of the module action, including
when the field discriminant is $4h$.

\paragraph{6. Central trace and period normalization.}
First use (E.1) to remove the unused tame primes. The central
specialization of (K8) is the corestriction of $z_D\otimes e_\chi$
from $\mathbf Q(\zeta_D)$ to $\mathbf Q$, with omitted support
$\operatorname{prime}(2Nd)$. For $D=1$ use the trivial character
and the same omitted support. This follows from the same repeated-two
norm identity even when $D=4|h|$. One must do this before using a
Gauss basis; the Gauss sum of an imprimitive character modulo a
larger cyclotomic level can be zero.

Theorem~9.7, p.~189, identifies the dual exponential of $z_D$ with
the retained-smoothed de Rham zeta element. Dual exponential commutes
with corestriction, whose de Rham map is field trace. Thus the
coefficient of $e_\chi$ is the exact unnormalized sum
\[
                    \sum_{g\in\operatorname{Gal}(\mathbf Q(\zeta_D)/\mathbf Q)}
                       \chi(g)\,g(\exp^*z_D).             \tag{K12}
\]
Theorem~6.6, p.~163, evaluates precisely this sum. In weight two,
$r=r'=1$, its Betti sign is $\epsilon=\chi(-1)$ and its power of
$2\pi i$ is zero. If $e_B^\epsilon$ is a fixed rational generator
of that Betti line and
$\delta(f,1,\xi)^\epsilon=b_\epsilon e_B^\epsilon$, then
$b_\epsilon\in\mathbf Q^\times$ is fixed and
\[
     \operatorname{per}_f\bigl((\text{K12})\bigr)^\epsilon
          =s_F L_{\operatorname{prime}(2Nd)}(f,\chi,1)\,
                        b_\epsilon e_B^\epsilon.          \tag{K13}
\]
Choose the original de Rham generator so that its Betti period
on this line is $\Omega_E^\epsilon e_B^\epsilon$, absorbing a
fixed rational comparison in $b_\epsilon$. Formula (K13) says
that the scalar multiplying $e_\chi$ in (K12) is
$s_F b_\epsilon L_{\operatorname{prime}(2Nd)}(f,\chi,1)/
\Omega_E^\epsilon$. Its Galois character is $\chi$.

To descend it, take the invariant basis $\sqrt h\,e_\chi$ of the
quadratic character line. With the usual complex convention
$\Omega_E^-=i|\Omega_E^-|$ and $\sqrt h=i\sqrt{|h|}$ for $h<0$,
division by $\sqrt h$ gives, up to the already fixed rational
period choice, exactly
\[
             s_F b_\epsilon\,\frac{L_*(h)}{h}
       \prod_{p\mid 2N}\text{(the omitted local factor at $p$)}.
                                                               \tag{K14}
\]
The signed integer $h$ is odd, so $h^{-1}$ has valuation zero.
This computation also explains the usual Gauss-basis description:
for the primitive character, $\tau(\chi)^2=\Delta_h$ and
$\tau(\chi)e_\chi$ is invariant; replacing it by $\sqrt h e_\chi$
costs the fixed dyadic factor $\sqrt{\Delta_h/h}$. It does not cost
$\varphi(D)$ or any power of the number of prime divisors of $h$.
The discriminant-$4h$ case has ratio two in this last basis change,
and the direct basis calculation (K14) already includes it.
One can verify the normalization without choosing an idempotent.
Write $\beta=\exp^*z_D$ in the original de Rham line and put
$\beta_\chi=\sum_g\chi(g)g\beta$. Then
$u\beta_\chi=\chi(u)\beta_\chi$ for every $u$ in the finite
cyclotomic group. Thus $\beta_\chi\otimes e_\chi$ is invariant,
and its coordinate in the invariant basis
$\sqrt h\,e_\chi$ is $\beta_\chi/\sqrt h$.
Theorem~6.6 uses precisely $\beta_\chi$, rather than
$\beta_\chi/\varphi(D)$. Theorem~9.7 applies here since
$1\le r=r'=1\le k-1=1$; it has no restriction that $p$ be odd.
The dual form $f^*$ in Theorem~6.6 equals $f$, because the elliptic
newform has rational coefficients.

In the real-period convention,
$L_*(h)=\sqrt{|h|}L(E^h,1)/|\Omega_E^\epsilon|$, up to the
fixed rational period comparison already included in $b_\epsilon$.
For $h>0$,
$\Omega_E^+\sqrt h=|\Omega_E^+|\sqrt{|h|}$; for $h<0$,
$\Omega_E^-\sqrt h=-|\Omega_E^-|\sqrt{|h|}$.
Consequently
\[
 \frac{L(E^h,1)}{\Omega_E^\epsilon\sqrt h}
     =\frac{L_*(h)}{h}
\]
with the same fixed period comparison. This proves the signed
factor in (K14), not only its valuation. For a primitive quadratic
character, $\tau(\chi)^2=\Delta_h$ and
$\tau(\chi)/\sqrt h$ equals a fixed choice of
$\sqrt{\Delta_h/h}$. Changing from its Gauss basis to the displayed
invariant basis multiplies the coordinate by this ratio, which
is one or two up to sign. In particular the discriminant-$4h$
case adds no factor depending on the number of odd conductor primes.


Identifying this generator with the chosen local generator for
$E^h$ at two has bounded valuation: odd squarefree $h$ belongs to
one of finitely many classes in $\mathbf Q_2^\times/\mathbf Q_2^{\times2}$,
and changing a representative by a unit square scales the quadratic
twist differential by a two-adic unit. The local factors at $2N$
also form a fixed finite nonzero list. At a good active prime the
local $L$-factor is one; at an unused good prime (K9) at $t=0$ is
$\#E^h(\mathbf F_q)/q$. This proves (K3), including all its
uniformity claims and the assertion about zero values.

\paragraph{7. Every hypothesis of Kato~13.4(2).}
Apply the theorem to the integral system (K8), after multiplying
by the fixed lattice comparison if needed. Its hypotheses are:
\begin{itemize}
\item Each generic sign component is nonzero. For a finite-order
cyclotomic character $\rho$, corestriction and the $\rho$-sum combine
into the single $(\chi\rho)$-weighted sum on the original class.
Theorems~6.6 and~9.7 apply with the same $\xi,c,d_s$. Their smoothing
factors are nonzero, since $c,d_s>1$ cannot be roots of unity.
Theorem~13.5(2), p.~227, gives only finitely many zero $L$-values
among characters supported on the fixed finite set
$\operatorname{prime}(2Nd)$. There are infinitely many $\rho$ of
each parity. A torsion Iwasawa class can have a nonzero specialization
at only finitely many such characters, since a nonzero annihilating
power series has only finitely many cyclotomic zeros. Hence both
generic sign components are nonzero.
\item The two real eigenspaces of $T\chi$ have rank one.
\item The theorem uses inverse Frobenius in its purity condition.
On $T\chi$ its good eigenvalues have complex absolute value
$q^{-1/2}$, so its integer $w$ is $-1$.
\item $V\chi$ is irreducible; twisting preserves irreducibility of
the non-CM rational Tate representation.
\item The image over $\mathbf Q(\mu_{2^\infty})$ contains an open
subgroup of $\mathrm{SL}_2(\mathbf Z_2)$. Intersecting with
$\ker\chi$ keeps an open subgroup, which contains a nonidentity
unipotent. It fixes a one-dimensional subspace and gives condition~(v).
\end{itemize}
These are all rational hypotheses. The residual irreducibility,
integral cokernel and $p\ne2$ conditions of part~(3) are not used.
Theorem~13.4(1) gives torsion $\mathbf H^2(T\chi)$, and part~(2)
gives, at every height-one prime away from two,
\[
 \ell\mathbf H^2(T\chi)_{0,\mathfrak p}
       \leq\ell(A_{\mathfrak p}/J_{\mathfrak p}),\qquad
 \mathbf H^2_0=\ker(\mathbf H^2\to\mathbf H^2_{\rm loc}). \tag{K15}
\]

The finiteness and Euler characteristic (12.2.1)--(12.2.2), p.~220,
give generic rank one in $\mathbf H^1$. Its localization away
from two is torsion-free. Indeed a torsion element for a scalar
prime would come from invariants of the corresponding finite-dimensional
rational coefficient algebra at that prime; its coefficient action is
abelian, whereas the open commutator subgroup on $V\chi$ has no
fixed vectors. The coefficient exact sequence excludes these
invariants. The scalar specialization need not be a finite-order
character. Thus $\mathbf H^1_{\mathfrak p}$ is free of rank one.
Localization commutes with $\operatorname{Hom}$ for this finitely
presented module. Writing $y_{2^\infty}=a e$ in a local basis gives
$J_{\mathfrak p}=(a)$. Formula (K15) is exactly the index bound for
the supplied class, not for an arbitrarily rescaled generator.

\paragraph{8. Paying for the full global complex.}
At an odd prime put $W_q=H^1(I_q,V\chi)$. The rational cone from
$j_*$ to full cochains is
\[
 [\,W_q\otimes A\xrightarrow{1-\sigma_q^{-1}\operatorname{Fr}_{W_q}}
          W_q\otimes A\,],\qquad\deg=(1,2).              \tag{K16}
\]
Its determinant $D_q$ is nonzero, since the cyclotomic Frobenius
exponent is nonzero. The map is injective even after localization
at a prime dividing $D_q$; it is not being reduced modulo that
prime. Thus the cone contributes only degree-two length
$v_{\mathfrak p}(D_q)$. The long exact sequence shows that full
$H^2$ increases by at most their sum and that $H^1$ is unchanged.
This inequality does not require vanishing of the next $j_*$ group.

At a good inactive prime, $W_q=V\chi(-1)$, so $D_q=E_q(\chi)$
with exactly the convention of (K9). At a good active prime,
ramified quadratic inertia acts by $-1$; its rational cohomology
vanishes and $D_q=1$. This includes every active base prime,
regardless of how many there are. At an old bad prime retain
$H^1(I_q,V\chi)$ itself; replacing it by $V^{I_q}\chi(-1)$
would in general be false.

Choose a fixed integral $F_2$ divisible over $A$ by the
characteristic elements of $\mathbf H^2_{\rm loc}(T\chi)$
for the finitely many allowed local types at two and finite
sign components. Its central value can be chosen nonzero:
derived local control, and the absence of local Iwasawa $H^3$,
identify its rational central $H^2$ with $H^2(\mathbf Q_2,V\chi)=0$.
The last equality follows from local duality and finiteness of
local elliptic torsion. Hence the characteristic elements have
no factor $t$. Kato (12.2.3), p.~220, gives the required finiteness
and torsion of this local Iwasawa module.

At the finitely many old bad primes, choose stable lattices in
the $W_q$ and multiply all the determinants $D_q$ for the finitely
many local twisting types. Clear their fixed scalar denominators.
Their central values are nonzero, since the cokernel of
$1-\operatorname{Fr}_{W_q}$ is $H^2(\mathbf Q_q,V\chi)=0$
by the same local duality. Call this product $F_{\rm bad}$ and set
$F_E=F_2F_{\rm bad}$.

The quotient $\mathbf H^2/\mathbf H^2_0$ injects into the local
$\mathbf H^2_{\rm loc}$, so $F_2$ pays for it in (K15).
Multiplying a class in a free rank-one module by a nonzero scalar
adds that scalar's valuation to its index. Formula (E.1) already
pays for every unused good determinant in (K16); $F_{\rm bad}$
pays for the remaining old ones. Adding these length inequalities
proves (K2).
Here is the length calculation with all contributions separated.
Fix a height-one prime $\mathfrak p$ of $A$ and write
$R=A_{\mathfrak p}$. Let $M$ be the localized degree-one module
for the primitive $j_*$ system, let $y=y_{2^\infty}$, and write
$\operatorname{ind}_R(y)=\ell_R(M/Ry)$.
The preceding freeness argument gives $M\simeq R$.
Put $Q_\chi=\prod_{q\mid F/d}E_q(\chi)$.
All primes in $F/d$ are good: the family has identical local
characters at primes dividing $N$, so a bad prime cannot ramify
in one constituent and be unramified in another.
The cone (K16) has zero degree-one cohomology, hence identifies
$M$ with the localized full degree-one module. Its long exact
sequence and (K15) give
\[
\begin{split}
 \ell_R H^2(C_{S,x})_{\mathfrak p}
 &\le \ell_R\mathbf H^2(T\chi)_{0,\mathfrak p}
       +\ell_R\mathbf H^2_{\rm loc}(T\chi)_{\mathfrak p}
       +\sum_{q\in S,\ q\ne2}v_{\mathfrak p}(D_q)\\
 &\le \operatorname{ind}_R(y)
       +v_{\mathfrak p}(F_2)+v_{\mathfrak p}(F_{\rm bad})
       +v_{\mathfrak p}(Q_\chi)\\
 &=\operatorname{ind}_R(F_EQ_\chi y)
  =\operatorname{ind}_R\bigl(F_E(Z_F)_x\bigr).
\end{split}
\]
The second line uses $D_q=1$ at active good primes, and
$D_q=E_q(\chi)$ at inactive good primes. It does not charge an
extra fixed constant for each active prime.

For completeness, nonvanishing at the center of the fixed
correction factors follows before multiplying over local types.
Local derived control and $H^3_{\rm Iw}=0$ identify
$\mathbf H^2_{\rm loc}/t\mathbf H^2_{\rm loc}$ rationally with
$H^2(\mathbf Q_2,V\chi)$. Local duality identifies the latter
with the dual of $H^0(\mathbf Q_2,V\chi)$, which is zero because
local elliptic torsion is finite. Nakayama at $(t)$ therefore
excludes a factor $t$ from its characteristic element.
At an old odd prime the cokernel of
$1-\operatorname{Fr}_{W_q}$ is $H^2(\mathbf Q_q,V\chi)=0$,
by the same duality. Its determinant at $t=0$ is nonzero.
Taking products over the finite list of local types, followed by
one fixed power of two to clear scalar denominators, gives
$F_E\in\Lambda$ with $F_E(0)\ne0$. This operation is independent
of the binary dimension and all later good conductor primes.


\paragraph{9. Uniformity and the boundary of this proposition.}
The only lattice denominator is fixed before the family. The
smoothing integers may grow, but their central valuation is $2e$.
The multiplier $F_E$ depends only on the original local types.
Transfer, coefficient push, real quotient and norm compatibility
are integral. The proof never divides by a conductor degree,
$|G|$, a class number, or a Fourier projector. It keeps one class
through (K1), (K2), and (K3).
\end{proof}

We now use this same class in the positive complex. The preceding
argument supplies its integral construction, its rational height-one
divisibility and its exact central reciprocity normalization.

\begin{proposition}[One integral even scalar for the family]\label{ex:app-even-scalar}
Retain Proposition~\ref{ex:app-kato-system} and assume, as for the
present curve, that $E[2](\Q)\simeq(\Z/2\Z)^2$. Let $C^+$ be the
fiber of the ordinary global cochains to the ordinary real cochains,
with real degree zero retained. It has a compatible finite free model
under the coefficient changes used here. For
$B=\mathscr O[G]$, it satisfies $C^+\otimes B\simeq B^2[-1]$.
The fixed real invariant generator, the positive lift of the retained
class, and the original determinant basis define the single scalar
$u^{\rm ev}$ by (E.4), with
\[
 u^{\rm ev}\in\Lambda[1/2][G],\qquad
 2^{c_E}u^{\rm ev}\in\Lambda[G].
\]
Here $c_E$ may be enlarged for the fixed data but is independent of
$G$ and the later good conductor primes. The assertion requires the
common real sign and local squareclasses at $2N$ in
Proposition~\ref{ex:app-kato-system}.
\end{proposition}
\begin{proof}\leavevmode
Let $C^+$ be the fiber of ordinary global cochains to ordinary real
cochains. Real $H^0$ is included. Over a finite Artin quotient,
continuous cochains are filtered unions of finite free function
modules, hence flat and therefore free. A residual splitting into
finite cohomology and disks lifts by the nilpotent perturbation
series to a finite free model with actual homotopies. Minimal models
at successive quotients are related by invertible matrices:
a homotopy equivalence reduces to an isomorphism of minimal residual
complexes, and Nakayama applies. Lifting these matrices and comparison
homotopies makes the models compatible. Their acyclic, termwise
surjective comparison cones have acyclic inverse limit.
Here is the compatibility argument. At successive finite coefficient
quotients, reduction of a minimal finite free model is a minimal model
of the reduced cochains. Two such models are homotopy equivalent;
reducing a homotopy inverse modulo the maximal ideal shows that the
comparison is degreewise invertible. Lift these invertible matrices
to the next quotient and change bases. The model differentials then
reduce to the preceding differentials exactly. The comparison maps
to the original cochains agree after reduction up to chain homotopy.
Since reduction of continuous cochains is termwise surjective, lift
that homotopy and adjust the upper comparison by its boundary. The
comparison maps are now exactly compatible and remain
quasi-isomorphisms.

Let $\mathcal K_n$ be their acyclic comparison cones. Their transition
maps are termwise surjective. If $s$ denotes the transition shift on
the product, successive lifting gives the exact sequence of complexes
\[
 0\longrightarrow\varprojlim_n\mathcal K_n
 \longrightarrow\prod_n\mathcal K_n
 \xrightarrow{\,1-s\,}\prod_n\mathcal K_n
 \longrightarrow0.
\]
Products of complexes of modules preserve acyclicity, so the long
exact cohomology sequence makes $\varprojlim_n\mathcal K_n$ acyclic.
The inverse-limit comparison is therefore a quasi-isomorphism from
the bounded finite free limit model to the continuous coefficient
complex. This also identifies its derived coefficient changes with
the finite-quotient constructions.
This proves the needed perfection and coefficient change.

% Replace baseline lines 3310--3335; retains equation tag (E.3).
\paragraph{The residual positive complex.}
Write $K_n=\mathbf Q(\zeta_{2^{n+2}})^+$, and let $S_n$ consist of
the places above two and a fixed finite set of odd rational primes.
Put $r_n=2^n$ and $s_n=\#S_n$. We first compute with the scalar
coefficient $\mathbf F_2=\mu_2$.

The ordinary two-class group of $K_n$ is trivial. Indeed this holds
for $K_0=\mathbf Q$. Assuming it for $K_{n-1}$, let $H$ be the
Hilbert two-class field of $K_n$, with the real places split.
The canonical extension $H/K_{n-1}$ is Galois with a finite
two-group $G_H$. Only the unique dyadic prime can ramify, and its
inertia $I$ has order two. The quotient of $G_H^{\rm ab}$ by the image
of $I$ gives an everywhere unramified abelian two-extension of
$K_{n-1}$ with real places split, hence is trivial by induction.
Thus $G_H^{\rm ab}$ is cyclic, generated by an element of order at most
two. The Frattini quotient of $G_H$ has one generator, so $G_H$ is cyclic and
has order at most two. Since it contains $K_n$, we obtain $H=K_n$.
The $S_n$-class group is a quotient, so its two-part also vanishes.

The signatures of units span all $r_n$ real signs. For $n\ge1$ put
$m=n+2$, $\zeta=\zeta_{2^m}$, and
\[
 u=1+\zeta+\zeta^{-1}
   =\zeta^{-1}\frac{\zeta^3-1}{\zeta-1}.
\]
This is a unit: the numerator and denominator generate the same
dyadic prime ideal in the cyclotomic field. Its real conjugates are
$1+2\cos(2\pi a/2^m)$ for odd $0<a<2^{m-1}$.
The negative ones have $2^m/3<a<2^{m-1}$, and their number is
\[
 2^{m-2}-\left\lfloor\frac{2^{m-1}}3+\frac12\right\rfloor,
\]
an odd integer. Hence $N_{K_n/\mathbf Q}(u)=-1$.
The signature vector of $u$ has augmentation one in the local ring
$\mathbf F_2[\operatorname{Gal}(K_n/\mathbf Q)]$ and is therefore
a unit. Its Galois translates span all signatures. At $n=0$ use $-1$.

Kummer theory now gives
$\dim H^1(G_{K_n,S_n},\mathbf F_2)=r_n+s_n$, and its signature map
to $\mathbf F_2^{r_n}$ is onto. The map in degree zero is the diagonal
injection. The positive complex is the fiber of ordinary global
cochains to the sum of ordinary real cochains, so its long exact
sequence gives $\dim H^1(C_n^+)=r_n+s_n-1$.
The vanishing two-part of the $S_n$-class group identifies global
$H^2$ with $\operatorname{Br}(\mathcal O_{K_n,S_n})[2]$.
The local invariant description has $r_n+s_n$ coordinates with one
sum relation. Projection to the $r_n$ real coordinates is onto,
because $s_n\ge1$, and has kernel of dimension $s_n-1$.
Together with signature surjectivity this gives
\[
 \dim H^1(C_n^+)=2^n+s_n-1,\qquad
 \dim H^2(C_n^+)=s_n-1,
 \qquad H^i(C_n^+)=0\quad(i\ne1,2).
\]
The last assertion uses the usual global-to-real comparison in
degrees greater than two; degree zero vanishes by the diagonal injection.

For each fixed support, $s_n$ is bounded. The dyadic prime is unique;
each fixed odd prime has open decomposition group in the real
cyclotomic $\mathbf Z_2$-extension, since its image in
$\mathbf Z_2^\times/\{\pm1\}$ has nonzero two-adic logarithm.
Use now the finite free model and coefficient change already constructed.
Its scalar reduction is a bounded finite free $\mathbf F_2[[t]]$-complex,
and its finite cyclotomic specializations are derived quotients by
$t^{2^n}$. Smith decomposition shows that torsion disks and their Tor
terms contribute bounded dimensions, while a free cohomology generator
contributes $2^n$. The displayed dimensions therefore give generic
rank one in degree one and zero in all other degrees.

For the present curve, $E[2]$ is the direct sum of two trivial
$\mathbf F_2$-representations. Thus the residual positive complex has
rank two in degree one only. The residue field of the local ring
$B=\mathscr O[G]$ is $\mathbf F_2((t))$; the augmentation ideal of
the finite two-group lies in its maximal ideal. A minimal finite free
$B$-model has zero residual differential, so the just-computed dimensions
force exactly two generators in degree one and no others. Consequently
\[
                    C^+\otimes B\simeq B^2[-1].\tag{E.3}
\]
This uses the local ring $B$, not semisimplicity of $\mathbf F_2[G]$.

The common real sign is needed here. The real cyclotomic character
and $\varepsilon$ are trivial on complex conjugation, so the real
action is the fixed action on $T_2E\otimes\chi_0$, extended to the
group ring. Its invariant module is the fixed free rank-one
$\Z_2$-lattice extended to $\Lambda[G]$, and then to
$\mathscr O[G]$. Thus a generator can be chosen before $G$ varies.
Real $H^1$ is killed by two over these coefficient rings, so the
positive lift below costs the same single factor two for every
binary dimension.
Choose a generator of the fixed real invariant lattice and its
connecting class $a$. Since real $H^1$ is killed by two,
$2^{c_E+1}F_EZ_F$ has an integral positive lift $b$.
For a basis $\theta$ of the original inverse determinant line, define
\[
                         u^{\rm ev}\theta=2^{-c_E}a\wedge b.
 \tag{E.4}
\]
The wedge does not depend on the choice of the positive lift.
Indeed, the difference of two lifts is in the image of real
degree-zero cohomology by the cone exact sequence. That module
is generated by the chosen real invariant generator, so the
difference is a multiple of $a$ and wedges to zero with $a$.
After rationalization set $z=F_EZ_F$ and choose any lift
$\widetilde z$ of $z$. Formula (E.4) then reads
\[
              u^{\rm ev}\theta=2a\wedge\widetilde z.
\]
Thus the factor $2^{c_E}$ used to make the lift integral cancels
exactly. Increasing the fixed clearing exponent does not change
the rational scalar. Character evaluation and derived central
specialization preserve this equality and its cone maps. The
class whose dual exponential occurs in (E.7) is consequently
the same $F_E(Z_F)_h$ as in (K3), and the single remaining factor
two is exactly the $+1$ in (E.7).

\paragraph{From the full cohomology bound to the same determinant scalar.}
Here is the rational divisibility step explicitly. Fix a character
$x$ of $G$, a height-one prime $\mathfrak p$ of $A$, and
$R=A_{\mathfrak p}$. Put $z=F_E(Z_F)_x$, without changing its
normalization. The full global degree-one module is free of rank one
over $R$ by the argument following (K15) and the injectivity of the
maps in (K16). Write it as $Re$, and write $z=\alpha e$.
The higher real cohomology vanishes over $R$, since two is invertible.
Global degree-zero cohomology is zero: an invariant would restrict
to an invariant of the open nonabelian subgroup on the elliptic
representation, on which the scalar coefficient character is trivial.
Thus the positive-cone exact sequence gives
\[
 0\longrightarrow R a\longrightarrow
 H^1(C^+_x\otimes R)\longrightarrow Re\longrightarrow0,
 \qquad
 H^2(C^+_x\otimes R)=H^2(C_{S,x}\otimes R).
\]
The element $a$ here is the image of the fixed real invariant generator
in (E.4), and it is primitive in this positive degree-one module,
because the quotient $Re$ is free. Choose a lift $\widetilde e$ of $e$;
then $(a,\widetilde e)$ is a basis. The equality already proved from
(E.4) becomes
\[
                 u_x^{\rm ev}\theta_x
                  =2\alpha\,a\wedge\widetilde e.
\]
Let $d_1,\ldots,d_r$ be the nonzero Smith entries of a finite free
model of $C^+_x\otimes R$. All its cohomology outside degrees one and
two is zero, its degree-one cohomology is free, and its degree-two
cohomology is torsion. Its inverse determinant lattice, under the
rational identification with $\bigwedge^2H^1$, is therefore generated,
up to an $R$-unit, by
$(\prod_i d_i)a\wedge\widetilde e$. It follows that
\[
 \begin{split}
 v_{\mathfrak p}(u_x^{\rm ev})
 &=v_{\mathfrak p}(\alpha)-\sum_i v_{\mathfrak p}(d_i)\\
 &=\ell_R(H^1(C_{S,x}\otimes R)/Rz)
       -\ell_R H^2(C_{S,x}\otimes R)\ \geq\ 0,
 \end{split}                                                   \tag{E.4a}
\]
where two is a unit and the inequality is exactly (K2).
This uses the full $H^2$ corrected in (K16), rather than only Kato's
$\mathbf H^2_0$.

The intersection of the height-one valuation rings of the normal
domain $A$ gives $u_x^{\rm ev}\in A$ for every $x$. Since $|G|$ is
a power of two, the character idempotents belong to $A[G]$ and these
memberships give $u^{\rm ev}\in A[G]$. This establishes membership
only; it is not a bound for the denominator from character inversion.
The uniform bound comes separately from (E.3): in its free rank-two
$\mathscr O[G]$-module the integral vectors $a,b$ of (E.4) have integral
wedge, so $2^{c_E}u^{\rm ev}\in\mathscr O[G]$ with the fixed exponent
already used for the original lattice. Coefficientwise intersection
in the original group basis gives


\[
              A\cap\mathscr O=\Lambda,\qquad
                  2^{c_E}u^{\rm ev}\in\Lambda[G].                 \tag{E.5}
\]
Indeed an element of $A$ has bounded powers of two as denominators
in its power-series coefficients; membership in $\mathscr O$ makes
their minimum valuation nonnegative. The coefficientwise intersection
retains $c_E$, without a term $v_2|G|$.
\end{proof}

\subsection{Central saturation and the numerical lower bound}
\begin{proposition}[Central calibration of the even scalar]\label{ex:app-central-calibration}
Use the scalar and determinant basis of
Proposition~\ref{ex:app-even-scalar}. At a corank-zero constituent
$E^h$ with finite $\Sha[2^\infty]$, its nonzero central value satisfies
\[
 v_2u_h^{\rm ev}(0)=v_2L_*(h)-2\omega(h)-s(h)+\epsilon_h,
 \qquad |\epsilon_h|\le C_{\rm loc},
\]
with one constant for the fixed local types. At such a constituent
$u_h^{\rm ev}(0)=0$ if and only if $L_*(h)=0$.
Consequently every admissible $h$ with $L_*(h)\ne0$ satisfies
$v_2L_*(h)\ge2\omega(h)+s(h)-c_E-C_{\rm loc}$: the forward
rank-zero theorem in Lemma~\ref{lem:forward-rank} supplies the
preceding corank and finiteness hypotheses in this case. The central saturation corrections are those
in (E.6)--(E.7); no primitivity after inverting $t$ is assumed to
survive central specialization.
\end{proposition}
\begin{proof}\leavevmode
At a corank-zero constituent put
\[
 g=\ell_{\Z_2}E^h(\Q)[2^\infty],\quad
 g_v=\ell_{\Z_2}E^h(\Q_v)[2^\infty],\quad
 r_i=\ell_{\Z_2}H^i(\mathbf R,T_h)\quad(i=1,2).
\]
Let $i$ be the index length of the image of global $H^1$ in the
free rank-one quotient $L_2=H^1(\Q_2,T_h)/H_f^1(\Q_2,T_h)$.
Orthogonal Kummer duality, with its real terms, gives
\[
 \ell H^2_{\rm glob}=s(h)-i+\sum_v g_v+r_2-g.                     \tag{E.6}
\]
Write $K=\ker(H^1_{\rm glob}\to H^1_{\mathbf R})$, let $k$ be its
torsion length and $m$ its free index. If $P=H^1(C_h^+)$, $p_0$ is
its torsion length and $a_s$ the saturation index of the real image
in $P/P_{\rm tors}$, exactness gives
\[
 \ell\operatorname{coker}(H^1_{\rm glob}\to H^1_{\mathbf R})
       =r_1-g+k-m,\qquad k=p_0+a_s.
\]
For the numerical calculation suppose also $L_*(h)\ne0$.
The wedge in (E.4) adds $a_s+1+v_{\rm free}(z)-m$.
Conversion to the inverse determinant adds $p_0-\ell H^2(C_h^+)$.
Substitution of these equalities cancels $m$ and $a_s$ exactly:
\[
 v_2u_h^{\rm ev}(0)
 =v_2\exp^*z-s(h)-\sum_v g_v+2g+1-r_1-e_2,                      \tag{E.7}
\]
where $e_2$ is the valuation of $\exp^*$ on a generator of $L_2$.
Primitivity of the real image after inverting $t$ has not been
specialized to $t=0$.

Here are the real local convention and the exact length sequence used
in (E.6). Put $T=T_2(E^h)$ and $D=E^h[2^\infty]$. Since $E^h[2]$ is
rational, complex conjugation $J$ satisfies $J\equiv1\pmod2$.
The integral projectors $(1\pm J)/2$ give $T=T^+\oplus T^-$.
At the real place use the bounded Kummer complex
\[
 U_{\mathbf R}=\tau_{\le1}R\Gamma(\mathbf R,T)
   =[\,T^+\oplus T^-\xrightarrow{\,J-1\,}T^-\,]
\]
in degrees zero and one, with its map to ordinary real cochains.
Modulo $2^n$ its degree-zero cohomology is
$T^+/2^n\oplus(T^-/2^n)[2]=E^h(\mathbf R)[2^n]$, and its degree-one
image is the real Kummer line $T^-/2$. The latter can also be checked
using a rectangular period lattice: dividing the nonidentity real
component, represented by half the imaginary period, produces the
minus-line cocycle. Thus this line is its exact annihilator under
finite real Tate duality. The ordinary completion
$E^h(\mathbf R)^\wedge_2[-1]$ alone would lose the first summand of
degree-zero cohomology.

At a finite place take the derived Kummer map from
$E^h(\Q_v)^\wedge_2[-1]$. Its reduction has full local degree-zero
cohomology and injective Kummer degree-one map. If $C_{\rm Sel}$
denotes the resulting ordinary global Selmer cone, these assertions
identify
\[
 H^1(C_{\rm Sel}\otimes^{\mathbf L}\Z/2^n)
       =\operatorname{Sel}_{2^n}(E^h/\Q).
\]
At corank zero the discrete Selmer group is finite, the Mordell--Weil
rank is zero, and its order is that of $\Sha[2^\infty]$. The
orthogonal-Kummer specialization of finite Poitou--Tate duality
\cite[5.7.1.4--5.7.1.7]{nekovar}, followed by the compact/discrete
limits, gives
\[
\begin{split}
0\longrightarrow L_2/J_2
 \longrightarrow\operatorname{Sel}_{2^\infty}(E^h/\Q)^\vee
 \longrightarrow H^2_{\rm glob}
 \longrightarrow \bigoplus_{v\in S_f}H^2(\Q_v,T)\oplus H^2(\mathbf R,T)
 \longrightarrow H^0(\Q,D)^\vee\longrightarrow0 ,
\end{split}
\]
where $J_2$ is the image of global $H^1$. There is no real
degree-one quotient: the compact real Kummer image is all of
$H^1(\mathbf R,T)$, and its discrete orthogonal Kummer group is zero.
At finite $v\ne2$ the compact Kummer image is likewise all of local
$H^1$, while at two its quotient is $L_2$. The real degree-two term
is retained. The finite-level images satisfy Mittag--Leffler, and
high-degree localization at the real place terminates the sequence.

This also makes $H^2_{\rm glob}$ finite. The cone sequence then
makes $H^2(C_{\rm Sel})$ finite, and universal coefficients give
\[
 0\to E^h(\Q)/2^n
 \to\operatorname{Sel}_{2^n}(E^h/\Q)
 \to H^2(C_{\rm Sel})[2^n]\to0.
\]
Comparison with the Kummer sequence gives
$\ell H^2(C_{\rm Sel})=s(h)$, with no extra real torsion. Equivalently
the cone sequence is
\[
0\to L_2/J_2\to H^2(C_{\rm Sel})\to H^2_{\rm glob}
 \to \bigoplus_{v\in S_f}H^2(\Q_v,T)\oplus H^2(\mathbf R,T)
 \to H^0(\Q,D)^\vee\to0.
\]
Its alternating lengths are exactly (E.6).

For the positive cone, $H^2_{\rm glob}\to H^2(\mathbf R,T)$ is onto:
any real component's image in $H^0(\Q,D)^\vee$ can be canceled by
a component at two, since restriction
$H^0(\Q,D)\to H^0(\Q_2,D)$ is injective. Thus its exact sequences are
\[
0\to T^+\xrightarrow{a}P\to K\to0,\qquad
0\to\operatorname{coker}(H^1_{\rm glob}\to H^1_{\mathbf R})
 \to H^2(C_h^+)\to H^2_{\rm glob}\to H^2_{\mathbf R}\to0.
\]
If $k$ and $p_0$ are the torsion lengths of $K$ and $P$, $m$ is the
free index of $K$ in global $H^1$, and $a_s$ is the saturation index
of the real line in $P/P_{\rm tors}$, these sequences give
\[
\ell\operatorname{coker}(H^1_{\rm glob}\to H^1_{\mathbf R})
 =r_1-g+k-m,\qquad k=p_0+a_s.
\]
The wedge $2^{-c_E}a\wedge b$ contributes
$a_s+1+v_{\rm free}(z)-m$, and the inverse determinant comparison
contributes $p_0-\ell H^2(C_h^+)$. Substitution leaves
$1+v_{\rm free}(z)+i-s(h)-\sum_vg_v+2g-r_1$.
Finally $v_2\exp^*z=v_{\rm free}(z)+i+e_2$, proving (E.7)
without specializing a primitivity assertion from the punctured
family.
Apply (K3) to the retained class $z=F_E(Z_F)_h$ used in the wedge
above. The primitive-conductor trace and the invariant
$\sqrt h\,e_{\chi_h}$ basis have already been computed in
paragraph~6, including discriminant $4h$. Absorb the fixed
smoothing and old-support factors into $\epsilon_{\rm rec}$.
Then
\[
 v_2\exp^*z
 =v_2L_*(h)+\sum_{q\ {\rm unused}}v_2\#E^h(\mathbf F_q)
                         +\epsilon_{\rm rec},\qquad
 |\epsilon_{\rm rec}|\le C_E.
\]
At unused good primes $g_q=v_2\#E^h(\mathbf F_q)$, so these
terms cancel. At an active good prime the inertia invariants in
$E^h[2^\infty]$ are $E[2]$, of length two for the present curve.
The remaining local types have fixed bounds, including global
torsion by its injection at two.
Thus
\[
 v_2u_h^{\rm ev}(0)=v_2L_*(h)-2\omega(h)-s(h)+\epsilon_h,\qquad
 |\epsilon_h|\le C_{\rm loc}.                                  \tag{E.8}
\]
The central rational positive complex has only rank-two $H^1$, so
this is specialization of (E.4). Corank-zero duality also shows
$u_h^{\rm ev}(0)=0$ if and only if $L_*(h)=0$.
At $L_*(h)\ne0$, Lemma~\ref{lem:forward-rank} supplies the
corank-zero and finite-$\Sha$ hypotheses just used.
Equations (E.5) and (E.8), already with $G=1$, yield
\[
          v_2L_*(h)\ge2\omega(h)+s(h)-c_E-C_{\rm loc}.             \tag{E.9}
\]
No binary-family dimension is needed for this lower bound.
\end{proof}

% Replaces baseline manuscript.tex lines 3416--3510 inclusive.
% Baseline SHA256 525c8c05e7a5acc83d6b1f89c1056fd1321bd7490fd227a63d1a520e941b5197.
\subsection{The actual even form and its phase}
\begin{proposition}[An actual even form with rational normalized phase]\label{ex:app-even-phase}
Fix a sign $\sigma$, unit squareclasses at a finite support $S_0$
containing $2N$, and a positive squarefree reference $D_0$ in that
filter with $L_*(\sigma D_0)\ne0$. For every fixed finite extra
exclusion set $B$ prime to $D_0$, there is an algebraic holomorphic
cuspidal form of weight $3/2$, normalized by $a(D_0)=1$, with the
scalar characters $\psi,\eta$ constructed below and
$\psi/\eta=\chi_{-\sigma}$. It has the prescribed coefficient
support. For every retained positive squarefree $D$ in the original
$S_0$ filter,
\[
 a(D)=z_Dr_D,\qquad z_D\in\mu_{2^\infty},\quad r_D\in\Q,
 \qquad r_D^2=L_*(\sigma D)/L_*(\sigma D_0).
\]
It satisfies the all-index good-prime recurrence (E.12) outside
$S_0\cup B$. Nonzero retained squarefree coefficients have integral
$v_2a(D)$, and $v_2a(D)-\omega(D)$ has a finite attained minimum.
The form and the initial sieve are fixed before this minimum is
chosen. Additional exclusions do not authorize a new minimum.
\end{proposition}
\begin{proof}\leavevmode
Fix a sign $\sigma$ and unit squareclasses at a finite support $S_0$
containing two and every prime dividing $N$. Let $D_0>0$ be squarefree
in these prescribed classes, with $L_*(\sigma D_0)\ne0$.
In (E.10)--(E.11), $D$ ranges over positive squarefree parameters in
these same unit squareclasses at every prime of $S_0$.
Fix also a finite extra exclusion set $B$ with
$\gcd(D_0,\prod_{p\in B}p)=1$.

\paragraph{A nonzero reference functional and its scalar local type.}
Use the global Shimura--Waldspurger packet with underlying
$\operatorname{PGL}_2$ representation $E^{(\sigma)}$ in the
normalization of Waldspurger's Corollary~2. The nonzero central value
at $\sigma D_0$ supplies the member with a nonzero $D_0$-Whittaker
functional. Equivalently, take the nonzero theta lift associated with
$E^{(\sigma D_0)}$ and the additive character of frequency $D_0$;
the change of additive character identifies its underlying packet
with the preceding one. The nonvanishing and the local identification
are the content of Waldspurger's global theta theorem and its
Fourier-functional calculation
\cite[Theorem~1 and Proposition~1]{waldspurgerICM}
and \cite[Proposition~18 and its proof, pp.~466--467]{waldspurger}. The positive real frequency selects the
holomorphic member of lowest weight $3/2$. The lift is in the
unary-orthogonal cuspidal spectrum.

Here is a local construction of a scalar type that retains the
reference functional. At $p\in S_0$, write $\ell=\ell_{D_0}$ and
choose $v$ with $\ell(v)\ne0$. Normalize Haar measure in
\[
 P_Rv=\frac1{\operatorname{vol}(p^{-R}\Z_p)}
       \int_{p^{-R}\Z_p}
          e_p(-D_0b)\pi(n(b))v\,db .
\]
Then $\ell(P_Rv)=\ell(v)$. Let $M(u)$ be the unit action normalized
by the Weil multiplier. Explicitly, if the chosen torus lifts have
cocycle $(u,v)_p$, their theta-cubed multiplier $\gamma$ has
two-power root-of-unity values and satisfies
$\gamma(u)\gamma(v)=(u,v)_p\gamma(uv)$; put
$M(u)=\gamma(u)^{-1}\pi(\widetilde m(u))$.
It is therefore an action of
$\Z_p^\times$ and
$M(u)\pi(n(b))M(u)^{-1}=\pi(n(u^2b))$.
Changing $u$ to its inverse changes none of the following support
conditions. Direct substitution gives
\[
 \ell(M(u)P_Rv)=
 {\bf1}_{\,u^2\equiv1\;(\mathrm{mod}\;p^R)}
                    \ell(M(u)v).
\]
The integral is a finite sum on this smooth vector.

The normalized $M(-1)$ acts by a scalar $\delta\in\{1,-1\}$.
Choose a continuous unit character $\alpha$ of two-power order
with $\alpha(-1)=\delta$. Such a character exists: for odd $p$,
$-1$ lies in the two-primary factor of the finite cyclic group
of residue-field units; for $p=2$ use the factor
$\{1,-1\}$ of $\Z_2^\times$.
Take $R$ large enough that both $v$ and $\alpha$ are constant on
the small neighborhoods appearing above. The projection
\[
 w=\int_{\Z_p^\times}\alpha(u)^{-1}M(u)P_Rv\,d^\times u
\]
satisfies $M(u)w=\alpha(u)w$. Its $\ell$-value is nonzero:
the neighborhoods of $1$ and $-1$ contribute the same nonzero
value. For odd $p$ the solution set is
$1+p^R\Z_p$ together with $-1+p^R\Z_p$.
For $p=2$ and $R\ge3$ it is
\[
 (1+2^{R-1}\Z_2)\ \cup\
 (-1+2^{R-1}\Z_2).
\]
Thus the four residue classes modulo $2^R$ form two neighborhoods;
there are not four independently assignable phases.
For completeness, the noncancellation in this projection is an exact
calculation. Normalize multiplicative Haar measure to give
$\Z_p^\times$ volume one, and put
$K_R=1+p^R\Z_p$ for odd $p$, and $K_R=1+2^{R-1}\Z_2$
for $p=2$. Increase $R$ until $M(k)v=v$ and $\alpha(k)=1$
for $k\in K_R$. On $K_R$ the integrand after applying $\ell$
is $\ell(v)$. On $-K_R$ it is
$\alpha(-1)^{-1}\ell(M(-1)v)=\delta^{-1}\delta\ell(v)$.
It vanishes elsewhere by the displayed additive projection formula.
Consequently
\[
                 \ell(w)=2\operatorname{vol}(K_R)\ell(v)\ne0.
\]
This calculation takes place in characteristic zero; reduction at two
is not involved. For the later support projection choose $m$ so that
the required unit squareclass is a union $A$ of residue classes modulo
$p^m$. Its operator is
\[
 \Pi_A=\frac1{p^m}\sum_{b\bmod p^m}
       \left(\sum_{a\in A}e_p(-ab/p^m)\right)\pi(n(b/p^m)).
\]
On the Fourier functional of integral frequency $D$ it acts as
$\mathbf1_A(D)$. It therefore preserves the reference value and
annihilates every excluded frequency. Conjugation by $M(u)$ permutes
the summands and replaces $A$ by $u^2A=A$, proving the asserted
commutation. Its denominators are fixed before the form is normalized.


The vector $P_Rv$, and then $w$, is fixed by upper unipotents
$n(\Z_p)$, because $D_0$ is a unit. Smoothness makes $w$
fixed by lower unipotents at sufficiently deep level. Together with
the unit eigencharacter, the elementary decomposition of
$K_0(p^m)$ gives a scalar $K_0$ type for sufficiently large $m$.
The Weil normalization contributes only two-power roots of unity,
so the resulting classical character has two-power order.
The forced values at $-1$ give the global parity required for a
weight-$3/2$ scalar character: the product of the finite central
actions is the inverse of the holomorphic archimedean action,
and the theta-cubed multiplier has the same product relation
on rational $-I$. Thus the quotient character satisfies
$\psi(-1)=1$. At $2$ one may, if desired, choose
$\alpha$ of arbitrarily large two-power conductor before increasing
$R$; this does not enlarge the prime support.

Apply the finite additive Fourier idempotent for the prescribed
unit-squareclass subset modulo a sufficiently high power of $p$.
This idempotent keeps the $D_0$ coefficient. Its support condition
is stable under multiplication by unit squares, so it commutes
with the unit action and retains a scalar type, at a possibly
deeper level. At every prime outside $S_0$ take the spherical vector.
The packet there is unramified; its standard squarefree coefficient
$c_p^0$ is $1$ at both valuations zero and one
\cite[Section~I.3, pp.~377--378]{waldspurger}.
This includes the primes dividing $D_0$. Local Whittaker uniqueness
and the tensor-product factorization show that the resulting
global reference coefficient is still nonzero.

\paragraph{The exact full-space comparison.}
Let the resulting scalar character be $\psi$ and put
$\eta=\psi\chi_{-\sigma}$. Both have two-power order.
Let $V$ be the full good-Hecke eigenspace of this scalar character,
at the level just constructed. Keep all old multiplicities. In
Waldspurger's notation its integral-weight eigenform and character are
\[
 \phi_0=f_E\otimes\eta,\qquad
 \chi_0=\psi\chi_{-1},\qquad
 \phi_0\otimes\chi_0^{-1}=f_E\otimes\chi_\sigma .
\]
This also fixes the signs in the packet convention: Waldspurger's
positive index $D$ twists the last representation by $\chi_D$.
The Hecke eigenvalues in $V$ are $\eta(p)a_p(E)$.

The nonzero space implies hypothesis (H1) by Waldspurger's
Proposition~2 \cite[Section~I.2, p.~377]{waldspurger}. In explicit terms, at a principal-series
place its two inducing characters take value $1$ at $-1$.
Corollary~2 \cite[pp.~483--484]{waldspurger} applies to every $f\in V$ and any
positive squarefree pair whose ratio is a square in every
$\Q_p$ at a prime dividing the level. Our $D,D_0$ satisfy
this condition because the level has prime support $S_0$.
The final paragraph of that proof removes (H2) for this
same-$\Q_2$-squareclass comparison. No assumption of
non-supercuspidality at two is required.

\paragraph{The printed coefficient identity and its dyadic hypothesis.}
For clarity, write $\mathcal L_D=L(E^{(\sigma D)},1)$ and let
$N_{\rm W}$ denote the level in Waldspurger's notation, not the
conductor of $E$. In \cite[Corollary~2, pp.~483--484]{waldspurger}
take
\[
 k=3,\quad \chi=\psi,\quad \chi_0=\psi\chi_{-1},\quad
 \rho_0=\rho_E\otimes\eta,\quad n_1=D_0,\quad n_2=D.
\]
Then $\rho_0\otimes\chi_0^{-1}=\rho_E\otimes\chi_\sigma$.
The corollary, in its printed order, is
\[
 a_f(D_0)^2\mathcal L_D\,\psi(D/D_0)D^{1/2}
   =a_f(D)^2\mathcal L_{D_0}D_0^{1/2}.
\]
The passage from the automorphic central argument $1/2$ to the
classical argument $1$ is the same for both terms. Dividing only
by the nonzero reference value gives (E.10). In particular this
is a literal square identity, with neither complex conjugation
nor an unspecified absolute-value normalization.

Every prime of $N_{\rm W}$ belongs to $S_0$ at this point of the
construction. The prescribed local classes make $D/D_0$ a square
in $\mathbb Q_p^\times$ at each such prime, including $p=2$.
In the proof of the corollary, the local coefficient functions
at these primes have equal values at $D,D_0$. Outside the level,
the squarefree arguments have valuation zero or one and the
unramified coefficient function has value one in both cases.
The real coefficient contributes exactly the power
$(D/D_0)^{(k-2)/4}$ before squaring. The quotient of epsilon
factors is computed on pp.~483--484: its Hilbert-symbol product
disappears by reciprocity and its remaining character is
$\chi(D/D_0)=\psi(D/D_0)$. Thus these local terms do not leave
an additional prime-dependent factor in the displayed identity.

The final paragraph on p.~484 explicitly explains the use of
(H2): it is needed in that argument only to compare Fourier
indices in different squareclasses of $\mathbb Q_2^\times$.
Our pair is in one squareclass. Hence the corollary as stated
requires (H1) here, but does not impose (H2), a nonsupercuspidal
condition at two, or a new dyadic conductor restriction. The
nonzero scalar eigenspace supplies (H1) by
\cite[Proposition~2, p.~377]{waldspurger}. Moreover the proof on
p.~483 extends its linear relation to every vector in the full
eigenspace before squaring. The resulting identity therefore
applies simultaneously to old multiplicities, and later to both
Galois-conjugated and character-twisted vectors at a common level.
At that later level the same squareclass test is repeated; its
prime support is still $S_0$.
Putting $n_1=D_0,n_2=D$ in the corollary, with weight $3/2$, gives
\[
 a_f(D)^2=
 \psi(D/D_0)\,
 \frac{\sqrt D\,L(E^{(\sigma D)},1)}
      {\sqrt{D_0}\,L(E^{(\sigma D_0)},1)}
 a_f(D_0)^2.                                      \tag{E.10}
\]
The period cancels, so the quotient in this formula is the quotient
of the specified $L_*$-values. No product over the primes dividing
$D$ or $D_0$ is missing: the unramified factors at squarefree
valuations are exactly those checked above.

For clarity, (E.10) gives proportional coefficient functionals on
the whole space. If $\ell_0(f)=a_f(D_0)$ and $\ell_D(f)=a_f(D)$,
then $\ell_D^2=C\ell_0^2$ as a polynomial identity on $V$.
If $C=0$, then $\ell_D=0$. Otherwise, choose a square root $c$ of
$C$. The polynomial ring on the finite-dimensional vector space
$V$ is an integral domain, so
$(\ell_D-c\ell_0)(\ell_D+c\ell_0)=0$ implies
$\ell_D=c\ell_0$ or $\ell_D=-c\ell_0$.
The scalar is common to every vector. This argument does not
assume that $V$ is one-dimensional.

\paragraph{An algebraic vector and the Galois comparison.}
Ramsey's section model \cite[Definition~3.1 and Theorem~4.2]{ramsey} realizes the half-integral-weight space as
sections of the indicated divisor sheaf on $X_1(4N_{\rm aux})$.
Here $N_{\rm aux}$ is supported on $S_0$ and includes the chosen
local depths; it need not equal the conductor $N$ of $E$.
Over characteristic-zero fields, flat field extension commutes
with these sections. Cuspidality, scalar character, support
conditions and prescribed good-Hecke eigenvalues are algebraic
linear conditions. Although infinitely many coefficients or
primes can occur in their description, a finite subset cuts out
the same subspace, by finite dimensionality.

One need not impose Petersson orthogonality as an additional
transcendental linear condition. The elliptic good-Hecke eigenspace
is automatically orthogonal to unary theta series: at a good prime,
the unary eigenvalues have absolute value $p+1$, whereas
$|\eta(p)a_p(E)|\le2\sqrt p<p+1$. The unary space and its
orthogonal complement are stable under the Hecke operators and
their adjoints. Thus their elliptic eigenspaces are disjoint.
The nonzero reference functional on the complex space is consequently
nonzero on its algebraic model. Choose such an algebraic $g$ and
scale it to $a(D_0)=1$.
This uses flat field extension only, not Ramsey's arbitrary-ring
base-change assertion in weight $3/2$; see \cite[Remark~4.4]{ramsey}.
Here the field-extension assertion can be checked without a
cohomology-vanishing theorem. Work over a number field $K$ containing
the character values and the roots of unity specifying the cusps and
the support idempotents. A finite affine cover of the modular curve
computes global sections as the kernel of the difference of restriction
maps to pairwise intersections. Tensoring with a field extension is
flat and preserves that kernel. Thus the sections in Ramsey's
Definition~3.1 commute with $K\subset\overline\Q\subset\mathbf C$,
also in weight $3/2$. Impose vanishing at the cusps, the diamond
character, the support idempotents, and the kernels of
$T(p^2)-\eta(p)a_p(E)$ inside this finite-dimensional space.
The span in its dual of all the defining linear equations has a finite
basis. These conditions therefore commute with the same field
extensions. Since the reference functional is nonzero after extension
to $\mathbf C$, it is nonzero before extension; a vector in a $K$-basis has
nonzero reference value and may be divided by that value.
This also places every coefficient of the chosen form in one number
field. No assertion about base change to a residue field, nor about
vanishing of $H^1$ in weight $3/2$, is used.


Let $\tau\in\operatorname{Aut}(\overline{\Q}/\Q)$.
Choose an odd integer $e$ describing its action on all values of
$\eta$ and $\psi$, and put $\nu=\eta^{(e-1)/2}$.
Coefficientwise conjugation gives a form $g^\tau$ with character
$\psi^e$ and good eigenvalues $\eta(p)^e a_p(E)$.
One can reduce this conjugation assertion to integral weight by
multiplying by the ordinary theta series: it has rational
coefficients, and its fourth power has trivial scalar character.
The algebraic divisibility condition at the cusps is precisely the
one in Ramsey's section model, so division after conjugation is
again a half-integral-weight form.
The coefficientwise twist $g\otimes\nu$ has character
$\psi\nu^2=\psi^e$ and exactly these good eigenvalues.
The identities hold because $\psi/\eta$ is quadratic.
The twist is modular at a level still supported on $S_0$; this
follows as well by writing it as the finite additive translation
sum using the Gauss sum of the primitive character inducing $\nu$.

Pass to a common such level. Both forms lie in its full
good-Hecke eigenspace, so the proportionality of its $D$ and $D_0$
functionals gives
\[
 \tau(a(D))
 =\frac{\nu(D)a(D)}{\nu(D_0)a(D_0)}
 =\nu(D/D_0)a(D).                                  
\]
There is no vanishing reference denominator: $D_0$ is prime to
$S_0$ and $a(D_0)=1$. No rationality of the central-value quotient
has been assumed in obtaining this equality.

Write the Dirichlet character as a product of its finite
prime-power factors $\eta_p$. For every $p$ dividing its conductor
choose $u_p\in\Z_p^\times$ with $u_p^2=D/D_0$, and set
\[
                 z_D=\prod_{p\mid\operatorname{cond}(\eta)}
                                      \eta_p(u_p).
\]
Then $z_D$ is a two-power root of unity,
$z_D^2=\eta(D/D_0)$, and
\[
 \tau(z_D)=z_D^e,\qquad
 \nu(D/D_0)=z_D^{e-1}.
\]
Changing the selected odd representative of $e$ has no effect on
the last identity, because $D/D_0$ is a local square at the
conductor. The preceding Galois identity now shows that $a(D)/z_D$ is fixed by
every $\tau$, hence lies in $\Q$.
Different choices of the local square roots change $z_D$ by a
sign, which changes only the sign of this rational number.

Finally, $\psi/\eta=\chi_{-\sigma}$ is trivial on $D/D_0$:
its conductor is supported at two and the ratio is a square there.
Thus $\psi(D/D_0)=z_D^2$. Dividing (E.10) by $z_D^2$ gives
\[
 a(D)=z_Dr_D,\qquad z_D\in\mu_{2^\infty},\qquad r_D\in\Q,
 \qquad r_D^2=\frac{L_*(\sigma D)}{L_*(\sigma D_0)}.       \tag{E.11}
\]
The argument also covers the zero coefficient case. In particular,
every nonzero coefficient in (E.11) has integral normalized
two-adic valuation.

\paragraph{Recurrence, later exclusions, and valuations.}
Only now impose the extra finite coprimality sieve $B$.
The zero-or-retain Fourier sieve is a finite additive translation
operator. It preserves cuspidality, keeps every coefficient whose
index is prime to $B$, and commutes with good $T(p^2)$ for
$p\notin S_1=S_0\cup B$. In particular it keeps $a(D_0)=1$.
The squarefree comparison (E.11) was proved before this enlargement,
so it requires no extra squareclass hypotheses at the primes of $B$.
After the sieve it is asserted at the retained indices.
The standard good $T(p^2)$ formula, with $\psi^2=\eta^2$, gives
\[
 a(p^2n)=\{\eta(p)a_p(E)-\psi(p)(-n/p)\}a(n)
                       -p\eta(p)^2a(n/p^2),\qquad p\notin S_1.
                                                        \tag{E.12}
\]
For squarefree $D$, write $a(Dp^{2r})=\eta(p)^r b_r$.
The first multiplier is
$a_p(E)-(\psi/\eta)(p)(-D/p)$; subsequent multipliers satisfy
$b_{r+1}=a_p(E)b_r-pb_{r-1}$. All are integers, including when
$p\mid D$. Iteration over the good primes gives the same assertion
for every retained square shell after removal of the $\eta$ phase.
More explicitly, for $p\notin S_1$ set
$A_{p,0}(D)=1$, $A_{p,1}(D)=a_p(E)-(\psi/\eta)(p)(-D/p)$,
and
$A_{p,r+1}(D)=a_p(E)A_{p,r}(D)-pA_{p,r-1}(D)$ for $r\ge1$.
Then, for every retained index $Dm^2$ with $D$ squarefree,
\[
 a(Dm^2)=z_D\eta(m)r_D
              \prod_{p^r\Vert m}A_{p,r}(D).
\]
The formula includes $p\mid D$, when the initial Legendre symbol
is zero. When applying the recurrence at one prime, the other
prime-square factors do not alter that Legendre symbol. This proves
the formula by induction on the prime factors of $m$, with
$a(n/p^2)=0$ when $p^2\nmid n$. All factors $A_{p,r}(D)$ are
integers. Thus every coefficient used by the extraction and shell
operators of Proposition~\ref{ex:app-coefficient-rules} is a rational
number times a two-power root of unity, including coefficients with
non-squarefree index. In the chosen residue field such a root reduces
to one. A rational coefficient of valuation zero reduces to the
element $1$ of its prime field $\F_2$. Consequently the fresh
coefficient-one reading of that proposition detects precisely equality
in the fixed minimum-depth bound. It uses the same $g$, $D_0$ and
normalization throughout.


Combining (E.9) and (E.11) makes the values
$v_2a(D)-\omega(D)$ at nonzero retained squarefree coefficients a
bounded-below subset of $\Z$. It is nonempty because $D_0$ is retained,
and therefore has an attained finite minimum. The form and the initial
finite sieve are fixed before this minimum is chosen. Restricting
later fresh primes by further exclusions does not authorize choosing
another minimum.
\end{proof}

% End 07-even-source.tex


% Begin 08-odd-source.tex
\section{Actual odd theta forms and their operators}\label{app:odd-source}
We now construct the forms in (3.8), with the same $h_*$, $D_w$, $L$
and actual Kummer stages as in Section~\ref{sec:rc}.
Write $m=|h_*|$. An admissible $D$ is positive, squarefree, prime to
the fixed support and in the prescribed unit classes. Then
$K_D=\Q(\sqrt{Dh_*})$ is imaginary, split at $2N$, has no extra
units, and $Dh_*$ is an odd fundamental discriminant.
The conductor-one point is the unaveraged sum
\[
 P_D=\sum_{[\mathfrak a]\in\operatorname{Pic}(\mathcal O_{K_D})}
       \chi_{h_*}(\Norm\mathfrak a)\,
                         \pi(x_D^{[\mathfrak a]}).               \tag{O.1}
\]
It is distinct from the conductor-$|h|$ ring-class point used in
Section~\ref{sec:limits}.

% Replace immutable baseline lines 3532--3591 inclusive.
% Keeps (O.2). This expands the existing non-CM argument.
\subsection{A fixed witness and one compatible sequence}
\begin{proposition}[One witness and a compatible detecting sequence]\label{ex:app-detector}
For the fixed $E=E_\ell$ under the full-corank-one hypothesis of
Theorem~\ref{hyp:RC}, use the prepared positive unit filter containing
$1$. There exist a negative odd fundamental discriminant $h_*$ with
$L(E^{h_*},1)\ne0$, a positive fundamental discriminant $D_w$ in
that filter with $L'(E^{D_w},1)\ne0$, and a fixed integer $L$ such
that $LP_D$ belongs to the $D$-twist rational line for every
admissible $D$ and $Q=LP_{D_w}$ is nontorsion.
There is one sequence of actual stages with $M\to\infty$, good
primes $\ell_0$, and additive maps
\[
 \lambda_M:E(\F_{\ell_0^2})\longrightarrow\Z/2^M\Z,
 \qquad v_2\lambda_M(\overline Q)=c_w<M,
\]
where $c_w$ is fixed. At these stages Frobenius squared is trivial
on $E[2^M]$, $(D_w/\ell_0)=1$, $(h_*/\ell_0)=-1$, and
$\lambda_M$ has Frobenius sign $+1$.
Any specified finite prime exclusion is permitted. Additional
field prescriptions require a compatible common lift; no preservation
of $c_w$ under arbitrary later field prescriptions is asserted.
\end{proposition}
\begin{proof}\leavevmode
We use the positive unit-squareclass filter containing one, with its
fixed support enlarged before choosing the witness. The rational
two-isogeny parity theorem \cite[Theorem 1.8]{dd} gives root number
$-1$ for $E$ under the full-corank-one assumption. Apply clause (ii)
of \cite[pp.~543--544]{bfh} to its weight-two newform, with every
old prime in the splitting set, to obtain a negative fundamental
discriminant $h_*$ with $L(E^{(h_*)},1)\ne0$. Including two makes
$h_*$ odd and excludes the extra-unit fields. Now add the primes
of $h_*$ to that set and apply clause (i) to obtain a positive
fundamental discriminant $D_w$ with $L'(E^{(D_w)},1)\ne0$.
Splitting at every old prime puts $D_w$ in the filter containing one
and makes it prime to $h_*$. These choices precede every detecting
prime and every theta precision.

The conductor-one genus character for
$K_D=\Q(\sqrt{Dh_*})$ cuts out
$F_D=\Q(\sqrt D,\sqrt{h_*})$. Thus the unaveraged sum $P_D$ in (O.1)
belongs to $E(F_D)$. Its rational character components are the
$D$- and $h_*$-components. In the degenerate case $D=1$ these are
the trivial and $h_*$ characters of the actual quadratic group.
The $h_*$-component has rank zero by the forward rank-zero theorem.
For $D_w$, the unaveraged height formula
\cite[Theorem 1.1]{cst} has Rankin derivative
$L'(E^{(D_w)},1)L(E^{(h_*)},1)\ne0$, so $P_{D_w}$ is nontorsion.
Here the conductor is one, all primes of $N$ split, and the
extra-unit fields have been excluded.

We make the fixed multiplier explicit. Choose distinct good rational
primes $p_1,p_2$ for $E$ and put
\[
 B_j=\prod_{f=1}^{4}\#E(\F_{p_j^f}),\qquad B=B_1B_2.
\]
For every number field of degree at most four, prime-to-$p_j$
torsion injects into a reduction group with residue degree at most
four. The two primes therefore show that $B$ kills its entire
torsion subgroup. Put $B_{\rm cusp}=C_\pi$, where $C_\pi$
annihilates the image of the cuspidal subgroup defined after (4.1).
This includes every Atkin--Lehner translation $\pi(W\infty)$,
which need not vanish merely because $\pi(\infty)=0$. For $\Delta_D=\operatorname{Gal}(F_D/\Q)$,
the integral character operator
\[
 4e_D=\frac4{|\Delta_D|}
                  \sum_{\gamma\in\Delta_D}\chi_D(\gamma)\gamma
\]
is defined also when $D=1$. The difference
$4P_D-4e_DP_D$ is torsion, since its only possible rational
component is the rank-zero $h_*$-component. Hence
$4BP_D=B(4e_DP_D)$ transforms exactly by $\chi_D$.
Choose $L$ to be a common multiple of the previously required fixed
multiplier and $4BB_{\rm cusp}$, for example their product. This
retains every earlier descent and local-condition factor.
Then $LP_D$ is a rational point of
$E^{(D)}$ under the standard twist identification, for every
admissible $D$. In particular
\[
                         Q=LP_{D_w}\in A(\Q),
                 \qquad A=E^{(D_w)},
\]
is nontorsion. No assertion that the original $P_D$ is
$E(K_D)$-rational is required.

Choose compatible division points $Q_n$ with
$2^nQ_n=Q$ and $2Q_{n+1}=Q_n$, and let
$c_Q(g)=(gQ_n-Q_n)_n\in T_2A$ be the Kummer cocycle.
Write $\rho_A$ for the Tate representation and $H=\ker\rho_A$.
On $H$, this cocycle is a homomorphism. Its image is compact,
and hence a closed $\Z_2$-submodule of $T_2A$. Moreover
\[
 c_Q(ghg^{-1})=\rho_A(g)c_Q(h)\qquad(h\in H),
\]
so its rational span is invariant under the Tate image.

The curve $A$ is non-CM. Open image \cite{serre} supplies both
irreducibility of its rational Tate representation and a central
scalar $z\ne1$ in its Tate image. If $c_Q(H)=0$, the rational
cocycle descends to that image. Commuting $z$ with $g$ gives
\[
 (z-1)c_Q(g)=(g-1)c_Q(z).
\]
Since $z-1$ is invertible over $\Q_2$, this is a coboundary.
It contradicts the Kummer injection
$A(\Q)\otimes\Q_2\hookrightarrow H^1(\Q,V_2A)$ and
nontorsion of $Q$. Thus the translation image is nonzero.
Irreducibility makes its rational span all of $V_2A$.
It consequently contains $2^{a_0}T_2A$ for some $a_0$.

Let $F_0/\Q$ be any fixed finite Galois auxiliary extension,
including $\Q(\sqrt{D_w},\sqrt{h_*})$ and any fixed ring-class
data to be retained with the action of complex conjugation.
Put $H_0=H\cap G_{F_0}$. It has finite index $d$ in $H$.
For $h\in H$, $h^d\in H_0$ and
$c_Q(h^d)=d\,c_Q(h)$. Hence $c_Q(H_0)$ still contains
$2^aT_2A$ for some fixed $a$.
This argument uses no independence between $F_0$ and the Tate
or division-point fields.

Fix actual complex conjugation $\tau$. Since $D_w>0$, the real
twist identification identifies its matrix on $T_2A$ with its
matrix $C$ on $T=T_2E$. We have $C^2=1$ and $\det C=-1$.
Its dual $+1$-eigenspace therefore contains a primitive integral
row $b:T\to\Z_2$, with $bC=b$ and $b(T)=\Z_2$.
Choose $t_0\in T$ with $b(t_0)=1$, and choose $h\in H_0$ with
$c_Q(h)=v=2^at_0$. Set $\sigma=h\tau$. If $d_\tau=c_Q(\tau)$,
then $(1+C)d_\tau=0$, because $\tau^2=1$. Consequently
\[
 \rho_A(\sigma)=C,\qquad
 c_Q(\sigma^2)=(1+C)(v+d_\tau)=(1+C)v,\qquad
 b\,c_Q(\sigma^2)=2^{a+1}.
\]
Put $c_w=a+1$. The element $\sigma$ and this finite valuation
are fixed before precision varies. Its restriction to $F_0$
is that of $\tau$: it fixes $\sqrt{D_w}$ and negates
$\sqrt{h_*}$. Since $h$ fixes $\sqrt{D_w}$, the matrix of
$\sigma$ on $T_2E$ is also $C$.

For each $M>c_w$, take the finite Galois extension generated by
$F_0$, $E[2^M]$, $A[2^M]$, and all $2^M$-division points of $Q$.
Restrict this same actual $\sigma$ to the extension. Finite
Chebotarev supplies infinitely many good rational primes $\ell_0$
in its conjugacy class. Choose a place realizing this representative,
and exclude two, all bad primes, the support of $D_wh_*$ and any
other finite set. Thus, with the torsion basis transported by this
place,
\[
 \operatorname{Fr}_{\ell_0}\equiv C\pmod{2^M},\qquad
 (D_w/\ell_0)=1,\qquad(h_*/\ell_0)=-1.
\]
The determinant also gives $\ell_0\equiv-1\pmod{2^M}$.
In particular $\operatorname{Fr}_{\ell_0}^2$ acts trivially on
$E[2^M]$.

Here is the finite-field map without a normalization ambiguity.
Put $n=2^M$ and write $F$ for Frobenius on the reduction of $E$.
For $P\in E(\F_{\ell_0^2})$, choose $R$ with $nR=P$ and set
\[
                    \kappa_M(P)=F^2R-R\in E[n]=T/nT.
\]
Changing $R$ by $n$-torsion does not change this value.
The map is additive, with kernel $nE(\F_{\ell_0^2})$.
It is onto: $F^2-1$ is a nonzero separable isogeny, so for
$S\in E[n]$ one can solve $(F^2-1)R=S$ and take $P=nR$.
Also $\kappa_M(FP)=F\kappa_M(P)$, by choosing $FR$ as a
division point. Therefore $\lambda_M=b\kappa_M$ has
Frobenius sign $+1$.

Good reduction identifies the reduction of the global division
cocycle with this exact map. The chosen place makes the twist
identification defined over $\F_{\ell_0}$, since
$(D_w/\ell_0)=1$. Thus the fixed witness satisfies
\[
 \lambda_M:E(\F_{\ell_0^2})\longrightarrow\Z/2^M\Z,\qquad
 v_2\lambda_M(\overline Q)=c_w<M.                         \tag{O.2}
\]
Choosing increasing $M$ and excluding previous primes gives one
sequence of actual stages. Subsequent finite prime exclusions
do not alter $c_w$. New field prescriptions require a compatible
common lift and, if added after this choice, are not asserted to
preserve this particular constant.
\end{proof}

% Completion audit replacement for baseline lines 3593--3748.
\subsection{Local functions and the theta sum}
\begin{proposition}[The actual odd theta form and its recurrence]\label{ex:app-theta-form}
Fix an actual stage of Proposition~\ref{ex:app-detector} with $M\ge2$.
Use $m=|h_*|$, the same $L$, the Eichler order of level $N$ in the
quaternion algebra ramified at $\ell_0,\infty$, and the local
Schwartz functions specified below, including (O.3).
With $c_0=|\operatorname{GL}_4(\F_3)|=24\,261\,120$, the sum
$H_M$ in (O.4) is well defined modulo $2^M$ in the character ring.
It gives an actual holomorphic form
\[
 g_M=\sum_n a_M(n)q^n,\qquad a_M(n)=H_M(mn),
 \qquad g_M\in M_{3/2}(\Gamma_0(K),\psi),
\]
with level supported only at the specified coefficient primes.
The characters are exactly
$\psi=\eta_B(m/\cdot)$ and $\eta=\eta_B\chi_{h_*}$, so
$\psi/\eta=\chi_{-1}$. For every prime outside that support,
the recurrence (O.6) holds at every integer index, with nonintegral
indices interpreted as zero. Changing the lifts of $P_M$ loses no
precision. A specified further local coefficient sieve retains this
conclusion outside the enlarged support.
\end{proposition}
\begin{proof}\leavevmode
In this appendix $M=M_\theta$; take $M\geq2$.  Retain the
same $m=|h_*|$, $L$, and detecting prime $\ell_0$ as in (O.2).
Since $\operatorname{Fr}_{\ell_0}^2=1$ on $E[2^M]$, evaluation at
arithmetic Frobenius squared identifies the finite-field Kummer
quotient with $T_2E/2^M T_2E$.  Equivariance and $bC=b$ imply
$\lambda_M(\operatorname{Fr}_{\ell_0} P)=\lambda_M(P)$.
The determinant gives $\ell_0\equiv-1\pmod4$; hence
$(m/\ell_0)=(-1/\ell_0)(h_*/\ell_0)=1$.
\paragraph{The finite-field evaluation map.}
Write $k_0=\mathbb F_{\ell_0^2}$ and let $F$ denote its arithmetic
Frobenius. For $P\in E(k_0)$, choose $Q_P\in E(\overline{k_0})$
with $2^M Q_P=P$. Then
\[
 \kappa_M(P)=FQ_P-Q_P\in E[2^M]
\]
is independent of the choice of $Q_P$, since $F$ is the identity
on $E[2^M]$. This is an additive map. It kills $2^ME(k_0)$,
and its kernel is exactly that subgroup: a zero value means that
$Q_P$ itself lies in $E(k_0)$. The finite group
$E(k_0)/2^ME(k_0)$ has order $|E(k_0)[2^M]|=2^{2M}$, as does
$E[2^M]$, so the induced injection is an isomorphism.
Composing $\kappa_M$ with the fixed row $b$ modulo $2^M$ gives
$\lambda_M$. The arithmetic Frobenius over $\mathbb F_{\ell_0}$
commutes with $F$ and acts on the torsion by $C$. Hence
$\kappa_M(\operatorname{Fr}_{\ell_0}P)=C\kappa_M(P)$ and
$bC=b$ proves the asserted Frobenius sign.
For the fixed witness, specialization of a compatible division
point identifies this evaluation with the cocycle value at
$\sigma^2$, whose $b$-value is $2^{a+1}$. Thus the visibility
valuation is fixed as the precision grows.

By \cite[Main Theorem 14.6.1]{voight-quaternion}, there is a quaternion
algebra $B/\mathbb Q$ ramified exactly at $\ell_0$ and infinity.
Let $R$ be an Eichler
order of level $N$.  Supersingular pairs are
$X=B^\times\backslash\widehat B^\times/\widehat R^\times$;
their automorphism groups are $R_z^\times$ for
$R_z=B\cap z\widehat Rz^{-1}$.  This description allows every
$N$ prime to $\ell_0$, including powers of two.  One may choose each
underlying supersingular curve over $\mathbb F_{\ell_0^2}$ with scalar
Frobenius; its cyclic subgroups of order $N$ are then stable.  Thus
all the corresponding coarse moduli points are defined over that
field.  The geometric and adelic descriptions are those of
\cite[Sections 1.1--1.2 and 2.5]{jetchevkane}.

Define $P_M(z)=\lambda_M(L\pi_{\ell_0}(z))$.  Here the specialization is
legitimate because $\ell_0\nmid N$: the modular curve has good
reduction and the map to the good-reduction elliptic curve extends
over the local base.  The good Hecke correspondences also extend.
Writing $\pi$ as the Abel--Jacobi map based at infinity followed by
the elliptic quotient gives, for $p\nmid N\ell_0$,
\[
 \sum_{z'\text{ a }p\text{-neighbor of }z}\pi_{\ell_0}(z')
       =a_p(E)\pi_{\ell_0}(z).
\]
Indeed the quotient intertwines $T_p$ with $[a_p(E)]$ and
$T_p(\infty)=(p+1)\infty$.  For an Atkin--Lehner involution
$W_{p^e}$ the analogous identity has eigenvalue $w_{p^e}$ and the
constant $\pi(W_{p^e}\infty)$, which $L$ kills.  Frobenius acts on
the supersingular set as the right division-uniformizer operation.
Consequently $P_M$ has the asserted good Hecke and Atkin--Lehner
eigenvalues modulo $2^M$, and uniformizer eigenvalue $+1$.
Choose arbitrary integral lifts of this finite function.

\paragraph{The finite Schwartz function.}
All local functions below are zero off their specified supports.
Set $V=B^0$ and $Q_B(x)=\operatorname{Nrd}(x)=-x^2$.
For $p\nmid2Nm\ell_0$ use $1_{R_p^0}$.
For odd $p^e\Vert N$, use the split integral unit support.  The
cyclic level quotient selects one of the two eigensummands and its
root $r\in\mathbb Z_p^\times$.  The projectors
$(x+r)/(2r)$ and $(r-x)/(2r)$ are integral, so reduction and Hensel
lifting show that the selected root modulo any fixed power of $p$
is locally constant.  Weight it by a unit character $\alpha_p(r)$
of finite $2$-power order with
\[
 \alpha_p(-1)=w_{p^e}\chi_{h_*}(p)^e. \tag{O.3}
\]
Such a character exists: at an odd prime the cyclic $2$-primary
part of $\mathbb F_p^\times$ admits a character taking $-1$ to
$-1$, and the trivial character supplies the other parity.
Conjugation by $R_p^\times$ preserves the cyclic quotient and its
selected summand.  Multiplication of $x$ by a unit $u$ preserves
the summands and replaces $r$ by $ur$, giving multiplier
$\alpha_p(u)$.

At two require $e_x=(1+x)/2\in R_2$ and
$x^2=r^2\in\mathbb Z_2^{\times2}$.  The eigenvalues
$a=(1+r)/2$ and $a'=(1-r)/2$ of $e_x$ differ by the unit $r$.
Its projectors $(e_x-a')/r$ and $(a-e_x)/r$ are integral and remain
distinct modulo two.  If $2^e\Vert N$, the cyclic quotient is a
free rank-one $\mathbb Z/2^e$-module.  Of these complementary
projectors exactly one acts as the identity on it, since the only
idempotents of $\mathbb Z/2^e$ are zero and one.  This defines the
selected root without a choice of its parity.  For every odd $u$,
\[
 e_{ux}=u e_x+(1-u)/2.
\]
This affine transformation preserves the two projectors
individually, so the selected root changes to $ur$, even when
the two eigenvalue parities switch.  The same root character rule
therefore applies at two.  A character with either prescribed
parity exists, using the trivial character or the character modulo
four.  If $2\nmid N$ there is no root character at two.

For $p\mid m$, require $x\in M_2(\mathbb Z_p)^0$ and $v_pQ_B(x)=1$.
The reduction of $x$ is a nonzero nilpotent: zero reduction would
force $p^2\mid Q_B(x)$.  All nonzero nilpotents are conjugate to
$X_0=\left(\begin{smallmatrix}0&1\\0&0\end{smallmatrix}\right)$.
The centralizer of $X_0$ in $\operatorname{GL}_2(\mathbb F_p)$ is
$\{\left(\begin{smallmatrix}a&b\\0&a\end{smallmatrix}\right):a\ne0\}$
and has square determinant.  Thus
\[
 \epsilon_p(gX_0g^{-1})=(\det g/p)
\]
is well-defined.  Conjugation by $u$ multiplies it by $(\det u/p)$,
and scalar multiplication by $u\in\mathbb Z_p^\times$ multiplies
it by $(u/p)$, as follows from conjugation of $X_0$ by
$\operatorname{diag}(u,1)$.

At $\ell_0$ require $x\in R_{\ell_0}^0$ with nonzero reduction
$\bar x=t\beta$ in the trace-zero line of $\mathbb F_{\ell_0^2}$,
and use $(t/\ell_0)$.  Order units act trivially on the residue
field.  A division uniformizer acts by Frobenius, which sends
$\beta$ to $-\beta$, giving multiplier $-1$.
Unit scaling gives $(u/\ell_0)$.  Moreover $\beta^2$ is a nonsquare
in $\mathbb F_{\ell_0}$: otherwise $\beta=\pm a$ for
$a\in\mathbb F_{\ell_0}$, contrary to its nonzero trace-zero condition.
Both $-1$ and $\beta^2$ are nonsquares, so $Q_B(x)=-t^2\beta^2$
is a nonzero square.  Since $(m/\ell_0)=1$, extracted indices have
$(D/\ell_0)=1$.

These supports are compact and open, and all weights are locally
constant.  Additional fixed unit-squareclass conditions multiply
them by locally constant characteristic functions.  They are
preserved by $x\mapsto ux$, which multiplies the norm by $u^2$.
Let $\mathcal O$ be the valuation ring at a place above two of the
field generated by the fixed root-character values at primes
dividing $N$.  All other weights are quadratic, so this field is
independent of $M$ and $\ell_0$.
With $f=\prod f_p$, the preceding calculations give
\[
 f(u^{-1}xu)=\chi_{h_*}(\operatorname{Nrd}u)^{-1}f(x)
       \quad (u\in\widehat R^\times),
 \qquad f(vx)=\alpha(v)f(x)\quad(v\in\widehat{\mathbb Z}^{\times}).
\]
The ramified local quadratic character on units at $p\mid m$ is
$(\cdot/p)$; all other order-unit character factors are trivial.
The scaling character has
\[
 \alpha(-1)=
 \Bigl(\prod_{p^e\Vert N}w_{p^e}\chi_{h_*}(p)^e\Bigr)
             \Bigl(\prod_{p\mid m}(-1/p)\Bigr)(-1)=1.
\]
For the last equality, $m\equiv3\pmod4$ and the first product is
$w(E^{h_*})=+1$.  More explicitly,
$w(E)=-\prod w_{p^e}$ and, since $(h_*,N)=1$,
$w(E^{h_*})=w(E)\chi_{h_*}(-N)$.
The latter twist formula also follows by multiplying the local
unramified-twist factors $\chi_{h_*}(p)^{v_pN}$ at $p\mid N$
and the factors $(-1/p)$ at the newly ramified good primes.

Let $\eta_B$ denote the classical character associated with $\alpha$
in the inverse-idele convention.  The formula
$\eta_B(d)=\prod_p\alpha_p(d)$ for positive integers $d$ prime to
the support is verified in the transformation calculation below.

Set $c_0=|\operatorname{GL}_4(\mathbb F_3)|=24\,261\,120$,
so $v_2(c_0)=9$, and define
\[
 H_M(n)=c_0\sum_{z\in X}
 \frac{P_M(z)\chi_{h_*}(\operatorname{Nrd}z)}{|R_z^\times|}
       \sum_{x\in V(\mathbb Q),\ Q_B(x)=n}f(z^{-1}xz). \tag{O.4}
\]
Left multiplication embeds each $R_z^\times$ into
$\operatorname{GL}_4(\mathbb Z)$; reduction modulo three is
injective on finite subgroups.  Indeed, a prime-order element
$1+3^aA$ in its kernel, with $a\geq1$ and $A$ nonzero modulo three,
has nonzero first term in its prime-th power modulo $3^{a+1}$,
or modulo $3^{a+2}$ when its order is three, a contradiction.
Thus $|R_z^\times|\mid c_0$.
The right-conjugation identity proves independence of the right
representative.  On the left, rational reduced norms are positive
and the global quadratic idele character is trivial on them.
The transformation law below shows that each summand is an actual
holomorphic theta series with the same level and character;
$c_0/|R_z^\times|$ is an integer.
Consequently the coefficients lie in the fixed character ring
$\mathcal O$.  Changing any lift of $P_M$ by $2^M$ changes every
coefficient by an element of $2^M\mathcal O$.  In particular
there is no loss of precision from an automorphism denominator.

\subsection{Transformation law and recurrence}
Use the standard global additive character, self-dual measures,
and the Gaussian $e^{-2\pi Q_B}$ at infinity.  Weil's operator
formulas \cite[I.13]{weil} give multiplication by $\mathbf e_p(bQ_B)$ for
$n(b)$, scaled pullback $f_p(x)\mapsto\gamma_{Q_B}(u)f_p(ux)$ for
$m(u)$ with $u$ a unit, and a constant times Fourier transform
for the Weyl element.  The upper unipotents $n(\mathbb Z_p)$ fix
$f_p$, because $Q_B$ is integral on its support.  If $L'_p$ is any
lattice of translations fixing $f_p$, its Fourier transform has
support in $(L'_p)^\vee$.  Choose $A_p$ with
$p^{A_p}Q_B((L'_p)^\vee)\subseteq\mathbb Z_p$.  Fourier conjugation
then proves that $n^-(p^{A_p}\mathbb Z_p)$ fixes $f_p$.
This includes $p=2$. At two take $A_2\ge2$ and use the comparison
with the ordinary theta multiplier given below; no splitting of
the metaplectic cover over all diagonal units is asserted.
For clarity, the local normalizations and the lower-unipotent
argument are as follows. Take
$\mathbf e_\infty(t)=e^{2\pi it}$ and
$\mathbf e_p(t)=e^{-2\pi i\{t\}_p}$, so that
$\prod_v\mathbf e_v(r)=1$ for $r\in\mathbb Q$.
For $Q(x)=x^{\mathsf T}Ax$ put
$B_Q(x,y)=Q(x+y)-Q(x)-Q(y)=2x^{\mathsf T}Ay$ and set
\[
 \mathcal F_Q\phi(y)=\int\phi(x)\mathbf e_p(B_Q(x,y))\,dx_Q,
 \qquad dx_Q=|\det(2A)|_p^{1/2}\,dx.
\]
Here $dx$ is the coordinate product measure with
$\operatorname{vol}(\mathbb Z_p^3)=1$.
At the real place use the same formula
$dx_Q=|\det(2A)|_\infty^{1/2}dx$ with coordinate Lebesgue measure.
Thus $\mathcal F_Q^2\phi(x)=\phi(-x)$.
If $\phi(x+\lambda)=\phi(x)$ for $\lambda\in L'$, translation
in this integral shows that its Fourier transform vanishes unless
$\mathbf e_p(B_Q(\lambda,y))=1$ for every $\lambda\in L'$.
This is precisely the asserted support in $(L')^\vee$.
Since this dual lattice is compact, an integer $A_p$ with
$p^{A_p}Q((L')^\vee)\subset\mathbb Z_p$ exists. The operator
for $n^-(c)$ is a Fourier conjugate of multiplication by
$\mathbf e_p(-cQ)$; consequently it fixes $\phi$ whenever
$c\in p^{A_p}\mathbb Z_p$. A scalar used to normalize the Weyl
operator cancels with its inverse in this conjugation.

All the specified local functions have such a translation lattice.
At an odd level prime the unit-root condition and its character
depend on a finite residue quotient. At two the square-unit
condition is a congruence modulo eight; the integral projectors
selecting the cyclic quotient, and the finite root character,
again depend on a finite residue quotient. At a prime dividing
$m$, valuation one and the nonzero-nilpotent sign depend on finite
residue data. The division-prime function is already a residue
function. Each support is contained in a lattice and is a finite
union of cosets of a sufficiently small lattice. This proves the
Schwartz--Bruhat and translation assertions for the actual weights.

For the reference $Q_0=\sum X_i^2$, the dyadic dual of
$\mathbb Z_2^3$ is $\tfrac12\mathbb Z_2^3$ and
$\operatorname{vol}_{dx_{Q_0}}(\mathbb Z_2^3)=2^{-3/2}$.
In particular its Fourier transform is
$2^{-3/2}1_{\frac12\mathbb Z_2^3}$, not $1_{\mathbb Z_2^3}$.
The lower-unipotent conclusion nevertheless holds at depth two,
because $4Q_0(\frac12\mathbb Z_2^3)\subset\mathbb Z_2$.
These positive measure factors affect neither the diagonal
eigenvalue quotient nor the counting coefficients of the theta sum.


The diagonal factor can be computed directly.  For
$Q_B=\sum_{i=1}^3 a_iX_i^2$, the one-dimensional Weil-index identity
\[
 \frac{\gamma(ua_i)}{\gamma(a_i)}
       =\frac{\gamma(u)}{\gamma(1)}(u,a_i)_p
\]
and multiplication under orthogonal sums give, relative to the
ordinary theta multiplier cubed, the factor
$(u,a_1a_2a_3)_p$.  These identities are \cite[II.25 and II.28]{weil};
the latter obtains the Hilbert symbol from the four-dimensional
quaternion norm form.  For $B=(-a,-b)$ the norm on $B^0$ is
$aX^2+bY^2+abZ^2$, whose determinant is $(ab)^2$.
There is therefore no additional determinant character for $Q_B$.
\paragraph{The normalized diagonal factor.}
The Weil-index identity used here can be derived without choosing
Gauss-sum signs. The norm form of the quaternion algebra $(u,a)$
is $\langle1,-u,-a,ua\rangle$. By multiplicativity under orthogonal
sums and $\gamma(-q)=\gamma(q)^{-1}$, its index is
\[
 \frac{\gamma(1)\gamma(ua)}{\gamma(u)\gamma(a)}.
\]
For a split quaternion algebra the norm form is hyperbolic and its
index is one; for the division algebra its index is minus one by
\cite[II.28, Proposition 4]{weil}. These are exactly the two values
of $(u,a)_p$. Solving this equality gives
\[
 \frac{\gamma(ua)}{\gamma(a)}
   =\frac{\gamma(u)}{\gamma(1)}(u,a)_p.
\]
Applying it to three diagonal entries gives the ordinary
one-dimensional theta factor cubed times $(u,\det Q_B)_p$.
This derivation includes $p=2$. A change of diagonalization changes
the determinant only by a square, so it changes no symbol.
For $Q_B/m$ the determinant becomes $m^{-3}\det Q_B$, whose
squareclass is $m$. Thus the extra factor is precisely $(u,m)_p$.
This establishes the determinant-character contribution in (O.5)
with the same additive character as the original theta multiplier.

\paragraph{An operator construction of the common double cover.}
Here is a direct justification that the same lifts may be used in
the comparison. It also fixes the signs of the diagonal operators.
Work over $k=\mathbb Q_p$, including $p=2$, or $k=\mathbb R$,
with the additive character just specified. For a nondegenerate
ternary $Q$ write $\gamma_Q=\gamma(Q)$ and define
\[
 \begin{split}
 U_Q(b)\phi(x)&=\mathbf e_k(bQ(x))\phi(x),\\
 D(a)\phi(x)&=|a|_k^{3/2}\phi(ax),\\
 W_Q&=\gamma_Q\mathcal F_Q,\qquad
 M_Q(a)=r_Q(a)D(a),\qquad
 r_Q(a)=\frac{\gamma(aQ)}{\gamma(Q)}.
 \end{split}
\]
The Weyl matrix in this convention is
$w=\left(\begin{smallmatrix}0&1\\-1&0\end{smallmatrix}\right)$.
The opposite Fourier-sign convention changes both compared
representations together.

The Fourier formula of \cite[I.14, Theorem~2]{weil}, or its
completion-of-the-square calculation, says for $b\ne0$ that
\[
 \int_{k^3}
 \mathbf e_k\bigl(bQ(z)+B_Q(z,v)\bigr)\,dz_Q
   =|b|_k^{-3/2}\gamma(bQ)\mathbf e_k(-b^{-1}Q(v)).
\]
The integral has the oscillatory meaning of that theorem.
Applying it to a Schwartz function gives the operator identity
\[
 \mathcal F_Q U_Q(b)\mathcal F_Q
 =\gamma(bQ)U_Q(-b^{-1})D(b^{-1})
                   \mathcal F_Q^{-1}U_Q(-b^{-1}).
\]
Thus no choice of an individual square root of a Gauss sum is
needed. The index identities proved above also give
\[
 \gamma(Q)^2=(\det Q,-1)_k\gamma(1)^6,
 \qquad \gamma(1)^4=(-1,-1)_k,
 \qquad r_Q(a)r_Q(b)=(a,b)_k r_Q(ab).
\]
For the middle equality apply the quaternion-norm formula to
$\langle1,1,1,1\rangle$; the first follows by applying the same
formula to $\langle1,-a_i,1,-a_i\rangle$ for each diagonal
entry $a_i$ of $Q$.

Consequently the displayed operators satisfy
\[
 \begin{gathered}
 U_Q(b)U_Q(c)=U_Q(b+c),\qquad
 M_Q(a)U_Q(b)M_Q(a)^{-1}=U_Q(a^2b),\\
 M_Q(a)M_Q(b)=(a,b)_kM_Q(ab),\qquad
 W_QM_Q(a)W_Q^{-1}=M_Q(a^{-1}),\\
 W_Q^2=(-1,-1)_kM_Q(-1),\\
 W_QU_Q(b)W_Q
 =(-b,-1)_kU_Q(-b^{-1})M_Q(-b^{-1})W_QU_Q(-b^{-1})
 \quad(b\ne0).
 \end{gathered}
\]
For the last equality, the preceding Fourier identity leaves
the scalar $\gamma(Q)^2\gamma(bQ)^2$; the index-square formula
reduces it to $(-b,-1)_k$. Every scalar in these relations is
independent of the ternary form. Interpreting a scalar $-1$ as
the central element $\varepsilon$ gives one common double
cover of $\operatorname{SL}_2(k)$. Indeed these relations reduce
every word to one of
$\varepsilon^eU(x)M(a)$ or
$\varepsilon^eU(x)M(a)WU(y)$, with $e\in\{0,1\}$.
The ordinary Bruhat decomposition makes their projections
unique, apart from $\varepsilon^e$; the operators send
$\varepsilon$ to $-\operatorname{id}$, so the central kernel
has exactly two elements. This realizes the rank-three pullback
of Weil's metaplectic operators, rather than assuming a section
over the diagonal units.

For $Q_0=\langle1,1,1\rangle$ these operators are the tensor
products of the three ordinary one-dimensional theta operators.
For either $Q$ or $Q_0$, changing a lift multiplies its operator
by the same $-1$. Hence the quotient of their eigenvalues is
independent of the lift. At two the relation
$M_Q(-1)^2=(-1,-1)_2\operatorname{id}=-\operatorname{id}$
explicitly retains the order-four lift of $-I$.

On a unit $a$ the quotient of the two diagonal index factors is
\[
 \frac{r_Q(a)}{r_{Q_0}(a)}=(a,\det Q)_k.
\]
Together with $f_p(ax)=\alpha_p(a)f_p(x)$, this is exactly
the eigenvalue quotient $\alpha_p(a)(a,\det Q)_p$ used below.
There is no remaining eighth-root factor at two.

\paragraph{The metaplectic comparison at two and the extracted character.}
We specify how the preceding index calculation fixes the multiplier,
including its dyadic part. Let
\[
 Q_0=X_1^2+X_2^2+X_3^2,\qquad
 \phi_{0,p}=1_{\mathbb Z_p^3},\qquad
 \theta(\tau)=\sum_{n\in\mathbb Z}e^{2\pi i n^2\tau}.
\]
Use the same additive character and metaplectic lifts for the
rank-three Weil representations of $Q_B$ and $Q_0$. The latter
theta series is literally $\theta(\tau)^3$; this fixes the meaning
of the ordinary theta multiplier cubed, including its eighth roots
of unity. For either representation the nontrivial central element
of the metaplectic cover acts by $-1$. Consequently the ratio of
two eigenvalues computed using the same lift is independent of
that lift. In particular a central factor in multiplying the
three lifts in the displayed triangular factorization cancels
in this ratio. We compare eigenvalues, and do not choose a group
section of the cover over $\Gamma_0(2^{A_2})$.

Here is the comparison on the three factors. Both upper unipotents
act trivially on the indicated vectors. Choose $A_p$ large enough
that both lower unipotents do also. This is the Fourier-support
argument above for $f_p$; for $\phi_{0,p}$ it requires no positive
depth at an odd prime and depth two at two, since the bilinear dual
of $\mathbb Z_2^3$ for $Q_0$ is $\tfrac12\mathbb Z_2^3$ and
$4Q_0(\tfrac12\mathbb Z_2^3)\subset\mathbb Z_2$.
The Weyl scalar cancels in conjugating an upper unipotent to a
lower one, so it introduces no additional eigenvalue. On a
diagonal unit $u$, the scaling part acts by $\alpha_p(u)$ on
$f_p$ and by one on $\phi_{0,p}$. The quotient of the Weil-index
parts, by the preceding calculation, is
\[
 \frac{\gamma(uQ_B)/\gamma(Q_B)}
      {(\gamma(u)/\gamma(1))^3}
       =(u,\det Q_B)_p.                                      \tag{O.5a}
\]
If the inverse Weil-index convention is used for the diagonal
operator, both numerator and denominator are inverted; the
quotient is unchanged because the Hilbert symbol has order two.
Thus this computation does not leave an undetermined inverse
or a dyadic root of unity. These operator comparisons are precisely
the scaled pullback, Fourier conjugation and scalar relations of
\cite[I.13--I.14 and II.25--II.28]{weil}.

Here $\det Q$ means the determinant of the matrix in
$Q(x)=x^{\mathsf T}Ax$. The associated bilinear Gram matrix is
$2A$. If that convention is used instead, both $Q$ and $Q_0$
acquire the same factor $2^3$ in their determinants, which cancels
in the comparison. In particular no extra squareclass of two
is introduced at the dyadic place.

For a general ternary form $Q$ in this argument, at each prescribed
prime take $A_p\ge1$ and increase $A_p$
also to the conductor of $u\mapsto(u,\det Q)_p$. The eigenvalue
relative to the ordinary theta multiplier, on a matrix whose
upper-left entry is $a$, is then
\[
                    \alpha_p(a)(a,\det Q)_p.                  \tag{O.5b}
\]
Indeed the triangular factorization has just the two trivial
unipotent factors and this diagonal factor, and its lift ambiguity
cancels as explained above. This also verifies multiplication:
the upper-left entry of a product is congruent to the product
of the upper-left entries modulo $p^{A_p}$, so both local
characters in (O.5b) multiply. The argument therefore applies
to every matrix in the indicated congruence subgroup at a prescribed
prime, not only to a choice of generators. At an unramified prime
outside the support, the full spherical invariance established
below treats matrices with nonunit upper-left entry as well.

For $Q_B$ the determinant squareclass is one at every place,
including two. Thus (O.5b) has no determinant factor. For the
extracted form $Q_B'=Q_B/m$ it is $m^{-3}\equiv m$.
At an odd prime $p\mid m$ and a unit $u$, and at two, respectively,
the additional factors are explicitly
\[
 (u,m)_p=(u/p)^{v_p(m)},\qquad
 (u,m)_2=(-1)^{((u-1)/2)((m-1)/2)}.
\]
Here $m$ is odd and $m\equiv3\pmod4$, so the second factor is
$(-1)^{(u-1)/2}$. It is already trivial on $1+4\mathbb Z_2$.
At every odd prime outside $m$ the factor on units is one. There
is consequently no new level prime, and there is no omitted
dyadic character.

Finally take a positive integer $d$ prime to the complete level.
The finite operator used in the classical transformation is the
one at $\gamma^{-1}$, with upper-left entry $d$. The root-character
factor is therefore $\prod_p\alpha_p(d)=\eta_B(d)$.
Hilbert reciprocity gives
\[
 \prod_{p\mid 2m}(d,m)_p
   =\prod_{q\mid d}(m/q)^{v_q(d)}=(m/d).                       \tag{O.5c}
\]
Indeed the infinite symbol is one because $m,d>0$, all other
finite symbols are one, and at an odd $q\mid d$ the symbol is
$(m/q)^{v_q(d)}$. This proves the displayed product without
discarding its dyadic factor. Both sides extend to the same
even Dirichlet character. Since $h_*=-m$, division by the genus
character now gives
\[
 \frac{\eta_B(d)(m/d)}{\eta_B(d)\chi_{h_*}(d)}=(-1/d).
\]
Thus the characters in (O.5), with the stated inverse-idele
convention for $\eta_B$, are exactly the characters supplied
by these local Weil operators.

Take $A_p$ also at least the conductor exponent of $\alpha_p$.
For $\gamma=\left(\begin{smallmatrix}a&b\\c&d\end{smallmatrix}\right)$
with $c\in p^{A_p}\mathbb Z_p$ and $a$ a unit, use
\[
 \gamma=n^-(c/a)m(a)n(b/a).
\]
The scalar is the diagonal theta factor times $\alpha_p(a)$.
In passing to the classical transformation law one applies the
finite operator at $\gamma^{-1}$, whose upper-left entry is $d$.
Thus the classical Dirichlet character is
$\eta_B(d)=\prod_p\alpha_p(d)$, not its inverse.  Its value at
$-1$ is one, as required by the weight-$3/2$ theta multiplier.

There are no additional level primes hidden in this argument.
Outside the prescribed support, $p$ is odd, $B_p$ is split, and
$R_p^0=\mathfrak{sl}_2(\mathbb Z_p)$ is self-dual for the bilinear
form of $Q_B=-a^2-bc$.  Its determinant is a unit.  Hence the
Fourier transform of $1_{R_p^0}$ equals itself; the unramified
Weil constant is one, as also follows from \cite[II.27]{weil} by integrating
over this self-dual lattice.  Upper unipotents, diagonal units,
and the Weyl element fix it.  These generate
$\operatorname{SL}_2(\mathbb Z_p)$, so the exponent there is zero.
At the finitely many prescribed primes the preceding $A_p$ suffice.
Conjugation by any $z\in\widehat B^\times$ is an isometry of $Q_B$,
commutes with the Weil representation, and preserves the
self-dual measure.  The same level and character therefore work
for every summand indexed by $z$.

\cite[III, Theorem 6]{weil} applies to this Schwartz--Bruhat function on
$V(\mathbb A)$ and proves invariance of its rational theta sum.
The normalized archimedean expression is
$\sum_x f(x)e^{2\pi i\tau Q_B(x)}$: the factor $y^{3/4}$ in the
Weil action is removed in passing to the classical form.
It has holomorphic weight $3/2$.  At a rational cusp, the slash
operation replaces the finite Schwartz function by its finite
Weil transform.  Its support meets $V(\mathbb Q)$ in a finite
union of rational lattice cosets, with bounded denominators.
Positive definiteness gives nonnegative rational exponents of
bounded denominator and normal convergence for $\Im\tau>0$.
This proves holomorphy at every cusp.  Theorem 6 alone asserts
automorphic invariance; this last argument supplies the
holomorphy conclusion used here.

\paragraph{Extraction.}
The support at every $p\mid m$ has $v_pQ_B=1$.  All nonzero
indices are integral and divisible by $m$.  Repeat the preceding
construction for $Q_B'=Q_B/m$ with the same finite weights and the
Gaussian for $Q_B'$.  This quadratic form is integral on the
support.  Away from $m$ it differs by a unit, so the good-prime
self-duality argument remains valid.  Its determinant squareclass
is $m^{-3}\equiv m$, giving the extra Dirichlet character
$(m/\cdot)$.  Hence
\[
 g_M(\tau)=\sum_n H_M(mn)q^n\in
      M_{3/2}(\Gamma_0(K),\psi).
\]
Here $K$ has only the stated prime support, and
\[
 \psi=\eta_B(m/\cdot),\qquad
 \eta=\eta_B\chi_{h_*},\qquad
 \psi/\eta=\chi_{-1}. \tag{O.5}
\]
This directly constructs the extracted form without a
half-integral-weight extraction operator.
The inverse in this classical passage can also be checked before
any character notation is introduced. Put
$g_\tau=n(x)m(\sqrt y)$ for $\tau=x+iy$, so $g_\tau i=\tau$.
For $\gamma\in\operatorname{SL}_2(\mathbb Q)$ write
$\gamma g_\tau=g_{\gamma\tau}k_\infty$ with
$k_\infty\in\operatorname{SO}_2(\mathbb R)$. In the adelic group,
\[
 (g_{\gamma\tau},1_f)
   =\gamma\,(g_\tau,\gamma_f^{-1})(k_\infty^{-1},1_f).
\]
This equality also holds with consistent metaplectic lifts.
The global rational lifts are compatible for the two forms:
\cite[II.30, Proposition~5]{weil} gives
$\prod_v\gamma_v(Q)=\prod_v\gamma_v(aQ)=1$ for rational $Q,a$,
and hence the local normalizing factors in $W_Q$ and $M_Q(a)$
have product one. The product formula also gives
$\prod_v|\det(2A)|_v^{1/2}=1$, so the product of the stated
self-dual measures is the measure for which the rational lattice
has covolume one. Apply the rational invariance of
\cite[III.41, Theorem~6]{weil} to the displayed adelic identity.
The archimedean Gaussian for either positive ternary form has
the same weight $3/2$; therefore its factor from $k_\infty^{-1}$
is the same as that for $\theta^3$. The relative finite factor
is evaluated at $\gamma_f^{-1}$, whose upper-left entry is $d$.
Thus the relative classical character is
$\prod_p\alpha_p(d)(d,\det Q)_p$, over the prescribed primes.
This proves the inverse-idele convention directly. It compares
the transformation identities and does not divide by a value
of a theta function.

For the extracted form there is likewise no measure or degree
normalization in its coefficients. With exactly the same finite
weights, replace only the quadratic form and archimedean Gaussian
by $Q_B/m$ and $e^{-2\pi Q_B/m}$. The summand indexed by $x$ now
has exponent $Q_B(x)/m$. The original support makes $Q_B(x)$ an
integer with valuation one at every prime dividing $m$, so these
exponents are nonnegative integers. Coefficient $n$ is exactly
the original sum over $Q_B(x)=mn$, namely $H_M(mn)$.
The self-dual measure changes by $|m|_v^{-3/2}$ locally, but a
theta coefficient is a discrete sum and contains no Haar
measure; this change is used only in the Fourier operators.
Outside the original support it is a unit rescaling of the same
unimodular local form, and at primes dividing $m$ the
lower-unipotent depth can be increased within the existing
support. The determinant changes by $m^{-3}$, whose squareclass
is $m$, so (O.5c) supplies exactly $(m/\cdot)$.
In particular no new level prime, factor of $m$, or dyadic
root of unity is introduced by extraction.


\paragraph{All-index recurrence and its phases.}
Let $p$ avoid the full support and put $s=\eta_B(p)$.
Locally use $R_p=M_2(\mathbb Z_p)$ and $V_p=\mathfrak{sl}_2$.
The $p+1$ neighbors are indexed by lines of $\mathbb F_p^2$.
For integral $x$ with nonzero reduction and $Q_B(x)=n$, the number
of preserved lines is $1+(-n/p)$, including the nilpotent case.
If $x\in pR_p^0$, all $p+1$ lines are preserved; the excess
over the preceding formula is $p$.

A nonintegral $x$ in a neighbor lies in $p^{-1}R_p^0$.
Write $x=y/p$ with $y$ primitive and $Q_B(y)=p^2n$.
Then $\bar y$ is nonzero nilpotent.  In a basis with
$\bar y=X_0$, write
$y=\left(\begin{smallmatrix}a&b\\c&-a\end{smallmatrix}\right)$.
Here $p\mid a,c$, $b$ is a unit, and
$p^2\mid-a^2-bc$ forces $p^2\mid c$.  Exactly the neighbor
belonging to the kernel line makes $y/p$ integral.  The integral
and nonintegral parts therefore exhaust the sum with no genericity
assumption on $n$.

Scaling by $p$ at all other places multiplies the weight by $s$.
Thus the primitive nonintegral contribution is
$s^{-1}H_M(p^2n)-H_M(n)$, and the integral contribution is
$(1+(-n/p))H_M(n)+psH_M(n/p^2)$.  Their sum is
\[
 s^{-1}H_M(p^2n)+(-n/p)H_M(n)+psH_M(n/p^2).
\]
Here and below a coefficient at a nonintegral index is zero.
For completeness, if $N_{ij}$ counts neighbor kernels from
$i$ to $j$ and $w_i=|\operatorname{Aut}(i)|$, counting actual
isogenies and dualizing gives
$N_{ij}w_j=N_{ji}w_i$, or $N_{ij}/w_i=N_{ji}/w_j$.
This proves the self-adjointness with the mass $1/w_i$ appearing
in $H_M$.  A neighbor changes the norm character by
$\chi_{h_*}(p)$.  Moving the neighbor operator onto $P_M$ and
using its actual modular-parametrization identity consequently
gives
\[
 H_M(p^2n)+s(-n/p)H_M(n)+ps^2H_M(n/p^2)
     =s\chi_{h_*}(p)a_p(E)H_M(n)\pmod {2^M\mathcal O}.
\]
This computation explains both powers of $s$ even when $s$ is
not quadratic.

Put $a_M(n)=H_M(mn)$ and
$\eta=\eta_B\chi_{h_*}$.  Since
$(-mn/p)=(m/p)(-n/p)$ and $\eta(p)^2=s^2$, one obtains
\[
 a_M(p^2n)=
 \{\eta(p)a_p(E)-\psi(p)(-n/p)\}a_M(n)
        -p\eta(p)^2a_M(n/p^2)\pmod {2^M\mathcal O}. \tag{O.6}
\]
The proof covers $p\mid n$ and $p^2\mid n$ explicitly.
At zero the coefficients vanish by the local support, and negative
indices vanish by positive definiteness.
An added finite coprimality or unit-squareclass sieve is realized
by the corresponding local characteristic function for $Q_B'$;
the same modularity proof applies.  At primes outside the enlarged
support the three indices $p^2n,n,n/p^2$ have the same retain-or-zero
status.  Hence these sieves preserve the displayed recurrence.
\end{proof}

\subsection{The unaveraged point comparison}
\paragraph{The actual weighted coefficient comparison.}
\begin{proposition}[The unaveraged weighted theta comparison]
\label{ex:app-theta-comparison}
Put $m=|h_*|$ and $F=\mathbf Q(\sqrt{Dh_*})$, where $D$ is a
positive squarefree parameter in the fixed odd filter. Assume
$(D/\ell_0)=1$ at the chosen actual stage, so that $\ell_0$ is
inert in $F$. Keep the local Schwartz functions, the actual function
$P_M(z)=\lambda_M(L\pi_{\ell_0}(z))$, and the integer $c_0$
of (O.4). Then a $2$-power root of unity $\zeta_M(D)$ satisfies
\[
 H_M(mD)=c_0\,2^{\omega(N)}\zeta_M(D)
                   \lambda_M(\overline{LP_D})\pmod {2^M}.\tag{O.7}
\]
The root of unity belongs to the field generated by the fixed
characters at the primes of $N$. No class number is divided out.
\end{proposition}


\begin{proof}
The filter makes $Dh_*$ an odd fundamental discriminant, with
$F$ split at every prime dividing $2N$. Its unit group is
$\{\pm1\}$. The condition $Q_B(x)=mD$ is equivalent to
$x^2=Dh_*$. Together with $(1+x)/2\in R_{z,2}$ and the other
integrality conditions, this says exactly that
\[
 \mathcal O_F\longrightarrow R_z,
 \qquad \sqrt{Dh_*}\longmapsto x
\]
is an optimal embedding. Optimality follows because the source
order is already maximal. All nonzero local weights on such an
embedding are roots of unity. The valuation-one condition at each
prime of $h_*$ holds because $(D,h_*)=1$; at $\ell_0$ the residue
of $x$ is nonzero. The fixed unit filters hold by admissibility.

First classify the local embeddings. Away from $N\ell_0$ there
is one local orbit: an $\mathcal O_{F,p}$-stable lattice is locally
a fractional ideal, hence principal, and this also covers ramified
primes dividing $Dh_*$. At $p^e\Vert N$, the algebra is split.
The two integral idempotents decompose both lattices of the
Eichler lattice pair. Its quotient is cyclic of order $p^e$,
so its elementary exponents lie wholly in one summand. The two
possibilities give exactly two orientations. This applies at two
using the integral spectral projectors specified in the construction.
At $\ell_0$, the maximal order of the division algebra has residue
field $\mathbf F_{\ell_0^2}$. The two residue-field embeddings
of the unramified quadratic order give the two orientations.
Consequently there are $2^{\omega(N)+1}$ local orientation tuples.

Fix one tuple. Choose a rational embedding
$\iota:F\hookrightarrow B$ and write $x_0=\iota(\sqrt{Dh_*})$.
Such an embedding exists because $F$ is a field at the two ramified
places of $B$, namely $\ell_0$ and infinity
\cite[Proposition 14.6.7, $(v)\Rightarrow(iii)$]{voight-quaternion}.
Choose $z_0$ realizing
the selected integral local embeddings. The adelic representatives
with this tuple are exactly
\[
       \widehat F^{\times}z_0\widehat R^{\times}.
\]
Their quotient by the rational centralizer $F^{\times}$ is
\[
 F^{\times}\backslash\widehat F^{\times}/
      (\widehat F^{\times}\cap z_0\widehat R^{\times}z_0^{-1})
 =F^{\times}\backslash\widehat F^{\times}/\widehat{\mathcal O}_F^{\times}
 =\operatorname{Pic}(\mathcal O_F).
\]
Thus this is an actual class-group torsor, not an average or a
quotient of a class-group orbit.

For an ideal-class representative $a\in\widehat F^{\times}$,
the element $a$ centralizes $x_0$, and therefore
\[
 f((az_0)^{-1}x_0(az_0))=f(z_0^{-1}x_0z_0).
\]
On the other hand,
\[
 \chi_{h_*}(\operatorname{Nrd}(az_0))
 =\chi_{h_*}(\operatorname{N}_{F/\mathbf Q}a)
                    \chi_{h_*}(\operatorname{Nrd}z_0).
\]
The first factor is exactly the genus character in (O.1).
To see that it is well-defined on the class group, for $D>1$
the biquadratic field $\mathbf Q(\sqrt D,\sqrt{h_*})$ has
discriminant $D\,h_*\,(Dh_*)=(Dh_*)^2$. Its relative
discriminant over $F$ is therefore one. The character obtained
by restricting $\chi_{h_*}$ to $G_F$ is unramified and, under
class field theory, is $a\mapsto
\chi_{h_*}(\operatorname{N}_{F/\mathbf Q}a)$. For $D=1$ it
is the trivial character. Ideal representatives can be chosen
prime to $h_*$, giving exactly the elementary norm formula in
(O.1). Reversing the Artin convention has no effect because this
character is quadratic.

Here is the precise specialization of the reduction theorem, including
its orbit and orientation data. In the notation of
\cite[\S\S3.1--3.3]{cornutjetchev}, take
\[
 F_{\rm CJ}=\Q,\quad K_{\rm CJ}=F=\Q(\sqrt{Dh_*}),\quad
 B_{\rm CJ}=M_2(\Q),\quad S=\{\ell_0\}.
\]
The field $F$ is imaginary because $D>0$ and $h_*<0$.
It embeds into $M_2(\Q)$ by its regular representation. The source
algebra has no finite ramification, so H.1 holds and the integer
in H.2 is $1+0+1=2$. Moreover $\ell_0\nmid Dh_*N$ and
\[
  (Dh_*/\ell_0)=(D/\ell_0)(h_*/\ell_0)=1\cdot(-1)=-1
\]
by the present support condition and (O.2). Thus H.3 holds in its
inert, rather than ramified, case. The resulting target algebra is
the definite algebra $B$ ramified precisely at $\ell_0$ and infinity.

Take a source Eichler order $R^{\rm src}=R'\cap R''$ of level $N$.
In \cite[\S3.2]{cornutjetchev} the source subgroup is
$H=\widehat{R^{\rm src}}^{\times}$ and the target subgroup is
$H_S=\widehat R^{\times}$. Indeed, away from $\ell_0$ the local
orders are transported by the fixed algebra isomorphisms. At $\ell_0$
the source subgroup is $\operatorname{GL}_2(\Z_{\ell_0})$,
whose determinant is $\Z_{\ell_0}^{\times}$; its inverse image
under the target reduced norm is exactly the unit group of the
maximal division order. The target level is the prime-to-$S$ part
of $N$, which is $N$ itself. This construction allows arbitrary
prime-power exponents in $N$, including its exponent at two.

Both fine conductors are one. To check this assertion using the
definition in \cite[\S3.3]{cornutjetchev}, an embedding of the maximal
order $\mathcal O_F$ into $R'\cap R''$ embeds it into each maximal
order. The intersection of either maximal order with $F$ is an order
containing $\mathcal O_F$, hence is $\mathcal O_F$. Conversely,
if both intersections are $\mathcal O_F$, their intersection with
$R'\cap R''$ is $\mathcal O_F$. The same argument applies on the
target. Thus $c'=c''=c=c_S=(1)$ in that theorem. Its parity
condition at $\ell_0$ is $0-0\equiv v_{\ell_0}(N)=0\pmod2$.
The local lattice pair there has distance zero. Since $F_{\ell_0}$
is the unramified quadratic extension, its maximal order fixes a
unique vertex of the local tree. Consequently
\[
 N_{\ell_0}(0,0,0)=1,\qquad
 |\operatorname{Gal}(F[c]/F[c_S])|=1,
\]
which also follows directly from
\cite[Lemma~2.1(ii)]{cornutjetchev}. The multiplicity $\kappa$ in
\cite[Theorem~3.1]{cornutjetchev} is therefore exactly one.

It is essential to retain the sign fiber in that theorem. At $\ell_0$
the quotient
\[
 F_{\ell_0}^{\times}\backslash B_{\ell_0}^{\times}/R_{\ell_0}^{\times}
       \simeq\Z/2\Z
\]
is measured by the valuation of the reduced norm modulo two:
norms from $F_{\ell_0}^{\times}$ have even valuation, whereas a
division uniformizer has norm valuation one.
By \cite[\S3.6]{cornutjetchev} this sign is precisely the choice of
residue-field embedding of the unramified quadratic order, with
origin fixed by the chosen reference embedding. Theorem~3.1 gives
a bijection onto one such sign, not onto both. Away from $\ell_0$
its map on local orbit data is a bijection, so we may also restrict
to the chosen orientations at the primes of $N$. The resulting
map is an equivariant bijection of the source class-group torsor
onto the target torsor for that full orientation tuple. Corollary~3.2
gives the same bijection after fixing any supersingular pair in the
target.

Choose the reduction place compatibly with the place used in (O.2).
Sections~4.9--4.10 of \cite{cornutjetchev} identify the projection
of this adelic map with geometric reduction of the corresponding
enhanced elliptic curve on $X_0(N)$. Their local formula is the
identity under the chosen identifications away from $\ell_0$, and
at $\ell_0$ replaces a matrix by a division-uniformizer power
whose exponent is the valuation of its determinant. Thus it is
the same sign and the same reduction map as above. Choose $z_0$
to be the image of the actual reference Heegner point $x_D$ in
(O.1). Writing
\[
 \chi_F([a])=\chi_{h_*}(\operatorname{N}_{F/\Q}a),\qquad
 z_a=[az_0],
\]
equivariance gives the exact finite-sum identity
\[
 \sum_{[a]\in\operatorname{Pic}(\mathcal O_F)}
       \chi_F([a])P_M(z_a)
 =\lambda_M\!\left(
       \overline{L\sum_{[a]\in\operatorname{Pic}(\mathcal O_F)}
                         \chi_F([a])\pi(x_D^{[a]})}\right)
 =\lambda_M(\overline{LP_D})\quad\text{in }\Z/2^M\Z.
\]
Here the supersingular pairs and their parametrization values are
defined over $\mathbb F_{\ell_0^2}$, so every term is in the domain
of the same $\lambda_M$. Additivity, specialization, and multiplication
by the fixed integer $L$ commute with this finite sum. If the Artin
action is written with inverses, reindexing $[a]\mapsto[a]^{-1}$
leaves it unchanged since $\chi_F$ is quadratic. Changing the
reference point within its torsor multiplies the sum only by
$\chi_F$ of that reference class.

The weights are constant on each embedding orbit used in the
following count. Indeed a unit of the definite order $R_z$ has
positive integral reduced norm with integral inverse, so its norm
is one. The conjugation law for $f$ therefore gives weight one
for its action. Along the class-group torsor the Schwartz factor
is constant because $a$ centralizes $x_0$, and the remaining
norm factor is exactly $\chi_F([a])$. Hence the finite-sum identity
just proved applies with the single phase
$\chi_{h_*}(\operatorname{Nrd}z_0)f(z_0^{-1}x_0z_0)$.


For each fixed $z$, the stabilizer of an embedding in $R_z^\times$
is $\mathcal O_F^\times=\{\pm1\}$. Each embedding orbit thus
has $|R_z^\times|/2$ elements. The denominator $|R_z^\times|$
in (O.4) leaves the factor $1/2$ per embedding class. The
vectors $x$ and $-x$ are not identified by this stabilizer:
at $\ell_0$ their residues define the two different embeddings
of the residue field, whereas conjugation by an order unit acts
trivially on that field. They already lie in the two different
$\ell_0$ orientation tuples counted above. The
contribution of the reference orientation tuple is
\[
 \frac{c_0}{2}\,\zeta_M(D)
                         \lambda_M(\overline{LP_D}),\qquad
 \zeta_M(D)=\chi_{h_*}(\operatorname{Nrd}z_0)
                               f(z_0^{-1}x_0z_0),
\]
up to the harmless quadratic reference-class factor just described.
All nonzero factors in this displayed phase are fixed-character
values or signs. Thus it is a $2$-power root of unity; its field
does not grow with $D$, $M$, or $\ell_0$.

It remains to account for the other orientations; they are not
all in the image of the one sign component of the reduction map.
At $p^e\Vert N$, use the local Atkin--Lehner normalizer. It
reverses the selected root, so multiplies the Schwartz weight by
\[
              \alpha_p(-1)=w_{p^e}\chi_{h_*}(p)^e.
\]
It multiplies the norm-weighted parametrization value by the same
sign: its reduced norm has valuation $e$, and the actual
parametrization transforms by $w_{p^e}$ after $L$ kills the fixed
cuspidal translation. The two signs multiply to one. This
normalizer has trivial components at the other primes, so it
changes no other orientation or local weight.

At $\ell_0$, a division-algebra uniformizer exchanges the two
residue embeddings. It sends the trace-zero residue $\beta$ to
$-\beta$, multiplying its weight by $(-1/\ell_0)=-1$. Its
reduced norm has valuation one and hence contributes
$\chi_{h_*}(\ell_0)=-1$. Its action on supersingular pairs is
Frobenius. Since $\lambda_M$ has Frobenius sign $+1$ and the
parametrization is defined over $\mathbf Q$, this introduces no
further sign. Again the two factors multiply to one.
The uniformizer operation used here may be described with arithmetic
or geometric Frobenius: they have the same action on the supersingular
pairs and their parametrization values over $\mathbb F_{\ell_0^2}$.
The identity $\lambda_M(FP)=\lambda_M(P)$ from (O.2) consequently
applies to the very operation that exchanges the residue embeddings.
Its norm has valuation one, and $\chi_{h_*}$ is unramified at
$\ell_0$, so its norm-character value is exactly
$\chi_{h_*}(\ell_0)=-1$, independently of the uniformizer's unit part.
Likewise, at $p^e\Vert N$, the norm of the Atkin--Lehner normalizer
has valuation $e$ and $\chi_{h_*}$ is unramified at $p$. Its
norm-character value is exactly $\chi_{h_*}(p)^e$. These observations
fix the norm signs in both orientation changes.


It follows that every one of the $2^{\omega(N)+1}$ tuples has
the same contribution as the reference tuple. Multiplying by
their number gives precisely $c_0 2^{\omega(N)}$, as asserted.
This counts individual embeddings with their stabilizer; a
conjugate-pair formulation of the geometric correspondence must
not be used to insert or remove another factor of two.

Finally, $c_0/|R_z^\times|$ is an integer for every $z$. Thus a
change of the chosen integer lift of $P_M(z)$ by $2^M$ changes
the actual coefficient by an element of $2^M\mathcal O$.
Every preceding equality is consequently valid at the original
retained precision. No denominator depending on the class number,
the discriminant, or the stage has entered.
\paragraph{Precision in the orientation count.}
The factor $1/2$ in the orbit count does not require dividing a
congruence modulo $2^M$ by two. Indeed, $-1\in R_z^\times$ implies
$2\mid |R_z^\times|$, and the chosen integer $c_0$ is divisible by
$|R_z^\times|$. Hence $c_0/2$ is an integer. Choose lifts of all
$P_M(z)$ before summing. Each orbit of embeddings contributes
$(c_0/|R_z^\times|)(|R_z^\times|/2)P_M(z)=(c_0/2)P_M(z)$,
with its prescribed unit weight. Replacing any lift by a multiple
of $2^M$ changes this expression by a multiple of $2^M$.
The comparison with the reduced point may therefore be made at
precision $M$, separately for each orientation. Summing the
$2^{\omega(N)+1}$ orientations gives the stated multiplier
$c_0 2^{\omega(N)}$ without any loss of precision.
In particular, write $c_\theta=9+\omega(N)$, the valuation of
the fixed integer $c_0 2^{\omega(N)}$. At the fixed witness $D_w$,
(O.2) and (O.7) imply
\[
 v_2a_M(D_w)=c_\theta+c_w
       \qquad\text{whenever }M>c_\theta+c_w.
\]
The phase is a unit and the congruence error has strictly larger
valuation, so this is an equality, not merely a lower bound.
Discarding the finitely many earlier stages is sufficient. Thus the
visibility used in choosing the odd minimum retains a fixed witness
and a fixed valuation; neither the class number nor the changing
prime $\ell_0$ introduces a new loss of precision.

\end{proof}

\begin{proposition}[Visibility and the native genus index]\label{ex:app-visibility}
Use the actual odd form and unaveraged point of
Propositions~\ref{ex:app-theta-form} and~\ref{ex:app-theta-comparison}.
There is one constant $C_{\rm det}$, independent of the actual stage
and admissible $D$, for which (O.8) holds. Explicitly, for a known
simple-zero $D$,
$v_2a_M(D)\ge\min\{M,j_{\rm gen}(D)-C_{\rm det}\}$.
For every admissible $D$, the strict inequality $v_2a_M(D)<M$
implies $\operatorname{ord}_{s=1}L(E^D,s)=1$ and
$j_{\rm gen}(D)\le v_2a_M(D)+C_{\rm det}$.
If an admissible odd-sign $D$ is not a simple zero, its coefficient
vanishes modulo $2^M$. No genus index of a torsion point is defined
or used in this assertion.
\end{proposition}
\begin{proof}\leavevmode
For a known simple-zero $D$, write $P_D=m_DG+T$ in its native
rank-one free lattice. The same $L$ kills $T$.
An integral finite-field Kummer map cannot decrease
$v_2(m_D)$, so (O.7) gives the lower bound by the native index.
The rational free line has index at most two in that native
lattice, by conjugation and norm. Absorb this and the fixed
initial convention in $C_{\rm det}$. The first inequality below
holds for a known simple-zero $D$; the visibility implication holds
for every admissible $D$:
\[
 \begin{split}
 v_2a_M(D)&\ge\min\{M,j_{\rm gen}(D)-C_{\rm det}\},\\
 v_2a_M(D)<M&\ \Longrightarrow\
  \ord_{s=1}L(E^D,s)=1,\quad
       j_{\rm gen}(D)\le v_2a_M(D)+C_{\rm det}.
 \end{split}                                                    \tag{O.8}
\]
If an odd-sign $D$ is not a simple zero, the height formula gives a
torsion genus point, killed by $L$, and its coefficient is zero
modulo $2^M$. Combined with the first inequality on simple-zero
parameters, this gives the all-$D$ truncated bound used below,
without assigning a rank-one lattice index to a torsion point.
\end{proof}

\subsection{Integer divisions and actual trace operators}
\begin{proposition}[Arithmetic realization of the coefficient rules]\label{ex:app-coefficient-rules}
Use the rational full-two-torsion graph conventions of
Definitions~\ref{ex:sec7-typed-graph} and~\ref{ex:sec7-symbol}.
For an even source of Proposition~\ref{ex:app-even-phase}, use its
fixed minimum depth. For an odd source, first fix a finite operator
budget $|U|\le W$, take $R\ge2$, and choose an actual stage satisfying
the depth inequalities of Proposition~\ref{ex:app-odd-minimum}.
Then the integer divisions (O.9)--(O.11) are valid at every Fourier
index, with the stated truncation for the odd source.
The fresh coefficient-one reading is
$\overline{(t(\operatorname{Fr}_q)\theta H_U)[1]}=\mathcal A(Uq)$.
On all tested nonempty configurations, these actual symbols depend
only on their typed graph and satisfy the rules of
Lemma~\ref{ex:sec7-coefficient-rules}: a pair difference
contracts the endpoints and adds their types and every external
incidence; whole-prime omission uses the same source, phase, depth
and stage; and an isolated primary vertex of neutral type and
neutral external incidences makes the symbol zero.
Every contraction has nonempty output. Stationary character tests
are carried along in the same product rule. This statement asserts
neither values at untested odd-stage weights nor an identity at an
empty contraction output.
\end{proposition}
\begin{proof}\leavevmode
All new primes below avoid the full coefficient support $S$,
including every late sieve and the current $\ell_0$.
Here $S$ is the actual coefficient support.  The coefficients of
$g$ vanish at all indices not prime to $S$.
For the odd source this is forced at $2,N,h_*,\ell_0$ by the
local Schwartz support and the extraction at $|h_*|$.
At every additional coefficient-sieve prime it is imposed by an
exact zero-or-retain Fourier sieve, which commutes with the good
recurrence outside the enlarged support.  This is the vanishing
used in the condition $(m,S)=1$ in the unary series below.
Later stationary prime-selection tests may enrich the common
label alphabet without resieving the intrinsic odd source;
such tests do not enlarge this coefficient-support condition.
Use the minimum-depth inequalities proved in Appendix~\ref{app:endpoints}.
The general bookkeeping assigns weight one to a split prime of
$\Q(E[2])$ and one-half to a prime with transposition Frobenius, called
simple. In our actual full-two-torsion family every prime is split,
$w(U)=|U|$, and the depth $\delta$ is an integer. Thus only matching
parity, meaning $e_U=0$, is needed here; the displayed general formula
also records where the unused simple-prime term would enter.
For a finite set $U$, put
\[
 n_U=\lceil\delta+w(U)\rceil,\quad e_U=n_U-\delta-w(U),\qquad
 B_U=2^{-n_U}\prod_{p\in U}(U_p-\eta(p)V_p)g,
\]
where $U_p(\sum a(n)q^n)=\sum a(pn)q^n$ and $V_pg(\tau)=g(p\tau)$.
These are operations on the actual form $g$. We verify the level
at a new odd prime using the classical theta multiplier. Write the
old level as $4A$, with $p\nmid4A$, and put
$\beta=\left(\begin{smallmatrix}p&0\\0&1\end{smallmatrix}\right)$.
For $\gamma=\left(\begin{smallmatrix}a&b\\c&d\end{smallmatrix}\right)
\in\Gamma_0(4Ap)$ one has
\[
 \beta\gamma\beta^{-1}=
       \begin{pmatrix}a&pb\\c/p&d\end{pmatrix}\in\Gamma_0(4A).
\]
The square-root factor in the theta multiplier is unchanged on
evaluating the latter matrix at $p\tau$. The quotient of its
arithmetic factor by that of $\gamma$ is
$(c/p\,/d)/(c/d)=(p/d)$. Cubing the multiplier leaves this
quadratic factor unchanged. Hence $V_pg$ has level $4Ap$ and
character $\psi(p/\cdot)$.

At this level the normalized sum
\[
 U_{p^2}h(\tau)=\frac1{p^2}\sum_{r\bmod p^2}
                             h((\tau+r)/p^2)
\]
preserves the same weight and character. Indeed, put
$\alpha_r=\left(\begin{smallmatrix}1&r\\0&p^2\end{smallmatrix}\right)$.
For $\gamma\in\Gamma_0(4Ap)$ choose
$r'\equiv(b+rd)/(a+rc)\pmod{p^2}$. Then
\[
 \alpha_r\gamma=\gamma'_r\alpha_{r'},\qquad
 \gamma'_r=
 \begin{pmatrix}
  a+rc&(b+rd-(a+rc)r')/p^2\\
  p^2c&d-cr'
 \end{pmatrix}\in\Gamma_0(4Ap).
\]
The map $r\mapsto r'$ is a permutation, since $a+rc$ is a unit
modulo $p$ and the fractional linear map has determinant one.
Writing $d'=d-cr'$, its nebentypus factor is unchanged since
$d'\equiv d\pmod{4Ap}$. Its theta arithmetic factor is unchanged
as well: $d'\equiv d\pmod c$, $d'\equiv d\pmod4$, and
$(p^2c/d')=(c/d')=(c/d)$. These are the usual Kronecker-symbol
identities for the theta multiplier on $\Gamma_0(4)$.
Its square-root factor, evaluated at $(\tau+r')/p^2$, is
$(c\tau+d)^{1/2}$. Thus the sum transforms with the same
theta multiplier cubed and the same character. The formulas also
hold for the triangular matrices, and holomorphy at the cusps is
preserved by the finite rational translates.
On Fourier expansions $U_{p^2}V_p=U_p$. Consequently both terms
of $U_p-\eta(p)V_p$ have level exponent at most one at $p$ and
character $\psi(p/\cdot)$. Iteration introduces no higher exponent
at previously added primes.

The same computation applies to the unary corrections. Start with
the odd-character unary theta series at the old support, including
its fixed coprimality sieve. Its $J$-series is its transform by
$V_J$, for squarefree $J\subset U$. A prime $p\mid J$ therefore
has exponent at most one and the extra character $(p/\cdot)$;
a prime $p\nmid J$ is absent from its original level. Multiplication
by the ordinary theta series introduces no odd level prime. This
proves the precise level and oldform alternatives used in (O.12).


Write an index as $Jrm^2$ with $J\subset U$ and $r$ squarefree,
prime to $U$. Products of prime sets are understood here.
The originating squarefree parameter in the coefficient of $g$ is
$D=(U\setminus J)r$, not the squarefree part $Jr$ of the tested index.
Remove the phase $\eta(p)^s$ in (O.6) and put
\[
 b_{-1}=0,\quad b_0=1,\quad b_1=a_p-\kappa,\quad
 b_{s+1}=a_pb_s-pb_{s-1},\quad
 \kappa=(\psi/\eta)(p)(-D/p).
\]
For $p\in U\setminus J$ the factor $b_s-b_{s-1}$ is odd;
for $p\in J$ the factor $b_{s+1}-b_s$ is even.
Consequently, for the even source,
\[
 v_2B_U[Jrm^2]\ge w(r)
              +\tfrac12\#\{p\in J:p\text{ is simple}\}-e_U.      \tag{O.9}
\]
The right side is an integer on the allowed filter and is
nonnegative. In our full-two-torsion application all primes are
split and $w=\omega$.
For the odd source, fix an integer $R\ge1$ and a weight bound $W$
first. At one sufficiently advanced actual stage, simultaneously for
all $U$ with $w(U)\le W$ and every Fourier index, (O.9) holds with
its right side replaced by its minimum with $R$. Thus the untruncated
bound on a prescribed bounded range of squarefree weights follows by
choosing $R$ and then the stage. For the unary-removal and two-trace
tests below, choose $R\ge2$ after fixing the operator budget.
For a split prime $a_p\equiv p+1\pmod4$. The recurrence gives the
same residue on each divided shell,
\[
        (b_{s+1}-b_s)/2\equiv
     c_p(D):=\{a_p-(\psi/\eta)(p)(-D/p)-1\}/2\pmod2.              \tag{O.10}
\]

\paragraph{The two integral shell calculations.}
Here every prime is split in the rational two-torsion field, so
$a_p\equiv p+1\pmod4$. For $p\mid D$ one has $\kappa=0$.
Reducing the recurrence modulo two gives $b_{s+1}=b_{s-1}$,
with $b_0=1$ and $b_1=0$. Thus $b_s-b_{s-1}$ is odd for
all $s\ge0$, including $s=0$ with $b_{-1}=0$.
For $p\nmid D$, one has $\kappa\in\{1,-1\}$, so every $b_s$
is odd. Write $d_s=b_{s+1}-b_s$. The recurrence modulo four gives
$d_s\equiv p d_{s-1}\pmod4$ for $s\ge1$.
Each $d_s$ is even, and multiplication of an even residue modulo
four by an odd number fixes that residue. Consequently
\[
 d_s\equiv d_0=a_p-\kappa-1\pmod4\qquad(s\ge0).
\]
This proves (O.10) without dividing a residue-field identity by two.
For an index $Jrm^2$, each $p\in J$ therefore supplies an actual
factor two, whereas the factors for $p\in U\setminus J$ are integral
units modulo two. The squarefree input has
$\omega(D)=|U|-|J|+\omega(r)$. Its depth bound plus these $|J|$
factors gives
$\delta+\omega(D)+|J|=n_U+\omega(r)$ before division by
$2^{n_U}$. All removed character factors are units.
For the odd source the same argument is a congruence modulo
$2^M\mathcal O_{\rm coef}$ before that division. It therefore
retains precision $M-n_U$ afterward; the subsequent division
in (O.11) retains precision $M-n_U-1$.
Choose $i\in\{0,1\}$ with $\eta\chi_{-4}^i$ odd and let
\[
 \Theta_J=\sum_{\substack{m>0\\(m,S)=1}}
                      \eta(m)\chi_{-4}(m)^i m q^{Jm^2}.
\]
Poisson summation gives this genuine unary weight-$3/2$ form.
Its scalar character is $\eta\chi_{-4}^i(-J/\cdot)$ and every
$p\mid J$ enters the level at exponent one.
In matching parity the shell recurrence gives
$B_U[Jm^2]-B_U[J]\eta(m)\in2\mathcal O_{\rm coef}$.
Thus
\[
                 H_U=\frac{B_U-\sum_JB_U[J]\Theta_J}{2}          \tag{O.11}
\]
is integral; in opposite parity take $H_U=B_U$.
This is division by the integer two, not reduction modulo the
maximal ideal of a ramified coefficient ring. For odd-stage
forms, the coarse bound and bounded-weight stabilization proved
below pay every fixed denominator before choosing the stage.

Write $\mathfrak m_{\rm coef}$ for the maximal ideal of the chosen
coefficient valuation ring $\mathcal O_{\rm coef}$. In the following
trace argument a reduced expansion means coefficientwise reduction
modulo $\mathfrak m_{\rm coef}$. The preceding divisions by integer
powers of two have already been justified in $2^a\mathcal O_{\rm coef}$;
reduction modulo the maximal ideal alone would not justify them.
No saturation of the integral Fourier module is assumed, and no
identification of its abstract quotient with all reduced expansions
is needed.
Multiply (O.11) and the undivided corrections by
$\theta=1+2\sum_{m>0}q^{m^2}$. They are actual weight-two forms.
For a fixed stage and $U$, choose a finite complete coefficient-field
extension containing all constituents and retain their old
multiplicities and diamond translates in the finite-dimensional
characteristic-zero space. Use modular Galois representations on
cuspidal constituents and the two character summands on Eisenstein
constituents. Holomorphy alone does not exclude the latter.
Diamond translates remain integral:
their correction coefficients have the form
\[
 \tfrac12 B_U[J]\{\psi_U(u)-\eta(u)\chi_{-4}(u)^i(-J/u)\},
\]
whose character ratio is quadratic. Good Hecke operations preserve
integrality by
$(T_qf)[n]=f[qn]+q(\langle q\rangle f)[n/q]$.
Finitely many Fourier coordinates inject this space into a finite
coordinate space. The module generated by these translates is
therefore finite over the coefficient DVR, and closed.
Chebotarev and continuity make scalar trace and determinant
operators limits of good Hecke and diamond operations, so they
preserve this module. No integral eigenspace projector is used.

\paragraph{Reading an added vertex.}
For the actual full-two-torsion family, all primes are split,
$\delta\in\mathbb Z$, and $n_U=\delta+|U|$.  Put $F_U=\theta H_U$
and let $q\notin S\cup U$ be a fresh prime, where $S$ includes
every late sieve and the current moving theta prime.
Write $u=\prod_{p\in U}p$ and let $\mathcal A(D)$ denote the
normalized valuation symbol on the tested range: it is one exactly when
$v_2a(D)=\delta+\omega(D)$ on the filter. For an odd source, that
range has a fixed weight bound chosen before the valid actual stage,
as specified before (F.7). No finite-stage value at larger weights is
asserted; the compatible limit tables used later are a separate object.
Expanding the operators, every term containing a $V_p$ has zero
coefficient at $q$, and no unary index $Jm^2$ with $J\subset U$
equals $q$.  Therefore
\[
 B_U[q]=2^{-n_U}a(uq),\qquad
 H_U[q]=2^{-n_U-1}a(uq).
\]
At coefficient one the diamond term in the good Hecke formula
vanishes.  Since $F_U\equiv H_U\pmod{2\mathcal O_{\rm coef}}$,
\[
 \overline{(t(\operatorname{Fr}_q)F_U)[1]}
   =\overline{H_U[q]}
   =\overline{\,2^{-\delta-|U|-1}a(uq)\,}
   =\mathcal A(uq)\in\mathbb F_2.
\]
Here residues lie in the coefficient residue field, with
$\mathbb F_2$ its prime subfield.  Every permitted phase is a
two-power root of unity and reduces to one.  For an odd source,
fix the finite operator budget $|U|\le W$ and $R\ge2$ first;
the stage choice in (O.9) may be taken with
$M_\theta\ge\max(0,\delta+W)+R$.
Thus $M_\theta-n_U-1\ge1$, so the normalized finite-stage error
is still divisible by the integer two.  The same reading holds
with value zero off the fixed or moving coefficient support.

The diamond correction above is a sum of integral multiples of
$\theta\Theta_J$; its fresh coefficient $q$ is zero after reduction.
It introduces neither an extra phase nor a denominator in this
reading.  All readings and subsequent omissions use this same
$g,\eta,\psi,\delta$ and actual stage; only the set $U$ changes.

For any fixed Fourier index $n$, Chebotarev approximation lets one
test a finite product of trace operators using distinct fresh primes
$q_1,\ldots,q_a$ avoiding $n$ and the level. The good Hecke formula
then reads the coefficient at $nq_1\cdots q_a$ of the original vector.
Different indices may use different fresh primes; the conclusion is
an identity of coefficientwise reduced expansions.

The determinant identity uses the actual diamond formula, not merely
the reduction of its eigenvalues. For a good prime $r$,
$d(\operatorname{Fr}_r)=r\langle r\rangle$. The formula above shows
that $(d(\operatorname{Fr}_r)-1)F_U$ reduces to a sum of integral
multiples of the undivided corrections $\theta B_U[J]\Theta_J$:
the scalar $r\psi_U(r)-1$ lies in $\mathfrak m_{\rm coef}$.
A fresh trace kills every such correction. Indeed, to test at $n$
choose $q\nmid nJ$; then $nq$ cannot have squarefree part $J$, and
$\theta\equiv1\pmod{\mathfrak m_{\rm coef}}$. This remains true
after further fresh traces. Continuity and the closed integral
module therefore give, on the reduced trace orbit of $F_U$,
\[
 d(u)t(v)=t(v),\qquad
 t(u)t(v)=t(uv)+t(uv^{-1}),\qquad t(1)=0.
\]
The middle equality follows from the characteristic-zero matrix
identity
$\operatorname{tr}(A)\operatorname{tr}(B)
=\operatorname{tr}(AB)+\det(B)\operatorname{tr}(AB^{-1})$,
using the determinant identity just proved; the last follows from
$\operatorname{tr}(1)=2$.

Two fresh primes outside $S\cup U$ contribute two to the squarefree
weight in (O.9). Thus $B_U[nq_1q_2]$ is divisible by four, no unary
correction occurs there, and $H_U[nq_1q_2]$ remains divisible by two.
Here $R\ge2$ and the actual-stage budget were fixed before the test.
It follows that $t(g)t(h)F_U=0$ after reduction. Applying the middle
identity to $t(gh)t(h)$ gives $t(gh^2)F_U=t(g)F_U$.
Invariance under all squares, and continuity, give factorization
through the elementary quadratic quotient unramified outside
$S\cup U$. Over $\Q$ this quotient is finite: its squareclasses
are generated by $-1$, $2$, and the primes in that support.
This does not assert abelianity of the characteristic-zero
Galois representations.

At $p\in U$, first keep the summands of $F_U$ separately.
The block $\theta B_U$ and each unary block with $p\in J$ have
level exponent at most one and primitive quadratic character at $p$.
The unary blocks with $p\notin J$ are degeneracies of forms at a
level omitting $p$. This distinction is preserved when retaining all
old multiplicities. In particular, the latter blocks cannot acquire a
Steinberg constituent merely by being viewed at a larger level.

Enlarge the characteristic-zero coefficient field to realize all the
cuspidal newform constituents and the Eisenstein character pairs of
these blocks. On a constituent and all its degeneracies, define $t(g)$
and $d(g)$ to be the scalar trace and determinant of its two-dimensional
Galois representation, acting as the identity on its multiplicity
space. Distinct characteristic-zero constituents with the same good
eigenvalues have the same trace and determinant by Chebotarev, so
these definitions agree on coinciding summands. In particular they
give well-defined operators on their sum. No decomposition modulo
two, and no residual irreducibility assertion, enters this construction.

For a cuspidal eigenform $f$, let $\Pi_f$ denote its unitary
adelic normalization: at a good prime $q$ its Satake parameters
$\alpha_q,\beta_q$ satisfy
$a_q=\sqrt q(\alpha_q+\beta_q)$ and
$\alpha_q\beta_q=\chi(q)$.
Apply \cite[\S\S0.3--0.7, Theorem~A, pp.~409--411]{carayol}
to $\Pi_f^\vee$, and set $\rho_f=\sigma(\Pi_f^\vee)^\vee$.
Here $F=\Q$, the real weight is $2$, and the coefficient place
lies above $2$. The degree of $F$ is odd, so the additional
finite discrete-series-place hypothesis for even degree does
not apply. The local prime $p$ is odd.
Carayol's reciprocity sends geometric Frobenius to a uniformizer,
and his $\omega_p$ is $|\cdot|_p$.
His normalized-induction formula is
\[
 \sigma(\operatorname{Ind}(\xi_1,\xi_2))
 =\xi_1^{-1}\omega_p^{-1/2}\oplus
   \xi_2^{-1}\omega_p^{-1/2}.
\]
Thus if $\Pi_{f,p}=\operatorname{Ind}(\mu_1,\mu_2)$,
$\rho_f$ on arithmetic Frobenius has eigenvalues
$\sqrt p\,\mu_1(p),\sqrt p\,\mu_2(p)$ whenever these
characters are unramified. In particular its good polynomial
is $X^2-a_qX+q\chi(q)$. Both the automorphic contragredient and
the dual on the Galois side are part of this convention.

For a block with ramified character, its local conductor is
exactly one: it is at most the level exponent and at least the
conductor of its determinant.
Casselman's newvector theorem and conductor calculation
\cite[Theorem~1, pp.~302--306]{casselman} show that a special
representation has conductor one only for an unramified twist;
a ramified twist has conductor at least two. A supercuspidal
representation has conductor at least two. Consequently
$\Pi_{f,p}=\operatorname{Ind}(\mu_1,\mu_2)$ with $\mu_1$
unramified and $\mu_2$ of conductor one, after interchanging
the characters if necessary. Its newvector space is a line.
On inertia the ramified character is the prescribed quadratic
character.

We check the Fourier normalization of the operator on this line.
Use normalized induction, so that
\[
 v\left(\begin{pmatrix}a&x\\0&d\end{pmatrix}g\right)
 =\mu_1(a)\mu_2(d)|a/d|_p^{1/2}v(g).
\]
Put $K_0(p)=\{g\in\mathrm{GL}_2(\Z_p):g_{21}\in p\Z_p\}$.
Use its character $k\mapsto\mu_2(k_{22})$.
This is the lower-right-character convention. Casselman's
upper-left-character convention is carried to it by
$h_p=\left(\begin{smallmatrix}0&1\\-p&0\end{smallmatrix}\right)$,
which normalizes $K_0(p)$ and interchanges its diagonal entries
under conjugation. We label the inducing characters so that
$\mu_1$ is unramified, using the interchange isomorphism for
irreducible principal series.
Normalize Haar measure by $\operatorname{vol}(K_0(p))=1$.
The character-weighted double-coset operator for
$K_0(p)\operatorname{diag}(p,1)K_0(p)$, with value one at
$\operatorname{diag}(p,1)$, is
\[
 H_p=\sum_{b\in\Z/p\Z}
       \Pi_{f,p}\left(\begin{pmatrix}p&b\\0&1\end{pmatrix}\right).
\]
Take the newvector supported on $B K_0(p)$ with $v(1)=1$.
Evaluating at $1$ gives
\[
 (H_pv)(1)=p\mu_1(p)p^{-1/2}=\sqrt p\,\mu_1(p).
\]
This is the Fourier operator $U_p^{(2)}$. Indeed, in the
Whittaker model use an additive character of conductor $\Z_p$.
Writing $a(n)=\operatorname{diag}(n,1)$ and
$n(x)=\left(\begin{smallmatrix}1&x\\0&1\end{smallmatrix}\right)$,
\[
 a(n)\begin{pmatrix}p&b\\0&1\end{pmatrix}=n(nb)a(np)
 \quad(n\in\Z_p),\qquad
 (H_pW)(a(n))=pW(a(np)).
\]
Normalize the archimedean Whittaker function by
$W_\infty(a(y))=y\exp(-2\pi y)$ for $y>0$.
Comparison of its positive-$n$ term with $y f(iy)$ gives the
finite Whittaker coefficient $a_n/n$. The displayed equality
therefore sends $a_n$ to $a_{pn}$. Hence
$U_p^{(2)}=\sqrt p\,\mu_1(p)$ on the new line, exactly the
arithmetic-Frobenius eigenvalue on the unramified Galois line.
This scalar is algebraic; the square roots only express the
unitary normalization. Old multiplicities at primes other than
$p$ carry the same scalar. Casselman's oldspace dimension formula
\cite[Corollary, p.~306, and Theorem~4, p.~313]{casselman}
does not introduce an additional multiplicity at this conductor-one
prime. For Eisenstein constituents, the two Galois characters and
the divisor-sum Fourier coefficients give the same equality directly;
Carayol's cuspidal theorem is not applied to them.

Choose an inertia element $\tau$ on which the quadratic character
is $-1$, and a Frobenius lift $\phi$. On a ramified block, a basis
of the two character lines gives
\[
 \rho(\tau)=\begin{pmatrix}1&0\\0&-1\end{pmatrix},\qquad
 \rho(\phi)=\begin{pmatrix}\lambda&0\\0&\mu\end{pmatrix},
 \qquad U_p^{(2)}=\lambda,\quad d(\phi)=\lambda\mu.
\]
For $\rho(g)=\left(\begin{smallmatrix}a&b\\c&e\end{smallmatrix}\right)$,
the two inertia differences are
\[
 t(g\tau)-t(g)=-2e,\qquad
 t(g\tau\phi)-t(g\phi)=-2e\mu.
\]
It follows in characteristic zero that
\[
 U_p^{(2)}\{t(g\tau\phi)-t(g\phi)\}
       =d(\phi)\{t(g\tau)-t(g)\}.                       \tag{O.12}
\]
On a block coming from level prime to $p$, the representation is
unramified and both differences vanish. This proves the identity
on its entire old multiplicity space even if $U_p^{(2)}$ is not
scalar or semisimple on that space. It therefore proves (O.12)
on the actual sum containing $F_U$.

These are identities on characteristic-zero forms before any
reduction. To relate them to the integral trace module above, choose
finitely many Fourier coordinates that inject its ambient space.
Every generator of that module has integral coordinates; hence the
module is a submodule of a finite free module over the complete
coefficient DVR. It is finite and closed. Simultaneous Chebotarev
approximation in the finitely many Galois representations makes
$t(g)$ and $d(g)$ limits of good Hecke and diamond operators on this
space. They preserve this integral module, so the identity can be
evaluated on the integral $F_U$ and then reduced. This does not
require integral projectors onto constituents. In particular no
division by the two in the displayed inertia differences is used.


Apply (O.12) at coefficient one.  Put $V=U\setminus\{p\}$,
$v=\prod_{r\in V}r$, so $n_U=n_V+1$.  For a fresh prime $r$ set
\[
 C_p(vr)=\frac{a_p-(\psi/\eta)(p)(-vr/p)-1}{2}\in\Z,
 \qquad c_p(vr)=C_p(vr)\bmod2.
\]
The two surviving operator terms at $pr$ and the good recurrence give
\[
 H_U[pr]\equiv
   \eta(p)C_p(vr)\,2^{-n_V-1}a(vr)\pmod{2\mathcal O_{\rm coef}}.
\]
For the even source this is an equality before reduction.  For the
odd source the error is in $2^{M_\theta-n_U-1}\mathcal O_{\rm coef}$,
which is contained in $2\mathcal O_{\rm coef}$ at our chosen stage.
Choose fresh $q,q_\tau,r_0,r_1$ representing respectively
$g,g\tau,g\phi,g\tau\phi$ in one sufficiently precise joint quotient
for the $U$ and $V$ forms and all retained characters.
The omitted form uses the same $g,\eta,\psi,\delta$ and stage,
and its level omits the entire rational prime $p$.
Thus the two $V$ symbols at $vr_0,vr_1$ agree.  Their Legendre
signs at $p$ are opposite, so $c_p(vr_0)+c_p(vr_1)=1$ in $\F_2$.
The coefficient-one reading and the reduced trace identity now give
\[
             \mathcal A_U(Uq_\tau)+\mathcal A_U(Uq)=\mathcal A_V(Vr_0)\qquad\text{in }\F_2.
\]
Here the subscript records the operator set, not a new coefficient
source or minimum. The factor $\eta(p)$ reduces to one; no factor
two is cancelled in (O.12). For every retained quadratic character
$\chi$ unramified at $p$, the representative of the omitted vertex
satisfies
\[
 \chi(g\phi)=\chi(g)\chi(\phi)
            =\chi(\operatorname{Fr}_q)\chi(\operatorname{Fr}_p).
\]
For a fixed-support character this is the product of the two vertex
labels. For the character of a third vertex $r\in V$, it adds the
two bits $(r/q)$ and $(r/p)$ in $\F_2$. Reciprocity then gives the
reverse bits from the same endpoint types. All additional stationary
tests were included in the joint quotient, so they obey this formula
as well. This proves the full product-label omission rule, including
every external incidence.

A fresh vertex with trivial label in the elementary extension for
$S\cup U$ has the same reduced reading as $t(1)$, hence gives zero.
An isolated primary vertex has precisely this fixed label and these
zero incidence bits. A contracted cluster of primary vertices has
the same property when its external incidences are neutral. The
result is still a nonempty configuration, including the one-vertex
case; no contraction deletes the last vertex.

Finally the fixed-$U$ factorization gives dependence on the entire
finite typed graph. For two disjoint prime realizations of the same
graph, choose a fresh bridge copy with the prescribed types and
matching incidences to both realizations. CRT and Dirichlet provide
the primes successively, and reciprocity determines their reverse
bits. Replace the vertices of the first copy by bridge vertices one
at a time, then the bridge vertices by those of the second copy.
At each step the other vertices form the same fixed $U$, and the
changed vertex has the same elementary Frobenius data for $S\cup U$;
its symbol therefore agrees. If the original realizations overlap,
first replace one by a fresh copy avoiding both original prime sets,
with incidences matching that first realization. A vertex and its
own replacement are never tested together, so their mutual bit can
be chosen freely. The resulting two disjoint realizations admit the
preceding bridge argument without conflicting prescriptions.
Only graphs of the original bounded size are evaluated; extra bridge
incidences are prime-selection conditions. Thus these are rules for
the actual arithmetic graph symbols, at the already chosen valid
stage. The simultaneous finite graph budget and stabilization are
treated below.
\end{proof}

% End 08-odd-source.tex


% Begin 09-normalization-endpoints.tex
\section{Nonvanishing, normalization and the final family}\label{app:endpoints}
The uses of the coefficient and address propositions have the following
order. First apply only their even-source case, whose depth is already
fixed by Proposition~\ref{ex:app-even-phase}. This gives the rough
companion and the genus lower bound. These in turn establish the odd
minimum and its finite-stage inequalities. Only then apply the odd
case of Proposition~\ref{ex:app-coefficient-rules} and the two-symbol
address construction. In particular the proof of the rough companion
does not use the later odd minimum.
We give the choice order underlying (3.3)--(3.9). This also separates
the two valuation arguments: the even upper transfer uses the common
Kato scalar, while the final contradiction uses the common Heegner
coordinate of Section~\ref{sec:limits}.

\subsection{An absolute bound of five new primes}
The theorem of Hoffstein--Luo \cite[p.~436]{hl} supplies fundamental
discriminants with nonzero central value and at most four prime
factors. Its finite splitting set is arbitrary. The bound four is
absolute; no uniform size threshold for those discriminants is needed.

\begin{proposition}[Nonvanishing with five primes and finite local prescriptions]\label{ex:app-five-prime}
Let $A/\Q$ be a fixed modular elliptic curve for this application,
let $S$ contain $2,3$ and all its bad primes, and let $B$ be a finite
prime exclusion set. In every nonempty signed unit-squareclass filter
at $S$ of even root number, there is a signed odd squarefree parameter
$u$, prime to $S\cup B$, with $L(A^u,1)\ne0$ and $\omega(u)\le5$.
Only the intermediate parameter supplied by Hoffstein--Luo is required
to be fundamental; the final field discriminant may be $4u$.
The number five is independent of the prescribed finite set. No
uniform bound on the sizes of the selected primes is asserted.
\end{proposition}
\begin{proof}\leavevmode
Let $A/\Q$ be any fixed modular elliptic curve for this application,
$S$ contain two, three and all its bad primes, and $B$ be any finite
exclusion set. Fix a nonempty signed unit filter of even root number.
By CRT and Dirichlet choose one signed prime $q_0$ in its classes,
outside $S\cup B$. Apply the theorem to the actual curve $A^{q_0}$,
putting
\[
                 S'=S\cup B\cup\{|q_0|\}\cup
                               \{p:p\mid N_{A^{q_0}}\}
\]
in its splitting set. The theorem gives a fundamental $d$ with
$\chi_d(p)=1$ for all $p\in S'$ and at most four prime factors.
The functional equation forces $d>0$, since $A^{q_0}$ has even
root number and every conductor prime splits. Splitting at two
gives $d\equiv1\pmod8$. Thus $q_0d$ has all original unit
classes, is prime to $B$, is squarefree, and has at most five
prime factors. If $q_0d\equiv3$ or $7\pmod8$, its field
discriminant is $4q_0d$; only $d$ was required to be fundamental
in the theorem. This is compatible with the coefficient
construction (E.2).
\end{proof}

Here $h$ is restricted to the prepared positive odd-sign unit
filter containing one, as in (3.3), so $h\equiv1\pmod8$.
The multiset of dummy labels and its size are fixed before $h$;
its realizing primes are chosen after $h$ and avoid its support.

For the rough companion the curve itself varies:
write $a$ for the bounded dummy product and take $A=E^{ha}$.
Choose $q_0$ so $k_0=aq_0$ is negative, $1$ modulo eight and
a square at every prime of $2Nh$, with the required exclusions.
Then $k_0$ is odd fundamental, prime to $N_{E^h}$, and
\[
          w(E^{hk_0})=w(E^h)\chi_{k_0}(-N_{E^h})=+1.
\]
This proves the sign directly for the actual application.
Apply Hoffstein--Luo to $E^{hk_0}$ with every conductor,
old, $h$, dummy and forbidden prime in its splitting set.
The resulting $k'=k_0d$ adds at most five new primes.
The constants and sizes in the analytic theorem may depend on
this varying curve and the whole finite set; the bound five does not.

\subsection{The even upper transfer}
\begin{proposition}[A uniformly bounded partitionable companion]\label{ex:app-rough-companion}
Fix $E$ and the finitely many prepared even sources, local types and
unit seeds used for (3.3). There are constants $A,C_E$, fixed before
$h$ and before any additional finite prime exclusion set, with the
following property. Let $h$ be positive squarefree in the prepared
odd-sign unit filter containing $1$, prime to its fixed support; in
particular $h\equiv1\pmod8$. There is a negative odd fundamental
discriminant $k'$, prime to $2Nh$ and the additional exclusions,
square at every place dividing $2Nh$, for which $hk'$ is partitionable
in the sense of Definition~\ref{ex:sec7-seed-partition},
$L_*(hk')\ne0$, and
\[
 \omega(k')\le A,\qquad
 v_2L_*(hk')\le2\omega(h)+2\omega(k')+s(hk')+C_E.
\]
Here $A$ is the fixed dummy-label count plus five. Dummy types and
their number are fixed first, but their distinct prime realizations
are selected after $h$ and the exclusions. The proof uses the
one-symbol stationary-address conclusion of
Lemma~\ref{ex:sec7-one-symbol-address}, whose proof is given in
Proposition~\ref{ex:app-simultaneous-address}.
\end{proposition}
\begin{proof}\leavevmode
For one fixed even source write (E.11) as
\[
             v_2L_*(\sigma D)=2v_2a(D)+\kappa_e.
\]
The lower bound (E.9) and a nonzero reference coefficient imply
that the nonempty set of integers
$v_2a(D)-\omega(D)$ has a finite attained minimum $\delta_e$.
A unit of the normalized symbol means precisely equality to this
minimum. It therefore gives
\[
 L_*(\sigma D)\ne0,\qquad
 v_2L_*(\sigma D)=2\omega(D)+C_{e,\rm unit},
 \quad C_{e,\rm unit}=2\delta_e+\kappa_e.                         \tag{F.1}
\]
This constant belongs to the fixed source, not to a curve later
introduced solely to obtain nonvanishing.

For the negative types used in the rough companion, the unsigned
index is $7$ modulo eight. Its total $(-1/D)$ label is nonzero, so
every unit seed is nonempty. A seed of least positive size has no
dependence on internal edge bits: a nonzero edge derivative
would be a smaller nonempty unit. Its labels therefore suffice.
If a list with the same total label contains the first $r-1$
seed anchors and a nonempty remainder, contract that remainder
into the last anchor. Its total label is forced to be the
remaining seed label. This proves partitionability regardless
of how many old $h$-primes are in the remainder.

For a partitionable $d$ already known to have Selmer corank zero,
the one-symbol address construction keeps every prime of $d$
stationary and gives arbitrarily large cubes $d_x$, $d_0=d$,
with unit symbol at each $x\ne0$. The scalar in (E.5) is common
to that cube. By (E.8) and (F.1),
\[
      v_2u_x^{\rm ev}(0)\le B_e:=C_{e,\rm unit}+C_{\rm loc}
                                      \quad(x\ne0).
\]
Choose $M-c_E>B_e$. The constant coefficient of
$2^{c_E}u^{\rm ev}\in\Lambda[G_b]$ has sign-value bits of total
Boolean degree at most $2^{M-1}-1$. For a larger cube dimension,
the Boolean comparison in Section~\ref{sec:limits} gives
$x\ne0$ with
$v_2(u_x^{\rm ev}(0)-u_0^{\rm ev}(0))\ge M-c_E$.
If the base were zero or had valuation greater than $B_e$, this
would contradict the upper bound at $x$. Therefore
\[
       v_2L_*(d)-2\omega(d)-s(d)
                       \le C_{e,\rm unit}+2C_{\rm loc}.          \tag{F.2}
\]
The required dimension may depend on $d$; this bound does not.
It precedes the odd minimum and uses no Heegner transfer.

Prepare these sources and nonempty seeds for every total type at
the fixed support. There are finitely many. A bounded supply of
dummy anchors and reserves works for all of them. For each $h$
realize that supply by distinct primes avoiding the entire fixed
support and $h$. Put every dummy and every other already fixed
forbidden prime in the exclusion set of the preceding
nonvanishing construction. It gives $k'$ with
$\omega(k')\le A$, where $A$ is the fixed dummy count plus five.
For the total type that actually occurs, $hk'$ is partitionable.
Nonvanishing gives corank zero by the forward theorem, so (F.2)
applies. Taking the maximum over the finitely many prepared
constants proves
\[
 v_2L_*(hk')\le2\omega(h)+2\omega(k')+s(hk')+C_E.                 \tag{F.3}
\]
No new even minimum is chosen after $h$ or its exclusion set.
Those excluded primes are not added as relaxed places in a
Selmer complex.
\end{proof}

\subsection{The odd minimum is chosen after its lower bound}
\begin{proposition}[The odd minimum and every finite operator budget]\label{ex:app-odd-minimum}
Fix the original odd filter, the detecting sequence of
Proposition~\ref{ex:app-detector}, and one nonprincipal ultrafilter
on its stages. The arithmetic lower bound (3.5), the native comparisons
(3.4), and Proposition~\ref{ex:app-rough-companion} give (F.4)--(F.5)
before any odd minimum is chosen. The bounded-weight visible depths
therefore have an attained finite integer minimum $\delta_o$.
For every fixed weight cutoff, the simultaneous inequality (F.6)
holds on a set in this same ultrafilter. For every finite operator
budget and $R\ge2$, one may choose a stage where all the integer
divisions and residual errors required in
Proposition~\ref{ex:app-coefficient-rules} are valid at every Fourier
index.
On a fixed squarefree-output range the unit symbol defined before
(F.7) is exhaustive, and its value $1$ implies (F.7), with the one
constant $C_o=\delta_o+C_{\rm det}$. No minimum is changed on
omitting a vertex or adding stationary tests, and no single finite
stage is claimed to serve every unbounded weight range.
\end{proposition}
\begin{proof}\leavevmode
The arithmetic result (3.5) and the two native-lattice height
comparisons (3.4) give, using (F.3),
\[
 2j_{\rm gen}(h)\ge2\omega(h)+s(h)-C_{\rm gen},\qquad
 C_{\rm gen}=2A+C_R+C_E-v_2L_*(h_*)+C_H.                         \tag{F.4}
\]
The variable $s(hk')$ cancels. This estimate precedes the theta
prime, precision, odd minimum, final field and binary dimension.

Combining (F.4) with (O.8) gives on all admissible $D$
\[
            v_2a_M(D)\ge\min\{M,\omega(D)-C_{\rm coarse}\},
 \qquad C_{\rm coarse}=C_{\rm gen}/2+C_{\rm det}.                 \tag{F.5}
\]
At visible coefficients (O.8) first gives a simple zero, allowing
(F.4); at the other odd-sign parameters the coefficient is zero
modulo $2^M$. This avoids a circular use of the future minimum.

Fix one nonprincipal ultrafilter on the actual stages with
$M=M_\theta\to\infty$. Consider the integers $r$ for which a
bounded-weight sequence of visible $D_M$ has
$v_2a_M(D_M)-\omega(D_M)=r$ on a set in the ultrafilter.
The fixed witness $D_w$ and (O.2),(O.7) make this set nonempty.
Equation (F.5) bounds it below. Let its attained minimum be
$\delta_o$. For every fixed $W$, on a set in the same ultrafilter,
\[
 v_2a_M(D)\ge\omega(D)+\delta_o
             \quad\hbox{for every visible }D,\ \omega(D)\le W.
 \tag{F.6}
\]
Otherwise choose a violating $D_M$ at each bad stage. One of the
finitely many integers between $-C_{\rm coarse}$ and
$\delta_o-1$ would occur on an ultrafilter-large set, a contradiction.
Invisible coefficients satisfy the required truncated inequality
once $M>W+\delta_o$.

For a fixed finite operator budget $|U|\le W_U$, put
$N_0=\max(0,\delta_o+W_U)$ and fix $R\ge2$.
Choose an integer $T\ge W_U+1$ with
$T-C_{\rm coarse}\ge N_0+R$, then choose a stage in the
ultrafilter-large set where (F.6) holds through weight $T$, with
$M\ge N_0+R$ and $M>T+\delta_o$.
For the index $Jrm^2$ put $D=(U\setminus J)r$.
The integral shell recurrence has an error in $2^M\mathcal O$.
If $\omega(D)>T$, (F.5) gives
\[
 v_2\bigl(2^{n_U}B_U[Jrm^2]\bigr)\ge N_0+R.
\]
If $\omega(D)\le T$, (F.6) for visible coefficients and
$v_2a_M(D)\ge M$ for invisible ones instead give
\[
 v_2\bigl(2^{n_U}B_U[Jrm^2]\bigr)
       \ge\min\{M,n_U+\omega(r)\}.
\]
Consequently $v_2B_U[Jrm^2]\ge\min\{R,\omega(r)\}$ for every
Fourier index and every $U$ in the budget at this one stage.
At $r=1$ the normalized originating coefficient
$2^{|J|-n_U}a_M(U\setminus J)$ is integral, so (O.10) proves
the unary congruence modulo the integer two.
The additional division in (O.11) leaves error valuation at least
$R-1\ge1$.  In particular $M-n_U-1\ge1$ for every reading and
first-shell omission in the budget.
It does not select one stage for infinitely many budgets.

More precisely, fix a squarefree-output cutoff $L$. On a set of
stages in the same ultrafilter, (F.6) holds through $L$ and
$M_\theta>\delta_o+L$. Define $\mathcal A_{M_\theta,L}(D)$ only
for positive squarefree $D$ with $\omega(D)\le L$: it is zero off the full coefficient filter,
one when $v_2a_{M_\theta}(D)=\delta_o+\omega(D)$ on that filter,
and zero otherwise. This is exhaustive on the tested range:
visible coefficients satisfy (F.6), while invisible ones have
valuation at least $M_\theta>\delta_o+\omega(D)$.
In particular a unit in this range is visible and satisfies
\[
 \ord_{s=1}L(E^D,s)=1,\qquad
        j_{\rm gen}(D)\le\omega(D)+C_o,\quad C_o=\delta_o+C_{\rm det}.
 \tag{F.7}
\]
For the preceding operator budget take $L\ge W_U+1$ and include
its valid-stage requirements; all fresh-vertex and omission outputs
then lie in this range. The constant $C_o$ is independent of $L$
and of the stage. No value is assigned by this definition at
unbounded weights of one finite stage. No minimum is changed when
a vertex is omitted or a stationary condition is later added.
\end{proof}

\subsection{Nonempty rational seed and a common finite stage}
\begin{proposition}[Nonempty seeds after fixing the complementary twist]\label{ex:app-nonempty-seeds}
Choose the negative even unit $k$ with the fresh ramification and
exclusions specified in Section~\ref{sec:rc}, and fix
$K=\Q(\sqrt{k})$ before the address dimension. Then $E/K$ has full
two-primary Selmer corank one. Transport the even coefficient source
to $E^k$, normalize at index $1$, and use its positive square-unit
filter. On that filter its coefficients have
$v_2a_{\rm new}(D)\ge\omega(D)$, and the source has a nonempty
$K$-split unit seed.
If the intrinsic odd limit symbol is a unit at the empty
configuration, $E$ already has a simple zero. Otherwise it too has
a nonempty unit seed. In this remaining case both seeds can be
read through the actual common alphabet, whose projections onto
the two intrinsic alphabets are separately surjective. No
surjectivity onto their product is required.
\end{proposition}
\begin{proof}\leavevmode
Choose the final negative even unit $k$ as in
Section~\ref{sec:rc}, with fresh ramification outside the old odd
label and torsion fields. Fix $K=\Q(\sqrt k)$ before the binary
dimension. The forward theorem gives corank zero for $E^k$;
the sum-and-difference isogeny, on compact Selmer modules tensored
with $\Q_2$, gives full corank one for $E/K$.
The intrinsic odd symbol omits the primes of $k$.
Later $\ell_0$ avoids them. Fresh ramification implies that every
old odd label still lifts to a $K$-split label.
Additional stationary tests can have relations; use their actual
common alphabet and its separate surjective projections.

\paragraph{Transport to the fixed complementary twist.}
Put $d_k=|k|$, so that $k=-d_k<0$. The old source has sign $-1$;
its Fourier coefficient at the positive index $d$ tests $E^{-d}$.
Choose the new source for $E^k$ with sign $+1$, normalize its
coefficient at $1$ to be $1$, and impose a support containing the
old support and all primes dividing $k$. For the following comparison
$D$ is positive and squarefree, is prime to that support, and lies
in the square unit classes of $1$ required by the new filter.
In particular $D$ is a square in every old support completion, so
$d_kD$ and $d_k$ lie in the same old filter. Waldspurger comparison
on the two sources gives
\[
 \begin{split}
  v_2 a_{\rm new}(D)
    &=v_2 a_{\rm old}(d_kD)-v_2 a_{\rm old}(d_k)\\
    &\ge \omega(D).
 \end{split}
\]
The equality is a valuation equality; the two sources may have
different root-of-unity phases. The normalized central-value ratios
agree, since both are
\[
 \sqrt D\,\frac{L(E^{kD},1)}{L(E^k,1)}.
\]
Fixed period constants cancel in each ratio. The old unit condition
at $d_k$, and $\omega(d_kD)=\omega(d_k)+\omega(D)$, prove the
inequality. The signed twist is $kD$ throughout; no Fourier
coefficient at a negative index is used. The final family already
satisfies these unit classes. In the nonempty-seed argument the
fresh prime $p'$ is chosen to match all fixed labels of $p$, so
$pp'$ satisfies the same filter, including its dyadic condition.

For a unit of this normalized source equality holds, so
$v_2L_*(h_xk)\le2\omega(h_x)+C_{e,\rm final}$ with one fixed
period-normalization constant.
Its $K$-split seed must be nonempty. Suppose otherwise. The integral
coefficient valuations then improve the preceding lower bound to
$\omega(D)+1$ on every nonempty product of allowed $K$-split primes.
For a fresh good $K$-split prime $p$ outside the support, apply the
trace identity to $f=\theta H_{\{p\}}$:
\[
 U_p^2t(\phi)=U_p^3+d(\phi)U_p,\qquad
 (t(\phi)f)[p^2]=f[p^3]+(d(\phi)f)[p].
\]
Here and in this calculation all congruences are modulo the ideal
$2\mathcal O_{\rm coef}$ before passage to the residue field.
Write $a_p=a_p(E^{(k)})$.
We check every theta convolution term in the three coefficients.
Put
\[
 B=B_{\{p\}}=\frac{(U_p-\eta(p)V_p)g}{2},\qquad
 b_1=B[1]=\frac{a(p)}2,\quad b_p=B[p],\qquad
 H=H_{\{p\}}=\frac{B-b_1\Theta_1-b_p\Theta_{\{p\}}}{2}.
\]
The inherited depth is zero, and $H$ is integral by (O.11).
Also $H[0]=0$, since the cuspidal source and the unary series have
zero constant term.
The ordinary lower bound gives $b_1\in\mathcal O_{\rm coef}$.
Under the assumed absence of a nonempty $K$-split unit seed one
has the stronger assertion
\[
                         b_1\in2\mathcal O_{\rm coef}.          \tag{G.4a}
\]
Indeed, if $p$ lies in the coefficient filter, the improved bound is
$v_2a(p)\ge2$; otherwise the actual filter gives $a(p)=0$.
No condition that $p$ itself belong to that filter is imposed.

Let $\delta_B,\delta_1,\delta_p$ be the scalars by which $d(\phi)$
acts on the full weight-two blocks
$\theta B,\theta\Theta_1,\theta\Theta_{\{p\}}$, respectively.
They are integral units.  The determinant is the cyclotomic character
times the weight-two nebentypus on every cuspidal constituent and
on the Eisenstein character pair, independently of old multiplicity.
For this positive-twist source $\psi/\eta=\chi_{-1}$ and $i=0$.
The half-integral-weight characters of $B$ and $\Theta_{\{p\}}$
are therefore both
\[
                 \psi(p/\cdot)=\eta(-p/\cdot).
\]
Multiplication by the same $\theta$ preserves this equality.
Consequently $\delta_B=\delta_p$ for the actual Frobenius lift
$\phi$, including its action on the ramified character.  The ratio
of the characters of $\theta\Theta_{\{p\}}$ and
$\theta\Theta_1$ is quadratic.  Their cyclotomic determinant factors
are identical, so $\delta_p/\delta_1\in\{1,-1\}$ and
$\delta_p-\delta_1\in2\mathcal O_{\rm coef}$.  We must not simply
equate these two scalars.  On the actual characteristic-zero
module the exact identity is
\[
 d(\phi)f=\delta_p f+
       \frac{(\delta_p-\delta_1)b_1}{2}\,\theta\Theta_1.
\]
Now $H[p]=0$, since $B[p]=b_p$, $\Theta_{\{p\}}[p]=1$ and
$\Theta_1[p]=0$.  Integrality of $H$ implies $f[p]\in2\mathcal O_{\rm coef}$.
Also $\theta\Theta_1[p]=2\sum_{0<j^2<p}\Theta_1[p-j^2]$ is even.
The displayed determinant identity therefore already proves
$(d(\phi)f)[p]\in2\mathcal O_{\rm coef}$ using only the ordinary
integrality of $b_1$.  The following direct expansion checks the
individual cross terms under the no-seed assumption in force here.

For $n_j=p-j^2>0$, one has $p\nmid n_j$ and
$\Theta_{\{p\}}[n_j]=0$.  We claim that
\[
                           B[n_j]\in2\mathcal O_{\rm coef}.    \tag{G.4b}
\]
If $n_j$ meets the old support, the coprimality sieve makes this
coefficient zero.  In particular $2\in S$, so every odd $j$ gives
an even $n_j$ and vanishes in this way.  Otherwise write
$n_j=r m^2$, with $r$ squarefree.
Since $p\nmid n_j$, the $V_p$ term is zero and
$B[n_j]=a(prm^2)/2$.  The integral good-prime shell recurrences
express $a(prm^2)$ as an integral multiple, up to unit character
factors, of $a(pr)$.  If $r>1$, the ordinary coefficient lower
bound gives $v_2a(pr)\ge\omega(pr)=1+\omega(r)\ge2$, or the
coefficient is zero off the filter.  This uses no assertion that
the primes dividing $r$ split in $K$.  If $r=1$, the same shell
recurrences express $B[m^2]$ as an integral multiple of $b_1$,
and (G.4a) applies.  This proves (G.4b).  Also every coefficient
of $b_1\Theta_1$ is even by (G.4a).

Now expand $\theta=1+2\sum_{j>0}q^{j^2}$ before discarding any
term.  Linearity of the determinant operator on the three actual
blocks gives the exact expression
\[
\begin{split}
 (d(\phi)f)[p]
 ={}&\frac{\delta_B B[p]-\delta_p b_p}{2}\\
 &+\sum_{0<j^2<p}
 \bigl(\delta_B B[p-j^2]
       -\delta_1 b_1\Theta_1[p-j^2]
       -\delta_p b_p\Theta_{\{p\}}[p-j^2]\bigr).
\end{split}                                                     \tag{G.4c}
\]
There is no constant-theta $\Theta_1[p]$ term since $p$ is not
a square, and $\Theta_{\{p\}}[p]=1$.
The first line is zero because $B[p]=b_p$ and
$\delta_B=\delta_p$.  Every summand in the second line is even
by (G.4a)--(G.4b) and the vanishing of its last term.  In particular
representations $p=j^2+m^2$, which can occur when $p\equiv1\pmod4$,
are covered by the case $r=1$ and are not omitted.  Thus
$(d(\phi)f)[p]\in2\mathcal O_{\rm coef}$.

The other two readings have no determinant operator.  Since $H$
is integral, for every positive integer $N$,
\[
 f[N]=H[N]+2\sum_{0<j^2\le N}H[N-j^2]
                 \equiv H[N]\pmod{2\mathcal O_{\rm coef}}.     \tag{G.4d}
\]
This includes $N=p^3$ and every term in its theta convolution.
Choose a fresh $K$-split prime $p'$ matching $\phi$ in a sufficiently
precise finite Galois quotient for the retained integral module,
and matching all fixed rational characters.  The good Hecke formula
at $p'\ne p$ identifies $(t(\phi)f)[p^2]$ modulo two with
$f[p^2p']$, hence with $H[p^2p']$ by (G.4d).
Neither unary series has support at $p^2p'$, and the recurrence at
$p\mid pp'$ has $\kappa=0$.  Consequently
\[
 H[p^2p']=\frac{a(p^3p')-\eta(p)a(pp')}{4}
        =\frac{\eta(p)(a_p-1)a(pp')}{4}.
\]
The product $pp'$ lies in the square filter even if $p$ alone does
not.  Its two primes split in $K$, so the no-seed assumption gives
$v_2a(pp')\ge3$.  This last coefficient is therefore even.
The trace identity and (G.4c)--(G.4d) now imply
$H[p^3]\in2\mathcal O_{\rm coef}$.  The following computation
uses precisely this unary-subtracted coefficient.

\paragraph{The coefficient behind the mod-eight contradiction.}
Write $a_p=a_p(E^k)$. The transported source tests positive twists of $E^k$. Thus
$\psi/\eta=\chi_{-1}$ and, at the squarefree index $1$,
$\kappa=(\psi/\eta)(p)(-1/p)=1$. Also $\eta$ is odd, so
the unary correction uses $i=0$. With $a(1)=1$ the recurrence gives
\[
 a(p^2)=\eta(p)(a_p-1),\qquad
 a(p^4)=\eta(p)^2\{a_p(a_p-1)-p\}.
\]
The inherited depth is zero. Hence
\[
 B_{\{p\}}[p]=\frac{\eta(p)(a_p-2)}2,\qquad
 B_{\{p\}}[p^3]=\frac{\eta(p)^2((a_p-1)^2-p)}2.
\]
Only $\Theta_{\{p\}}$ contributes at $p^3$, with coefficient
$\eta(p)p$. Therefore the actual integer division gives
\[
 H_{\{p\}}[p^3]
   =\frac{\eta(p)^2(a_p-1)(a_p-p-1)}4.
\]
This expression is integral because full rational two-torsion gives
$a_p\equiv p+1\pmod4$. The positive-twist character convention is
essential in this computation. Under the assumed absence of a
nonempty unit seed, the trace identity makes this coefficient zero
in the residue field. Since its phase is a unit and $a_p-1$ is odd,
it follows that $a_p\equiv p+1\pmod8$. The contradiction is obtained
only after this conditional deduction.
Choose $p$ also with $(k/p)=1$, $p\equiv1\pmod4$ and
$(\ell/p)=-1$, outside the finite support.  These prescriptions
are compatible: $K$ has a ramified prime outside the ramification
of $\mathbf Q(i,\sqrt\ell)$, so
$K\cap\mathbf Q(i,\sqrt\ell)=\mathbf Q$.  The corresponding
quadratic residue classes are therefore nonempty and Dirichlet's
theorem supplies infinitely many such primes.  Since $k$ is a
square modulo $p$, the reductions of $E^{(k)}$ and $E_\ell$ are
isomorphic.  A half of a nonzero two-torsion point with root $e_i$
on a split cubic requires both differences $e_i-e_j$ to be
squares.  For the roots $0,\ell,-3\ell$, one of these required
differences is, respectively, $-\ell,\ell,-4\ell$.  All three are
nonsquares under the chosen conditions.  Thus none of the three
nonzero points of order two has a half over $\mathbf F_p$.
If the two-primary subgroup had order greater than four, its
structure as a finite abelian group with exactly four elements
killed by two would give an element of order four, and hence such
a half.  Its order is consequently exactly four.  This gives
$v_2(p+1-a_p)=2$, contradicting the congruence modulo eight above.
This proves the required nonempty even seed. If the intrinsic odd
symbol already has a unit at the empty configuration, its detection
formula gives a simple zero for $E$ and the converse is finished.
Otherwise take a least nonempty odd unit seed. Thus the remaining
two-symbol address argument has two nonempty seeds.
\end{proof}

\begin{proposition}[Compatible finite graph tables on one ultrafilter]\label{ex:app-odd-stabilization}
Use $\delta_o$ and the ultrafilter of
Proposition~\ref{ex:app-odd-minimum}. First fix the abstract type of
the bounded elementary label groups and their projection data on
a set in this same ultrafilter. For every finite cutoff $W$,
the valid odd stages give finitely many possible symbol tables on
the compatible typed graphs with at most $W+1$ vertices.
One table occurs on a set in that ultrafilter. These selected tables
agree on their common restrictions, and their union defines the
intrinsic odd symbol on all compatible finite typed graphs.
It retains its bounded-weight witness and the nonempty-output
coefficient rules. After fixing $k$, its stationary tests may be
included in a common alphabet with separately surjective
projections on the same ultrafilter, without changing $\delta_o$.
Every specified finite collection of tests has a common actual
stage. No stage common to all cutoffs is selected.
\end{proposition}
\begin{proof}\leavevmode
Here the finite graph argument has a definite arithmetic order.
The moving theta prime contributes only a bounded number of quadratic
label coordinates. The root characters at the old level were fixed
before $M_\theta$; growing Kummer precision introduces no growing-order
character on the variable primes. Thus the elementary label groups,
their projections and their finite addition tables have bounded size.
Choose their abstract type on the same ultrafilter as $\delta_o$.
For each fixed weight cutoff $W$, restrict first to the
ultrafilter-large set of stages where the preceding budget and
precision conditions hold for every $|U|\le W$, with $R=2$.
On those stages the coefficient reading, factorization, contraction
and isolation rules hold on all configurations of at most $W+1$
vertices, with the stated nonempty-output restriction.  They therefore define actual symbol truth tables on
the compatible typed graph data.  There are finitely many such
tables; select the one occurring on a set in the same ultrafilter.
Two selected tables agree on their common restriction by taking
the finite intersection of their valid-stage sets. Their union defines
the odd limit symbol on compatible finite typed graphs of every size;
it is not a total valuation symbol at any one finite stage.
Every seed or contraction identity uses only finitely many such graphs
and is inherited at a cutoff containing all of them.
The bounded-weight witness remains a witness in these tables.
No stage common to all cutoffs is selected.
After fixing $k$, include its tests and choose the common projection
data on the same filter, without redefining the intrinsic odd minimum.
Stabilize the finite alphabet and its two projections in this way.
\end{proof}

\begin{proposition}[Stationary and simultaneous address constructions]\label{ex:app-simultaneous-address}
Use the definitions and realization freedoms in
Definitions~\ref{ex:sec7-typed-graph}, \ref{ex:sec7-symbol},
\ref{ex:sec7-address} and Lemma~\ref{ex:sec7-realization},
and the coefficient rules of Proposition~\ref{ex:app-coefficient-rules}.
Fix $b\ge1$. Either take one symbol on a nonempty set of always
active terminals with a nonzero polynomial in their pair bits, or
take one or two symbols on terminals with nonzero linear activations,
including a terminal of every nonzero activation, such that each
symbol separately has a nonzero terminal polynomial at every
$x\ne0$. In the latter case the symbols may read different
compatible projections of the common alphabet and fixed stationary
external data. Require the active data to satisfy each prescribed
filter at every nonzero address.
Starting from fixed compatible typed terminal and stationary
data, one can adjoin finitely many primary auxiliaries with nonzero
linear activations and choose only incidences involving them so that
all the symbols equal $1$ at every $x\ne0$. All original data stay
fixed, no auxiliary is present at zero, and the required neutral
external incidences and splitting conditions are retained.
Only nonempty configurations and nonempty contraction outputs are
used. A nonzero iterated difference, hence a partitionability
certificate, suffices for the required nonzero terminal polynomial.
For the two final symbols the zero-total packets below supply these
terminal hypotheses using only separate projection surjectivity.
For the even stationary application, the original prime representatives
may be fixed at the outset. For a stabilized odd symbol, first fix the
abstract data and the finite budget, then choose an actual stage
carrying them, and only then choose prime representatives. Previously
fixed representatives may be retained when compatible with the full
type at that stage; arbitrary earlier primes are not asserted to
satisfy a later moving condition.
\end{proposition}
\begin{proof}\leavevmode
\paragraph{Total labels of the final seeds.}
Both final symbols have projected total label zero. The intrinsic
odd filter contains $1$: its fixed unit squareclasses, positive sign,
and moving condition $(D/\ell_0)=1$ are those of $1$. The even source
for $E^k$ is likewise normalized in the unit filter containing $1$.
Every nonzero coefficient, and hence every unit seed, lies in its
respective filter. Additional stationary tests at primes dividing
$k$ refine the common alphabet but are forgotten by the intrinsic
odd projection. They do not add conditions to that odd symbol.

Take repeated packets with total zero in the common elementary
label group, including sufficiently many copies of every label.
For either projection separately, keep the first $r-1$ anchors of
its unit seed and contract the nonempty remainder into the last
anchor. Since the total projected label is zero, the remainder has
the last seed label. This argument uses surjectivity of each
projection separately. It does not require that a fixed old odd
seed admit lifts whose total is zero in the common alphabet.
The whole packet already has that total, and contractions preserve
it. Thus the final twists satisfy the stationary tests as well as
both coefficient filters.
Concretely, choose a packet with an even number of copies of every
common label, and enough copies to provide the anchors of either
seed under its respective projection.  Its total label is zero.
For every nonzero activation $\lambda$, take a disjoint copy of
this whole packet and give all its vertices activation $\lambda$.
At every $x\ne0$ there are exactly $2^{b-1}$ active packets.
The active list is therefore nonempty, has total common label zero,
and contains enough anchors for each projected seed separately.
Its terminal polynomial for each symbol is nonzero by the
contraction certificate, regardless of the initially fixed bits.
It also contains a terminal of each required activation.  This
constructs the terminal input of the address lemma explicitly.
No simultaneous lifting of a prescribed pair of seed labels is
needed.  Every added network vertex is primary, so the total
common label remains zero at each address.

\paragraph{Network operators and their degree.}
We spell out the simultaneous argument, including the product over
all addresses.  Write $\partial_e$ for the finite difference in a
terminal pair bit $e$.  Products of these operators commute and
$\partial_e^2=0$.  The contraction rule identifies a forest product
with its contracted graph.  A cycle product is zero: after a
spanning path is contracted its last edge is internal and changes
neither a type nor any external incidence.  No contraction removes
the last vertex.

A network of primary auxiliary vertices acts on any of the symbols
by an operator in the $\partial_e$.  To see this directly, expand in
the bits incident to the auxiliary vertices.  Contract the chosen
forest edges.  A component with no terminal becomes an isolated
primary vertex, and its contribution vanishes by primary isolation.
A component meeting several terminals gives their forest difference.
If $Q$ is the grounded auxiliary Laplacian and $c_a$ is its attachment
column at terminal $a$, the degree-zero and degree-one terms are
\[
 \det(Q)\,1+
 \sum_{a<b}c_a^{\mathsf t}\operatorname{adj}(Q)c_b\,
                    \partial_{ab}.
\]
All other terms have difference degree at least two.  The displayed
coefficients are the rooted-forest minors, by Cauchy--Binet on the
incidence matrix.  They are polynomial identities, so apply to
singular $Q$ as well.  Distinct networks have no edges between their
auxiliary vertices.  Their operators multiply: expanding both
forests gives their union, and unions containing a terminal cycle
give zero by the preceding observation.  This also shows that the
operator is the same for both projected symbols.

Fix a target $x_*\ne0$.  Choose a basis
$\lambda_1,\ldots,\lambda_b$ of the dual address space with
$\lambda_i(x_*)=1$ for all $i$.  Such a basis is obtained from one
functional taking value one by adding it to a basis of the space
of functionals vanishing at $x_*$.  Group $i$ has three primary vertices, each activated
by $\lambda_i$.  In its physical three-dimensional coordinate space
let $e_i=(1,1,1)$ and choose $u_i,v_i$ spanning its orthogonal
complement.  The vector $e_i$ has norm one.  Prescribe the symmetric
matrix $Q$ by its bilinear matrix $I_b$ on the $e_i$ and
\[
 H(T,B)=\begin{pmatrix}0&T\\T^{\mathsf t}&B\end{pmatrix}
\]
on the remaining coordinates.  In physical coordinates each
$h_i=Qe_i=e_i$ is supported on its own group.  Attach this column to
a terminal activated by $\lambda_i$, which exists by hypothesis.
For an active set $J$ the sum of these active attachment columns is
$Q_J\mathbf1_J$.  Thus the prescribed principal matrix $Q_J$ really
is the grounded graph Laplacian on that address, including its
diagonal entries.

For a selected terminal pair $a,b$ active at $x_*$, let their
activations be $\alpha,\beta$.  Attach an extra column $c$, dual to
$u_1$, to both.  If $\alpha\ne\beta$, attach the same column to a
terminal of activation $\alpha+\beta$; it is absent at $x_*$.  The
extra columns sum to zero at every address since
$\alpha+\beta+(\alpha+\beta)=0$.  If $\alpha=\beta$, the two
columns already cancel.  Hence they do not alter $Q_J$.

For the singular network set $T=I+P$, where $P$ is the cyclic shift,
and give $B$ one diagonal entry one.  On every proper group subset,
the restricted cyclic shift is a union of paths, so $T_J$ is
invertible and $\det Q_J=\det(T_J)^2=1$.  On all groups, the kernels
of $T$ and $T^{\mathsf t}$ are generated by the all-one vector.
The one nonzero diagonal entry of $B$ forces the lower component of
a vector in $\ker H$ to vanish.  Thus $\ker Q$ is precisely the line
generated by $r=\sum_i u_i$, and
$\operatorname{adj}(Q)=rr^{\mathsf t}$.  The anchor columns pair
to zero with $r$, while $c^{\mathsf t}r=1$.  Only $a,b$ receive
active extra columns at $x_*$.  Consequently its network operator
$S_{x_*,e}$, for $e=ab$, has the expansions
\[
 S_{x_*,e}|_{x_*}=\partial_e+\text{terms of degree at least two},
 \qquad
 S_{x_*,e}|_x=1+\text{terms of positive degree}\quad(x\ne x_*).
\]
For $b=1$, $T=0$ and $B=(1)$ give exactly the same calculation.

For a varying network set $T=I+L$, with $L$ the consecutive lower
shift, and $B_\varepsilon=\varepsilon E_{bb}$ for
$\varepsilon\in\{0,1\}$.  Every $T_J$ is invertible and both
operators $V^\varepsilon_{x_*,e}|_x$ have constant coefficient one.
The upper-left block of $H^{-1}$ is
$T^{-\mathsf t}B_\varepsilon T^{-1}$.  Its change on the selected
columns is $((T_J^{-1})_{b1})^2$, interpreted as zero when an endpoint
is absent.  This is one exactly when the entire path is present,
which means $x=x_*$.  Anchor columns have zero change.  Therefore
\[
 \begin{aligned}
 (V^0_{x_*,e}+V^1_{x_*,e})|_{x_*}
       &=\partial_e+\text{terms of degree at least two},\\
 (V^0_{x_*,e}+V^1_{x_*,e})|_x
       &=\text{terms of degree at least two}\quad(x\ne x_*).
 \end{aligned}
\]

\paragraph{One choice for both symbols at every nonzero address.}
For each $x\ne0$ let $P_{1,x},P_{2,x}$ be the two nonzero terminal
polynomials, of degrees $d_{1,x},d_{2,x}$, and let $H_{s,x}$ be their
highest homogeneous parts.  Work in the ordinary polynomial ring
in the terminal pair variables for this address.  Since it is a
domain, $H_{1,x}H_{2,x}\ne0$.  Choose a coefficient-one monomial
\[
 M_x=\prod_e X_e^{n_{x,e}},\qquad n_{x,e}\in\{0,1,2\},\qquad
 \sum_e n_{x,e}=d_{1,x}+d_{2,x}.
\]
For $n_{x,e}=2$ install one singular network $S_{x,e}$; it acts on
both symbol factors.  For $n_{x,e}=1$ install one varying network
$V^{\varepsilon_{x,e}}_{x,e}$.  Use disjoint auxiliary sets for all
these networks.  Let $\Gamma_y(\varepsilon)$ be the resulting graph
at address $y$, and consider the finite sum in $\mathbf F_2$
\[
 \sum_{\varepsilon}
       \prod_{y\ne0}\prod_{s=1}^2 A_s(\Gamma_y(\varepsilon)).
                                                               \tag{G.1a}
\]
The total degree available to all terminal polynomials is
$D=\sum_{y\ne0}(d_{1,y}+d_{2,y})$.

Each singular network at $x$ has zero constant at $x$, so its action
costs at least one difference in each of the two factors there, hence
at least two in (G.1a).  For a varying network, summing over its
$\varepsilon$ cancels the all-constant term.  Its first-degree
terms are independent of $\varepsilon$ away from its target, and
their difference at the target is $\partial_e$.  Thus every
surviving term costs at least one difference, and a cost-one term
acts at its target on exactly one of the two factors.  The sum of
all these minimum costs is exactly $D$.  Any higher-degree action,
or any additional action at a wrong address, exceeds $D$ and
annihilates at least one terminal polynomial.  Hence only the
minimum-cost terms survive.

At address $x$, a variable with exponent two is therefore
differentiated once in each factor; a variable with exponent one
is differentiated in exactly one factor, summed over the two
possibilities.  The resulting scalar is precisely
$[M_x](H_{1,x}H_{2,x})=1$.  Lower-degree parts cannot contribute.
Taking the product over all $x\ne0$ proves that (G.1a) equals one.
At least one choice of the finitely many $\varepsilon$ consequently
makes every factor in that product one.  This is one graph for both
symbols at every nonzero address.  Boolean reduction of
$H_{1,x}H_{2,x}$ before selecting $M_x$ would destroy this
exponent-two accounting and is not used.
Here is the stationary and one-symbol specialization, so that the
application in Proposition~\ref{ex:app-rough-companion} uses the same
proved construction. If all terminals are always active, attach
each column $h_i=Qe_i$ to any one fixed terminal. Since $h_i$ is
supported on its own group, their restricted sum is
$Q_J\mathbf1_J$ at every address. The two additional columns $c$
at a selected terminal pair cancel directly; no balancing terminal
of activation $\alpha+\beta$ is needed. The same singular program
therefore isolates that pair difference at its target address.
For one symbol, choose a coefficient-one top monomial of each
nonzero $P_x$ and put a separate singular program at $x$ for each
variable in that monomial. Its own programs already spend
$\deg P_x$ contraction degrees at $x$; any extra degree kills the
term. Thus exactly the selected top coefficient survives, and it
is one. A constant nonzero $P_x$ needs no program. This proves the
stationary one-symbol case as well as the moving one-symbol case.
In all cases the primary realization of
Lemma~\ref{ex:sec7-realization} changes only incidences with
the newly adjoined vertices. These are finite constructions for
each $b$, with no assertion of one infinite prime array.

\end{proof}

\begin{proposition}[One actual stage for the final family]\label{ex:app-final-family}
In the nonempty-seed case of Proposition~\ref{ex:app-nonempty-seeds},
fix the final $k,K$ and the stabilized common alphabet. For every
sufficiently large $b$, the address construction is realized at one
actual theta stage by a finite set $Q$ of distinct good odd primes
split in $K$, avoiding the fixed support and the specified finite
exclusions. With nonzero activations $\lambda_q$ it gives exactly
the family in (3.9),
\[
 h_x=\prod_{q\in Q}(q^*)^{\lambda_q(x)},\qquad h_0=1,
 \qquad q^*=(-1)^{(q-1)/2}q.
\]
All fixed local types are retained, and for every $x\ne0$,
\[
 \operatorname{ord}_{s=1}L(E^{h_x},s)=1,\quad L_*(h_xk)\ne0,
 \qquad j_{\rm gen}(h_x)\le\omega(h_x)+C_o,\quad
 v_2L_*(h_xk)\le2\omega(h_x)+C_e.
\]
The constants precede $b$. The finite graph budget and then the
actual theta precision are chosen after $b$. The theta prime remains
a coefficient prime and exclusion, not a relaxed Selmer prime.
\end{proposition}
\begin{proof}\leavevmode
For each chosen $b$, this construction has a finite budget $W_b$.
Enlarge it to bound every graph product and squarefree output in its
finite auxiliary tests, as well as every operator set. Contractions
only decrease the number of vertices. Thus every final $h_x$ lies in
the tested range of (F.7); this enlargement changes no fixed constant.
Only now take one actual theta stage in the finite intersection
of the selected table sets through $W_b+1$, the valid-stage set
for the full budget $|U|\le W_b$ with $R=2$, and the required
precision and finite exclusion sets.  In particular (F.6) holds
through the threshold $T$ used for that budget.
Realize its prescribed forward residue bits by distinct primes using CRT and Dirichlet;
quadratic reciprocity supplies the reverse bits.
There is no uniform bound on $W_b$, no intersection over all
budgets, and no product-surjectivity assumption on the alphabets.
The resulting one family satisfies (3.9), with constants fixed
before $b$. The theta prime $\ell_0$ is a coefficient prime and
an exclusion, not an active relaxed Selmer prime or a moving
arithmetic prime $r_i$.
\paragraph{Finite intersection and the upper bound.}
The required stage conditions form a finite list once the address
program is fixed. Let $\mathcal U$ be the chosen ultrafilter and
let $A_1,\ldots,A_t\in\mathcal U$ be the selected truth-table
sets used by that program, including all intermediate graphs.
Let $B\in\mathcal U$ impose (F.6) through the threshold $T$
for its operator budget and impose
$M\ge\max(0,\delta_o+W_b)+2$ and $M>T+\delta_o$.
Finite prime exclusions are included in $B$.
Then $B\cap\bigcap_{i=1}^t A_i\in\mathcal U$, hence is nonempty.
Choose a stage there before realizing the prime labels. The
coefficient rules hold for every realization at this stage, not
merely for a previously selected finite list of Fourier indices.
The address construction makes each nonzero address a unit of
both sources. For the odd source this says exactly
\[
 v_2a_M(h_x)=\delta_o+\omega(h_x)<M.
\]
The strict inequality follows from the tested-range precision
condition. Applying (O.8) gives the simple zero and
$j_{\rm gen}(h_x)\le\omega(h_x)+\delta_o+C_{\rm det}$.
The normalized even source gives
$v_2L_*(h_xk)\le2\omega(h_x)+C_{e,\rm final}$.
Thus one may take $C_o=\delta_o+C_{\rm det}$ and
$C_e=C_{e,\rm final}$, both fixed before $b$. Only the number
of tests, the threshold $T$, and the chosen stage depend on the
address dimension.
\end{proof}

\subsection{One coordinate in the final contradiction}
The family in Proposition~\ref{ex:app-final-family} is now used with
the same coordinate of Lemma~\ref{ex:arith-common-coordinate}
and its central calibration in
Proposition~\ref{ex:arith-central-calibration}. The denominator and
coefficient comparisons are those of
Propositions~\ref{ex:limits-fixed-field-clearing}
and~\ref{ex:limits-coefficient-transfer}; the zero-base assertion is
Lemma~\ref{ex:arith-base-zero}.
For this fixed $K$, (3.4), (F.7), and the even bound just proved give
\[
 2j_{\rm rc}(h_x,k)\le4\omega(h_x)+C_J,\qquad
 C_J=2C_o+C_{e,\rm final}-v_2L_*(h_*)+C_{H,\rm final}
\]
at every $x\ne0$.
The finite support comparison uses the same class, duality map,
bases and coefficient stages. Its integral correction
\[
 \eta=\prod_q(1+v_q(1+g_q)),\qquad l_q/e_q^2=1+2v_q,
 \qquad \widetilde u=\eta u_{\rm rel}
\]
is formed before characters. The actual integral determinant
triangle gives $\widetilde u_x=\epsilon_xu_{{\rm prim},x}$ with
$\epsilon_x\in\Lambda^\times$, rather than independently chosen
units in $\mathscr O^\times$. Therefore the central identity is
\[
 \begin{split}
 v_2\widetilde u_x(0)={}&2j_{\rm rc}(h_x,k)-4\omega(h_x)
       -s(h_x)-s(h_xk)\\
       &-\sum_{p\mid k}t_p+\epsilon'_x,\qquad
 0\le t_p\le2,\quad |\epsilon'_x|\le C_{\rm cal}.
 \end{split}                                                    \tag{F.8}
\]
The nonnegative finite $\Sha$ terms give one upper bound
$v_2\widetilde u_x(0)\le B$ for all $x\ne0$, independent of $b$.
At the unknown base no finiteness of $\Sha$ is assumed.
The retained central class and corank-one argument instead give
$\widetilde u_0(0)=0$ if the base derivative vanishes.

The fixed-field construction of Section~\ref{sec:limits} has
\[
 \begin{gathered}
 c_{\rm det}=a+4(h_K+2c_E+2t_K),\quad
 L_T(M)=16(d_K+5)M,\\
 f_x\widetilde u_x=d_x,\qquad
 c=v_2f_0(0)\le c_{\rm det}.
 \end{gathered}
\]
Choose $M>\max(c_{\rm det},B+2c_{\rm det},0)$, then
$b>(L_T(M)+2)(2^M-1)$ large enough for the address construction,
then its finite graph budget and actual theta stage $M_\theta$.
The residual comparison tests coefficients of $f$ from
$-L_T(M)$ through zero and only the constant coefficient of $d$.
It gives $x\ne0$ with
\[
       v_2(\widetilde u_x(0)-\widetilde u_0(0))\ge M-2c>B.
\]
The estimate uses both $v_{\mathscr O}(f_x)\le c$ and
$\widetilde u_x\in\Lambda[1/2]$; neither a bound on the whole
numerator nor a finite-Artin inverse of $t$ is asserted.
This contradicts (F.8) under the supposed base zero.
The Rankin factorization and $L(E^k,1)\ne0$ give the asserted
simple zero of $E$.

This completes the proof of the restricted converse,
Theorem~\ref{hyp:RC}, from the displayed classical inputs and
the arithmetic, coefficient and determinant constructions.
The conversion from that converse to pole witnesses, the integer
model, the natural compiler and finally the ordinary rational value
is the separate argument in the main text.

% End 09-normalization-endpoints.tex


\section{De Rham comparison and finite descent in the height argument}
\label{app:de-rham-descent}

\begin{lemma}\label{lem:de-rham-descent}
Let $K/\mathbb Q_\ell$ and $F/K$ be finite extensions, let $A/F$
be an abelian surface, and let $C_\ell/\mathbb Q_\ell$ be a finite
coefficient field. Suppose a continuous representation
$r:G_K\longrightarrow\mathrm{GL}_2(C_\ell)$ satisfies
\[
 V_\ell(A)\otimes_{\mathbb Q_\ell}C_\ell
       \simeq r|_{G_F}\oplus r|_{G_F}.
\]
Then $r$ is de Rham. At every coefficient embedding its
Hodge--Tate weights are $0,-1$, with
$\mathbb Q_\ell(1)$ assigned weight $-1$.
\end{lemma}

\begin{proof}
The smooth proper variety $A$ has a smooth proper rigid
analytification. Apply \cite[Corollary 1.8]{scholze2013}, with the
corrections in \cite{scholze-erratum}. The passage from algebraic
to analytic cohomology uses \cite[Theorems 9.1 and 9.3]{scholze2013}:
Theorem 9.1 identifies coherent cohomology, and hence the filtered
de Rham hypercohomology, while Theorem 9.3 identifies the
$\mathbb Z/\ell^n\mathbb Z$-cohomology, compatibly in $n$.
Its invertibility condition holds because $F$ has characteristic
zero. Taking inverse limits and then inverting $\ell$ identifies
the algebraic and analytic $\ell$-adic cohomology, compatibly with
the action of $G_F$. Thus
$H^1_{\mathrm{et}}(A_{\overline F},\mathbb Q_\ell)$ is de Rham.

The canonical, Galois-equivariant identification
\cite[Part I, Theorem 12.1(a)]{milne-av}
\[
 V_\ell(A)\simeq
 H^1_{\mathrm{et}}(A_{\overline F},\mathbb Q_\ell)^\vee
\]
gives the same property for the Tate module. De Rham representations
are closed under duals and subquotients
\cite[Theorem 5.2.1 and Proposition 6.3.3]{brinon-conrad}.
Extension of the coefficient field is, on underlying
$\mathbb Q_\ell$-representations, a finite direct sum.
The displayed decomposition therefore makes $r|_{G_F}$ de Rham.

Here finite-extension descent requires no semistable descent
assertion. To see it explicitly, pass to a finite Galois extension
$F'/K$ containing $F$, and let $W$ be the underlying
$\mathbb Q_\ell$-representation of $r$. The vector space
\[
 D_{F'}=(B_{\mathrm{dR}}\otimes_{\mathbb Q_\ell}W)^{G_{F'}}
\]
has its natural semilinear $\mathrm{Gal}(F'/K)$-action, and its
invariants are
$D_K=(B_{\mathrm{dR}}\otimes_{\mathbb Q_\ell}W)^{G_K}$.
Finite Galois descent gives $F'\otimes_K D_K\simeq D_{F'}$,
also on each filtration step. Since $W|_{G_{F'}}$ is de Rham,
$\dim_K D_K=\dim_{F'}D_{F'}=\dim_{\mathbb Q_\ell}W$.
This proves the assertion over $K$; it is the finite-extension
case of \cite[Proposition 6.3.8]{brinon-conrad}.

The two Hodge numbers of $H^1_{\mathrm{dR}}(A/F)$ are $2,2$.
Consequently the Tate module has weights $0,0,-1,-1$.
At each coefficient embedding, its two isomorphic summands
have the same weight multiset, so each summand has weights $0,-1$.
The filtration-compatible descent preserves these weights.
\end{proof}

The Hodge--Tate comparison in \cite[Chapter III, Theorem 4.1]{faltings1988}
alone would not prove the de Rham assertion. The preceding argument
uses the de Rham comparison itself and does not infer de Rhamness
from Hodge--Tate weights.

% Begin references.tex
\begin{thebibliography}{99}
\bibitem{scholze2013} P. Scholze,
\emph{$p$-adic Hodge theory for rigid-analytic varieties},
Forum of Mathematics, Pi 1 (2013), e1, 1--77,
Corollary 1.8 and Theorems 9.1, 9.3.
\url{https://doi.org/10.1017/fmp.2013.1}.
\bibitem{scholze-erratum} P. Scholze,
\emph{$p$-adic Hodge theory for rigid-analytic varieties---corrigendum},
Forum of Mathematics, Pi 4 (2016), e6, 1--4.
\url{https://doi.org/10.1017/fmp.2016.4}.
\bibitem{brinon-conrad} O. Brinon and B. Conrad,
\emph{CMI summer school notes on $p$-adic Hodge theory},
preliminary version, June 24, 2009,
Theorem 5.2.1 and Propositions 6.3.3, 6.3.8.
\url{https://math.stanford.edu/~conrad/papers/notes.pdf}.
\bibitem{milne-av} J. S. Milne, \emph{Abelian varieties},
version 2.00 (2008), Part I, Theorem 12.1(a).
\url{https://www.jmilne.org/math/CourseNotes/AV.pdf}.
\bibitem{babagranath} S. Baba and H. Granath,
\emph{Genus 2 curves with quaternionic multiplication}, 2005 preprint,
Sections 2.1, 3 and Proposition 10.
\url{https://webdoc.sub.gwdg.de/ebook/serien/e/mpi_mathematik/2005/18.pdf}.
\bibitem{buzzard} K. Buzzard, \emph{Integral models of certain Shimura
curves}, 1996, corrected 2002, Theorem 2.1 and Corollary 2.3.
\url{https://www.ma.imperial.ac.uk/~buzzard/maths/research/papers/shimura.pdf}.
\bibitem{linyang} Y.-H. Lin and Y. Yang, \emph{Quaternionic loci in
Siegel's modular threefold}, arXiv:1807.00466v1, 2018, equations (2), (14).
\url{https://arxiv.org/abs/1807.00466v1}.
\bibitem{berger2002} L. Berger, \emph{Repr\'esentations $p$-adiques et
\'equations diff\'erentielles}, Inventiones Mathematicae 148 (2002),
219--284, $p$-adic monodromy theorem.
\url{https://doi.org/10.1007/s002220100202}.
\bibitem{pan2022} L. Pan, \emph{The Fontaine--Mazur conjecture in the
residually reducible case}, arXiv:1901.07166v2,
Theorem 1.0.4, the version used here. Published in Journal of the
American Mathematical Society 35 (2022), 1031--1169.
\url{https://arxiv.org/abs/1901.07166v2}.
\bibitem{conrad-local-characters} B. Conrad,
\emph{Lifting global representations with local properties},
author preprint, December 11, 2011, Appendix B, pp.~32--33.
\url{https://math.stanford.edu/~conrad/papers/locchar.pdf}.
\bibitem{li-newforms} W.-C. W. Li,
\emph{Newforms and functional equations}, Mathematische Annalen
212 (1975), 285--315, p.~296.
\bibitem{shimura-factors} G. Shimura,
\emph{On the factors of the jacobian variety of a modular function field},
Journal of the Mathematical Society of Japan 25 (1973), 523--544,
Theorem 1.
\url{https://doi.org/10.2969/jmsj/02530523}.
\bibitem{javanpeykar} A. Javanpeykar, with an appendix by P. Bruin,
\emph{Polynomial bounds for Arakelov invariants of Belyi curves},
Algebra \& Number Theory 8 (2014), 89--140,
Theorem 1.1.1 and Section 2.3.
\url{https://doi.org/10.2140/ant.2014.8.89}.
\bibitem{gr2014} \'E. Gaudron and G. R\'emond,
\emph{Polarisations et isog\'enies}, Duke Mathematical Journal
163 (2014), 2057--2108, Theorem 1.4 and Section 9.
\url{https://doi.org/10.1215/00127094-2782528}.
\bibitem{fowler} G. Fowler, \emph{Some uniform effective results on
Andr\'e--Oort for sums of powers in $\mathbb C^n$}, Mathematical
Proceedings of the Cambridge Philosophical Society 180 (2026),
607--641, Proposition 2.7.
\url{https://doi.org/10.1017/S0305004125101825}.
\bibitem{cmheight} A. Barquero-Sanchez, L. Cadwallader, O. Cannon,
T. Genao and R. Masri, \emph{Faltings heights of CM elliptic curves
and special Gamma values}, Research in Number Theory 3 (2017),
article 13, Section 6, equation (6.2).
\url{https://doi.org/10.1007/s40993-017-0077-7}.
\bibitem{silverbergzarhin} A. Silverberg and Yu. G. Zarhin,
\emph{Semistable reduction and torsion subgroups of abelian varieties},
Annales de l'Institut Fourier 45 (1995), 403--420, Theorem 3.5.
\url{https://www.numdam.org/article/AIF_1995__45_2_403_0.pdf}.
\bibitem{serrelocal} J.-P. Serre, \emph{Local class field theory},
in J. W. S. Cassels and A. Fr\"ohlich (eds.), \emph{Algebraic Number Theory},
Academic Press, 1967, Chapter VI, 128--161, Sections 4.1--4.3.
\bibitem{faltings1983} G. Faltings, \emph{Endlichkeitss\"atze f\"ur
abelsche Variet\"aten \"uber Zahlk\"orpern}, Inventiones Mathematicae
73 (1983), 349--366; erratum, 75 (1984), 381.
\bibitem{faltings1988} G. Faltings, \emph{$p$-adic Hodge theory},
Journal of the American Mathematical Society 1 (1988), 255--299.

\bibitem{lny} K. Lauter, M. Naehrig and T. Yang, \emph{Hilbert theta
series and invariants of genus 2 curves}, Journal of Number Theory,
2015, Section 2.2, doi:10.1016/j.jnt.2015.02.020.
\url{https://people.math.wisc.edu/~tonghaiyang/2015-LNYFinal-JNT.pdf}.
\bibitem{pazuki} F. Pazuki, \emph{Theta height and Faltings height},
Bulletin de la Soci\'et\'e Math\'ematique de France, 2012,
Sections 2.2--2.3 and Corollary 1.3.
\url{https://www.numdam.org/item/10.24033/bsmf.2623.pdf}.
\bibitem{poonenstoll} B. Poonen and M. Stoll, \emph{The Cassels--Tate
pairing on polarized abelian varieties}, Annals of Mathematics 150
(1999), 1109--1149. \url{https://math.mit.edu/~poonen/papers/sha.pdf}.
\bibitem{greentao} B. Green and T. Tao, \emph{Linear equations in primes},
Annals of Mathematics 171 (2010), 1753--1850.
\url{https://annals.math.princeton.edu/2010/171-3/p08}.
\bibitem{moebius} B. Green and T. Tao, \emph{The M\"obius function is
strongly orthogonal to nilsequences}, Annals of Mathematics 175 (2012).
\url{https://annals.math.princeton.edu/2012/175-2/p03}.
\bibitem{gtz} B. Green, T. Tao and T. Ziegler, \emph{An inverse theorem
for the Gowers $U^{s+1}[N]$-norm}, arXiv:1009.3998v5, 2026 corrected version.
\url{https://arxiv.org/abs/1009.3998v5}; author erratum:
\url{https://people.maths.ox.ac.uk/greenbj/papers/GI(s)-erratum.pdf}.
\bibitem{bmor} M. A. Bennett, G. Martin, K. O'Bryant and A. Rechnitzer,
\emph{Explicit bounds for primes in arithmetic progressions},
Illinois Journal of Mathematics 62 (2018), 427--532,
arXiv:1802.00085v3, Theorem 1.2.
\url{https://arxiv.org/abs/1802.00085v3}.
\bibitem{milne} J. S. Milne, \emph{Arithmetic duality theorems},
second edition, 2006, Chapter I.
\url{https://www.jmilne.org/math/Books/ADTnot.pdf}.
\bibitem{stackscohom} The Stacks Project,
\emph{Stalks of higher direct images} and \emph{The Leray spectral sequence},
Theorem 59.53.1, Tag 03Q9, and Proposition 59.54.2, Section Tag 03QA,
HTML snapshots of November 5, 2026.
\url{https://stacks.math.columbia.edu/tag/03Q9};
\url{https://stacks.math.columbia.edu/tag/03QA}.
\bibitem{nekovar} J. Nekov\'a\v r, \emph{Selmer complexes},
Ast\'erisque 310, 2006, Proposition 5.2.4, Theorems 6.3.4 and 7.8.8, and Sections 5.7 and 6.4.
\url{https://www.numdam.org/item/AST_2006__310__1_0/}.
\bibitem{afp} Archive of Formal Proofs, \emph{DPRM\_Theorem},
\texttt{Alpha\_Sequence}, \texttt{Exponential\_Relation}, and
\texttt{Exponentiation}, directory
\texttt{thys/DPRM\_Theorem/Diophantine/}, AFP development git mirror
revision \texttt{b4f3b892d0ae25cea27412ba4007cdb63e44b0eb}
(May 10, 2026). The three source files match this revision byte for byte.
\url{https://isa-afp.org/entries/DPRM_Theorem.html}.
\bibitem{edixhoven}
B. Edixhoven, \emph{N\'eron models and tame ramification},
Compositio Mathematica 81 (1992), 291--306,
Theorem~5.3 and Remark~5.4.1.
\url{https://www.numdam.org/item/CM_1992__81_3_291_0/}.
\bibitem{sutherland} A. V. Sutherland,
\emph{Isogeny kernels and division polynomials}, MIT 18.783,
Fall 2025, Lecture 5, September 18, 2025, \S\S5.5--5.6.
\url{https://ocw.mit.edu/courses/18-783-elliptic-curves-fall-2025/mit18_783_f25_lec05.pdf}.
\bibitem{silverman} J. H. Silverman, \emph{The arithmetic of elliptic
curves}, second edition, Springer, 2009, Chapters IV and VII.
\bibitem{dd} T. Dokchitser and V. Dokchitser, \emph{Root numbers and parity
of ranks of elliptic curves}, Journal f\"ur die reine und angewandte
Mathematik 658 (2011), 39--64, Theorem 1.8.
\url{https://arxiv.org/abs/0906.1815v1}.
\bibitem{bcdt} C. Breuil, B. Conrad, F. Diamond and R. Taylor,
\emph{On the modularity of elliptic curves over $\Q$: wild 3-adic exercises},
author manuscript, Theorem A and Introduction, conditions (1)--(6).
\url{https://virtualmath1.stanford.edu/~conrad/papers/tswfinal.pdf}.
Published in Journal of the American Mathematical Society 14 (2001), 843--939.
\bibitem{milnemf} J. S. Milne, \emph{Modular functions and modular forms},
v1.31, March 22, 2017, Remark 11.18 and Theorem 11.19.
\url{https://www.jmilne.org/math/CourseNotes/MF.pdf}.
\bibitem{elkik} R. Elkik, \emph{Le th\'eor\`eme de Manin--Drinfeld},
Ast\'erisque 183 (1990), 59--67, Theorem, p.~59.
\url{https://www.numdam.org/item/AST_1990__183__59_0.pdf}.
\bibitem{gz} B. Gross and D. Zagier, \emph{Heegner points and derivatives
of $L$-series}, Inventiones Mathematicae 84 (1986), 225--320.
\bibitem{kolyvagin} V. A. Kolyvagin,
\emph{On the Mordell--Weil group and the Shafarevich--Tate group
of modular elliptic curves}, MPIM preprint MPI/90-69 (1990),
Theorems 1--5, pp.~4--5.
\url{https://archive.mpim-bonn.mpg.de/id/eprint/325/1/preprint_1990_69.pdf}.
\bibitem{kato} K. Kato, \emph{$p$-adic Hodge theory and values of zeta
functions of modular forms}, Ast\'erisque 295 (2004), 117--290.
\url{https://www.numdam.org/article/AST_2004__295__117_0.pdf}.
\bibitem{bfh} D. Bump, S. Friedberg and J. Hoffstein,
\emph{Nonvanishing theorems for $L$-functions of modular forms and their
derivatives}, Inventiones Mathematicae 102 (1990), 543--618.
\url{https://wstein.org/papers/bib/bump-friedberg-hoffstein-nonvanishing.pdf}.
\bibitem{hl} J. Hoffstein and W. Luo, \emph{Nonvanishing of $L$-series and
the combinatorial sieve}, Mathematical Research Letters 4 (1997), 435--444.
\url{https://intlpress.com/site/pub/files/_fulltext/journals/mrl/1997/0004/0003/MRL-1997-0004-0003-a012.pdf}.
\bibitem{desjardins} J. Desjardins, \emph{Root number of twists of an
elliptic curve}, Journal de Th\'eorie des Nombres de Bordeaux
32 (2020), 73--101. \url{https://jtnb.centre-mersenne.org/item/10.5802/jtnb.1112.pdf}.
\bibitem{cst} L. Cai, J. Shu and Y. Tian, \emph{Explicit Gross--Zagier
and Waldspurger formulae}, Algebra \& Number Theory 8 (2014).
\url{https://arxiv.org/abs/1408.1733}.
\bibitem{ccl} L. Cai, Y. Chen and Y. Liu, \emph{Heegner points on modular
curves}, Transactions of the American Mathematical Society 370 (2018),
3721--3743. \url{https://archive.ymsc.tsinghua.edu.cn/pacm_download/632/12220-Heegner_points_on_modular_curves.pdf}.
\bibitem{jetchevkane} D. Jetchev and B. Kane,
\emph{Equidistribution of Heegner Points and Ternary Quadratic Forms},
arXiv:0908.3905 (2009), Sections 1.1--1.2 and 2.5.
\url{https://arxiv.org/abs/0908.3905}.
\bibitem{cornutjetchev} C. Cornut and D. Jetchev,
\emph{Liftings of reduction maps for quaternion algebras}, 2010 preprint.
\url{https://webusers.imj-prg.fr/~christophe.cornut/papers/liftings.pdf}.
\bibitem{waldspurger} J.-L. Waldspurger, \emph{Sur les coefficients de
Fourier des formes modulaires de poids demi-entier},
Journal de Math\'ematiques Pures et Appliqu\'ees 60 (1981), 375--484,
Proposition 18 and Corollary 2.
\bibitem{ramsey} N. Ramsey, \emph{Geometric and $p$-adic modular forms
of half-integral weight}, Annales de l'Institut Fourier 56 (2006),
599--624. \url{https://www.numdam.org/item/10.5802/aif.2195.pdf}.
\bibitem{casselman}
W. Casselman, \emph{On some results of Atkin and Lehner},
Mathematische Annalen 201 (1973), 301--314,
Theorem~1 and its proof, pp.~302--306, and Theorem~4, p.~313.
\bibitem{carayol} H. Carayol, \emph{Sur les repr\'esentations
$\ell$-adiques associ\'ees aux formes modulaires de Hilbert},
Annales Scientifiques de l'\'Ecole Normale Sup\'erieure 19 (1986), 409--468,
Theorem A. \url{https://www.numdam.org/item/ASENS_1986_4_19_3_409_0/}.
\bibitem{waldspurgerICM} J.-L. Waldspurger,
\emph{Correspondances de Shimura}, Proceedings of the International
Congress of Mathematicians, 1983, 525--531.
\url{https://www.imj-prg.fr/wp-content/uploads/2020/prix/waldspurger1983.pdf}.
\bibitem{weil} A. Weil, \emph{Sur certains groupes d'op\'erateurs unitaires},
Acta Mathematica 111 (1964), 143--211, I.13, II.25, II.28, III, Theorem 6.
\url{https://webhomes.maths.ed.ac.uk/~v1ranick/papers/weil2.pdf}.
\bibitem{serre} J.-P. Serre, \emph{Propri\'et\'es galoisiennes des points
d'ordre fini des courbes elliptiques}, Inventiones Mathematicae 15
(1972), 259--331.
\bibitem{mazur1995} B. Mazur,
\emph{Speculations about the topology of rational points: an up-date},
Ast\'erisque 228 (1995), 165--181, Conjecture 3, p.~165.
\url{https://www.numdam.org/article/AST_1995__228__165_0.pdf}.
\bibitem{matiyasevich1971} Yu. V. Matiyasevich,
\emph{Diophantine representation of enumerable predicates},
Mathematics of the USSR-Izvestiya 5 (1971), 1--28,
Basic Theorem, p.~3; integer-witness conversion, p.~2.
\url{https://doi.org/10.1070/IM1971v005n01ABEH001004}.
\bibitem{openaihtp} OpenAI,
\emph{Hilbert's tenth problem over the rational numbers},
OpenAI Math Release preprint, 6 October 2026.
\url{https://github.com/openai/math/blob/main/preprints/Hilberts-tenth-problem-over-the-rational-numbers-October-6-2026/main.pdf}.
\bibitem{howard2012} B. Howard, \emph{The Heegner point Kolyvagin system},
Compositio Mathematica 140 (2004), 1439--1472; arXiv:1202.6340v1.
\bibitem{lang1956} S. Lang, \emph{Algebraic groups over finite fields},
American Journal of Mathematics 78 (1956), 555--563.
\bibitem{williams-mertens}
K. S. Williams, \emph{Mertens' theorem for arithmetic progressions},
J. Number Theory \textbf{6} (1974), 353--359, Theorem 1.
\url{https://people.math.carleton.ca/~williams/papers/pdf/057.pdf}.

\bibitem{voight-quaternion}
J. Voight, \emph{Quaternion algebras}, corrected online edition
v1.0.7u, August 5, 2026, Main Theorem 14.6.1 and Proposition 14.6.7,
pp. 223--224.
\url{https://jvoight.github.io/quat-book.pdf}.


% Insert inside the existing thebibliography environment.
% This is the only new bibliography key used by assembled-original-extension.tex.
\bibitem{koenigsmann-defining} J. Koenigsmann,
\emph{Defining $\mathbb Z$ in $\mathbb Q$},
Annals of Mathematics \textbf{183} (2016), 73--93.
\href{https://doi.org/10.4007/annals.2016.183.1.2}{doi:10.4007/annals.2016.183.1.2}.

% Proposed additions only. No existing bibliography entry is replaced.
\bibitem{GonzalezRotger2006} J. Gonz\'alez and V. Rotger,
\emph{Non-elliptic Shimura curves of genus one},
Journal of the Mathematical Society of Japan \textbf{58} (2006), 927--948,
Theorem 4.2 and Table 2, pp.~938--939; Proposition 5.2, pp.~940--941.
\url{https://doi.org/10.2969/jmsj/1179759530}.

\bibitem{Yuan2024} X. Yuan,
\emph{Explicit Kodaira--Spencer maps over Shimura curves},
Acta Mathematica Sinica, Chinese Series \textbf{67} (2024), 227--249.
\url{https://doi.org/10.12386/A20220154}.
We use the English translation dated April 30, 2024, identified by
the author as a translation of the published Chinese article:
\url{https://faculty.bicmr.pku.edu.cn/~yxy/preprints/explicit_KS.pdf}.
The author identifies it at
\url{https://faculty.bicmr.pku.edu.cn/~yxy/publications.html}.
The relevant locators in that translation are Theorem 1.1, p.~3;
metric definitions, pp.~2--3; Sections 2.1--2.2 and Lemma 2.1,
pp.~4--10; and the analytic and metric calculations in Sections
4.2--4.3, pp.~20--26.

\bibitem{Deligne1979} P. Deligne,
\emph{Vari\'et\'es de Shimura: interpr\'etation modulaire, et techniques
de construction de mod\`eles canoniques},
in \emph{Automorphic Forms, Representations and $L$-Functions},
Proceedings of Symposia in Pure Mathematics \textbf{33}, Part 2 (1979),
247--289, Sections 2.2.5 and 2.3.1.
English translation by J. S. Milne, corrected September 6, 2020, p.~29:
\url{https://www.jmilne.org/math/Documents/DeligneSV.pdf}.

\bibitem{GaudronRemond2023} \'E. Gaudron and G. R\'emond,
\emph{Nouveaux th\'eor\`emes d'isog\'enie},
M\'emoires de la Soci\'et\'e Math\'ematique de France \textbf{176} (2023),
vi+123 pp., Section 1.5 and Theorem 1.9(1), pp.~7--8.
\url{https://doi.org/10.24033/msmf.484}.

\bibitem{vonKaenel2021} R. von K\"anel,
\emph{The effective Shafarevich conjecture for abelian varieties of
$\mathrm{GL}_2$-type}, Forum of Mathematics, Sigma \textbf{9} (2021), e39,
1--29, equations (2.2) and (5.9), pp.~5 and 13.
\url{https://doi.org/10.1017/fms.2021.29}.

\bibitem{voight2021} J. Voight, \emph{Quaternion algebras},
Graduate Texts in Mathematics \textbf{288}, Springer, 2021,
Chapter 43, equations (43.1.2)--(43.1.3), p.~800, and
Main Theorem 43.6.14, p.~819; proof pp.~822--823.
\url{https://doi.org/10.1007/978-3-030-56694-4_43}.

\bibitem{calegari-icm2022} F. Calegari,
\emph{Reciprocity in the Langlands program since Fermat's Last Theorem},
in \emph{Proceedings of the International Congress of Mathematicians 2022},
Vol.~2, EMS Press, 2023, 610--651, Section 7, p.~624;
definition of geometric representations, pp.~611--612.
\url{https://doi.org/10.4171/ICM2022/101}.

\bibitem{diamond-im} F. Diamond and J. Im,
\emph{Modular forms and modular curves}, in Seminar on Fermat's Last Theorem,
CMS Conference Proceedings 17 (1995), 39--133, \S12.5.
\bibitem{baker-et-al-finiteness} M. H. Baker, E. Gonz\'alez-Jim\'enez,
J. Gonz\'alez and B. Poonen, \emph{Finiteness results for modular curves
of genus at least 2}, American Journal of Mathematics 127 (2005),
1325--1387, \S\S3.1, 3.2.1.
\bibitem{mascot-moduli} N. Mascot, \emph{Moduli-friendly Eisenstein series
over the p-adics and the computation of modular Galois representations},
Research in Number Theory 8 (2022), article 48,
\href{https://doi.org/10.1007/s40993-022-00329-6}{doi:10.1007/s40993-022-00329-6},
\S2.1, (2.1.8).
\end{thebibliography}

% End references.tex

\clearpage
\section*{AI usage statement}
This manuscript was developed primarily by OpenAI Codex agents (Astra 6 and Sol 6.1)
and Anthropic Claude agents (Opus 5.5), with auditing from Fable 5.1.
Their contributions included mathematical
exploration, construction and analysis of the polynomial, proof drafting,
literature searches, computational checks, internal review, and drafting
and revision of the exposition. Reviews by multiple AI agents formed part
of this process; they do not constitute independent human peer review.
\end{document}
